\documentclass[11pt,openany]{book}
\setcounter{secnumdepth}{3}
\usepackage[margin=1.1in]{geometry}
\usepackage[T1]{fontenc}
\usepackage[utf8]{inputenc}
\usepackage{lmodern}
\usepackage[english]{babel}
\usepackage{amsmath,amssymb,amsthm}
\usepackage{graphicx}
\usepackage{tikz}
\usetikzlibrary{arrows.meta,decorations.markings,decorations.pathmorphing,decorations.pathreplacing,calc}
\input{figures/icons}
\usepackage[unicode,colorlinks=true,linkcolor=blue!60!black,citecolor=blue!60!black]{hyperref}
\usepackage{microtype}
\setlength{\emergencystretch}{3em} % long names (licences, references) in narrow lines
\usepackage{empheq}
\usepackage{array}
\usepackage{booktabs}
\usepackage{multirow}
\usepackage{longtable}
\usepackage{circuitikz}
\definecolor{keyyellow}{rgb}{1.0,0.97,0.78}
% key equations: \begin{empheq}[box=\keybox]{equation} ... \end{empheq} -- light-yellow box, number stays outside
\newcommand{\keybox}[1]{{\setlength{\fboxsep}{7pt}\colorbox{keyyellow}{#1}}}
\usepackage[most]{tcolorbox}
% key statements (propositions etc.): the same light-yellow background, breakable across pages
\newtcolorbox{keyblock}{enhanced,breakable,colback=keyyellow,colframe=keyyellow,boxrule=0pt,arc=0pt,left=6pt,right=6pt,top=2pt,bottom=2pt,boxsep=0pt}
\renewcommand{\chaptermark}[1]{\markboth{\small\chaptername~\thechapter.\ #1}{}}
\renewcommand{\sectionmark}[1]{\markright{\small\thesection.\ #1}}
\renewcommand{\topfraction}{0.95}
\renewcommand{\textfraction}{0.05}
\renewcommand{\floatpagefraction}{0.75}

% part pages: title at the top third, and text may follow on the same page (used for the part introductions)
\makeatletter
\renewcommand\part{\if@openright\cleardoublepage\else\clearpage\fi\thispagestyle{plain}\if@twocolumn\onecolumn\@tempswatrue\else\@tempswafalse\fi\null\vspace{0.18\textheight}\secdef\@part\@spart}
\renewcommand\@endpart{\par\vspace{3em}\if@tempswa\twocolumn\fi}
\makeatother
\newtheorem{theorem}{Theorem}[chapter]
\newtheorem{proposition}{Proposition}[chapter]
\newtheorem{lemma}{Lemma}[chapter]
\theoremstyle{remark}
\newtheorem{remark}{Remark}[chapter]
% exercises: \begin{exercise}[\lvlB] ... \end{exercise}, numbered "Exercise 4.3"; difficulty marks
\theoremstyle{definition}
\newtheorem{exercise}{Exercise}[chapter]
\newcommand{\lvlA}{$\star$}
\newcommand{\lvlB}{$\star\star$}
\newcommand{\lvlC}{$\star\star\star$}
% answer to an exercise in the Answers appendix: \answer{ex:NN:k}{text}
\newcommand{\answer}[2]{\par\noindent\textbf{\ref{#1}.}~#2\par\smallskip}
\newcommand{\dd}{\mathrm{d}}
% neuron type code: first four letters of the primary mediator
\newcommand{\nt}[1]{\textsf{\textbf{#1}}}
\newcommand{\mV}{\,\text{mV}}
\newcommand{\ms}{\,\text{ms}}
\newcommand{\mM}{\,\text{mM}}
\newcommand{\um}{\,\text{\textmu m}}
\newcommand{\ion}[2]{\mathrm{#1}^{#2}}
\newcommand{\Na}{\ion{Na}{+}}
\newcommand{\K}{\ion{K}{+}}
\newcommand{\Cl}{\ion{Cl}{-}}
\newcommand{\Ca}{\ion{Ca}{2+}}
\newcommand{\conc}[1]{[\mathrm{#1}]}

\title{\Huge The 20-Watt Computer\\[16pt] \Large How the Brain Computes: A Quantitative Approach for Engineers}
\hypersetup{pdftitle={The 20-Watt Computer. How the Brain Computes: A Quantitative Approach for Engineers},pdfauthor={Konstantin Kladko}}
\author{\Large Konstantin~Kladko}
\date{}

\begin{document}
\frontmatter
\maketitle
\thispagestyle{empty}
\vspace*{\fill}
\noindent{\small \copyright~2026 Konstantin~Kladko.\\[4pt]
The text, figures, exercises and answers are licensed under the Creative Commons Attribution-NonCommercial-ShareAlike 4.0 International License (CC BY-NC-SA 4.0): \url{https://creativecommons.org/licenses/by-nc-sa/4.0/}. You may copy the book, teach with it, translate and adapt it for non-commercial purposes, provided you credit the author and share derivative works under the same license.\\[4pt]
The model and build programs are licensed under the MIT License.\\[4pt]
Translated from the Russian edition.\\[4pt]
Sources, models and new versions: \url{https://github.com/kladkogex/20-watt-computer}.}
\clearpage

\tableofcontents
\chapter*{How to Read This Book}
\addcontentsline{toc}{chapter}{How to Read This Book}

You do not have to read this book from cover to cover. Chapters~1--4 are the common foundation: the neuron types, what \nt{GLUT} and \nt{GABA} compute, and the language they speak. After that the book branches. The map in Fig.~\ref{fig:roadmap} shows which chapters build on which: an arrow from chapter $A$ to chapter $B$ means that $B$ uses concepts and formulas from $A$. The map shows only the main dependencies; minor back-references do not get in the way.

\begin{figure}[h]
\centering
\begin{tikzpicture}[>=Stealth,scale=0.95,
  ch/.style={draw,thick,rounded corners=3pt,align=center,font=\scriptsize,text width=1.85cm,minimum height=0.62cm,inner sep=2pt},
  base/.style={ch,fill=blue!8}, bus/.style={ch,fill=orange!14}, sum/.style={ch,fill=gray!18},
  io/.style={ch,fill=green!10}, sort/.style={ch,fill=cyan!8}, lab/.style={ch,fill=violet!10},
  dep/.style={->,thick,gray!60!black}]
% foundation
\node[base] (c1) at (0,0) {1 Neuron types};
\node[base] (c2) at (2.3,0) {2 \nt{GLUT}};
\node[base] (c3) at (4.6,0) {3 \nt{GLUT} and \nt{GABA}};
\node[base] (c4) at (6.9,0) {4 Morse code};
% broadcast channels and the synthesis
\node[bus] (c5) at (4.6,-1.5) {5 \nt{SERO}};
\node[bus] (c6) at (7.3,-1.5) {6 \nt{DOPA}};
\node[bus] (c7) at (9.2,-0.9) {7 Dopamine theories};
\node[bus] (c8) at (9.2,-2.1) {8 \nt{NORA}};
\node[bus] (c9) at (11.5,-2.1) {9 \nt{ACET}};
\node[sum] (c10) at (13.8,-0.9) {10 The whole machine};
% inputs and outputs
\node[io] (c12) at (2.3,-4.1) {12 Recog-\\nition};
\node[io] (c11) at (4.6,-4.1) {11 Inputs};
\node[io] (c13) at (6.9,-4.1) {13 Muscles};
\node[io] (c14) at (9.2,-3.05) {14 Pain};
\node[io] (c15) at (9.2,-4.1) {15 Motor\\learning};
\node[io] (c16) at (9.2,-5.15) {16 Body\\chemistry};
% lab, sorting, sleep
\node[lab] (c17) at (11.5,-4.1) {17 Debugging};
\node[lab] (c21) at (13.8,-4.1) {21 Living machines};
\node[lab] (c22) at (13.8,-5.15) {22 DishBrain};
\node[sort] (c18) at (11.5,-5.15) {18 Neurons as\\instruments};
\node[sort] (c19) at (13.8,-6.2) {19 Dendrite\\count};
\node[bus] (c20) at (11.5,-6.2) {20 Sleep};
% dependencies
\draw[dep] (c1) -- (c2); \draw[dep] (c2) -- (c3); \draw[dep] (c3) -- (c4);
\draw[dep] (c1) -- (c5); \draw[dep] (c4) -- (c6); \draw[dep] (c5) -- (c6);
\draw[dep] (c6) -- (c7); \draw[dep] (c6) -- (c8); \draw[dep] (c8) -- (c9);
\draw[dep] (c7) -- (c10); \draw[dep] (c9) -- (c10);
\draw[dep] (c4) -- (c11); \draw[dep] (c11) -- (c12); \draw[dep] (c11) -- (c13); \draw[dep] (c5) -- (c13);
\draw[dep] (c13) -- (c14); \draw[dep] (c13) -- (c15); \draw[dep] (c13) -- (c16);
\draw[dep] (c15) -- (c17); \draw[dep] (c9) -- (c17); \draw[dep] (c17) -- (c21); \draw[dep] (c21) -- (c22);
\draw[dep] (c16) -- (c18); \draw[dep] (c18) -- (c19); \draw[dep] (c16) -- (c20);
% legend
\begin{scope}[yshift=-7.25cm]
\foreach \s/\t [count=\i] in {base/foundation, bus/channels and modes, sum/synthesis, io/inputs and outputs, sort/neuron classifications, lab/laboratory}{
  \pgfmathtruncatemacro{\col}{mod(\i-1,3)}\pgfmathtruncatemacro{\row}{(\i-1)/3}
  \node[\s,text width=0.3cm,minimum height=0.32cm] at ({\col*4.8+0.2},{-\row*0.55}) {};
  \node[anchor=west,font=\scriptsize] at ({\col*4.8+0.55},{-\row*0.55}) {\t};}
\end{scope}
\end{tikzpicture}
\caption{Chapter map. An arrow leads from a chapter to the one that builds on it. Other cross-references between chapters are pointers, not dependencies.}
\label{fig:roadmap}
\end{figure}

\section*{Paths Through the Book}

\begin{center}
\footnotesize
\setlength{\tabcolsep}{4pt}
\renewcommand{\arraystretch}{1.2}
\begin{tabular}{>{\raggedright}p{3.0cm}>{\raggedright}p{3.9cm}>{\raggedright\arraybackslash}p{6.4cm}}
\toprule
Path & For whom & Chapters in order \\
\midrule
one-quarter course & instructor, 10 weeks & week 1: ch.~1--2; 2: ch.~3; 3: ch.~4; 4: ch.~5--6; 5: ch.~7--8; 6: ch.~9; 7: ch.~10; 8: ch.~11--12; 9: ch.~13 and~15; 10: ch.~17 \\
AI and machine learning & someone who knows neural networks and wants to understand how the brain learns & 1--4, 5 (skim), 6, 7, 8, 9, 11 (skim), 12, 13 (skim), 15 \\
electronics and hardware & circuit designer, neuromorphic chip developer & 1--4, 5, 11, 13, 16, 18, 19, 17 (skim chapters 9 and 15) \\
neural interfaces and living computers & someone who builds brain--computer interfaces and chips with living neurons & 1, 2, 4, 11, 13, 17 (skim chapters 9 and 15), 21, 22 (skim chapter 12) \\
quick overview & an engineer with a free weekend & 1, 2, 3, 6, 10; chapter~6's references to chapters~4--5 are clear from context \\
\bottomrule
\end{tabular}
\end{center}

``Skim'' means: it is enough to read the beginning of the chapter, its statements in yellow boxes, and its summary table.

Each chapter ends with exercises: estimates, derivations, simulations, and design problems; short answers are collected in Appendix~\ref{sec:answers}. All notation is in Appendix~\ref{sec:notation}, and the experiments behind each chapter's main claims are in Appendix~\ref{sec:supplement}. The numbers in the book come from small programs in the \texttt{models/} folder; you can run and modify them.

\mainmatter

\chapter{Neuron Types}
\label{sec:neurontypes}

\section{The neuron as a black box}

Suppose you are handed a sack of eighty-six billion neurons~\cite{Azevedo2009} and asked to sort them into piles. How would you do it?

An engineer handed a sack of unfamiliar chips would not sort them by the shape of the package. He would ask: how many inputs does the part have, how many outputs, and what does it put out? Let us draw the neuron the same way, as a block in a schematic (Fig.~\ref{fig:blackbox}).

\begin{figure}[t]
\centering
\begin{tikzpicture}[>=Stealth,x=1cm,y=1cm]
% the box
\node[draw,very thick,minimum width=2.6cm,minimum height=2.6cm,fill=blue!6] (N) at (0,0) {};
\node at (0,0.55) {\small neuron};
\node at (0,-0.15) {$\displaystyle u=\sum_i w_i x_i$};
\node at (0,-0.8) {\small $y=\Theta(u-\theta)$};
% inputs
\foreach \k/\y/\lab in {1/1.05/{x_1},2/0.55/{x_2},3/0.05/{x_3}}{
  \draw[->,thick,blue!60!black] (-4.2,\y) -- (-1.3,\y);
  \node[left] at (-4.2,\y) {$\lab$};
  \node[above,blue!60!black] at (-2.1,\y-0.06) {\scriptsize $w_{\k}$};}
\node at (-4.45,-0.4) {$\vdots$};
\node at (-2.1,-0.4) {$\vdots$};
\draw[->,thick,blue!60!black] (-4.2,-1.05) -- (-1.3,-1.05);
\node[left] at (-4.2,-1.05) {$x_n$};
\node[above,blue!60!black] at (-2.1,-1.11) {\scriptsize $w_{n}$};
\draw[decorate,decoration={brace,amplitude=5pt},gray!60!black] (-4.9,-1.2) -- (-4.9,1.2);
\node[left,align=right,gray!40!black] at (-5.1,0) {\small $n\sim10^3$--$10^5$\\[-2pt]\small inputs};
% single output wire, then fan-out
\draw[very thick,red!70!black] (1.3,0) -- (3.2,0);
\foreach \x in {1.6,2.2,2.34,2.9}{\draw[thick,red!70!black] (\x,0) -- (\x,0.4);}
\node[above right,font=\tiny\itshape,gray!30!black,inner sep=1pt] at (1.42,0.45) {click\dots\ click-click\dots};
\node[below,red!70!black,align=center] at (2.25,-0.05) {\scriptsize one output $y$};
\fill[red!70!black] (3.2,0) circle (0.06);
\foreach \y in {1.3,0.65,0,-0.65,-1.3}{
  \draw[->,thick,red!70!black] (3.2,0) -- (5.0,\y);}
\draw[decorate,decoration={brace,amplitude=5pt},gray!60!black] (5.3,1.45) -- (5.3,-1.45);
\node[right,align=left,gray!40!black] at (5.5,0) {\small copies of $y$\\[-2pt]\small to $10^3$--$10^4$\\[-2pt]\small other neurons};
\end{tikzpicture}
\caption{The neuron through an engineer's eyes. Many inputs $x_i$, each with its own weight $w_i$; inside, a weighted sum $u$ and a comparison with the threshold $\theta$ ($\Theta$ is the step function: $1$ if its argument is positive and $0$ otherwise); one output $y$, a train of spikes, copied by branching to thousands of recipients.}
\label{fig:blackbox}
\end{figure}

The output of the block is not a number but a train of short voltage spikes: ``click,'' a pause, ``click-click,'' a long pause. Every spike has the same height, so all the information is in \emph{when} they occur. Inside the block, to a first rough approximation, sit an adder and a comparator: the inputs are summed with weights, and if the sum crosses a threshold, the block fires a spike. In the next chapter we will replace this with a model that works honestly in continuous time.

There is one output, but the wire branches, and every branch carries \emph{the very same} copy of the signal. In electronics this is called fan-out. For a cortical neuron it is on the order of several thousand; a logic gate on a chip typically drives a handful of other gates.

\emph{What}, then, travels along the output wire to the recipient? The spike runs to the end of each branch and there turns into a chemical signal: the branch releases a pinch of a particular substance, a \emph{neurotransmitter}, into the narrow gap between the cells, and it changes the recipient's input current. That gives us a sorting criterion: \emph{which transmitter does the neuron's output release?}

\section{One output, one signal type}

Sorting makes sense only if each neuron has a single output type. Is that so? Experiment says: almost.

\begin{keyblock}
\begin{proposition}[One neuron, one output type]
\label{prop:dale}
Almost every neuron has one principal transmitter, and it releases it from every branch of its output. The type of a neuron is therefore set by its principal transmitter~\cite{Strata1999}.
\end{proposition}
\end{keyblock}

Let us agree to label a neuron type by \emph{the first four letters of the name of its principal transmitter}. A neuron whose principal transmitter is glutamate is a \nt{GLUT} neuron. All the types are collected in Table~\ref{tab:types}.

\begin{table}[t]
\centering
\footnotesize
\setlength{\tabcolsep}{3pt}
\renewcommand{\arraystretch}{1.15}
\begin{tabular}{>{\raggedright}p{1.0cm}>{\raggedright}p{2.05cm}>{\raggedright}p{1.75cm}>{\raggedright}p{1.6cm}>{\centering}p{0.7cm}>{\raggedright}p{1.55cm}>{\raggedright}p{1.8cm}>{\raggedright\arraybackslash}p{3.2cm}}
\toprule
Type & Principal transmitter & Bus & Speed & Sign & Neurons in the brain & Outputs per neuron & Role in computation \\
\midrule
\nt{GLUT} & glutamate & address & fast and slow & $+$ & $\sim8\cdot10^{10}$ & $\sim10^4$ & main positive weight \\
\nt{GABA} & GABA & address & fast and slow & $-$ & $\sim3\cdot10^{9}$ & $10^3$--$10^4$ & main negative weight \\
\nt{GLYC} & glycine & address & fast & $-$ & $10^6$--$10^7$ & $10^2$--$10^3$ & negative weight in the spinal cord and brainstem; second key for one of the glutamate receptors \\
\nt{ACET} & acetylcholine & both & fast and slow & $\pm$ & $\sim10^6$ & muscle: $10$--$10^3$; brain: $\sim10^5$ & output to muscle; in the brain, learning rate (hypothesis) \\
\nt{DOPA} & dopamine & broadcast & slow & $\pm$ & $\sim5\cdot10^{5}$ & $10^5$--$10^6$ & reward-prediction error $\delta$ \\
\nt{SERO} & serotonin & broadcast & slow (one receptor type is fast) & $\pm$ & $\sim3\cdot10^{5}$ & $\sim10^5$ & planning horizon (hypothesis) \\
\nt{HIST} & histamine & broadcast & slow & $+$ & $\sim6\cdot10^{4}$ & no estimate & global ``stay awake'' signal \\
\nt{NORA} & noradrenaline & broadcast & slow & $\pm$ & $\sim5\cdot10^{4}$ & $\sim10^5$ & gain (hypothesis) \\
\nt{ADRE} & adrenaline & broadcast & slow & $\pm$ & $\sim10^4$ & no estimate & few neurons; role in the brain unclear \\
\nt{ASPA} & aspartate & address & fast & $+$ & no estimate & no estimate & weak positive weight; role disputed \\
\bottomrule
\end{tabular}
\caption{Neuron types, in decreasing order of their number in the brain. \emph{Bus}: address, the transmitter is released into a narrow gap toward a specific recipient; broadcast, from a small group of source neurons onto large regions (Section~\ref{sec:buses}). \emph{Speed}: fast, through ionotropic receptors, milliseconds; slow, through metabotropic receptors, hundreds of milliseconds and longer. \emph{Sign}: the usual one; the recipient has the final say (Section~\ref{sec:receiver}), and $\pm$ means both signs occur. \emph{Neurons in the brain}: how many neurons of this type the whole human brain has, to order of magnitude; the total is about $8.6\cdot10^{10}$~\cite{Azevedo2009}. Nearly all \nt{GLUT} neurons, about $7\cdot10^{10}$, are the small input neurons of the cerebellum; the \nt{GLUT} and \nt{GABA} numbers come from this count and the fraction of \nt{GABA} neurons in the cortex~\cite{Hendry1987}. Of the sparse types, only \nt{HIST}~\cite{Airaksinen1991} and the main one of the two \nt{DOPA} groups~\cite{Pakkenberg1991} have been counted directly, so the \nt{DOPA} figure is a lower bound; the \nt{SERO} figure is the sum of two partial counts~\cite{Baker1990,Baker1991}. The numbers for \nt{NORA}, \nt{GLYC}, \nt{ACET}, and \nt{ADRE} are rough order-of-magnitude estimates with no direct count. \emph{Outputs per neuron}: how many synaptic terminals a single neuron has, that is, release sites on the branches of its output wire (Fig.~\ref{fig:blackbox}), to order of magnitude. For the broadcast types the terminals have not been counted directly: for \nt{DOPA} the number is inferred from what fraction of its target the branching of a single output wire covers~\cite{Matsuda2009}; for the others the estimates are rougher still. \nt{ACET} has two cases: the neuron that drives a muscle and the broadcast neuron in the brain.}
\label{tab:types}
\end{table}

% the pie is a table-type float with a figure caption, so it can never float ahead of Table~\ref{tab:types}
\begin{table}[tb]
\makeatletter\def\@captype{figure}\makeatother
\centering
\definecolor{vizglut}{HTML}{2A78D6}
\definecolor{vizgaba}{HTML}{E34948}
\definecolor{vizrest}{HTML}{8A8986}
\begin{tikzpicture}[>=Stealth]
% pie: GLUT 96.4 %, GABA 3.6 % (13.0 degrees), the rest 0.006 % (a hairline at 0 degrees)
\fill[vizglut] (0,0) -- (13:1.9) arc (13:360:1.9) -- cycle;
\fill[vizgaba] (0,0) -- (0:1.9) arc (0:13:1.9) -- cycle;
\draw[white,line width=1pt] (0,0) -- (13:1.9);
\draw[vizrest,line width=1.2pt] (0,0) -- (0:1.9);
\node[white,align=center,font=\small] at (190:0.95) {\nt{GLUT}\\$96.4\,\%$};
\draw[gray!60!black] (6.5:1.75) -- (2.05,2.1);
\node[anchor=west,font=\small] at (2.05,2.1) {\nt{GABA} $3.6\,\%$};
% zoom box for the remaining types
\draw[densely dashed,vizrest] (1.9,0) -- (3.1,1.65);
\draw[densely dashed,vizrest] (1.9,0) -- (3.1,-2.2);
\draw[vizrest,rounded corners=3pt] (3.1,-2.2) rectangle (10.1,1.65);
\node[anchor=west,font=\scriptsize] at (3.2,1.4) {all others together: $\approx0.006\,\%$; log scale};
\foreach \code/\lg/\y/\val/\lx in {GLYC/6.477/1.0/{$10^6$--$10^7$}/4.05,ACET/6.0/0.6/{$\sim10^6$}/3.0,DOPA/5.699/0.2/{$\sim5\cdot10^5$}/2.699,SERO/5.477/-0.2/{$\sim3\cdot10^5$}/2.477,HIST/4.778/-0.6/{$\sim6\cdot10^4$}/1.778,NORA/4.699/-1.0/{$\sim5\cdot10^4$}/1.699,ADRE/4.0/-1.4/{$\sim10^4$}/1.0}{
  \node[anchor=east,font=\scriptsize] at (4.35,\y) {\nt{\code}};
  \fill[vizrest] (4.45,\y-0.13) rectangle ({4.45+\lg-3},\y+0.13);
  \node[anchor=west,font=\scriptsize] at ({4.5+\lx},\y) {\val};}
\draw[vizrest,thick] (7.45,1.0) -- (8.45,1.0);
\draw[vizrest,thick] (7.45,0.92) -- (7.45,1.08);
\draw[vizrest,thick] (8.45,0.92) -- (8.45,1.08);
\draw[gray!60!black] (4.45,-1.72) -- (8.45,-1.72);
\foreach \e in {3,4,5,6,7}{
  \draw[gray!60!black] ({4.45+\e-3},-1.72) -- ({4.45+\e-3},-1.66);
  \node[below,font=\scriptsize] at ({4.45+\e-3},-1.72) {$10^{\e}$};}
\end{tikzpicture}
\caption{Shares of neuron types in the brain, from the estimates of Table~\ref{tab:types}. \nt{GLUT} fills almost the entire circle; \nt{GABA} is a narrow sector, and all the other types together are a hairline at the edge. On the right, that hairline is magnified on a logarithmic scale of neuron count; for \nt{GLYC} the bar reaches the middle of the $10^6$--$10^7$ range, and the segment shows the whole range. \nt{ASPA} has no estimate and is not shown.}
\label{fig:typepie}
\end{table}

Figure~\ref{fig:typepie} shows how unequal these piles are: out of every thousand neurons, $964$ are \nt{GLUT} and $36$ are \nt{GABA}, while all the other types together account for about six neurons in a hundred thousand.

GABA is already a four-letter name. Substances that are almost never a principal transmitter, namely neuropeptides, gases, lipids, and purines, get no code of their own; they are covered in Appendix~\ref{sec:othermediators}. The word ``almost'' in Proposition~\ref{prop:dale} matters. Many neurons release auxiliary substances along with the principal transmitter~\cite{Yao2023}. Some also release a second fast transmitter, even one of the opposite sign: some \nt{DOPA} neurons also release GABA~\cite{Tritsch2012}. And some neurons release different transmitters from different branches of one and the same output~\cite{Zhang2015}. All of this has been measured, so ``one neuron, one transmitter'' is a rule with exceptions, not a law. But as a first division it holds up: in the cell atlas of the mouse brain, where neurons are sorted into more than five thousand groups by the genes active in them, the neurons still split into large classes primarily by principal transmitter~\cite{Yao2023}.

This rule has a consequence that immediately sets the brain apart from an artificial neural network. In an artificial network one and the same neuron can have a positive weight to one recipient and a negative weight to another: the weights $w_{ij}$ are just numbers in a matrix. In the brain the sign of the action is set mainly by the transmitter, and a neuron has one transmitter. So if neuron $j$ excites one recipient, it excites all the others as well. The entire \emph{column} of the weight matrix leaving neuron $j$ has one sign:
\begin{empheq}[box=\keybox]{equation}
\operatorname{sign} w_{ij} \;=\; s_j \quad\text{for all recipients } i, \qquad s_j\in\{+1,-1\}.
\label{eq:dalesign}
\end{empheq}
Neurons divide into excitatory ($s_j=+1$) and inhibitory ($s_j=-1$), and this is a property of the neuron itself, not of the connection. If a network needs a negative connection from an excitatory neuron, it has to route it through an intermediary: the excitatory neuron excites an inhibitory one, which inhibits the target: a negative weight with a one-link delay.

More than a hundred transmitters are known, but almost all of the brain's work rests on half a dozen, and they fall into two completely different classes.

\section{Two buses: address and broadcast}
\label{sec:buses}

Computer networks have two ways of delivering a message. The first is \emph{addressed}, unicast: a packet goes from one sender to a specific recipient, quickly, and no one else sees it. The second is \emph{broadcast}: the sender transmits one message to every node of the network at once. Neurons use both, and the transmitters divide exactly along this line (Fig.~\ref{fig:phone_radio}).

\begin{figure}[t]
\centering
\begin{tikzpicture}[>=Stealth,scale=0.95]
% ---- unicast: point-to-point wiring
\node[above] at (2.0,3.3) {\small \textbf{Address}: glutamate and GABA};
\foreach \i in {0,...,4}{
  \fill[blue!60!black] (0.2+0.9*\i,2.5) circle (0.1);
  \fill[blue!60!black] (0.2+0.9*\i,0.6) circle (0.1);}
\foreach \a/\b in {0/0,0/1,1/1,1/3,2/2,3/2,3/4,4/4,2/0}{
  \draw[->,blue!60!black] (0.2+0.9*\a,2.38) -- (0.2+0.9*\b,0.72);}
\fill[red!70!black] (1.1,1.55) circle (0.09);
\pic[scale=1.2] at (1.75,1.9) {letter};
\draw[->,red!70!black,thick] (1.1,1.55) -- (0.28,0.7);
\draw[->,red!70!black,thick] (1.1,1.55) -- (1.95,0.72);
\node[align=center,below] at (2.0,0.2) {\small billions of neurons,\\[-2pt]\small each with its own addressees,\\[-2pt]\small time: milliseconds};
% ---- broadcast: one small source, global targets
\begin{scope}[xshift=7.2cm]
\node[above] at (1.8,3.3) {\small \textbf{Broadcast}: dopamine, \dots};
\foreach \i in {0,...,8}{\fill[blue!60!black] (-0.2+0.5*\i,2.75) circle (0.08);}
\fill[orange!85!black] (1.8,0.35) ellipse (0.35 and 0.18);
\pic[scale=1.5] at (2.7,0.75) {mast};
\foreach \x in {-0.2,0.3,0.8,1.3,1.8,2.3,2.8,3.3,3.8}{
  \draw[->,orange!85!black] (1.8,0.5) .. controls (1.8,1.3) and (\x,1.6) .. (\x,2.62);}
\node[align=center,below] at (1.8,0.1) {\small thousands to hundreds of thousands of neurons,\\[-2pt]\small one message to all,\\[-2pt]\small time: hundreds of milliseconds and longer};
\end{scope}
\end{tikzpicture}
\caption{Two ways of transmitting. Left, the address bus, like mail with an address on the envelope; one neuron and two of its recipients are shown in red. Right, the broadcast bus, like a radio station: a small group of source neurons, and everyone receives the same message.}
\label{fig:phone_radio}
\end{figure}

The \emph{address bus} is glutamate and GABA. The overwhelming majority of neurons release them. Each output leads to specific cells, the transmitter acts only within the narrow gap of the contact, and it acts fast: the recipient's ion channels open within a millisecond, and it is all over within a few milliseconds. This is the \emph{data} bus: it carries what the brain computes.

The \emph{broadcast bus} is dopamine, serotonin, noradrenaline, and partly acetylcholine. They are released by tiny groups of neurons, but the output of each of these neurons branches incredibly widely. The transmitter is often released not into a narrow gap but into the extracellular space, where it spreads out and acts on everyone nearby. It acts slowly: rather than opening channels directly, it triggers a chain of chemical reactions inside the recipient. This bus carries not data but \emph{control parameters}, shared by large parts of the network, that change its operating mode: the gain, the learning rate, the signal ``that was good.'' That is why such transmitters are called \emph{modulators}. For a long time each such channel was thought to be a single number for the whole brain. Modern measurements have refined the picture: a channel is not one wire but a small bundle of parallel lines. The source neurons of one transmitter divide into subsystems with their own recipients and their own message, and how much transmitter to release on site is also adjusted by local signals~\cite{Ren2018,Poe2020}. The number of such lines is in the units or tens, not billions, so the channel remains a broadcast channel.

If you need one picture from this chapter, here it is. The brain is a huge network with addressed data transfer, above which hang a few broadcast channels carrying control signals. A programmer will recognize a familiar design: computation plus a small set of control variables, each shared by a large part of the network.

\section{\nt{GLUT}: the positive weight}

Glutamate is the brain's main excitatory transmitter. When a neuron's output releases glutamate, channels open in the recipient through which sodium flows in, the voltage across the recipient's membrane rises, and its chance of firing a spike of its own increases. In the language of Fig.~\ref{fig:blackbox}, this is an input with a \emph{positive} weight, $w_i>0$.

Who puts out glutamate? Almost everyone who transmits data over long distances: three-quarters to four-fifths of cortical neurons and most of the neurons that connect different brain regions. The brain's network is first and foremost a network of glutamate connections.

One engineering detail. After release, the glutamate concentration in the gap shoots up a thousandfold, to about a millimolar, and decays with a time constant of about a millisecond~\cite{Clements1992}: glutamate diffuses out of the gap, and dedicated pump proteins carry it back into cells. Why? For the same reason a communication line returns the signal to zero after each symbol. If the transmitter were not cleared, the next spike would arrive on a ``busy'' line, and spikes a few milliseconds apart would merge.

\section{\nt{GABA}: the negative weight}

Imagine a network in which all weights are positive. If on average each spike gives rise to more than one next spike, activity grows exponentially and within a few steps engulfs the whole network. An electronics engineer will recognize a positive feedback loop with gain above one: the circuit saturates. To prevent this, negative weights are needed. Their main source is gamma-aminobutyric acid, GABA for short.

GABA opens channels in the recipient that pass chloride ions. The chloride equilibrium potential lies near the resting potential or slightly below it, so open chloride channels hold the membrane near rest and keep the voltage from rising to threshold. This is an input with a \emph{negative} weight, $w_i<0$. GABA neurons make up a fifth to a quarter of cortical neurons~\cite{Hendry1987}. Their outputs are short, and they inhibit their nearest neighbors.

Inhibition in the brain does much more than subtract. It works as negative feedback that holds the whole network at its operating point. In a normally working cortex, the excitatory and inhibitory currents arriving at a neuron are thought to balance each other almost exactly: this regime was predicted by theory~\cite{vanVreeswijk1996} and is seen in current measurements in living cortex~\cite{Haider2006}, and the neuron sits near threshold all the time. A circuit stabilized by strong negative feedback around an unstable point responds quickly and sensitively to small signals, an old engineering trick.

\section{\nt{DOPA}: the error signal}

The first broadcast channel is dopamine. Humans have on the order of half a million dopamine neurons, roughly one in a hundred seventy thousand; that is the count in the main one of their two groups~\cite{Pakkenberg1991}, so it is a lower bound. Their outputs spread to the regions that select actions.

What does this channel transmit? The answer comes from experiments that record the spikes of dopamine neurons in an animal receiving a reward~\cite{Schultz1997}. From time to time the animal is unexpectedly given a drop of juice, and the dopamine neurons respond with a short burst of spikes. Then a lamp is lit before the juice, always one second ahead of it. Once the animal has learned this association, the neurons start responding with a burst to the \emph{lamp}, and to the juice itself, arriving right on time, they do not respond at all. And if the juice is withheld after the lamp, then at the moment it should have arrived the neurons \emph{fall silent}, and their rate drops below the usual level.

The rule that explains all three cases (Fig.~\ref{fig:rpe}): the neurons report not how good things are but how much \emph{better or worse than expected} they are.

\begin{figure}[t]
\centering
\begin{tikzpicture}[>=Stealth,x=3.4cm,y=1cm]
\foreach \y/\lab in {2.8/{unexpected juice},1.4/{juice after lamp},0/{lamp, but no juice}}{
  \draw[gray!70!black] (0,\y) -- (3,\y);
  \draw[densely dotted,gray!60!black] (0,\y+0.3) -- (3,\y+0.3);
  \node[left,align=right,font=\small] at (-0.03,\y+0.35) {\lab};}
% row 1: juice without a cue -> burst at the juice
\draw[very thick,orange!85!black] plot[domain=0:3,samples=120] (\x,{2.8+0.3+0.75*exp(-((\x-2.08)/0.07)^2)});
\pic[scale=0.8] at (2,2.58) {juice};
% row 2: cue then juice -> burst at the cue, nothing at the juice
\draw[very thick,orange!85!black] plot[domain=0:3,samples=120] (\x,{1.4+0.3+0.75*exp(-((\x-1.08)/0.07)^2)});
\pic[scale=0.85] at (1,1.15) {lamp}; \pic[scale=0.85] at (2,1.15) {juice};
% row 3: cue, juice omitted -> burst at the cue, dip when the juice was due
\draw[very thick,orange!85!black] plot[domain=0:3,samples=120] (\x,{0.3+0.75*exp(-((\x-1.08)/0.07)^2)-0.28*exp(-((\x-2.1)/0.12)^2)});
\pic[scale=0.85] at (1,-0.25) {lamp}; \pic[scale=0.85] at (2,-0.25) {nojuice};
\node[right,font=\scriptsize] at (2.36,0.1) {dip};
\node[right,font=\scriptsize] at (2.13,3.75) {surprise!};
\node[left,font=\scriptsize] at (0.97,2.25) {burst at the lamp};
\node[above,font=\scriptsize] at (2,1.75) {nothing};
\draw[->,gray!70!black] (0,-0.75) -- (3.05,-0.75) node[right,black,font=\small] {$t$};
\draw[gray!60!black] (1,-0.8) -- (1,-0.7) node[below=2pt,font=\scriptsize] {lamp, $1\,\text{s}$};
\draw[gray!60!black] (2,-0.8) -- (2,-0.7) node[below=2pt,font=\scriptsize] {juice, $2\,\text{s}$};
\end{tikzpicture}
\caption{Firing rate of \nt{DOPA} neurons in three experiments (schematic, after~\cite{Schultz1997}). The dotted line is the steady background rate. Lamp: the cue; drop: the juice; crossed-out drop: juice that was promised but not given. The response is the prediction error~\eqref{eq:rpe}.}
\label{fig:rpe}
\end{figure}

A programmer familiar with reinforcement learning~\cite{SuttonBarto2018} will recognize this quantity at once. An agent learning to choose actions keeps an estimate $V(s)$ of how much reward it expects in state $s$. After each step it compares the expectation with what it got:
\begin{empheq}[box=\keybox]{equation}
\delta_t \;=\; r_t \;+\; \gamma\,V(s_{t+1}) \;-\; V(s_t) ,
\label{eq:rpe}
\end{empheq}
and corrects the estimate: $V(s_t)\leftarrow V(s_t)+\alpha\,\delta_t$. Here $r_t$ is the reward received, $\gamma<1$ is the discount factor (how much less future reward is worth than immediate reward), and $\alpha$ is the learning rate. The quantity $\delta_t$ is called the \emph{temporal-difference error}.

It behaves exactly like the dopamine neurons. Unexpected juice: $V$ is still zero, $r>0$, $\delta>0$. Learned juice: the lamp predicted the reward, $V(s_t)$ before the juice already equals $r$, and $\delta=0$. The juice did not come: $V(s_t)=r$ but $r_t=0$, $\delta<0$. And the burst moves to the lamp because it is at the moment of the lamp that the expectation jumps up: $\gamma V(s_{t+1})-V(s_t)>0$. The steps $t,t+1,\dots$ here are a convention of the algorithm: the body has no clock that ticks out steps, and we will obtain the same quantity in continuous time in Chapter~\ref{sec:dofa}.

So dopamine is a broadcast of a scalar error signal $\delta$ to everyone who must learn from it. To a first approximation it is one wire for the whole brain and one number at each moment; how much of a bundle this wire really is, is taken up in Chapter~\ref{sec:dofaalt}.

\section{The other broadcast channels}

For the other modulators there are only working hypotheses that translate their roles into the same language of global parameters~\cite{Doya2002}. Hold them as exactly that: hypotheses.

\emph{\nt{NORA}, noradrenaline: gain.} There are even fewer noradrenaline neurons than dopamine neurons, by a rough estimate several tens of thousands, and their outputs spread through practically the entire brain. They become sharply active when something unexpected or important happens. Noradrenaline changes the steepness of the transfer characteristic of the recipient neurons: weak inputs become weaker, strong ones stronger. This is measured as follows: the recipient's background spikes are suppressed more than its response to the input, and the signal-to-noise ratio increases~\cite{Foote1975NE}. A simple model describes this with one number: if a neuron responds with rate $f=F\bigl(g\cdot u\bigr)$ to input current $u$, noradrenaline increases the gain $g$~\cite{AstonJones2005} (more in Chapter~\ref{sec:nora}). A network with large $g$ works with ``more contrast'' and decides faster; with small $g$ it works ``blurred'' and tries out more options, like a reinforcement-learning agent in exploration mode.

\emph{\nt{ACET}, acetylcholine: learning rate.} Acetylcholine controls how much the network listens to the current input and how much to its own stored data. When the acetylcholine level is high, inputs from the sense organs are amplified, the internal, ``recalling'' connections are weakened, and at the same time weight changes are made easier~\cite{Hasselmo2006}. Roughly speaking, it is a ``write/read'' switch and, together with it, the learning rate $\alpha$.

\emph{\nt{SERO}, serotonin: planning horizon.} What serotonin transmits is the least understood. One hypothesis ties it to the discount factor $\gamma$ in~\eqref{eq:rpe}: a high serotonin level corresponds to a willingness to wait for a large reward rather than grab a small one now. Serotonin neurons work steadily and slowly, a few spikes per second during waking, and fall almost silent in one of the phases of sleep~\cite{JacobsAzmitia1992}.

All six transmitters are collected in Table~\ref{tab:transmitters}.

\begin{table}[t]
\centering
\small
\begin{tabular}{>{\raggedright}p{2.4cm}>{\raggedright}p{3.0cm}>{\raggedright}p{2.0cm}>{\raggedright\arraybackslash}p{5.2cm}}
\toprule
Neuron type & Fraction of neurons & Bus & Role in computation \\
\midrule
\nt{GLUT} (glutamate) & majority & address & positive weight, $w>0$ \\
\nt{GABA} (GABA) & 20--25\% in the cortex & address & negative weight, $w<0$; stabilization \\
\nt{DOPA} (dopamine) & $\sim 6\cdot10^{-6}$ & broadcast & prediction error $\delta$ \\
\nt{NORA} (noradrenaline) & $\sim 6\cdot10^{-7}$ & broadcast & gain $g$ (hypothesis) \\
\nt{ACET} (acetylcholine) & small & both & learning rate $\alpha$; ``write/read'' (hypothesis) \\
\nt{SERO} (serotonin) & $\sim 3\cdot10^{-6}$ & broadcast & discounting $\gamma$ (hypothesis) \\
\bottomrule
\end{tabular}
\caption{The brain's main transmitters in the language of computation. The right-hand column is a strong simplification: reliable for glutamate, GABA, and dopamine; for the others these are hypotheses.}
\label{tab:transmitters}
\end{table}

\section{The recipient sets the sign}
\label{sec:receiver}

I misled you a little when I said glutamate is a positive weight and GABA a negative one. No molecule carries a sign within it. A molecule is just a key. What happens when it is caught is decided by the \emph{lock}: the receptor on the recipient cell.

The best example is dopamine. It has receptors of two families~\cite{Beaulieu2011}. At one it makes the cell easier to excite, at the other harder. In the region that selects actions, some neurons carry receptors of the first type and others of the second, and one and the same dopamine signal simultaneously strengthens one group and weakens the other. It is like a single broadcast packet that different network nodes interpret differently.

\begin{keyblock}
\begin{proposition}[The recipient sets the meaning of the message]
\label{prop:receptor}
The sign and speed of a transmitter's action are set not by the transmitter itself but by the receptor it binds to. Receptors come in two kinds. \emph{Ionotropic} receptors are themselves ion channels and open within a fraction of a millisecond: this is direct control of current, like the gate of a transistor. \emph{Metabotropic} receptors trigger a chain of chemical reactions inside the cell and act over hundreds of milliseconds and longer: this is more like a write to a settings register that changes the cell's subsequent operation. The address bus speaks mainly through the former, the broadcast bus mainly through the latter.
\end{proposition}
\end{keyblock}

Why, then, can one say ``glutamate excites'' at all? Because almost all the glutamate receptors met in practice are excitatory, and almost all GABA receptors are inhibitory. This is not a law of nature but a very reliable convention, and rule~\eqref{eq:dalesign} rests on it.

\section{What to take away}

The overwhelming majority of neurons form the addressed data bus, with weights of a single sign per neuron; a tiny minority form broadcast channels that deliver a few shared control signals to large parts of the brain. About two neurons in every hundred thousand, and on them depends how the entire rest of the network learns and in what mode it works. To an engineer this is no surprise: every computer has vastly fewer control registers than memory cells, and yet the machine does not work without them.

In the next chapter we will test what a network of \nt{GLUT} neurons alone, the most numerous type, can compute while running in continuous time.

\section{Exercises}

\begin{exercise}[\lvlA]
\label{ex:01:1}
\emph{Piles and wires.} From Table~\ref{tab:types} (log-scale midpoint of each range; \nt{ACET}: $10^5$ outputs): (a)~if a neuron were a person, how many Earth populations would the \nt{GLUT} neurons make, and how big a city would all the broadcast neurons make? (b)~By what factor does the share of terminals of \nt{DOPA}, \nt{SERO}, \nt{NORA}, and \nt{ACET} exceed the share of their neurons?
\end{exercise}

\begin{exercise}[\lvlA]
\label{ex:01:2}
\emph{Return to zero.} Glutamate in the gap decays as $c_0e^{-t/\tau_c}$, $c_0=1\mM$, $\tau_c=1\ms$; the recipient detects concentrations above $10$~\textmu M. (a)~How long is the line ``busy'' after one release? (b)~The pumps are partly blocked, $\tau_c=10\ms$. Up to what release rate does the line still manage to clear?
\end{exercise}

\begin{exercise}[\lvlB]
\label{ex:01:3}
\emph{Where the burst moves.} In the experiment of Fig.~\ref{fig:rpe}, the lamp is followed by states $s_0\to\dots\to s_{K-1}$ with step $\Delta$, $T=K\Delta$; the reward $r$ arrives on the last transition, and the moment of the lamp is unpredictable. Find the learned $V(s_k)$ and the response $\delta$ to the lamp. Setting $\gamma=e^{-\Delta/\tau_\gamma}$, find $\tau_\gamma$ for $T=1$~s, $\Delta=0.1$~s, $\gamma=0.95$.
\end{exercise}

\begin{exercise}[\lvlB]
\label{ex:01:4}
\emph{Simulation: learning from error.} Simulate the chain of Exercise~\ref{ex:01:3} with $K=10$, $\gamma=0.95$, $r=1$ and the rule $V(s_t)\leftarrow V(s_t)+\alpha\delta_t$. On which trial $n^*$ does the response to the lamp first overtake the response to the juice for $\alpha=0.05$, $0.1$, $0.2$, $0.5$? How does $n^*$ depend on $\alpha$?
\end{exercise}

\begin{exercise}[\lvlB]
\label{ex:01:5}
\emph{Design: broadcasting one number.} The number $\delta$ must reach all $10^{10}$ neurons; every recipient, relays included, must receive $10$ terminals from the previous layer, and each link adds $1\ms$. How many neurons and links does the most economical relay tree need with $10^4$ outputs per neuron (\nt{GLUT}) and with $10^{5.5}$ (\nt{DOPA})?
\end{exercise}

\begin{exercise}[\lvlB]
\label{ex:01:6}
\emph{Why balance is sensitive.} A network is $80\,\%$ \nt{GLUT} and $20\,\%$ \nt{GABA}, all-to-all with weights $J_E/N$ and $-J_I/N$; $\tau\dot r=-r+Wr+h$. At $J_E=10$ the only nonzero eigenvalue of $W$ is $0.5$. Find the ratio of inhibitory to excitatory current, and by what percentage it suffices to weaken inhibition for the network to run away into an avalanche.
\end{exercise}

\begin{exercise}[\lvlC]
\label{ex:01:7}
\emph{Research: is \nt{SERO} the $\gamma$?} By Exercise~\ref{ex:01:3}, the response of \nt{DOPA} neurons to the lamp is $re^{-T/\tau_\gamma}$. Delays $T=1,\dots,5$~s, $n$ presentations each, and $\ln\delta$ is measured with scatter $0.5$. How large must $n$ be to tell $\tau_\gamma=2$~s from $2.4$~s (level $0.05$, power $0.8$)? What mechanism other than $\gamma$ would give the same decline with $T$?
\end{exercise}

\chapter{What \nt{GLUT} Neurons Compute}
\label{sec:glutcomputer}
\begin{chaptergoals}
\item how a leaky integrate-and-fire neuron makes OR and a coincidence detector, an AND in time;
\item why \nt{GLUT} neurons alone compute only monotone functions, so activity cannot stop activity;
\item how on/off wire pairs let \nt{GLUT} neurons compute any logical function and signal readiness;
\item how delay lines turn a time difference into a place, and how feedback makes memory and timers.
\end{chaptergoals}

\section{The rules of the game}

In Chapter~\ref{sec:neurontypes} we saw that the overwhelming majority of the brain's neurons are \nt{GLUT}: they only excite, and their weights are only positive. Let us take as many such neurons as we like and ask: what can such a network compute if we allow ourselves only what happens in a living organism?

A living organism has no clock generator that switches all the elements at once. Each neuron works on its own, in continuous time: spikes come and go whenever they happen to, with various delays. There are \emph{inputs}, sensory neurons that fire in response to light, sound, or pressure, and \emph{outputs}, neurons that drive the muscles. Between them lies the network, and no one tells it when to start computing and when to stop. To an electronics engineer this is an \emph{asynchronous} circuit whose elements have delays that are not known precisely in advance.

We take the simplest neuron model among those that work honestly in continuous time. The voltage $V$ across the neuron's membrane is zero at rest. Each spike arriving from neuron $j$ raises it by a jump of $w_j\ge0$, after which the voltage leaks back to zero on its own with time constant $\tau$:
\begin{empheq}[box=\keybox]{equation}
\tau\,\frac{\dd V}{\dd t} \;=\; -V \;+\; \tau\sum_j w_j \sum_k \delta\bigl(t-t_j^{(k)}-d_j\bigr), \qquad w_j\ge 0 .
\label{eq:glutneuron}
\end{empheq}
Here $t_j^{(k)}$ are the spike times of neuron $j$, $d_j$ is the delay on the path from it, and $\delta$ is the delta function, which stands for an instantaneous jump. When $V$ reaches the threshold $\theta$, the neuron fires a spike, $V$ is reset to zero, and for about a millisecond the neuron cannot fire again. Typical numbers: $\tau$ ranges from a fraction of a millisecond to twenty milliseconds in different neurons, and delays from a fraction of a millisecond to tens of milliseconds. This is the \emph{leaky integrate-and-fire neuron}: an RC circuit with a threshold. It is a deliberate simplification: in cortical neurons the input branches generate local spikes of their own~\cite{Smith2013}, and a single neuron behaves more like a small multilayer network (Chapter~\ref{sec:synthesis}). For the questions of this chapter one RC circuit is enough.

The main difference from a logic gate is the leak: the neuron remembers a spike not forever but for about $\tau$, and only what arrives within this window adds up.

\section{OR and AND}

Let us start with gates. If a single spike lifts the voltage above threshold, $w\ge\theta$, the neuron fires on every spike from any input. That is \emph{OR}.

AND is more interesting. Let $w<\theta<2w$: one spike is not enough, two are needed. But now the question ``did both arrive?'' makes no sense until we say \emph{within what time}. Let the first spike arrive at time $0$ and the second a time $\Delta$ later. By the time the second arrives, $we^{-\Delta/\tau}$ remains of the first, and the neuron fires if $we^{-\Delta/\tau}+w\ge\theta$. Hence the coincidence window:
\begin{empheq}[box=\keybox]{equation}
\Delta \;\le\; \Delta_{\max} \;=\; \tau\,\ln\frac{w}{\theta-w} .
\label{eq:window}
\end{empheq}
This is not a logical AND but a \emph{coincidence detector}: an AND in time (Fig.~\ref{fig:coincidence}). With $\theta=1.5w$ the window is $\tau\ln2\approx0.7\tau$. For a cortical neuron with $\tau=10\ms$ that is about seven milliseconds. In neurons of the auditory system, where $\tau$ is less than a millisecond, the window shrinks to a fraction of a millisecond.

\begin{figure}[t]
\centering
\begin{tikzpicture}[>=Stealth,x=0.9cm,y=1.6cm]
\foreach \dx/\lab/\c [count=\i] in {0.5/{coincident}/red!70!black, 2.6/{not coincident}/blue!60!black}{
 \begin{scope}[xshift={(\i-1)*7.2cm}]
  \draw[->,gray!70!black] (0,0) -- (6.3,0) node[below left,black] {\small $t$};
  \draw[->,gray!70!black] (0,0) -- (0,1.75) node[left,black] {\small $V$};
  \draw[densely dashed,gray!60!black] (0,1.5) -- (6.0,1.5) node[right,black] {\small $\theta$};
  \draw[very thick,\c] (0,0) -- (1,0) -- (1,1)
      plot[domain=1:1+\dx,samples=30] (\x,{exp(-(\x-1)/1.5)});
  \pgfmathsetmacro{\v}{exp(-\dx/1.5)+1}
  \draw[very thick,\c] (1+\dx,{exp(-\dx/1.5)}) -- (1+\dx,{min(\v,1.5)});
  \ifdim\v pt<1.5pt
    \draw[very thick,\c] plot[domain=1+\dx:6,samples=40] (\x,{\v*exp(-(\x-1-\dx)/1.5)});
  \else
    \draw[very thick,\c] (1+\dx,1.5) -- (1+\dx,0) -- (6,0);
    \draw[->,thick,\c] (1+\dx,1.55) -- (1+\dx,1.72) node[above,font=\scriptsize] {spike};
  \fi
  \draw[->,thick] (1,-0.35) -- (1,-0.08); \draw[->,thick] (1+\dx,-0.35) -- (1+\dx,-0.08);
  \draw[decorate,decoration={brace,amplitude=3pt,mirror}] (1,-0.42) -- (1+\dx,-0.42) node[midway,below=3pt] {\scriptsize $\Delta$};
  \node[\c] at (3.5,-0.95) {\small \lab};
 \end{scope}}
\end{tikzpicture}
\caption{Coincidence detector, threshold $\theta=1.5w$. Left: the second spike arrived after $\Delta<\Delta_{\max}$, the sum reached threshold, and the neuron fired. Right: $\Delta>\Delta_{\max}$; almost nothing was left of the first spike, and the neuron stayed silent.}
\label{fig:coincidence}
\end{figure}

Another way to get AND is through duration rather than precise timing. While a light is on, a sensory neuron fires frequently and steadily; if one such input is not enough for the threshold but two are, the neuron is active as long as both are. In this way the network computes on \emph{levels}, and the precise arrival times of spikes do not matter. Most of what follows will work this way.

\section{What cannot be done}

Now let us try to build NOT: a neuron that is active when its input is silent. It will not work: with no input spikes, \eqref{eq:glutneuron} has not a single jump, and the voltage stays at zero, below threshold. What about a network of many neurons? No again, and this can be proved for all networks at once.

It is most convenient to talk about rates. Let $r_i$ be the firing rate of neuron $i$, $I_i$ the drive from the inputs, and $f_i$ its transfer characteristic: how the rate depends on the total drive. For an integrating neuron this characteristic is nondecreasing: more drive, more frequent spikes. The rates in a network that has come to equilibrium satisfy the equations
\begin{equation}
r_i \;=\; f_i\Bigl(\sum_j w_{ij}\,r_j + I_i\Bigr), \qquad w_{ij}\ge 0 .
\label{eq:ratenet}
\end{equation}

\begin{keyblock}
\begin{proposition}[Monotonicity]
\label{prop:monotone}
Let network~\eqref{eq:ratenet} consist of \nt{GLUT} neurons only and start from complete silence. If the inputs are strengthened, the steady-state rate of no neuron decreases. Everything such a network computes is a \emph{monotone} function of its inputs.
\end{proposition}
\end{keyblock}

Proof. Start from silence, $r^{(0)}=0$, and keep substituting the rates into the right-hand side of~\eqref{eq:ratenet}: $r^{(k+1)}=F\bigl(r^{(k)}\bigr)$. The right-hand side $F$ is nondecreasing in every argument, because the weights are nonnegative and the $f_i$ are nondecreasing. Since $r^{(1)}\ge0=r^{(0)}$, by induction $r^{(k+1)}\ge r^{(k)}$: the rates only grow and converge to the smallest solution of~\eqref{eq:ratenet}. If the inputs $I$ are replaced by larger $I'$, then at every step $r^{(k)}(I)\le r^{(k)}(I')$, and so in the limit too. A real network reaches equilibrium not in steps but continuously, yet the same holds for it: with positive connections, an inequality between two runs that holds at the start is preserved for all time.

For individual spikes the statement is not literally true: an extra input spike can make a neuron fire slightly earlier, and because of the millisecond of refractoriness its next spike is lost. But the effect is weak and unreliable, and one cannot compute with it. For rates, monotonicity is strict.

The consequences are the same as for any circuit without inverters. No NOT. No exclusive OR: adding a second input cannot turn the output off. And the most unpleasant one: \emph{activity cannot be stopped by activity}; one can only remove what sustains it.

\section{Two wires: on and off}

The way out is suggested by the organism itself. In many sense organs a stimulus is transmitted by two sets of neurons: in vision, some cells respond to an increase in light, others to a decrease~\cite{Schiller1992}. Let us call them \emph{on} and \emph{off} neurons. Instead of one wire $x$ the network receives a pair $(x,\bar x)$, and the absence that it cannot compute is reported to it by the sensor itself.

With a pair of wires, NOT is free: just swap the wires. AND and OR each take two neurons:
\begin{empheq}[box=\keybox]{equation}
\begin{aligned}
x\wedge y \;&\to\; \bigl(\;x\wedge y,\;\; \bar x\vee\bar y\;\bigr),\\
x\vee y \;&\to\; \bigl(\;x\vee y,\;\; \bar x\wedge\bar y\;\bigr).
\end{aligned}
\label{eq:dualrail}
\end{empheq}
All operations on the right are monotone, so Proposition~\ref{prop:monotone} is not violated.

For an asynchronous circuit the dual-rail code has a second, even more important advantage. A pair of wires has three states: ``yes,'' ``no,'' and, when both are silent, \emph{``no answer yet.''} The recipient learns that the computation is finished not from a clock but from the signal itself: exactly one wire in each pair has become active. Between stimuli the sensors are silent, and the network returns to silence, ready for the next time. In asynchronous electronics this trick is used to build circuits that work correctly for any element delays~\cite{Sparso2001}.

\begin{keyblock}
\begin{proposition}[Asynchronous computation]
\label{prop:dualrail}
If each input arrives as an on/off pair of wires, a network of \nt{GLUT} neurons can compute any logical function of its inputs. The answer also arrives as pairs of wires, and its readiness is seen from the fact that one wire in each pair has become active. No clock is needed for this. The price is roughly twice as many neurons.
\end{proposition}
\end{keyblock}

An example is exclusive OR: ``exactly one of two stimuli has changed.'' The direct wire of the answer is $(x\wedge\bar y)\vee(\bar x\wedge y)$, the complementary one is $(x\wedge y)\vee(\bar x\wedge\bar y)$. That is six neurons, and none of them needs inhibition (Fig.~\ref{fig:dualxor}).

\begin{figure}[t]
\centering
\begin{tikzpicture}[>=Stealth,x=1cm,y=1cm,
  nrn/.style={circle,draw,very thick,fill=blue!6,minimum size=0.75cm,inner sep=0pt,font=\small},
  inp/.style={font=\small}]
\node[inp] (X)  at (0,3.3) {$x$};
\node[inp] (Xb) at (0,2.2) {$\bar x$};
\node[inp] (Y)  at (0,1.1) {$y$};
\node[inp] (Yb) at (0,0)   {$\bar y$};
\node[nrn] (A1) at (3.2,3.3) {$2$};
\node[nrn] (A2) at (3.2,2.2) {$2$};
\node[nrn] (A3) at (3.2,1.1) {$2$};
\node[nrn] (A4) at (3.2,0)   {$2$};
\node[nrn] (O)  at (7.2,2.75) {$1$};
\node[nrn] (Ob) at (7.2,0.55) {$1$};
\foreach \s/\t in {X/A1,Yb/A1,Xb/A2,Y/A2,X/A3,Y/A3,Xb/A4,Yb/A4}{\draw[->,thick,blue!60!black] (\s) -- (\t);}
\draw[->,thick,blue!60!black] (A1) -- node[pos=0.45,above,sloped,font=\scriptsize,black] {$x\wedge\bar y$} (O);
\draw[->,thick,blue!60!black] (A2) -- node[pos=0.45,below,sloped,font=\scriptsize,black] {$\bar x\wedge y$} (O);
\draw[->,thick,blue!60!black] (A3) -- node[pos=0.45,above,sloped,font=\scriptsize,black] {$x\wedge y$} (Ob);
\draw[->,thick,blue!60!black] (A4) -- node[pos=0.45,below,sloped,font=\scriptsize,black] {$\bar x\wedge\bar y$} (Ob);
\draw[->,very thick,red!70!black] (O) -- (8.6,2.75) node[right,black] {$x\oplus y$: ``yes''};
\draw[->,very thick,red!70!black] (Ob) -- (8.6,0.55) node[right,black] {$\overline{x\oplus y}$: ``no''};
\node[below,font=\small,gray!40!black] at (3.2,-0.55) {AND: threshold $2$};
\node[below,font=\small,gray!40!black] at (7.2,-0.55) {OR: threshold $1$};
\node[below,font=\small,gray!40!black] at (0,-0.55) {wire pairs};
\end{tikzpicture}
\caption{Exclusive OR in the dual-rail code, built from \nt{GLUT} neurons only. The number in a circle is the threshold in units of one input's weight. Four AND neurons recognize the four input combinations, and two OR neurons collect them into the answer's pair of wires.}
\label{fig:dualxor}
\end{figure}

\section{Delays instead of a clock}

For computations involving time, wire delays and coincidence detectors are enough. The best example is how the brain determines where a sound came from.

Sound from a source off to one side reaches the near ear earlier than the far one. The difference depends on direction: from zero when the source is straight ahead to $0.22\,\text{m}/343\,\text{m/s}\approx0.6\ms$ in humans when it is directly to the side. The brain distinguishes differences of about ten \emph{microseconds}~\cite{KlumppEady1956,Grothe2010}, although the neurons themselves fire with a precision no better than a fraction of a millisecond. How?

Run a wire from the left ear along a row of coincidence detectors from left to right, and one from the right ear from right to left (Fig.~\ref{fig:itd}). The signal travels along the wire at speed $v$. To the detector located at distance $x$ from the left end of a row of length $L$, the left signal takes time $x/v$, the right one $(L-x)/v$. If the sound reached the right ear later by $\delta t$, the two signals meet simultaneously at the point where
\begin{empheq}[box=\keybox]{equation}
\frac{x}{v} \;=\; \frac{L-x}{v} + \delta t
\quad\Longrightarrow\quad
x \;=\; \frac{L}{2} + \frac{v\,\delta t}{2} .
\label{eq:itd}
\end{empheq}
Only the detector at which both signals arrive within the window~\eqref{eq:window} fires. The index of the detector that fired is the direction to the source. A time difference has turned into a \emph{place}.

\begin{figure}[t]
\centering
\begin{tikzpicture}[>=Stealth,
  nrn/.style={circle,draw,very thick,fill=blue!6,minimum size=0.7cm,inner sep=0pt}]
\foreach \k in {0,...,6}{\node[nrn] (D\k) at (1.4*\k,0) {\scriptsize $2$};}
\node[nrn,fill=red!15] at (D4) {\scriptsize $2$};
\draw[->,thick,blue!60!black] (-1.4,0.9) node[left,black] {\small left ear} -- (9.2,0.9);
\draw[->,thick,orange!85!black] (9.8,-0.9) node[right,black] {\small right ear} -- (-0.8,-0.9);
\foreach \k in {0,...,6}{
  \draw[->,blue!60!black] (1.4*\k,0.9) -- (D\k.north);
  \draw[->,orange!85!black] (1.4*\k,-0.9) -- (D\k.south);}
\draw[->,thick,red!70!black] (D4.north east) ++(0.1,0.05) -- ++(0.5,0.45) node[above right,font=\small,black] {fired};
\node[font=\small,gray!40!black] at (4.2,-1.5) {delay grows $\longrightarrow$ \hspace{2.2cm} $\longleftarrow$ delay grows};
\pic[scale=1.4] at (-2.3,1.75) {speaker};
\node[font=\scriptsize,align=center] at (-2.3,2.35) {source on the left};
\end{tikzpicture}
\caption{Finding the direction of a sound. The signals from the two ears run toward each other; each circle is a coincidence detector with threshold $2$. If the sound reached the right ear later, the signals meet to the right of the middle, by formula~\eqref{eq:itd}. The output of the network is the index of the neuron that fired. Here the source is on the left: the sound reached the left ear first.}
\label{fig:itd}
\end{figure}

There is no clock here, no inhibition, not even a dual-rail code: the network does not answer ``yes'' or ``no'' but points to one neuron out of many. Such a circuit of delay lines and coincidence detectors apparently does operate in the auditory brainstem of birds~\cite{Carr1990}. The microsecond precision comes not from the precision of each neuron but from there being many detectors, each looking at its own delay. In mammals no such row of delays has been found: there the same task is apparently solved by coincidence detectors tuned by precisely timed inhibition from \nt{GLYC} neurons~\cite{Brand2002,Grothe2010}. How inhibition narrows the coincidence window we will work out in Chapter~\ref{sec:gaba} using \nt{GABA} as the example; glycine inhibits the same way, by opening chloride channels in the recipient.

\section{Memory}

A computer needs memory. In such a network, memory is activity that sustains itself. Take a group of \nt{GLUT} neurons strongly connected to one another and describe it by a single rate $r$ (Fig.~\ref{fig:recurrentgroup}a). At equilibrium $r=f(wr+I)$, where $w$ is the strength of the group's feedback onto itself and $I$ is the external drive. Let us solve this equation graphically by intersecting the curve $f(wr+I)$ with the line $r$ (Fig.~\ref{fig:bistable}).

\begin{figure}[t]
\centering
% (b): tau dr/dt = -r + f(w r + I), f as in fig:bistable, w=1, I=0.6; Euler dt=0.001 tau.
\begin{tikzpicture}[>=Stealth,
  nrn/.style={circle,draw,thick,fill=blue!6,minimum size=0.5cm,inner sep=0pt}]
\begin{scope}
\node[font=\small] at (-0.5,1.55) {(a)};
\foreach \k in {1,...,6}{\node[nrn] (n\k) at ({1.4+1.0*cos(60*\k)},{1.0*sin(60*\k)}) {};}
\foreach \a/\b in {1/2,2/3,3/4,4/5,5/6,6/1,1/4,2/5,3/6}{\draw[<->,blue!60!black] (n\a) -- (n\b);}
\foreach \k in {2,3,4}{\draw[->,thick,green!45!black] ($(n\k)+(-0.95,0)$) -- (n\k);}
\node[font=\small,green!45!black] at (-0.35,-0.45) {$I$};
\node[font=\Large] at (3.1,0) {$\approx$};
\node[nrn,minimum size=0.85cm,font=\small] (R) at (4.35,-0.1) {$r$};
\draw[->,thick,blue!60!black] (R) edge[loop above,looseness=6] node[above,font=\small] {$w$} (R);
\draw[->,thick,green!45!black] (3.55,-0.95) node[left,font=\small] {$I$} -- (R);
\end{scope}
\begin{scope}[xshift=6.3cm,y=2.6cm,x=1.05cm]
\node[font=\small] at (-0.6,1.25) {(b)};
\draw[->,gray!70!black] (0,0) -- (7.3,0) node[below left,font=\small,black] {$t/\tau$};
\draw[->,gray!70!black] (0,0) -- (0,1.12) node[left,font=\small,black] {$r$};
\foreach \t in {0,2,4,6}{\draw[gray!60!black] (\t,-0.02) -- (\t,0.02); \node[below,font=\scriptsize] at (\t,0) {\t};}
\draw[densely dotted,gray!60!black] (0,0.5) -- (7,0.5) node[right,font=\scriptsize,black,align=left] {write\\threshold};
\fill[green!45!black,opacity=0.25] (1,-0.2) rectangle (2,-0.13);
\fill[gray!60!black,opacity=0.35] (1,-0.29) rectangle (1.5,-0.22);
\draw[very thick,blue!60!black] plot coordinates {(0.0,0.002) (0.1,0.002) (0.2,0.002) (0.3,0.002) (0.4,0.002) (0.5,0.002) (0.6,0.002) (0.7,0.002) (0.8,0.002) (0.9,0.002) (1.0,0.002) (1.1,0.083) (1.2,0.165) (1.3,0.242) (1.4,0.313) (1.5,0.378) (1.6,0.437) (1.7,0.491) (1.8,0.539) (1.9,0.583) (2.0,0.623) (2.1,0.643) (2.2,0.665) (2.3,0.688) (2.4,0.71) (2.5,0.732) (2.6,0.753) (2.7,0.773) (2.8,0.792) (2.9,0.81) (3.0,0.826) (3.2,0.855) (3.4,0.88) (3.6,0.9) (3.8,0.917) (4.0,0.931) (4.4,0.953) (4.8,0.967) (5.2,0.977) (5.6,0.984) (6.0,0.988) (6.5,0.992) (7.0,0.994)};
\draw[very thick,gray!60!black,densely dashed] plot coordinates {(0.0,0.002) (0.1,0.002) (0.2,0.002) (0.3,0.002) (0.4,0.002) (0.5,0.002) (0.6,0.002) (0.7,0.002) (0.8,0.002) (0.9,0.002) (1.0,0.002) (1.1,0.083) (1.2,0.165) (1.3,0.242) (1.4,0.313) (1.5,0.378) (1.6,0.358) (1.7,0.336) (1.8,0.313) (1.9,0.291) (2.0,0.269) (2.2,0.228) (2.4,0.191) (2.6,0.159) (2.8,0.133) (3.0,0.11) (3.2,0.091) (3.4,0.076) (3.6,0.063) (3.8,0.052) (4.0,0.043) (4.4,0.03) (4.8,0.021) (5.2,0.015) (5.6,0.011) (6.0,0.008) (6.5,0.006) (7.0,0.004)};
\node[font=\scriptsize,blue!60!black,anchor=west] at (3.3,0.8) {drive for $1\tau$: written};
\node[font=\scriptsize,gray!50!black,anchor=west] at (2.3,0.3) {drive for $0.5\tau$: died out};
\end{scope}
\end{tikzpicture}
\caption{A group strongly connected to itself. (a)~Many \nt{GLUT} neurons exciting one another behave like a single neuron with rate $r$ that excites itself with strength $w$; an external drive $I$ comes in from outside. (b)~The group's rate over time for $w=1$ and the same curve $f$ as in Fig.~\ref{fig:bistable}. The drive $I=0.6$ is applied for the time marked by the bar under the axis. If it lasts $1\tau$, the group gets past the unstable point $r=0.5$, flares up, and stays on after the drive ends. If only $0.5\tau$, it does not reach that point and dies back down: to write a bit, the drive must last longer than about $0.7\tau$.}
\label{fig:recurrentgroup}
\end{figure}

\begin{figure}[t]
\centering
\begin{tikzpicture}[>=Stealth,x=5cm,y=3.2cm]
\draw[->,gray!70!black] (0,0) -- (1.1,0) node[below left,black] {\small $r$};
\draw[->,gray!70!black] (0,0) -- (0,1.1);
\draw[thick,gray!60!black] (0,0) -- (1.05,1.05) node[right,black] {\small $r$};
\draw[very thick,blue!60!black] plot[domain=0:1.05,samples=80] (\x,{1/(1+exp(-(\x-0.5)/0.08))});
\draw[very thick,red!70!black,densely dashed] plot[domain=0:1.05,samples=80] (\x,{1/(1+exp(-(\x+0.3-0.5)/0.08))});
\fill[blue!60!black] (0.002,0.002) circle (0.018);
\draw[blue!60!black,fill=white,thick] (0.5,0.5) circle (0.018);
\fill[blue!60!black] (0.998,0.998) circle (0.018);
\node[below right,font=\small] at (0.02,0.0) {off};
\node[right,font=\small] at (0.52,0.46) {unstable};
\node[below,font=\small] at (0.99,0.96) {on};
\node[anchor=east,font=\small,red!70!black] at (0.19,0.62) {$f(wr+I)$};
\node[anchor=west,font=\small,blue!60!black] at (0.45,0.1) {$f(wr)$};
\end{tikzpicture}
\caption{Memory from a group of \nt{GLUT} neurons with strong feedback. Solid curve: no external drive; two stable intersections with the line $r$, ``off'' and ``on,'' and an unstable one between them. A brief drive $I$ (dashed curve) leaves only the upper one.}
\label{fig:bistable}
\end{figure}

If the feedback is strong, there are three intersections: the lower one is silence, the upper one is the group firing at full rate, and the middle one is unstable. We have obtained an element with two stable states, a flip-flop. Writing a one into it is easy: a brief drive lifts the curve, the lower intersection disappears, and the group flares up and stays on after the drive ends (Fig.~\ref{fig:recurrentgroup}b). A drive that is too short fails to bring the group to the unstable point, and it dies back down. Such self-sustaining activity, which persists for seconds after the stimulus is gone, is indeed observed in the brain and is considered the basis of short-term memory~\cite{Wang2001}. But it is not the only way to remember for seconds. The record can also be kept in a slow variable of the synapses themselves: after a burst of spikes, facilitation, as in the ``dash'' synapse of Chapter~\ref{sec:code}, lasts about a second, and between brief flashes the network can be almost silent: not a flip-flop but a charged capacitor~\cite{Mongillo2008}. Such ``quiet'' intervals during memory maintenance are observed~\cite{Stokes2015}, and storage without continuous spiking is cheaper in energy. Which of the two ways the cortex uses in which case is under debate.

And how does one erase? By Proposition~\ref{prop:monotone}, not by activity; something has to be \emph{removed}. The first way is to remove the support: if the group stays on only together with the drive from another group, the memory dies out along with it. The second is fatigue. In real synapses the stock of transmitter is depleted during prolonged activity~\cite{Tsodyks1997}, and in neurons slow currents that lower excitability build up. The feedback weakens, the upper intersection disappears, and the group goes out by itself within seconds. The result is a memory that is switched on by a signal and switched off only by time.

\section{Sequences}

Instead of a clock one can have an \emph{event-triggered timer}. Connect groups of \nt{GLUT} neurons into a chain: the first excites the second, the second the third, and so on. An external signal ignites the first group, and a wave of activity runs along the chain, each link one or two delays after the previous one and only once: right after a spike the neurons are refractory, and the wave has already moved on.

Such a chain measures time from an event: link number $k$ has fired, which means $k$ delays have passed since the start. Its outputs can be fed to the muscles to produce a sequence of movements that unfolds on its own. In songbirds, during singing, individual \nt{GLUT} neurons in one brain region fire strictly in turn, each at its own moment of the song, much like such a chain~\cite{Hahnloser2002}.

Unlike a clock, the chain is silent until triggered, runs once, and measures time only for those connected to it. The network still has no common beat.

\section{What is missing}

A lot can be built from \nt{GLUT} neurons alone, without inhibition. But three things are missing, and all three because of Proposition~\ref{prop:monotone}.

\emph{Selection.} The network must choose one direction, one action. If two stimuli are nearly equal, both outputs of the network in Fig.~\ref{fig:itd} will fire, and there is nothing to suppress the weaker in favor of the stronger.

\emph{Reset.} A memory cannot be erased by a signal.

\emph{Stability.} In a network whose connections do not arise from an engineer's plan, positive feedback loops are inevitable, and activity in them can build up like an avalanche (Chapter~\ref{sec:neurontypes}). The only brake is the same fatigue, slow and crude.

All three troubles are cured by one remedy: a neuron that cancels a spike with a spike. That is \nt{GABA}, and it is needed not to compute but to select, to reset, and to keep the computation under control.

\section{Exercises}

\begin{exercise}[\lvlA]
\label{ex:02:1}
\emph{A row of detectors for sound.} The signals from the ears run along the wires of Fig.~\ref{fig:itd} at $5$~m/s, and the largest arrival difference is $0.64\ms$. The detectors, neurons~\eqref{eq:glutneuron} with $\tau=0.5\ms$, $\theta=1.5w$, are spaced so that neighbors differ by $10$~\textmu s. How many detectors are in the row, and how many of them fire for one sound?
\end{exercise}

\begin{exercise}[\lvlB]
\label{ex:02:2}
\emph{Unequal inputs shift the window.} A coincidence detector~\eqref{eq:glutneuron} with $\tau=0.5\ms$, $\theta=1.5$ receives one spike from each of two inputs with weights $1$ and $0.8$. Find the coincidence window and its center. By how many microseconds will the row of Fig.~\ref{fig:itd} err if the input from one ear is $20\,\%$ weaker than from the other?
\end{exercise}

\begin{exercise}[\lvlB]
\label{ex:02:3}
\emph{Which networks do not oscillate.} The network $\tau\dot r_i=-r_i+f_i\bigl(\textstyle\sum_jw_{ij}r_j+I_i\bigr)$ with nondecreasing $f_i$ and $w_{ij}\ge0$ for $i\ne j$ is monotone and does not oscillate stably. Prove that a network with an even number of inhibitory connections in every cycle reduces to it by the substitution $r_i\to\pm r_i$. Can mutual inhibition of two groups oscillate? A \nt{GLUT}--\nt{GABA} loop?
\end{exercise}

\begin{exercise}[\lvlB]
\label{ex:02:4}
\emph{Design: a dual-rail adder.} Each bit is carried by a pair of wires $(x,\bar x)$. Build a full adder of three bits that outputs the pairs $(S,\bar S)$ for the sum and $(C,\bar C)$ for the carry from \nt{GLUT} neurons, giving weights and thresholds. Use no more than eight neurons. How many would a circuit of AND and OR elements, as in Fig.~\ref{fig:dualxor}, require?
\end{exercise}

\begin{exercise}[\lvlB]
\label{ex:02:5}
\emph{Simulation: how long a write takes.} The group of Fig.~\ref{fig:recurrentgroup}: $\tau\dot r=-r+f(r+I)$, $f(u)=1/\bigl(1+e^{-(u-0.5)/0.08}\bigr)$, in the lower state before the write; the drive $I$ is applied for time $T$. Find numerically the smallest $T$ that writes the bit at $I=0.6$, and the threshold $I_c$ below which no $T$ writes the bit.
\end{exercise}

\begin{exercise}[\lvlA]
\label{ex:02:6}
\emph{A chain as a timer.} A link of the chain is $100$ \nt{GLUT} neurons with a delay of $2\ms$ and an independent scatter of $0.2\ms$. (a)~How many neurons are needed to measure $1$~s, and what is the relative error? How does it depend on the interval? (b)~A real timer has a $5\,\%$ error at any interval. What source of scatter produces this?
\end{exercise}

\begin{exercise}[\lvlC]
\label{ex:02:7}
\emph{Research: flip-flop or capacitor.} $100$ neurons hold a bit for $10$~s. Flip-flop: each fires $20$ spikes per second. Capacitor: facilitation $u$ decays with $\tau_F=1$~s, the bit is read when $u\ge u_{\max}/2$, and a refresh is a burst of three spikes from each neuron. How much more economical is the capacitor, and what happens to the bit if the group is silenced for $200\ms$?
\end{exercise}

\chapter{\nt{GLUT} and \nt{GABA} Together}
\label{sec:gaba}

\section{The Rules of the Game}

A network of \nt{GLUT} neurons alone cannot select, cannot reset memory, and cannot keep activity from running away in an avalanche, all because of monotonicity (Proposition~\ref{prop:monotone}). Let us add the second most numerous neuron type, \nt{GABA}. The rules are as before: only what happens in a living organism, and no clock signal.

The neuron model is the same as~\eqref{eq:glutneuron}, but a spike from a \nt{GABA} neuron lowers the recipient's voltage instead of raising it. Weights now come in two signs, and the sender sets the sign, rule~\eqref{eq:dalesign}.

A few parameters of the circuit. A fifth to a quarter of cortical neurons are \nt{GABA}~\cite{Hendry1987}. Their outputs are short: they inhibit their neighbors, not distant regions. Inhibition arrives within a millisecond and lasts a few milliseconds. With no input a \nt{GABA} neuron is silent: to inhibit someone, it must itself be excited, usually by \nt{GLUT} neurons of the same network.

So a negative connection from a \nt{GLUT} neuron $A$ to a neuron $B$ has to be made through an intermediary: $A$ excites a \nt{GABA} neuron, and that one inhibits $B$. Changing the sign costs one neuron and one delay, about a millisecond.

For inhibition it matters not only \emph{from whom} a connection comes but also \emph{where} it lands: the landing site largely determines what the connection does~\cite{Markram2004}. \nt{GLUT} outputs land mainly on the recipient's input branches, its \emph{dendrites}, and carry data: each such input is one term of the sum. An output onto the cell body acts on the sum as a whole: that is how many fast \nt{GABA} neurons divide the recipient's response and set when it fires. An output onto the place where the recipient's spike is born has the last word: a veto over the entire output at once. And an output onto the terminal of another neuron's wire changes the release probability $p$ of one input line without touching the others~\cite{Rudomin1999}; this is how, in particular, the pain gate of Chapter~\ref{sec:pain} works. Outside the brain, outputs land on muscle fibers (Chapter~\ref{sec:muscles}) and on organs (Chapter~\ref{sec:chemoutput}), and some land nowhere at all and release their transmitter into the volume of the tissue or into the blood (Chapters~\ref{sec:serocomputer} and~\ref{sec:chemoutput}).

\section{NOT and Veto}

Can we now build NOT? Almost. A neuron that is active while its input is silent has to get its excitation from somewhere. In the living brain it has it: every neuron constantly receives background activity from thousands of others. Let the background keep neuron $Y$ active, and let input $x$ inhibit it through a \nt{GABA} intermediary. That is NOT.

More useful, however, is the \emph{veto}: ``$a$, but not when $b$'':
\begin{empheq}[box=\keybox]{equation}
y \;=\; a \wedge \bar b .
\label{eq:veto}
\end{empheq}
Neuron $Y$ receives excitation from $a$ and inhibition from a \nt{GABA} neuron that $b$ excites. No background is needed here: $a$ itself supplies the excitation. Everything can be assembled from vetoes and OR. For example, exclusive OR: $x\oplus y=(x\wedge\bar y)\vee(\bar x\wedge y)$ takes two veto neurons, two \nt{GABA} intermediaries, and one OR neuron.

\begin{keyblock}
\begin{proposition}[Completeness of \nt{GLUT}+\nt{GABA}]
\label{prop:gabacomplete}
A network of \nt{GLUT} and \nt{GABA} neurons can compute any logical function of its inputs, with each bit carried by a single wire. The on/off sensors of Chapter~\ref{sec:glutcomputer} are no longer needed for this. If the function must output one when all inputs are silent, a source of background activity is needed.
\end{proposition}
\end{keyblock}

Proof: the veto~\eqref{eq:veto} together with OR is functionally complete given a constant one, because NOT $x$ is the veto ``one, but not when $x$.'' And functions that output zero when all inputs are silent are assembled from vetoes and OR even without the one.

\section{Direction of Motion}

A veto in time, like the AND in time of Chapter~\ref{sec:glutcomputer}, gives a detector of the direction of motion.

Let two neighboring sensors, $A$ and $B$, be a distance $\Delta x$ apart. The output neuron $Y$ receives excitation from $B$ and inhibition from $A$ through a \nt{GABA} intermediary, with delay $d$ (Fig.~\ref{fig:direction}). A spot of light moving at speed $v$ from $A$ to $B$ passes $A$ at time $0$, and the inhibition reaches $Y$ at time $d$; it passes $B$ at time $\Delta x/v$, and the excitation arrives at that moment. If these moments coincide to within the window~\eqref{eq:window}, the inhibition cancels the excitation, and $Y$ stays silent. This happens at a speed of about
\begin{empheq}[box=\keybox]{equation}
v^* \;=\; \frac{\Delta x}{d} .
\label{eq:vstar}
\end{empheq}
For motion from $B$ to $A$, the excitation from $B$ arrives at time $0$, and the inhibition from $A$ only at time $\Delta x/v+d$, when $Y$ has already fired.

\begin{figure}[t]
\centering
\begin{tikzpicture}[>=Stealth,
  nrn/.style={circle,draw,very thick,fill=blue!6,minimum size=0.9cm,inner sep=0pt},
  gab/.style={nrn,fill=red!10},
  inh/.style={-{Bar[width=7pt]},thick,red!70!black}]
\node[nrn] (A) at (0,2) {$A$};
\node[nrn] (B) at (3,2) {$B$};
\node[gab] (G) at (0.9,0.6) {\small \nt{GABA}};
\node[nrn] (Y) at (3,-0.6) {$Y$};
\draw[->,thick] (A) -- (G);
\draw[inh] (G) -- node[below left=-2pt] {\scriptsize delay $d$} (Y);
\draw[->,thick] (B) -- (Y);
\draw[->,very thick,red!70!black] (Y) -- (4.6,-0.6) node[right,black] {\small output};
\draw[->,thick,orange!85!black] (-0.6,2.9) -- (3.6,2.9) node[right,black] {\small silent};
\draw[<-,thick,orange!85!black] (-0.6,3.4) -- (3.6,3.4) node[right,black] {\small responds};
\foreach \yy/\dd in {2.9/-1,3.4/1}{
  \foreach \k in {1,2,3}{\draw[orange!70!yellow,line width=0.8pt] (1.5+\dd*0.12+\dd*0.13*\k,\yy+0.1-0.1*\k+0.1) -- ++(\dd*0.25,0);}
  \fill[yellow!70!orange,draw=orange!70!black] (1.5,\yy) circle (0.14);}
\node[font=\scriptsize,align=center] at (1.5,3.95) {spot of light};
\end{tikzpicture}
\caption{Direction-of-motion detector. $B$ excites $Y$; $A$ inhibits it through a \nt{GABA} intermediary with delay $d$. For motion from $A$ to $B$ at a speed of about $\Delta x/d$, the inhibition cancels the excitation; for motion in the opposite direction it arrives too late. The yellow spot with streaks is light moving over the sensors.}
\label{fig:direction}
\end{figure}

In the retina, selectivity for the direction of motion apparently arises in just this way: from asymmetric inhibition on the side from which motion should not evoke a response~\cite{Barlow1965}.

\section{Subtraction or Division}

So far our inhibition has subtracted. In fact it works more subtly.

A synapse opens a channel, that is, it switches on a conductance. For an excitatory synapse the equilibrium potential lies far above threshold, about $70\mV$ above rest: an open channel pulls the voltage up. For an inhibitory \nt{GABA} synapse the channel passes chloride, whose equilibrium potential lies near rest, and the open channel pulls the voltage not down but \emph{toward rest}. Measuring voltage from rest, the steady-state voltage of a membrane with leak conductance $g_L$, excitatory conductance $g_E$, and inhibitory conductance $g_I$ is
\begin{empheq}[box=\keybox]{equation}
V_\infty \;=\; \frac{g_E\,E_E}{g_L+g_E+g_I} ,
\label{eq:shunt}
\end{empheq}
where $E_E\approx70\mV$ is the equilibrium potential of the excitatory synapse. This is Ohm's law for a voltage divider: $g_E$ pulls toward $E_E$, while $g_L$ and $g_I$ pull toward zero.

$g_I$ appears not in the numerator with a minus sign but in the denominator: inhibition does not \emph{subtract} from excitation, it \emph{divides} it (Fig.~\ref{fig:shunt}). Such inhibition is called shunting inhibition: the open chloride channels short-circuit the membrane. For the network this is gain control. If $g_I$ is proportional to the total activity of the neighbors, each neuron responds not to the absolute strength of its input but to its share of the total activity. Such division by total activity occurs in many parts of the brain and is considered one of its canonical operations~\cite{Carandini2012}: one and the same network works with both weak and strong input.

A caveat: formula~\eqref{eq:shunt} is about a neuron that does not fire. In a firing neuron the voltage stays near threshold all the time, the current through the inhibitory conductance is nearly constant, and the shunt acts on the \emph{rate} of spikes mainly by subtraction~\cite{HoltKoch1997}. The rate is divided if the neuron simultaneously receives balanced excitatory and inhibitory background conductances, as in the jittery, balanced regime of living cortex~\cite{Chance2002}. So the brain gets division not from a shunt alone but from inhibition together with background noise~\cite{Carandini2012}.

\begin{figure}[t]
\centering
\begin{tikzpicture}[>=Stealth,x=1.25cm,y=0.05cm]
\draw[->,gray!70!black] (0,0) -- (5.4,0) node[right,black] {\small $g_E/g_L$};
\draw[->,gray!70!black] (0,0) -- (0,68) node[above,black] {\small $V_\infty$, mV};
\foreach \v in {20,40,60}{\draw[gray!60!black] (-0.05,\v) -- (0.05,\v) node[left=3pt,font=\scriptsize] {\v};}
\foreach \x in {1,2,3,4,5}{\draw[gray!60!black] (\x,-1.5) -- (\x,1.5) node[below=5pt,font=\scriptsize] {\x};}
\draw[very thick,blue!60!black] plot[domain=0:5,samples=60] (\x,{70*\x/(1+\x)});
\draw[very thick,red!70!black] plot[domain=0:5,samples=60] (\x,{70*\x/(3+\x)});
\draw[thick,densely dashed,gray!50!black] plot[domain=0.56:5,samples=60] (\x,{70*\x/(1+\x)-25});
\node[right,font=\small,blue!60!black] at (5.02,58.3) {no inhibition};
\node[right,font=\small,red!70!black] at (5.02,45.5) {divides: $g_I=2g_L$};
\node[right,font=\small,gray!40!black] at (5.02,32.5) {subtracts $25\mV$};
\end{tikzpicture}
\caption{Shunting inhibition divides the voltage rather than subtracting from it. Blue curve: the steady-state voltage~\eqref{eq:shunt} without inhibition; red: with inhibitory conductance $g_I=2g_L$. The red curve is the blue one stretched threefold horizontally: as if the input had been divided by $1+g_I/g_L=3$. Dashed: inhibition that subtracts a constant $25\mV$; the curve drops as a whole, and its slope does not change.}
\label{fig:shunt}
\end{figure}

There is a second consequence as well. The membrane time constant is $\tau_{\text{eff}}=C/(g_L+g_E+g_I)$, and inhibition shortens it. The coincidence window~\eqref{eq:window} is proportional to $\tau$, so inhibition that arrives together with excitation \emph{narrows the window}. A circuit in which an input excites a neuron directly and at the same time, with a millisecond delay, inhibits it through a \nt{GABA} intermediary is one of the most common in the brain: the neuron has time to respond only to what arrived in the first millisecond or two~\cite{Pouille2001}.

\section{Selection}

Now for the main thing the \nt{GLUT} network lacked: choosing one of several. Take two groups of \nt{GLUT} neurons with inputs $I_1$ and $I_2$. Each inhibits the other through its own \nt{GABA} intermediaries with strength $\beta$ (Fig.~\ref{fig:wta}). For the rates $r_1,r_2$ we get
\begin{equation}
\tau\,\frac{\dd r_1}{\dd t} = -r_1 + \bigl[I_1-\beta r_2\bigr]_+ ,\qquad
\tau\,\frac{\dd r_2}{\dd t} = -r_2 + \bigl[I_2-\beta r_1\bigr]_+ ,
\label{eq:wta}
\end{equation}
where $[u]_+=\max(u,0)$: a rate cannot be negative.

\begin{figure}[t]
\centering
\begin{tikzpicture}[>=Stealth,
  nrn/.style={circle,draw,very thick,fill=blue!6,minimum size=0.9cm,inner sep=0pt},
  gab/.style={nrn,fill=red!10,minimum size=0.75cm},
  inh/.style={-{Bar[width=7pt]},thick,red!70!black}]
\node[nrn] (E1) at (0,0) {$r_1$};
\node[nrn] (E2) at (4,0) {$r_2$};
\node[gab] (G1) at (1.4,1.1) {};
\node[gab] (G2) at (2.6,-1.1) {};
\draw[->,thick] (E1) -- (G1); \draw[inh] (G1) -- (E2);
\draw[->,thick] (E2) -- (G2); \draw[inh] (G2) -- (E1);
\draw[->,thick,blue!60!black] (-1.6,0) node[left,black] {$I_1$} -- (E1);
\draw[->,thick,blue!60!black] (5.6,0) node[right,black] {$I_2$} -- (E2);
\end{tikzpicture}
\caption{Choosing one of two. Two groups of \nt{GLUT} neurons (blue) inhibit each other through \nt{GABA} intermediaries (red).}
\label{fig:wta}
\end{figure}

While both neurons are active, system~\eqref{eq:wta} is linear, and its behavior is set by two eigenvalues: $-(1+\beta)/\tau$ for equal deviations of the two rates and $-(1-\beta)/\tau$ for opposite ones. With weak inhibition, $\beta<1$, both eigenvalues are negative: both neurons stay active, and inhibition merely damps the weaker one. With strong inhibition, $\beta>1$, the second eigenvalue becomes positive: any difference between $r_1$ and $r_2$ grows until one neuron falls completely silent.

\begin{keyblock}
\begin{proposition}[Winner-take-all]
\label{prop:wta}
If the mutual inhibition exceeds one, $\beta>1$, the state in which both groups are active is unstable. The network settles into a state in which one group is active and the other is silent. The winner, with rate $r_1=I_1$, keeps its victory as long as $I_2<\beta I_1$.
\end{proposition}
\end{keyblock}

We checked this on the model. With $\beta=0.5$ and inputs $1.0$ and $0.9$, both groups are active, with rates $0.73$ and $0.53$. With $\beta=2$ and the same inputs, the second group is completely silent. Such a choice has memory: if the inputs are then swapped, the previous winner remains the winner, because $1.0<2\cdot0.9$.

With a hundred groups it is the same: each inhibits all the others, and one survives.

\section{Resetting Memory}

The memory of Chapter~\ref{sec:glutcomputer} could be switched off only by fatigue. Let us add a \nt{GABA} neuron that is excited by a ``reset'' signal and inhibits the group. The inhibition lowers the curve $f(wr+I)$ in Fig.~\ref{fig:bistable}, the upper intersection disappears, and the group goes out, and it stays in the lower state after the reset ends. Now the memory is written by one signal and erased by another, like a flip-flop with set and reset inputs.

It is even more elegant to connect two such groups by mutual inhibition, as in Fig.~\ref{fig:wta}, and give each feedback onto itself. A signal at the input of one of the groups switches the cell into that group's state, and the cell remembers the last decision. This is a direct analog of a latch made of two cross-coupled NOR gates.

\section{Stability and Rhythm}

The third trouble of the \nt{GLUT} network remains: the avalanche. Take a group of \nt{GLUT} neurons with self-feedback $w_{EE}$ and a group of \nt{GABA} neurons that the first group excites with strength $w_{IE}$ and that inhibit the first group with strength $w_{EI}$. For the rates $E$ and $I$ we get
\begin{equation}
\tau_E\frac{\dd E}{\dd t} = -E + w_{EE}E - w_{EI}I + h, \qquad
\tau_I\frac{\dd I}{\dd t} = -I + w_{IE}E ,
\label{eq:ei}
\end{equation}
where $h$ is the external input~\cite{WilsonCowan1972}. Without inhibition, for $w_{EE}>1$ activity grows without bound: each spike gives rise to more than one next spike. With inhibition the steady-state rate
\begin{equation}
E^* \;=\; \frac{h}{1-w_{EE}+w_{EI}w_{IE}}
\label{eq:eistar}
\end{equation}
is finite if the denominator is positive. But that is not enough: the equilibrium must also attract. Linear analysis of~\eqref{eq:ei} gives two conditions.

\begin{keyblock}
\begin{proposition}[Stability conditions]
\label{prop:ei}
Equilibrium~\eqref{eq:eistar} is stable if the inhibition is
\begin{enumerate}
\item[(i)] \emph{strong enough}: $w_{EI}\,w_{IE} > w_{EE}-1$;
\item[(ii)] \emph{fast enough}: $\tau_I < \tau_E/(w_{EE}-1)$.
\end{enumerate}
If (i) is violated, activity builds up like an avalanche. If (i) holds but (ii) is violated, the inhibition lags, and activity begins to oscillate.
\end{proposition}
\end{keyblock}

Condition~(i) is the determinant of the matrix of system~\eqref{eq:ei}, condition~(ii) its trace. The meaning of the second is clear to anyone who has tuned a controller: delayed feedback does not damp a deviation but drives it into oscillation.

\begin{figure}[t]
\centering
\begin{tikzpicture}[>=Stealth]
\draw[->,gray!70!black] (0,0) -- (10.9,0) node[below left,black] {\small $t$, ms};
\draw[->,gray!70!black] (0,0) -- (0,3.3) node[left,black] {\small $E$};
\foreach \t in {0,50,100,150}{\draw[gray!70!black] (\t*0.07,0.06) -- (\t*0.07,-0.06) node[below,font=\scriptsize] {\t};}
\input{figures/ei_traces}
\node[font=\small,gray!40!black,anchor=west] at (0.9,3.15) {no inhibition: avalanche};
\node[font=\small,blue!60!black,anchor=west] at (7.4,1.25) {fast inhibition};
\node[font=\small,red!70!black,anchor=west] at (7.4,2.55) {slow inhibition};
\end{tikzpicture}
\caption{Rate of a \nt{GLUT} group with feedback $w_{EE}=1.5$, $\tau_E=10\ms$, after the input is switched on. Without inhibition (dashed) activity grows without bound. With inhibition of strength $w_{EI}w_{IE}=2$ and $\tau_I=2\ms$ (blue curve) it settles at $E^*=h/(1-1.5+2)=0.67h$. With the same inhibition but slow, $\tau_I=30\ms$ (red curve), activity oscillates.}
\label{fig:ei}
\end{figure}

The cortex is thought to operate in exactly this regime: its own feedback is strong enough to run away into an avalanche without inhibition, and only fast inhibition holds it at its operating point~\cite{Tsodyks1997isn}.

This regime has a paradoxical signature by which it can be recognized from outside~\cite{Tsodyks1997isn}. Add to~\eqref{eq:ei} an external input $h_I$ directly to the \nt{GABA} group. The steady-state rates become
\begin{equation}
E^* \;=\; \frac{h-w_{EI}\,h_I}{D},\qquad
I^* \;=\; \frac{w_{IE}\,h+(1-w_{EE})\,h_I}{D},\qquad
D \;=\; 1-w_{EE}+w_{EI}w_{IE} .
\label{eq:isn}
\end{equation}
If $w_{EE}>1$, the coefficient of $h_I$ in the second formula is negative: push the \nt{GABA} neurons, and their rate \emph{drops}. The reason is the loop. The extra inhibition damps the \nt{GLUT} group, which then excites the \nt{GABA} neurons less, and they lose more than they gained. With the numbers of Fig.~\ref{fig:ei} ($w_{EE}=1.5$, $w_{EI}w_{IE}=2$, $D=1.5$), each unit of added input lowers $I^*$ by a third of a unit. This signature has been found in the cortex: when the response of visual cortex neurons is suppressed by a surrounding stimulus, the inhibitory current they receive does not rise but falls together with the excitatory current~\cite{Ozeki2009}. Later the same signature was found in various areas and layers of the cortex~\cite{Sanzeni2020}.

Oscillations when condition~(ii) is violated also occur in the brain. The loop ``\nt{GLUT} excites \nt{GABA}, \nt{GABA} inhibits \nt{GLUT}'' with a delay of a few milliseconds generates a rhythm with a period on the order of $12$--$50\ms$ (frequency $20$--$80$\,Hz), and such fast rhythms in the cortex are well known~\cite{Cardin2009,Buzsaki2012}. This is not a computer's clock: the rhythm is local and arises only when the network is active. But over a few cycles it divides time into windows in which neurons can fire.

\section{Summary}

\begin{table}[t]
\centering
\small
\begin{tabular}{>{\raggedright}p{3.3cm}>{\raggedright}p{4.4cm}>{\raggedright\arraybackslash}p{4.4cm}}
\toprule
& \nt{GLUT} only & \nt{GLUT} + \nt{GABA} \\
\midrule
NOT & only through on/off sensors & veto $a\wedge\bar b$; NOT given background activity \\
wires per bit & two & one \\
direction of motion & no & veto with a delay \\
gain control & no & division: shunt together with background noise \\
coincidence window & $\sim\tau$ & narrowed by inhibition \\
choosing one of many & no & mutual inhibition, $\beta>1$ \\
memory reset & only by fatigue & by a signal through \nt{GABA} \\
stability & only by fatigue & fast negative feedback \\
rhythm & not needed & arises with slow inhibition \\
\bottomrule
\end{tabular}
\caption{What \nt{GABA} adds to a network of \nt{GLUT} neurons.}
\label{tab:gaba}
\end{table}

\nt{GABA} adds almost no new \emph{computations} to the network (all of this could have been computed with dual-rail sensors); what it adds is \emph{control} (Table~\ref{tab:gaba}): selection, reset, gain control, and stability. In the brain, excitation carries the data, and inhibition decides which data get through, when, and how loudly. For every fourth or fifth neuron of the cortex, that is a very sensible division of labor.

\section{Exercises}

\begin{exercise}[\lvlA]
\label{ex:03:1}
\emph{The shunt in numbers.} $E_E=70\mV$. Using~\eqref{eq:shunt}, find $V_\infty$ for $g_E=g_L$ and $4g_L$, without inhibition and with a shunt $g_I=2g_L$. What $V_\infty$ at $g_E=4g_L$ would a subtraction give that removes as much at $g_E=g_L$ as the shunt does? By what factor does the shunt narrow the coincidence window?
\end{exercise}

\begin{exercise}[\lvlB]
\label{ex:03:2}
\emph{Deriving the stability conditions.} Find the trace and determinant of the matrix of the linear system~\eqref{eq:ei} and derive from them both conditions of Proposition~\ref{prop:ei}. At the boundary of condition~(ii) the equilibrium loses stability through oscillation. Find the period of these oscillations for the numbers of Fig.~\ref{fig:ei}.
\end{exercise}

\begin{exercise}[\lvlB]
\label{ex:03:3}
\emph{Rhythm from a delay.} Replace the slow inhibition of model~\eqref{eq:ei} with fast inhibition after a delay $d$: $\tau_E\dot E=-E(t)-GE(t-d)+h$. For $\tau_E=10\ms$, find the smallest delay $d_c$ that produces oscillations, and their period, for $G=2$ and $G=5$. Do they fall within the fast cortical rhythms, $12$--$50\ms$, unlike in Exercise~\ref{ex:03:2}?
\end{exercise}

\begin{exercise}[\lvlB]
\label{ex:03:4}
\emph{Simulation: selection with memory.} Integrate~\eqref{eq:wta} with $\beta=2$, $\tau=1$. (a)~With $I_1=1$, slowly raise $I_2$ from $0$ to $3$ and lower it back: at what $I_2$ does the network switch? (b)~By how much does the decision time from silence grow when the input difference $\varepsilon$ is reduced tenfold?
\end{exercise}

\begin{exercise}[\lvlB]
\label{ex:03:5}
\emph{Design: a detector for a band of speeds.} The detector of Fig.~\ref{fig:direction} must stay silent for motion from $A$ to $B$ with a travel time of $5$--$20\ms$. A \nt{GABA} intermediary with delay $d_k$ vetoes $Y$ over the interval $[d_k,\,d_k+5\ms]$ after a spike from $A$; the excitation from $B$ arrives immediately. How many intermediaries are needed, and with what delays?
\end{exercise}

\begin{exercise}[\lvlB]
\label{ex:03:6}
\emph{Vetoes and background.} How many neurons does the parity $x\oplus y\oplus z$ require in the general scheme ``a \nt{GABA} intermediary per input, a veto for each input combination with $f=1$, a common OR''? Build it from two neurons, a \nt{GABA} and a \nt{GLUT}, that receive the inputs directly, giving weights and thresholds.
\end{exercise}

\begin{exercise}[\lvlC]
\label{ex:03:7}
\emph{Design: choosing one of a hundred.} Each of $100$ \nt{GLUT} groups must inhibit all the others equally, but not itself. Linear \nt{GABA} intermediaries are excited by the groups and inhibit them; the groups may excite one another, but not themselves. What is the smallest number of intermediaries? And if there are no direct excitatory connections at all?
\end{exercise}

\begin{exercise}[\lvlC]
\label{ex:03:8}
\emph{Research: when inhibition is paradoxical.} In the model of Fig.~\ref{fig:ei} ($w_{EI}=2$, $w_{IE}=1$), the \nt{GLUT} group is given the transfer function $f(u)=[u]_+^2$, while the \nt{GABA} group stays linear. Above what input $h_c$ does an extra drive to the \nt{GABA} group lower its own rate? Check by simulation.
\end{exercise}

\chapter{Morse Code: What We Know About the Language of \nt{GLUT} and \nt{GABA} Neurons}
\chaptermark{Morse Code}
\label{sec:code}

\section{An alphabet of one symbol}

Morse code has two symbols, the dot and the dash, and pauses of three lengths between them. A neuron's alphabet, as we saw in Chapter~\ref{sec:neurontypes}, is poorer still: it has one symbol. A spike lasts about a millisecond, and all the spikes of a neuron have the same height and shape. So the whole message is in the pauses, that is, in the sequence of times $t_1,t_2,t_3,\dots$ It is a code with no dashes, in which only the rhythm of the dots carries meaning (Fig.~\ref{fig:morse}).

\begin{figure}[t]
\centering
\begin{tikzpicture}[>=Stealth,x=0.08cm,y=1cm]
% Morse letter: dot, dash, dot
\node[left,font=\small] at (-6,2.55) {Morse:};
\fill[black] (0,2.45) rectangle (6,2.65);
\fill[black] (14,2.45) rectangle (32,2.65);
\fill[black] (40,2.45) rectangle (46,2.65);
\node[right,font=\small,gray!40!black] at (52,2.55) {dot, dash, dot: the meaning is in durations and pauses};
% stimulus onset
\node[left,font=\small] at (-6,0.4) {neuron:};
\draw[->,thick,green!45!black] (0,-0.55) -- (0,-0.05);
\node[below,font=\scriptsize,green!45!black] at (0,-0.55) {stimulus};
\draw[gray!70!black] (0,0) -- (152,0) node[right,black] {\small $t$};
% spikes: single ones blue, the burst red
\foreach \s in {14,26,41,88,112,131}{\draw[very thick,blue!60!black] (\s,0) -- (\s,0.8);}
\foreach \s in {60,63,66,69}{\draw[very thick,red!70!black] (\s,0) -- (\s,0.8);}
% latency
\draw[<->,thick] (0,1.1) -- node[above,font=\scriptsize] {$t_1$: first-spike time} (14,1.1);
% burst
\draw[decorate,decoration={brace,amplitude=4pt},red!70!black] (58,0.95) -- (71,0.95) node[midway,above=4pt,font=\scriptsize] {burst: a ``dash''};
% counting windows
\foreach \a/\b/\n in {0/50/3,50/100/5,100/150/2}{
  \draw[decorate,decoration={brace,amplitude=4pt,mirror},gray!50!black] (\a+1,-0.2) -- (\b-1,-0.2) node[midway,below=4pt,font=\small] {$n=\n$};}
\node[font=\scriptsize,gray!40!black] at (75,-1.05) {rate code: the number of spikes $n$ in each window};
\foreach \x in {0,50,100,150}{\draw[gray!60!black] (\x,-0.06) -- (\x,0.06);}
\end{tikzpicture}
\caption{Morse code and the neuron's code. The same spike train can be read in different ways: count the spikes in $50\ms$ windows (Section~\ref{sec:rate}), measure the latency $t_1$~\eqref{eq:latency}, or notice the burst, a ``dash''~\eqref{eq:burst}.}
\label{fig:morse}
\end{figure}

What pauses are possible at all? From below they are limited by refractoriness: after a spike a neuron cannot fire again for about a millisecond, so more than a few hundred spikes per second do not occur. From above nothing limits them: a neuron can stay silent for minutes. In cortical \nt{GLUT} neurons, rates range from a fraction of a spike to a few tens of spikes per second, and most neurons spend most of the time near the lower end.

In Chapters~\ref{sec:glutcomputer} and~\ref{sec:gaba} we adopted one or another way of writing numbers with spikes as we went along. Now let us ask directly: what language do \nt{GLUT} and \nt{GABA} neurons use to talk to each other? There is no complete answer; this language has not been deciphered. But some things about it are firmly known, and almost everything known can be obtained by reasoning like a communications engineer: channel capacity, noise, receiver, cost of transmission.

\section{How much can be transmitted}

Let us start with an upper bound. Suppose the recipient resolves spike times with precision $\Delta t$: shifts smaller than $\Delta t$ are indistinguishable to it. Cut time into bins of length $\Delta t$, shorter than the refractory time, so that each bin holds either one spike or none (Fig.~\ref{fig:cells}). At mean rate $r$, a spike falls into a bin with probability $p=r\Delta t$. The stream carries the most information when the bins are independent, and then each bin gives an entropy of $-p\log_2p-(1-p)\log_2(1-p)\approx p\log_2(e/p)$ bits for $p\ll1$. Dividing by $\Delta t$, we get the limiting transmission rate and its share per spike~\cite{MacKay1952}:
\begin{empheq}[box=\keybox]{equation}
R_{\max} \;\approx\; r\,\log_2\frac{e}{r\,\Delta t}\ \ \text{bit/s},
\qquad
\frac{R_{\max}}{r} \;\approx\; \log_2\frac{e}{r\,\Delta t}\ \ \text{bits per spike}.
\label{eq:spikeentropy}
\end{empheq}
For $r=10$ spikes per second and $\Delta t=1\ms$, this is about eight bits per spike and eighty bits per second.

\begin{figure}[t]
\centering
\begin{tikzpicture}[>=Stealth,x=0.5cm,y=1cm]
% time cut into cells of width Delta t; a cell holds one impulse or none
\foreach \k in {0,...,23}{\draw[gray!55] (\k,0) rectangle (\k+1,0.7);}
\foreach \k in {2,7,8,15,21}{
  \fill[red!10] (\k,0) rectangle (\k+1,0.7);
  \draw[very thick,red!70!black] (\k+0.5,0.05) -- (\k+0.5,0.65);}
\foreach \k in {0,...,23}{
  \pgfmathtruncatemacro{\bit}{(\k==2)||(\k==7)||(\k==8)||(\k==15)||(\k==21)}
  \ifnum\bit=1 \def\bitcol{red!70!black}\else \def\bitcol{gray!60!black}\fi
  \node[font=\scriptsize,text=\bitcol] at (\k+0.5,-0.25) {\bit};}
\draw[<->,gray!60!black] (0,0.95) -- (1,0.95) node[midway,above,font=\scriptsize,black] {$\Delta t$};
\node[font=\small,anchor=east] at (-0.3,0.35) {spikes};
\node[font=\small,anchor=east] at (-0.3,-0.25) {bits};
\draw[decorate,decoration={brace,amplitude=5pt,mirror},gray!60!black] (0,-0.55) -- (24,-0.55)
  node[midway,below=6pt,font=\scriptsize,black,align=center] {a recipient that only counts: ``$n=5$'';\\ a recipient that sees the bins: which $5$ of the $24$, that is $\log_2\binom{24}{5}\approx15$ bits};
\end{tikzpicture}
\caption{Where capacity~\eqref{eq:spikeentropy} comes from. Time is cut into bins of length $\Delta t$; each bin either has a spike ($1$) or not ($0$): the spike train is a binary string. A spike falls into a bin with probability $p=r\Delta t$. The finer the bins, the longer the string and the more bits; a spike counter loses almost everything recorded in \emph{where} the ones are.}
\label{fig:cells}
\end{figure}

Bits per spike depend on timing precision only logarithmically: at a precision of $10\ms$, about five bits per spike remain. But if the recipient only counts spikes per second, then at $r=10$ it gets less than two bits for the whole second (Section~\ref{sec:rate}).

\begin{keyblock}
\begin{proposition}[Capacity of a spike channel]
\label{prop:capacity}
A spike train with mean rate $r$, read with timing precision $\Delta t$, carries no more than~\eqref{eq:spikeentropy}. Almost all of this capacity lies in the \emph{timing} of the spikes: a recipient that only counts spikes extracts tens of times less from the same stream than one that resolves their times to within a millisecond.
\end{proposition}
\end{keyblock}

This is an upper bound, not a measurement. Measurements on sensory neurons whose stimulus is known give from one to several bits per spike; in the best cases this is up to half the bound computed for their precision~\cite{Rieke1997}. So at least at the input to the nervous system, spike timing is used rather than wasted.

\section{A noisy channel}

Capacity~\eqref{eq:spikeentropy} was computed without noise. About where noise arises, three facts are firmly established.

\emph{The spike generator itself is precise.} If a cortical neuron, taken out of the brain together with a slice of tissue, is given the same rapidly fluctuating current many times over, its spikes repeat with a precision of about a millisecond~\cite{Mainen1995}. If the current is constant, the spike times drift from trial to trial: precise timing is produced by fast transients of the input, not by the neuron itself.

\emph{The synapse is unreliable.} A spike that reaches a \nt{GLUT} synapse releases a packet of transmitter only with some probability $p$. At hippocampal synapses more than half of the spikes simply never reach the recipient, and the cause is precisely the randomness of release, not conduction failures~\cite{AllenStevens1994}. To a communications engineer this is an erasure channel; several synapses between two neurons partly save the situation.

\emph{In living cortex the spike train looks random.} The intervals between spikes are scattered roughly as in a random Poisson stream, and when the same stimulus is repeated, the spike count per window varies so that its variance is close to its mean~\cite{Softky1993}.

How does a precise element produce a random stream? A neuron receives thousands of inputs, each unreliable, while excitation and inhibition are balanced, as in Chapter~\ref{sec:gaba}. The membrane voltage jitters near threshold, and these jitters set the spike times. Whether this is noise or a message we cannot read is largely an open question. Part of the ``noise'' has already turned out to be a message: in the visual cortex a considerable share of the variability follows the animal's own movements~\cite{Stringer2019}. The difficulty here is fundamental: to someone who does not know the code, a well-compressed message is indistinguishable from a random stream.

\section{Rate: slow but reliable}
\label{sec:rate}

The simplest code is the \emph{rate code}: the number is written in how many spikes the neuron fires in a window $T$. Neither erasures nor timing jitter hurt such a code. But it has a price. If the stream is Poisson, the spike count $n$ in a window has mean $rT$ and spread $\sqrt{rT}$:
\begin{empheq}[box=\keybox]{equation}
\frac{\sigma_n}{\bar n} \;=\; \frac{1}{\sqrt{r\,T}} .
\label{eq:rateerror}
\end{empheq}
To learn a rate to within ten percent takes a hundred spikes, which at twenty spikes per second is five seconds. Neighboring levels are distinguishable if they are a couple of spreads apart, so up to rate $r$ in time $T$ about $\sqrt{rT}$ levels can be distinguished: at $r=10$ and $T=1\,\text{s}$ that is three or four levels, less than two bits.
The brain does not work that slowly. A person recognizes an animal in a picture flashed for an instant, and the response in the brain appears after about $150\ms$~\cite{Thorpe1996}. By a rough estimate, in that time the signal passes through about ten successive processing stages, $10$--$20\ms$ per stage. At the usual cortical rates, each neuron manages to fire zero or one spike per stage, and its rate in such a window is undefined. There are two ways out: take many neurons, or look at timing.

\emph{Many neurons.} Suppose $N$ neurons carry the same number, and the recipient adds up their spikes. If their noise is independent, $rT$ in~\eqref{eq:rateerror} is replaced by $NrT$: a hundred neurons at $r=20$ give thirty spikes in $15\ms$ and a precision of about twenty percent. Space substitutes for time. However, the noise of neighboring cortical neurons is not independent: the correlation coefficient of the spike counts of a pair of neighbors is typically about $0.1$~\cite{Zohary1994}. If it equals $c$, the variance of the average over $N$ neurons is
\begin{empheq}[box=\keybox]{equation}
\sigma^2_N \;=\; \sigma^2_1\,\frac{1+(N-1)\,c}{N}
\;\xrightarrow[N\to\infty]{}\; c\,\sigma^2_1 .
\label{eq:correlated}
\end{empheq}

\begin{keyblock}
\begin{proposition}[The limit of averaging]
\label{prop:pooling}
With correlated noise, averaging over a group reduces the error by no more than a factor of $1/\sqrt{c}$. For $c=0.1$ that is threefold, and this gain is reached already with a few tens of neurons; adding more is useless.
\end{proposition}
\end{keyblock}

Not every correlation is equally harmful. What limits precision is only the part of the shared noise that \emph{looks like the signal itself}, that is, that shifts the whole group the same way a change in the measured quantity would~\cite{MorenoBote2014}. How much such noise the brain actually has is still being worked out.

There is another way to use many neurons: write the number in \emph{which} neuron fired, as in the sound-direction finder (Fig.~\ref{fig:itd}). This is the \emph{place} code, the most reliable one known: for sensory neurons the index of the wire tells which point of the skin was touched or what pitch a sound has, and this has been established many times over. It is expensive in neurons but fast: a message takes one delay.

\section{Timing: fast, but measured from what?}

The second way out is to write the number in the time of a spike. Take neuron~\eqref{eq:glutneuron} and, from time $t=0$, apply a constant input that would raise the voltage to level $U$. The voltage grows as $V=U\bigl(1-e^{-t/\tau}\bigr)$ and reaches the threshold $\theta$ at time
\begin{empheq}[box=\keybox]{equation}
t_1 \;=\; \tau\,\ln\frac{U}{U-\theta}, \qquad U>\theta .
\label{eq:latency}
\end{empheq}
A strong input gives an early spike, a weak one a late spike, a subthreshold one none at all. With $\tau=10\ms$, input $U=4\theta$ gives the first spike after $2.9\ms$, $U=2\theta$ after $6.9\ms$, and $U=1.2\theta$ after $17.9\ms$. A single spike carries a number.

But from what moment is $t_1$ measured? There is no clock, and the recipient does not know when the stimulus began. There are three answers.

\emph{Relative latency.} Two neurons see the same stimulus, and the difference of their latencies $t_1^A-t_1^B$ depends only on their inputs, not on the moment of onset. Coincidence detectors with delays compare the times (Fig.~\ref{fig:itd}). If $N$ neurons each fire one spike, the very \emph{order} of their arrival carries up to $\log_2N!$ bits: for ten neurons about twenty-two bits, whereas the answer ``fired or not'' would give no more than ten. In the retina it has been shown that an image can be reconstructed quickly and reliably from the relative latencies of the first spikes of its output neurons~\cite{Gollisch2008}.

\emph{A network burst as the start of a word.} An abrupt change in the stimulus makes many neurons fire almost simultaneously. Such a network burst itself serves as a start mark, like the long pause between words in Morse code, and the recipient measures latencies from it.

\emph{A local rhythm.} In Chapter~\ref{sec:gaba} the \nt{GLUT}--\nt{GABA} loop generated a local rhythm with a period of $12$--$50\ms$. It is not a clock, but within its cycle a neuron with strong input manages to fire before a neuron with weak input, and~\eqref{eq:latency} turns into a \emph{phase} code: the number is written in which part of the cycle the spike arrived. Such a code has been observed: neurons responsible for a particular place in space fire earlier and earlier in the cycle of a slow local rhythm as the animal moves through ``its'' place~\cite{OKeefe1993}. How widely the brain uses phase is still unknown.

\section{The recipient chooses the code}

A spike train contains rate, timing, and order. Which of these gets read is decided by the recipient, and it decides with one parameter, its time constant $\tau$.

Average equation~\eqref{eq:glutneuron} over time. If independent streams with rates $r_j$ arrive at the neuron, the mean voltage and its relative jitter are
\begin{empheq}[box=\keybox]{equation}
\bar V \;=\; \tau\sum_j w_j\,r_j ,
\qquad
\frac{\sigma_V}{\bar V} \;\approx\; \frac{1}{\sqrt{2\,r_{\text{in}}\,\tau}} ,
\label{eq:reader}
\end{empheq}
where the second equality is written for equal weights, and $r_{\text{in}}$ is the total rate of all inputs. When $\tau$ is large, the voltage hardly jitters and follows the weighted sum of rates: the recipient is an \emph{integrator}, it reads rates and forgets timing. When $\tau$ is small, the voltage has time to fall between spikes, and threshold is reached only when several spikes arrive within the window~\eqref{eq:window}: the recipient is a \emph{coincidence detector}, it reads timing (Fig.~\ref{fig:reader}). A thousand inputs at five spikes per second with $\tau=10\ms$ give a jitter of about ten percent, almost pure integration. If inhibition, as in Chapter~\ref{sec:gaba}, has shortened $\tau$ to $2\ms$, the jitter grows to over twenty percent, and the neuron begins to notice coincidences.

\begin{figure}[t]
\centering
\begin{tikzpicture}[>=Stealth,x=0.115cm,y=1cm]
% common input: the same spike train for both receivers
\foreach \s in {6,21,22,38,52,55.5,59,62.5,66,69.5,73,76.5,80,83.5,87,90.5,94,97.5}{\draw[thick,gray!40!black] (\s,5.25) -- (\s,5.65);}
\draw[gray!70!black] (0,5.25) -- (105,5.25);
\node[left,font=\small] at (-1,5.45) {input};
\node[above,font=\scriptsize,gray!40!black] at (13.5,5.65) {sparse};
\node[above,font=\scriptsize,gray!40!black] at (21.5,6.0) {pair};
\node[above,font=\scriptsize,gray!40!black] at (75,5.65) {dense};
% axes and thresholds
\foreach \y/\th in {2.7/2.5,0/1.5}{
  \draw[gray!70!black] (0,\y) -- (106,\y) node[right,black] {\small $t$};
  \draw[densely dashed,gray!60!black] (0,\y+0.6*\th) -- (104,\y+0.6*\th) node[right,black,font=\small] {$\theta$};
}
\node[anchor=south west,font=\small] at (0,4.3) {$\tau=10\ms$: integrator};
\node[anchor=south east,font=\small] at (104,1.0) {$\tau=2\ms$: coincidence detector};
% integrator: tau=10 ms, theta=2.5w
\begin{scope}[yshift=2.7cm]
\foreach \a/\b/\v in {0/6/0.000,6/21/1.000,21/22/1.223,22/38/2.107,38/52/1.425,52/55.5/1.351,55.5/59/1.952,59/62.5/2.376,62.5/66/0.000,66/69.5/1.000,69.5/73/1.705,73/76.5/2.201,76.5/80/0.000,80/83.5/1.000,83.5/87/1.705,87/90.5/2.201,90.5/94/0.000,94/97.5/1.000,97.5/104/1.705}{\draw[very thick,blue!60!black] plot[domain=\a:\b,samples=14] (\x,{0.6*\v*exp(-(\x-\a)/10)});}
\foreach \s/\p/\q in {6/0.000/1.000,21/0.223/1.223,22/1.107/2.107,38/0.425/1.425,52/0.351/1.351,55.5/0.952/1.952,59/1.376/2.376,62.5/1.674/2.500,62.5/2.500/0.000,66/0.000/1.000,69.5/0.705/1.705,73/1.201/2.201,76.5/1.551/2.500,76.5/2.500/0.000,80/0.000/1.000,83.5/0.705/1.705,87/1.201/2.201,90.5/1.551/2.500,90.5/2.500/0.000,94/0.000/1.000,97.5/0.705/1.705}{\draw[very thick,blue!60!black] (\s,0.6*\p) -- (\s,0.6*\q);}
\foreach \s in {62.5,76.5,90.5}{\draw[->,thick,red!70!black] (\s,1.58) -- (\s,1.95);}
\end{scope}
% coincidence: tau=2 ms, theta=1.5w
\begin{scope}[yshift=0cm]
\foreach \a/\b/\v in {0/6/0.000,6/21/1.000,21/22/1.001,22/38/0.000,38/52/1.000,52/55.5/1.001,55.5/59/1.174,59/62.5/1.204,62.5/66/1.209,66/69.5/1.210,69.5/73/1.210,73/76.5/1.210,76.5/80/1.210,80/83.5/1.210,83.5/87/1.210,87/90.5/1.210,90.5/94/1.210,94/97.5/1.210,97.5/104/1.210}{\draw[very thick,blue!60!black] plot[domain=\a:\b,samples=14] (\x,{0.6*\v*exp(-(\x-\a)/2)});}
\foreach \s/\p/\q in {6/0.000/1.000,21/0.001/1.001,22/0.607/1.500,22/1.500/0.000,38/0.000/1.000,52/0.001/1.001,55.5/0.174/1.174,59/0.204/1.204,62.5/0.209/1.209,66/0.210/1.210,69.5/0.210/1.210,73/0.210/1.210,76.5/0.210/1.210,80/0.210/1.210,83.5/0.210/1.210,87/0.210/1.210,90.5/0.210/1.210,94/0.210/1.210,97.5/0.210/1.210}{\draw[very thick,blue!60!black] (\s,0.6*\p) -- (\s,0.6*\q);}
\foreach \s in {22}{\draw[->,thick,red!70!black] (\s,0.98) -- (\s,1.35);}
\end{scope}
\foreach \x in {0,50,100}{\draw[gray!60!black] (\x,-0.06) -- (\x,0.06) node[below,font=\scriptsize] {\x\,ms};}
\end{tikzpicture}
\caption{The same input, two different recipients. Each spike raises the voltage by $w$, after which it leaks away, as in neuron model~\eqref{eq:glutneuron}; red arrows are the recipient's spikes. The slow recipient ($\tau=10\ms$, $\theta=2.5w$) fires where spikes are \emph{many} and does not notice the pair. The fast one ($\tau=2\ms$, $\theta=1.5w$) fires only on the pair within window~\eqref{eq:window} and lets a dense but non-coincident stream pass.}
\label{fig:reader}
\end{figure}

Measurements show that in active cortex the membrane conductance increases severalfold compared with rest~\cite{Destexhe2003}. So $\tau$ there is shorter, and cortical neurons in their working regime are closer to coincidence detectors than to integrators.

\begin{keyblock}
\begin{proposition}[The recipient sets the code]
\label{prop:reader}
One and the same spike train carries rate for a slow recipient, timing for a fast one, and, for a volume into which transmitter is released, a concentration level smoothed over seconds (Chapter~\ref{sec:serocomputer}). Which code is used is determined not by the sender but by the recipient's time constant, and inhibition adjusts that constant.
\end{proposition}
\end{keyblock}
\section{Dots and dashes: the synapse as a filter}

So far our synapse has been a wire with a weight. In fact its response depends on the preceding spikes, and this gives the language of neurons a second symbol.

\emph{Depression: the synapse transmits changes.} Each release uses up a fraction $U$ of the stock of transmitter ready for release~\cite{Tsodyks1997}, and the stock recovers with a time constant $\tau_{\text{rec}}$ on the order of hundreds of milliseconds. At rate $r$ the amplitude of the response to a spike settles at approximately $A(r)=A_0/(1+U r\tau_{\text{rec}})$, where $A_0$ is the response of a rested synapse. The current the recipient receives per unit time is
\begin{empheq}[box=\keybox]{equation}
r\,A(r) \;=\; \frac{A_0\,r}{1+U r\,\tau_{\text{rec}}}
\;\xrightarrow[r\to\infty]{}\; \frac{A_0}{U\tau_{\text{rec}}} .
\label{eq:depression}
\end{empheq}
With $U=0.5$ and $\tau_{\text{rec}}=0.8\,\text{s}$, saturation begins already at $1/(U\tau_{\text{rec}})=2.5$ spikes per second. If the rate rises from five to twenty, the steady-state current grows by only a third. But the first spikes at the new rate arrive while the stock is still at its old level, and the current momentarily jumps fourfold, then decays within $\tau_{\text{rec}}/(1+Ur\tau_{\text{rec}})\approx90\ms$ (Fig.~\ref{fig:depression}).

Such a synapse transmits not the rate but its \emph{change}, essentially the relative change $\Delta r/r$~\cite{Abbott1997depression}. To an electronics engineer this is a high-pass filter built right into the wire.

\begin{figure}[t]
\centering
\begin{tikzpicture}[>=Stealth,x=12.5cm,y=1cm]
% input rate
\draw[->,gray!70!black] (0,2.6) -- (0.9,2.6) node[right,black] {\small $t$};
\draw[->,gray!70!black] (0,2.6) -- (0,4.3) node[above,black] {\small $r$};
\draw[very thick,blue!60!black] (0,2.9) -- (0.2,2.9) -- (0.2,3.8) -- (0.8,3.8);
\node[left,font=\scriptsize] at (0,2.9) {$5$};
\node[left,font=\scriptsize] at (0,3.8) {$20$};
\node[right,font=\small,blue!60!black] at (0.4,4.1) {rate at the synapse input};
% output current
\draw[->,gray!70!black] (0,0) -- (0.9,0) node[right,black] {\small $t$};
\draw[->,gray!70!black] (0,0) -- (0,2.45) node[left,black] {\small $rA$};
\draw[densely dashed,gray!60!black] (0,0.667) -- (0.8,0.667);
\draw[very thick,red!70!black] (0,0.5) -- (0.2,0.5) -- (0.2,2.0)
   plot[domain=0.2:0.8,samples=60] (\x,{0.667+(2.0-0.667)*exp(-(\x-0.2)/0.089)});
\node[left,font=\scriptsize] at (0,0.5) {$1.7$};
\node[left,font=\scriptsize] at (0,2.0) {$6.7$};
\node[right,font=\scriptsize] at (0.8,0.667) {$2.2$};
\node[right,font=\small,red!70!black] at (0.28,1.6) {current to the recipient: a jump, then decay within $\approx90\ms$};
\draw[gray!60!black] (0.2,-0.08) -- (0.2,0.08); \node[below,font=\scriptsize] at (0.2,0) {$0.2\,\text{s}$};
\draw[gray!60!black] (0.8,-0.08) -- (0.8,0.08); \node[below,font=\scriptsize] at (0.8,0) {$0.8\,\text{s}$};
\end{tikzpicture}
\caption{A depressing synapse according to formula~\eqref{eq:depression} with $U=0.5$, $\tau_{\text{rec}}=0.8\,\text{s}$; current in units of $A_0$ per second; the dashed line is the steady-state current: it rose only from $1.7$ to $2.2$, but at the moment of the step it jumped to $6.7$. The synapse tells the recipient that the rate has changed, not what it is.}
\label{fig:depression}
\end{figure}

\emph{Facilitation: a burst as a dash.} At other synapses the release probability is low but rises for tens of milliseconds after each spike. A single spike through such a synapse is most often lost. But a \emph{burst}, several spikes in a row a few milliseconds apart, gets through almost certainly. Even without facilitation, the probability that all $k$ spikes of a burst are lost is
\begin{empheq}[box=\keybox]{equation}
P_{\text{loss}} \;=\; (1-p)^k ,
\label{eq:burst}
\end{empheq}
which for $p=0.2$ is $0.8$ for a single spike and $0.41$ for a burst of four; facilitation lowers it further still. So the neuron acquires a second symbol: a single spike is a dot, a burst is a dash. A depressing synapse best transmits the first spike after a pause; a facilitating one passes bursts and suppresses isolated dots. That bursts are transmitted more reliably has been measured; that the brain actually uses them as a separate symbol is a plausible hypothesis~\cite{Lisman1997}.

\section{How \nt{GABA} neurons speak}

The most numerous of them in the cortex are the fast ones~\cite{Tremblay2016}. Their spikes are shorter than those of \nt{GLUT} neurons, and they can fire them steadily for long periods at hundreds per second, with almost no fatigue. Their synapses are more reliable: the connection to each recipient consists of many release sites, and a spike is almost never lost. Their outputs land near the site where the recipient's spike is born, so they control not how the recipient sums its inputs but whether and when it fires.

They themselves receive excitation from many neighboring \nt{GLUT} neurons and report the total activity around them, which is exactly what the division of Chapter~\ref{sec:gaba} needs. Finally, fast \nt{GABA} neurons are connected to one another and readily fire together; they are the ones that set the local rhythm discussed above~\cite{Cardin2009}.

In the language of Morse code, \nt{GABA} neurons are responsible not for the letters but for the \emph{pauses}. They cut the continuous stream of \nt{GLUT} spikes into windows within which the recipient can fire, narrow the coincidence windows, and turn down the whole stream when it gets too loud.

\begin{keyblock}
\begin{proposition}[Division of roles]
\label{prop:punctuation}
\nt{GLUT} neurons carry the content of the message: \emph{what} and \emph{how much}. Fast \nt{GABA} neurons carry its markup: \emph{when} one may speak and \emph{how loudly}; one more class, by inhibiting the inhibitors, decides \emph{for whom} the gates are open. Individual parts of this division have been measured; as a general rule it is a working hypothesis.
\end{proposition}
\end{keyblock}

There are also other, slower \nt{GABA} neurons whose outputs land on individual inputs of the recipient. They work as the veto~\eqref{eq:veto} of Chapter~\ref{sec:gaba}: they shut off one stream of \nt{GLUT} spikes without touching the others.

The division into fast and slow \nt{GABA} neurons is an old picture, and it is incomplete. Today at least three large classes are distinguished in the cortex~\cite{Tremblay2016}, and the third is arranged unexpectedly: it inhibits not \nt{GLUT} neurons but other \nt{GABA} neurons, above all those that close off inputs~\cite{Pfeffer2013}. Inhibition of inhibition is a double negation. Such a neuron \emph{opens} the gates without sending the recipient a single excitatory spike. It is switched on by signals from other cortical areas and by the broadcast channels, so it works as an enable input for a whole part of the network: while it is active, the vetoes are lifted and the stream gets through~\cite{Pi2013}. In Appendix~\ref{sec:othermediators} this trick is called disinhibition.
\section{An economical language}

The last thing firmly known about the language of neurons is that it is expensive. A spike, together with the synaptic work it causes, costs about $10^9$ ATP molecules (an estimate for rodent cortex~\cite{Attwell2001}), and one molecule yields about $10^{-19}$ joules. So a spike costs on the order of $E_{\text{spk}}\sim10^{-10}$ joules. The whole brain consumes about $20$ watts, and if half of that goes to signaling, $P\sim10$ watts, the mean rate of one neuron is
\begin{empheq}[box=\keybox]{equation}
\bar r \;\approx\; \frac{P}{N\,E_{\text{spk}}}
\;\sim\; \frac{10\ \text{W}}{8.6\cdot10^{10}\cdot10^{-10}\ \text{J}}
\;\approx\; 1\ \text{spike per second}.
\label{eq:energy}
\end{empheq}
More detailed estimates for the cortex give even less~\cite{Lennie2003}. The brain's code must be \emph{sparse}: only a small fraction of neurons can be firing fast.

From~\eqref{eq:spikeentropy} one sees that this is not so bad. At a precision of $1\ms$, a spike of a neuron firing at $100$ per second carries no more than five bits, while a spike of a neuron firing once a second carries up to eleven. If you pay per spike, it pays to speak rarely. The cortex does in fact trade precision for energy: when food is scarce, neurons keep their firing rate but spend less on synaptic currents, and their responses become less precise~\cite{Padamsey2022}.

In Morse code, frequent letters are short: the most frequent letter of English text, E, is a single dot. The neurons' code, as far as is known, follows the same rule in another form: what is expected costs few spikes~\cite{Barlow1961}. A neuron under a constant stimulus gradually lowers its rate, the sensors of Chapter~\ref{sec:glutcomputer} respond to changes rather than levels, and the \nt{DOPA} neurons of Chapter~\ref{sec:neurontypes} report only the reward-prediction error.

\begin{keyblock}
\begin{proposition}[Economy]
\label{prop:economy}
Energy limits a neuron's mean rate to about one spike per second. That is why the language of \nt{GLUT} neurons is sparse and spends its spikes on news: on what has changed or was not predicted. The energy limit has been measured; that the brain's code is tuned for the greatest number of bits per joule is a well-founded hypothesis.
\end{proposition}
\end{keyblock}

\section{Summary}

\begin{table}[t]
\centering
\footnotesize
\setlength{\tabcolsep}{4pt}
\renewcommand{\arraystretch}{1.15}
\begin{tabular}{>{\raggedright}p{2.6cm}>{\raggedright}p{2.3cm}>{\raggedright}p{3.0cm}>{\raggedright}p{2.0cm}>{\raggedright\arraybackslash}p{3.6cm}}
\toprule
Encoding & What it carries & Who reads it & Time per message & What is known \\
\midrule
index of the neuron that fired (place code) & which value & any recipient; coincidence detectors & one delay & established for sensory neurons and in sound localization \\
rate of one neuron & magnitude & integrator with large $\tau$; the volume (ch.~\ref{sec:serocomputer}) & hundreds of ms to seconds & established; too slow for a $10$--$20\ms$ processing stage \\
group rate & magnitude & neuron with thousands of inputs & $10$--$50\ms$ & established; gain limited by correlations~\eqref{eq:correlated} \\
first-spike time, order & magnitude, order & coincidence detectors with delays & one delay & established at the input of the nervous system; debated in the cortex \\
phase in a local rhythm & magnitude, position & neurons sharing a rhythm & cycle of $12$--$50\ms$ and slower & observed in some areas; general role is a hypothesis \\
burst (``dash'') & ``important'': reliable delivery & facilitating synapse & $\sim10\ms$ & reliability measured; a separate symbol is a hypothesis \\
rate change & news & depressing synapse & $\sim0.1\,\text{s}$ & established \\
spikes of fast \nt{GABA} neurons & when and how loudly & all neighboring \nt{GLUT} neurons & $1$--$30\ms$ & rhythm and division measured; ``punctuation'' as a rule is a hypothesis \\
\bottomrule
\end{tabular}
\caption{The ways \nt{GLUT} and \nt{GABA} neurons encode messages in spikes. The right-hand column separates what has been measured from what is assumed.}
\label{tab:code}
\end{table}

Table~\ref{tab:code} collects what we know; it also shows what we do not know. We do not know how much the cortex uses precise spike timing rather than only group rates. We do not know whether neurons have ``words,'' stable combinations of spikes from many neurons that are read as a whole. And we do not know whether the language is the same in different parts of the brain. Morse code can be learned in an evening because its table is known. The table of the neurons' language has yet to be compiled, and the ones who compile it will most likely be communications engineers.

\section{Exercises}

\begin{exercise}[\lvlA]
\label{ex:04:1}
\emph{Capacity in numbers.} $\Delta t=1\ms$. (a)~How many bits per joule does a neuron firing $1$ spike per second give at the spike cost of~\eqref{eq:energy}? (b)~At $r=10$, by what factor less than bound~\eqref{eq:spikeentropy} does a recipient extract if it only counts spikes per second?
\end{exercise}

\begin{exercise}[\lvlB]
\label{ex:04:2}
\emph{The optimal rate.} A neuron spends power $P_0+rE_{\text{spk}}$, where $P_0$ is the other half of the $20$~W divided among all neurons, and $E_{\text{spk}}=10^{-10}$~J. At what rate $r$ is the number of bits per joule $R_{\max}(r)/(P_0+rE_{\text{spk}})$ from~\eqref{eq:spikeentropy} with $\Delta t=1\ms$ greatest?
\end{exercise}

\begin{exercise}[\lvlB]
\label{ex:04:3}
\emph{Which correlation hurts.} $N=1000$ neurons, variance $\sigma^2=1$, pairwise correlation $c=0.1$. Compute the Fisher information $\mathcal I=f'^{\mathsf T}\Sigma^{-1}f'$ ($\Sigma$ is the covariance matrix) for equal slopes $f'=\mathbf 1$ and for slopes $\pm1$ with $\mathbf 1^{\mathsf T}f'=0$. By what factor do they differ?
\end{exercise}

\begin{exercise}[\lvlB]
\label{ex:04:4}
\emph{Simulation: a coincidence detector.} Neuron~\eqref{eq:glutneuron} receives $1000$ Poisson inputs at $5$ spikes per second each, $w=1$; the threshold $\theta$ is set so that the free voltage crosses it once a second. By simulation, find how many simultaneous inputs $m$ fire the neuron with probability $1/2$ for $\tau=10$ and $2\ms$.
\end{exercise}

\begin{exercise}[\lvlB]
\label{ex:04:5}
\emph{Design: a detector of relative changes.} Choose synapse~\eqref{eq:depression} so that the steady-state current at $40$ spikes per second is no more than $10\,\%$ above that at $10$, and after a step to $40$ the current settles to its new level with a time constant of no more than $50\ms$. What is the smallest $\tau_{\text{rec}}$, and at what $U$?
\end{exercise}

\begin{exercise}[\lvlA]
\label{ex:04:6}
\emph{First-spike time and bursts.} (a)~Using~\eqref{eq:latency}, find the input for which the first spike arrives exactly one time constant later. What does it depend on? (b)~With release probability $p=0.2$, how many spikes must a burst have for the probability of losing them all to be below $5\,\%$?
\end{exercise}

\begin{exercise}[\lvlC]
\label{ex:04:7}
\emph{Research: how many bits are in the arrangement.} The neuron is a set of independent bins~\eqref{eq:spikeentropy}, $r=10$, $\Delta t=1\ms$. (a)~Split the entropy of a ``word'' of $20$ bins into the entropy of the spike count and the remainder. (b)~Simulate estimating the entropy from word frequencies with $N=30$ and $10^4$ recordings. How much is it underestimated?
\end{exercise}

\chapter{What \nt{SERO} Neurons Compute}
\label{sec:serocomputer}

\section{The first broadcast channel}

So far every neuron in this book has talked over the address bus: \nt{GLUT} and \nt{GABA} sent spikes to specific recipients, and we worked out what such a network computes and what language it speaks. Now we turn to the second bus of Section~\ref{sec:buses}, the broadcast bus. Its first channel is \nt{SERO}. As before, we allow ourselves only what exists in a living organism.

This chapter introduces three devices that every broadcast channel will use later: a neuron that, given a steady background drive, runs on its own until something stops it; a wire that is really a volume of tissue; and a slow concentration in place of individual spikes. \nt{DOPA}, \nt{NORA} and \nt{ACET} in Chapters~\ref{sec:dofa}--\ref{sec:acet} are built from the same parts.

From an engineer's point of view, the \nt{SERO} neuron differs in four important ways.

\emph{First: the recipient chooses the sign.} Serotonin has receptors of both signs, and rule~\eqref{eq:dalesign} does not apply here: the sign is set not by the sender but by the recipient's receptor (Proposition~\ref{prop:receptor})~\cite{BarnesSharp1999}.

\emph{Second: a \nt{SERO} neuron runs on its own, given a steady background.} In the waking state it fires steadily, a few times per second, and it needs no push for each spike: it is a pacemaker. But it does need a constant background drive. In a brain slice, where this drive is absent, most of these cells are silent; the steady rhythm returns when noradrenaline is applied through $\alpha_1$ receptors, and blocking those receptors abolishes it~\cite{VandermaelenAghajanian1983}. In the organism this background comes from \nt{NORA}. In circuit terms it is an oscillator that needs a power supply: while the supply is on, it runs by itself until inhibition stops it.

\emph{Third: it is slow.} Almost all serotonin receptors are metabotropic: they do not open a channel themselves but trigger a chain of reactions inside the cell. The response is slow: the inhibitory current through the main receptor rises and decays within about a second~\cite{CourtneyFord2016}, tens of times longer than the response of a fast \nt{GLUT} synapse.

\emph{Fourth: it releases its transmitter not into a wire but into a volume.} In the regions that \nt{SERO} neurons project to, serotonin is often released not into a narrow cleft but into the extracellular space, where it spreads out~\cite{Bunin1998}. The recipient senses a concentration and cannot know which sender a molecule came from.

At these time scales individual spikes no longer matter, so we describe the network by rates $r_j$ and concentrations. The serotonin concentration $c$ in the volume into which neurons $j$ release rises with release and falls as the transmitter is pumped back:
\begin{empheq}[box=\keybox]{equation}
\tau_c\,\frac{\dd c}{\dd t} \;=\; -c \;+\; \kappa\sum_j r_j ,
\qquad
r_i \;=\; f_i\bigl(b_i + s_i\,g\,c\bigr), \quad s_i=\pm1 .
\label{eq:seroneuron}
\end{empheq}
This is one more way to read a spike train, complementing Proposition~\ref{prop:reader}: the volume sees neither individual spikes nor their order, only the rate averaged over the time $\tau_c$. The first equation is the same RC circuit as the membrane: $\tau_c$ is the lifetime of the transmitter in the volume (it has not been measured directly, and we take it to be of order a second), and $\kappa$ is how much transmitter one spike delivers. The second equation describes the recipient: $b_i$ is its own drive, positive for a pacemaker (it includes the steady background input from \nt{NORA}); $g$ is the receptor sensitivity; the sign $s_i$ is chosen by the recipient; $f_i$ is a nondecreasing transfer function.

\section{NOT with one neuron}

Take a pacemaker with an inhibitory receptor, $s=-1$. While there is no serotonin around, it runs on its own drive $b$. When serotonin appears, the drive drops by $gc$, and if $gc>b$ the neuron falls silent. This is NOT: there is an output when there is no input. It took one neuron (Fig.~\ref{fig:serogates}). Proposition~\ref{prop:monotone} does not apply here: it required positive weights.

If the same pacemaker is exposed to serotonin from two different volumes, it is silent when serotonin is present in at least one of them. This is NOR, and any logic function can be built from NOR gates. The two-wire code of the \nt{GLUT} network is not needed here: the absence of a signal is computed by neurons that run until something stops them.

\begin{figure}[t]
\centering
\begin{tikzpicture}[>=Stealth,
  nrn/.style={circle,draw,very thick,fill=blue!6,minimum size=0.95cm,inner sep=0pt},
  pace/.style={nrn,double,double distance=1.2pt},
  inh/.style={-{Bar[width=7pt]},thick,red!70!black}]
\node[pace] (N1) at (0,0) {};
\draw[inh] (-1.6,0) node[left,black] {$c$} -- (N1);
\draw[->,very thick,red!70!black] (N1) -- (1.3,0) node[right,black] {$\bar c$};
\node[below=0.55cm] at (0,0) {\small NOT};
\node[pace] (N2) at (5,0) {};
\draw[inh] (3.4,0.45) node[left,black] {$c_1$} -- (N2);
\draw[inh] (3.4,-0.45) node[left,black] {$c_2$} -- (N2);
\draw[->,very thick,red!70!black] (N2) -- (6.3,0) node[right,black] {$\overline{c_1\vee c_2}$};
\node[below=0.55cm] at (5,0) {\small NOR};
\end{tikzpicture}
\caption{Logic from \nt{SERO} neurons. A double circle is a pacemaker: given a steady background drive, it runs on its own until it is inhibited. A line ending in a bar is an inhibitory receptor, $s=-1$; $c$, $c_1$, $c_2$ are serotonin concentrations in the volumes around the recipient. Left: NOT; right: NOR.}
\label{fig:serogates}
\end{figure}

\begin{keyblock}
\begin{proposition}[Completeness of \nt{SERO} logic]
\label{prop:serocomplete}
If the network contains pacemakers with inhibitory receptors and releases into different volumes do not mix, network~\eqref{eq:seroneuron} can compute any logic function, with each bit carried by a single wire.
\end{proposition}
\end{keyblock}

Logically, a \nt{SERO} network is more powerful than a \nt{GLUT} network, but the condition ``releases into different volumes do not mix'' does not hold in the brain (Section~\ref{sec:volume}).

\section{Negative feedback}

Serotonin neurons carry inhibitory receptors on themselves~\cite{Sprouse1987}. The serotonin they release comes back to them and slows them down. An engineer knows this circuit: an amplifier with negative feedback (Fig.~\ref{fig:serofeedback}).

What is measured here, and what is a model? In the raphe nucleus itself, where the \nt{SERO} neurons sit, serotonin inhibits neighbors not through a common volume but point to point, through synapses: the transporter clears it quickly and does not let it accumulate outside the cleft. A single release causes only a brief pause in firing, and the mean rate does not change~\cite{CourtneyFord2016}. So the loop below, closed through a common volume, is a model: we compute what such feedback would do if it operated at the level of mean rates.

\begin{figure}[t]
\centering
\begin{tikzpicture}[>=Stealth,
  nrn/.style={circle,draw,very thick,fill=blue!6,minimum size=0.95cm,inner sep=0pt},
  pace/.style={nrn,double,double distance=1.2pt},
  blk/.style={draw,very thick,rounded corners=2pt,fill=orange!10,minimum height=0.9cm,minimum width=2.2cm,align=center,font=\small},
  inh/.style={-{Bar[width=7pt]},thick,red!70!black}]
\node[pace] (N) at (0,0) {};
\node[blk] (V) at (3.6,0) {volume\\$\tau_c$};
\draw[->,very thick] (N) -- node[above] {\small $r$} (V);
\draw[->,very thick] (V) -- (6.0,0) node[right] {\small $c$ to all recipients};
\draw[inh] (V.south) -- ++(0,-0.9) -| node[pos=0.25,below] {\small $-g\,c$} (N.south);
\draw[->,thick,blue!60!black] (-1.7,0) node[left,black] {\small $b$} -- (N);
\end{tikzpicture}
\caption{A \nt{SERO} neuron with feedback through its own transmitter: the concentration $c$ goes to all recipients and inhibits the neuron itself. In the model this is a level regulator; in the brain, such a role for the loop is a hypothesis.}
\label{fig:serofeedback}
\end{figure}

Let us compute what this loop does. Let the transfer function be linear, $f(u)=au$ for $u>0$. At equilibrium $c=\kappa r$ and $r=a(b-gc)$. Substituting one into the other gives
\begin{empheq}[box=\keybox]{equation}
r \;=\; \frac{a\,b}{1+G}, \qquad G \;=\; a\,g\,\kappa .
\label{eq:feedback}
\end{empheq}
The quantity $G$ is the loop gain. This is the same formula as for any amplifier with negative feedback, and it gives the same two benefits.

\emph{Robustness to variation.} Suppose the neuron's excitability $a$ changes by a fraction $\varepsilon$. Without feedback, the rate would change by the same fraction. With feedback and $G\gg1$, the rate $r\approx b/(g\kappa)$ hardly depends on $a$: its change is $1+G$ times smaller. At $G=9$ the variation is suppressed tenfold.

\emph{Speed.} Substitute $r=a(b-gc)$ into the first equation~\eqref{eq:seroneuron}: we get $\tau_c\,\dd c/\dd t=-(1+G)\,c+\kappa ab$. The concentration approaches its level with time constant $\tau_c/(1+G)$, which is $1+G$ times faster than without feedback.

Here is the first computation the \nt{SERO} system could perform: holding the overall serotonin level in the brain at a set point, the way a thermostat holds a temperature. This is a hypothesis: experimentally, autoinhibition is so far seen only as brief pauses that do not change the mean rate.

\section{From spikes to a level}

The second computation is hidden in the first equation~\eqref{eq:seroneuron}. Neurons emit individual spikes, but the recipient senses a concentration, which smooths them over the time $\tau_c$. If the source rate changes, the concentration follows it with a lag:
\begin{equation}
c(t) \;=\; \frac{\kappa}{\tau_c}\int_{-\infty}^{t} e^{-(t-t')/\tau_c}\,\sum_j r_j(t')\,\dd t' .
\label{eq:lowpass}
\end{equation}
This is a moving average of source activity over the last few $\tau_c$: a low-pass filter (Fig.~\ref{fig:lowpass}). In the same way, the simplest digital-to-analog converter passes pulses through an RC circuit and turns their rate into a level. The \nt{SERO} system does the same with hundreds of thousands of neurons at once.

This quantity is also a memory. After the sources fall silent, the concentration persists for a time of order $\tau_c$. It is not a bit but a smoothly fading trace.

\begin{figure}[t]
\centering
\begin{tikzpicture}[>=Stealth,x=1.45cm,y=1cm]
\draw[gray!70!black] (0,3.2) -- (8.1,3.2);
\node[left,font=\small] at (-0.05,3.4) {spikes};
\node[above,font=\scriptsize,gray!40!black] at (1,3.65) {$2$ per second};
\node[above,font=\scriptsize,gray!40!black] at (3.5,3.65) {$8$ per second};
\node[above,font=\scriptsize,gray!40!black] at (6.5,3.65) {$2$ per second};
\draw[->,gray!70!black] (0,0) -- (8.3,0) node[right,black] {\small $t$};
\draw[->,gray!70!black] (0,0) -- (0,2.75) node[left,black] {\small $c$};
\foreach \s in {0.136,0.529,0.760,1.223,1.714,1.748,1.755,2.663,2.701,2.734,3.414,3.493,3.719,3.800,3.928,3.948,4.074,4.327,4.420,4.589,4.728,4.736,4.914,5.025,5.205,5.220,6.224,6.544,7.178}{\draw[thick,blue!60!black] (\s,3.2) -- (\s,3.6);}
\foreach \a/\b/\v in {0.000/0.136/0.000,0.136/0.529/1.000,0.529/0.760/1.675,0.760/1.223/2.330,1.223/1.714/2.466,1.714/1.748/2.509,1.748/1.755/3.425,1.755/2.663/4.402,2.663/2.701/2.775,2.701/2.734/3.672,2.734/3.414/4.553,3.414/3.493/3.306,3.493/3.719/4.055,3.719/3.800/4.235,3.800/3.928/4.905,3.928/3.948/5.316,3.948/4.074/6.211,4.074/4.327/6.476,4.327/4.420/6.028,4.420/4.589/6.493,4.589/4.728/6.483,4.728/4.736/6.642,4.736/4.914/7.589,4.914/5.025/7.351,5.025/5.205/7.579,5.205/5.220/7.331,5.220/6.224/8.221,6.224/6.544/4.012,6.544/7.178/3.914,7.178/8.000/3.076}{\draw[very thick,orange!85!black] plot[domain=\a:\b,samples=10] (\x,{0.25*\v*exp(-(\x-\a))});}
\foreach \s/\p/\q in {0.136/0.000/1.000,0.529/0.675/1.675,0.760/1.330/2.330,1.223/1.466/2.466,1.714/1.509/2.509,1.748/2.425/3.425,1.755/3.402/4.402,2.663/1.775/2.775,2.701/2.672/3.672,2.734/3.553/4.553,3.414/2.306/3.306,3.493/3.055/4.055,3.719/3.235/4.235,3.800/3.905/4.905,3.928/4.316/5.316,3.948/5.211/6.211,4.074/5.476/6.476,4.327/5.028/6.028,4.420/5.493/6.493,4.589/5.483/6.483,4.728/5.642/6.642,4.736/6.589/7.589,4.914/6.351/7.351,5.025/6.579/7.579,5.205/6.331/7.331,5.220/7.221/8.221,6.224/3.012/4.012,6.544/2.914/3.914,7.178/2.076/3.076}{\draw[very thick,orange!85!black] (\s,0.25*\p) -- (\s,0.25*\q);}
\draw[thick,densely dashed,black] plot[domain=0:2,samples=30] (\x,{0.25*2*(1-exp(-\x))})
  plot[domain=2:5,samples=40] (\x,{0.25*(8-6.271*exp(-(\x-2)))})
  plot[domain=5:8,samples=40] (\x,{0.25*(2+5.688*exp(-(\x-5)))});
\foreach \x in {0,2,5,8}{\draw[gray!60!black] (\x,-0.05) -- (\x,0.05) node[below=3pt,font=\scriptsize] {\x\,s};}
\end{tikzpicture}
\caption{From spikes to a level. Top: spikes of \nt{SERO} neurons, sparse for two seconds, dense for three seconds, then sparse again. Bottom: serotonin concentration~\eqref{eq:lowpass} with $\tau_c=1\,\text{s}$. Dashed line: the same if the spikes came perfectly evenly.}
\label{fig:lowpass}
\end{figure}

\section{The wire is a volume}
\label{sec:volume}

Let us return to the condition of Proposition~\ref{prop:serocomplete}. Serotonin released nearby by different senders adds up to a single concentration.

\emph{The wire here is not a fiber but a volume of tissue.} Two signals stay distinct only if they are released into regions farther apart than the transmitter has time to spread. In time $\tau$ a molecule with diffusion coefficient $D$ travels
\begin{empheq}[box=\keybox]{equation}
\ell \;\sim\; \sqrt{D\,\tau}\, .
\label{eq:diffusionlength}
\end{empheq}
The tortuosity of the extracellular gaps slows the diffusion of a small molecule by a factor of about $2.6$~\cite{Sykova2008}, while the diffusion coefficient of serotonin itself in tissue has not been measured; from the size of the molecule we estimate it at a few units of $10^{-6}$~cm$^2$/s. If a signal lives $\tau\sim1$~s (also an estimate), then $\ell\sim\sqrt{3\cdot10^{-6}}$~cm $\approx20$~\textmu m. The address in such a network is a \emph{place}: a recipient knows where a signal came from only by where it itself is located.

Real \nt{SERO} neurons are built so that almost none of these separate addresses remain. The output of each one is spread over huge parts of the brain: the branches of a single cell reach many regions far apart from one another~\cite{GagnonParent2014}. The number of release sites per cell is estimated at about $10^5$ (Table~\ref{tab:types}); nobody has counted them directly. The ``wires'' of different \nt{SERO} neurons overlap and mix. They do not merge into one wire entirely: \nt{SERO} neurons fall into several subsystems that project to different regions and respond differently to the same event~\cite{Ren2018}. But there are a handful of such subsystems, not hundreds of thousands.

\begin{keyblock}
\begin{proposition}[Number of independent wires]
\label{prop:volumewires}
In a network with volume transmission, the number of independent signals is limited not by the number of neurons but by the number of non-overlapping regions of size $\ell\sim\sqrt{D\tau}$ into which they release transmitter. If the output of each neuron is spread over a large volume, the whole network carries only a few shared variables.
\end{proposition}
\end{keyblock}

So the brain has nothing from which to build the logic circuits of Proposition~\ref{prop:serocomplete}. The regulator~\eqref{eq:feedback} and the filter~\eqref{eq:lowpass}, however, need no separate wires: one shared quantity is enough for them. The filter follows directly from the measured release into the volume in the target regions~\cite{Bunin1998}; the level regulator is a hypothesis.

\section{The planning horizon}
\label{sec:horizon}

What could a single shared slow quantity be good for? A good candidate is a parameter that must be the same for the whole network and must change no faster than over seconds or minutes. Such a parameter has already appeared in Chapter~\ref{sec:neurontypes}: the coefficient $\gamma$ in the error formula~\eqref{eq:rpe}, which says how much less a future reward is worth than an immediate one. In continuous time its place is taken by the \emph{horizon} $\tau_\gamma$: a reward $R$ that is a wait $D$ away is worth $R\,e^{-D/\tau_\gamma}$ now.

The horizon decides whether waiting is worthwhile. Suppose one can take a small reward $r_0$ right away or wait for a large reward $R$. Waiting pays as long as
\begin{empheq}[box=\keybox]{equation}
D \;<\; \tau_\gamma\,\ln\frac{R}{r_0} .
\label{eq:patience}
\end{empheq}
With $\tau_\gamma=2\,\text{s}$ and a delayed reward three times larger, it is worth waiting up to $2.2\,\text{s}$; if the horizon is twice as long, up to $4.4\,\text{s}$. Patience is proportional to the horizon.

Formula~\eqref{eq:patience} rests on exponential discounting. Discounting measured at delays of seconds is closer to the hyperbola $R/(1+D/\tau_\gamma)$: this is how, in monkeys, both the value of a delayed reward in choice and the response of \nt{DOPA} neurons to its cue fall off~\cite{KobayashiSchultz2008}. For the hyperbola, condition~\eqref{eq:patience} becomes $D<\tau_\gamma\,(R/r_0-1)$: the dependence on the reward ratio is no longer logarithmic but linear, and with the same numbers the waiting limit shifts by almost a factor of two. The hyperbola follows from the same exponential if the delay is random (Exercise~\ref{ex:05:7}), and patience remains proportional to the time scale.

The hypothesis is that \nt{SERO} sets the horizon~\cite{Doya2002}. It has everything such a role requires: one shared level for the whole network that changes over seconds. There are also direct experiments: if serotonin neurons are switched on artificially, an animal waits longer for a delayed reward before giving up~\cite{Miyazaki2014}, and keeps trying longer to obtain reward where it has become scarce~\cite{Lottem2018}. How exactly the serotonin concentration becomes $\tau_\gamma$ inside the network is unknown. In the next chapter the horizon enters the error that \nt{DOPA} broadcasts.

\section{Summary}

\begin{table}[t]
\centering
\small
\begin{tabular}{>{\raggedright}p{3.6cm}>{\raggedright}p{4.3cm}>{\raggedright\arraybackslash}p{4.3cm}}
\toprule
& \nt{GLUT} & \nt{SERO} \\
\midrule
response time & milliseconds & a fraction of a second to a second \\
sign of action & $+$ only & $+$ and $-$, chosen by the recipient \\
without input & silent & runs on its own given a steady background drive \\
addressing & point to point & volume; the address is a place \\
NOT & only through ``off'' sensors & one neuron \\
memory & self-sustained activity, dies out through fatigue & concentration, fades over the time $\tau_c$ \\
main computations & coincidences, delays, logic, sequences & smoothing; level regulation and horizon $\tau_\gamma$ (hypotheses) \\
\bottomrule
\end{tabular}
\caption{What networks of \nt{GLUT} neurons alone and of \nt{SERO} neurons alone compute.}
\label{tab:glutvssero}
\end{table}

As a logic element (Table~\ref{tab:glutvssero}), the \nt{SERO} neuron is richer than \nt{GLUT}, but as a communication system it is a broadcast: slow and almost without addresses. So it does not try to be a processor; instead it smooths and broadcasts a few shared quantities of the kind discussed in Chapter~\ref{sec:neurontypes}, hypothetically including the planning horizon. The pacemaker, the volume and the slow concentration will meet us three more times: in \nt{DOPA}, \nt{NORA} and \nt{ACET}.

\section{Exercises}

\begin{exercise}[\lvlA]
\label{ex:05:1}
\emph{The wire is a volume.} Take $D=3\cdot10^{-6}$~cm$^2$/s for serotonin (Section~\ref{sec:volume}). Using~\eqref{eq:diffusionlength}, find $\ell$ for a signal that lives $1\,\text{s}$, and the time for the transmitter to spread over $5$~cm. What then makes \nt{SERO} a broadcast: diffusion, or a branching wire with $\sim10^5$ release sites?
\end{exercise}

\begin{exercise}[\lvlA]
\label{ex:05:2}
\emph{Level regulator.} Show that in~\eqref{eq:seroneuron} with $f(u)=au$ the relative sensitivity $S=(a/r)\,\partial r/\partial a$ equals $1/(1+G)$ exactly, not only for $G\gg1$. At $G=9$ the excitability $a$ rises by $20\,\%$. By what percentage does $r$ change, without linearizing?
\end{exercise}

\begin{exercise}[\lvlB]
\label{ex:05:3}
\emph{A slow receptor in the loop.} The \nt{SERO} receptor responds with a delay: $r(t)=a\bigl(b-gc(t-T)\bigr)$, and the volume obeys~\eqref{eq:seroneuron}. Show that for $G>1$ the loop is stable only if $T<T^*=\tau_c\arccos(-1/G)/\sqrt{G^2-1}$. Find $T^*$ at $\tau_c=1\,\text{s}$ for $G=2$ and $G=9$. What suppression of variation is available with a delay of hundreds of milliseconds?
\end{exercise}

\begin{exercise}[\lvlB]
\label{ex:05:4}
\emph{Shot noise of the level.} By~\eqref{eq:lowpass}, each spike adds an exponential with $\tau_c=1\,\text{s}$ to $c$. Find the coefficient of variation of $c$ if all $3\cdot10^5$ \nt{SERO} neurons release Poisson spikes at $2$ per second into one volume. What $\tau_c$ would give $\mathrm{CV}=10\,\%$ for a single neuron firing $4$ spikes per second?
\end{exercise}

\begin{exercise}[\lvlB]
\label{ex:05:5}
\emph{Simulation: feedback versus noise.} Simulate $N=50$ Poisson pacemakers with rates $r_i=a_i[b-gc]_+$, $a_i$ uniform in $[2.8,\,5.2]$, and the volume~\eqref{eq:seroneuron} with $\kappa=1/N$, $\tau_c=1\,\text{s}$. By what factor does feedback ($b=1$, $g=2.25$) reduce the CV of the level $c$ compared with $b=0.1$, $g=0$? Does it equalize the neurons' rates?
\end{exercise}

\begin{exercise}[\lvlB]
\label{ex:05:6}
\emph{Design: XOR in \nt{SERO} logic.} A gate is a pacemaker that emits $r_0$ if $b-g\sum_kc_k>0$ and $0$ otherwise, into its own volume $\tau_c\,\dd c/\dd t=-c+\kappa r$; it computes NOR. Build XOR from five such gates and find its settling time at $\tau_c=1\,\text{s}$ and $G'=g\kappa r_0/b=2$.
\end{exercise}

\begin{exercise}[\lvlB]
\label{ex:05:7}
\emph{Patience with an unknown delay.} (a)~An animal agrees to wait for a reward three times larger if the wait is no longer than $6\,\text{s}$. What horizon $\tau_\gamma$ does this imply? (b)~Let the delay $D$ be exponentially distributed with mean $\bar D$. Show that waiting pays if $\bar D<\tau_\gamma(R/r_0-1)$, and compare with a fixed delay at $\tau_\gamma=2\,\text{s}$, $R=3r_0$.
\end{exercise}

\begin{exercise}[\lvlC]
\label{ex:05:8}
\emph{How a concentration sets the horizon.} A \nt{DOPA} neuron receives $V$ directly with weight $k_1$ and through a \nt{GABA} intermediary that smooths with $\tau_d=0.1\,\text{s}$, with weight $-k_2$; for slow $V$ the sum equals $(k_1-k_2)V+k_2\tau_d\dot V$. Choose $k_1$, $k_2$ to match~\eqref{eq:ctd} with $\tau_\gamma=2\,\text{s}$. \nt{SERO} multiplies $k_1$ by $m$: how much must $m$ change to double the horizon?
\end{exercise}

\chapter{Neural Networks That Can Learn: Adding \nt{DOPA}}
\chaptermark{Networks That Learn}
\label{sec:dofa}

\section{The rules of the game}

In all the circuits of the previous chapters, the weights were set by an engineer. In the organism there is no engineer. The weights must set themselves, in the course of operation, and in such a way that the network does something useful. This is what we call learning.

One more rule joins the old ones. A synapse changes its own weight, and it can know only what happens \emph{next to it}: the activity of the sender $x$, the activity of the recipient $y$, and whatever substances reach it. Nobody can send a synapse a personal correction.

This rules out the main method of artificial neural networks, backpropagation of error. There, the output error travels back through the network over the same weights, in a separate phase after the forward pass, and every weight receives its own correction. That requires backward wires that exactly mirror the forward ones, and a common clock. Neither has been found in the brain, at least not in the form used in artificial networks~\cite{Lillicrap2020,WhittingtonBogacz2019}.

We will write the change of a weight as a slow continuous process:
\begin{equation}
\frac{\dd w}{\dd t} \;=\; \eta\,\Phi(\text{what the synapse knows}),
\label{eq:plasticity}
\end{equation}
where $\eta$ is the learning rate, so small that the weight hardly changes during a single spike. This whole chapter is about what the function $\Phi$ can be.

\section{Where the weight is stored}
\label{sec:weightstore}

What is a synapse that ``changes its weight''? A connection has two sides: the terminal of the sender's wire (the \emph{presynaptic} side) and a patch of the recipient's membrane with receptors (the \emph{postsynaptic} side), with a narrow cleft between them. In \nt{GLUT} connections the recipient's membrane usually sits on a \emph{spine}, a small protrusion of the recipient's input branch. The weight of a connection is conveniently factored into three terms~\cite{delCastillo1954}:
\begin{empheq}[box=\keybox]{equation}
w \;=\; N_s\,p\,A_1 .
\label{eq:quantal}
\end{empheq}
Here $N_s$ is the number of release sites between the two neurons, $p$ is the probability that a spike releases a packet of transmitter at one site (the same $p$ as in Chapter~\ref{sec:code}), and $A_1$ is the recipient's response to one packet, that is, how many receptors catch it. The factor $p$ lives in the sender, $A_1$ in the recipient, and $N_s$ requires both sides: new release sites grow and old ones disappear, and this goes on throughout life.

\emph{The recipient decides.} At a typical \nt{GLUT} synapse, one of the glutamate receptors conducts current only when two conditions coincide: glutamate has arrived from the sender, and the recipient's membrane is already excited~\cite{Nowak1984}. This is Hebb's rule built into a single molecule: the coincidence detector of Chapter~\ref{sec:glutcomputer}, placed at the input of the synapse itself. When the receptor fires, it lets calcium into the recipient, and calcium triggers the weight change. The recipient is the only place where both the input and the output are visible at once, so the decision is made there.

\emph{Either side can carry out the decision.} The best-studied long-term strengthening happens in the recipient: new receptors are inserted into the membrane~\cite{Malinow2002}, and the spine grows~\cite{Matsuzaki2004}, so $A_1$ grows. But the recipient can also change $p$: it releases a substance that travels back across the cleft to the sender (the reverse channel of Appendix~\ref{sec:othermediators}), and the sender starts releasing less transmitter~\cite{WilsonNicoll2001}. And the short-term weight changes of Chapter~\ref{sec:code}, depression and facilitation, are changes of $p$ that the sender manages on its own according to its recent activity.

For an engineer, the weight register is split between two chips. The recipient executes the learning rule, and it can write the result either locally or at the sender, through the reverse channel. So the word ``local'' in the rules of this chapter means not one side of the connection but the whole connection.

\section{Learning without a teacher}

The simplest thing a synapse can do is strengthen when the sender and the recipient are active together:
\begin{equation}
\frac{\dd w}{\dd t} \;=\; \eta\,x\,y .
\label{eq:hebb}
\end{equation}
This is Hebb's rule~\cite{Hebb1949}. It is local and needs no clock. But it has the same problem as the \nt{GLUT} network in Chapter~\ref{sec:glutcomputer}: positive feedback. A strong weight makes $y$ larger, a large $y$ strengthens the weight further, and the weights grow without bound. The cure is negative feedback on the weight: let the weight also decay in proportion to $y^2$. For a neuron with inputs $\mathbf x$ and a linear output $y=\mathbf w\cdot\mathbf x$ we get
\begin{empheq}[box=\keybox]{equation}
\frac{\dd\mathbf w}{\dd t} \;=\; \eta\,y\,\bigl(\mathbf x - y\,\mathbf w\bigr) .
\label{eq:oja}
\end{empheq}
The rule is still local: the synapse knows its input $x_i$, the output $y$ and its weight $w_i$.

\begin{keyblock}
\begin{proposition}[What Hebb's rule finds]
\label{prop:oja}
Rule~\eqref{eq:oja} drives the weight vector to unit length along the direction in which the inputs vary most, that is, to the principal eigenvector of the input correlation matrix $C=\langle\mathbf x\mathbf x^{\mathsf T}\rangle$~\cite{Oja1982}. The neuron learns to output the single most expressive quantity into which its inputs can be compressed.
\end{proposition}
\end{keyblock}

Proof. Average~\eqref{eq:oja} over the inputs: $\dd\mathbf w/\dd t=\eta\bigl(C\mathbf w-(\mathbf w^{\mathsf T}C\mathbf w)\,\mathbf w\bigr)$. The fixed points are the unit-length eigenvectors of $C$. If $\mathbf w$ is close to an eigenvector $\mathbf e_k$ with eigenvalue $\lambda_k$, a small perturbation along another eigenvector $\mathbf e_j$ grows as $e^{\eta(\lambda_j-\lambda_k)t}$. Only the point for which all $\lambda_j<\lambda_k$ is stable, that is, the principal eigenvector.

There is also a version of Hebb's rule that looks at spike timing. If the sender's spike arrives $s$ milliseconds \emph{before} the recipient's spike, the weight grows by $A_+e^{-s/\tau_+}$; if $s$ milliseconds \emph{after}, it decreases by $A_-e^{-s/\tau_-}$, with $\tau_\pm$ of order twenty milliseconds. This rule has been measured at synapses of \nt{GLUT} neurons~\cite{BiPoo1998}. It strengthens connections that \emph{could have caused} the response and weakens the latecomers (Fig.~\ref{fig:stdp}). Run one sequence of activations through a network many times, and connections from each link to the next will grow by themselves: the chain of Chapter~\ref{sec:glutcomputer}, assembled without an engineer.

\begin{figure}[t]
\centering
\begin{tikzpicture}[>=Stealth,x=0.07cm,y=1.3cm]
\draw[->,gray!70!black] (-66,0) -- (68,0) node[right,black] {\small $s$};
\draw[->,gray!70!black] (0,-1.2) -- (0,1.35) node[above,black] {\small weight change};
\draw[very thick,blue!60!black] plot[domain=0.5:60,samples=50] (\x,{exp(-\x/20)});
\draw[very thick,red!70!black] plot[domain=-60:-0.5,samples=50] (\x,{-0.6*exp(\x/20)});
\draw[blue!60!black,thick] (0.5,0) -- (0.5,1);
\draw[red!70!black,thick] (-0.5,0) -- (-0.5,-0.6);
\foreach \x in {-40,-20,20,40}{\draw[gray!60!black] (\x,-0.04) -- (\x,0.04);}
\node[below,font=\scriptsize] at (20,-0.04) {$20\ms$};
\node[above,font=\scriptsize] at (-20,0.04) {$-20\ms$};
\node[align=left,font=\small,blue!60!black,anchor=west] at (22,0.95) {sender first:\\ connection strengthens};
\node[align=right,font=\small,red!70!black,anchor=east] at (-12,-0.95) {sender later:\\ connection weakens};
\end{tikzpicture}
\caption{Hebb's rule in time. Horizontal axis: $s=t_{\text{post}}-t_{\text{pre}}$, how far the recipient's spike lags behind the sender's spike; $\tau_\pm=20\ms$, $A_-=0.6A_+$. The window, a few tens of milliseconds wide, is of the same order as the coincidence window~\eqref{eq:window}.}
\label{fig:stdp}
\end{figure}

But all the rules in this section learn what is \emph{frequent}, not what is \emph{useful}. A synapse that knows only $x$ and $y$ does not know whether an action brought juice or pain. To learn from outcomes, the synapse needs a third signal.

\section{The third factor}

The third signal is brought by \nt{DOPA}: dopamine neurons broadcast a single number, the reward-prediction error $\delta$~\eqref{eq:rpe} (Chapter~\ref{sec:neurontypes}). Let the weight change be proportional to the product of three factors: the activity of the sender, the activity of the recipient, and $\delta$.

The difficulty is that the reward arrives not when the network chose the action but hundreds of milliseconds or seconds later. By the time $\delta$ arrives, the product $xy$ has long been zero, and the synapse needs a memory that it took part. The simplest memory is our familiar RC circuit:
\begin{empheq}[box=\keybox]{equation}
\tau_e\,\frac{\dd e}{\dd t} \;=\; -e \;+\; x\,y ,
\qquad
\frac{\dd w}{\dd t} \;=\; \eta\,\delta(t)\,e(t) .
\label{eq:threefactor}
\end{empheq}
The quantity $e$ is called the \emph{eligibility trace}~\cite{Fremaux2016}. Joint activity leaves a tag in the synapse that fades over the time $\tau_e$. If dopamine arrives within that time, the weight changes with the sign of $\delta$; if not, the tag vanishes without a trace (Fig.~\ref{fig:trace}). This has been measured at \nt{GLUT} synapses: dopamine that arrives within a second or two after joint activity turns it into strengthening of the connection, while dopamine that arrives later does not~\cite{Yagishita2014}. Plasticity windows lasting seconds have also been found where the third signal is not dopamine but strong excitation of the recipient itself~\cite{Bittner2017}.

\begin{figure}[t]
\centering
\begin{tikzpicture}[>=Stealth,x=7cm,y=1cm]
\fill[orange!12] (0.7,-0.2) rectangle (0.8,5.2);
% rows: x*y, delta, traces, weights
\foreach \y/\lab in {4.4/{$x\,y$},3.1/{$\delta$},1.4/{trace $e$},0/{weight $w$}}{
  \draw[gray!70!black] (0,\y) -- (1.42,\y);
  \node[left,font=\small] at (-0.02,\y+0.3) {\lab};}
\draw[very thick,blue!60!black] (0,4.4) -- (0,4.9) -- (0.2,4.9) -- (0.2,4.4);
\node[right,font=\scriptsize,blue!60!black] at (0.21,4.8) {sender and recipient active together};
\draw[very thick,orange!85!black] (0,3.1) -- (0.7,3.1) -- (0.7,3.7) -- (0.8,3.7) -- (0.8,3.1) -- (1.4,3.1);
\node[right,font=\scriptsize,orange!85!black] at (0.81,3.6) {dopamine arrives};
% traces, each in units of its own maximum
\draw[very thick,blue!60!black] plot[domain=0:0.2,samples=20] (\x,{1.4+1.1*(1-exp(-\x))/(1-exp(-0.2))})
  plot[domain=0.2:1.4,samples=60] (\x,{1.4+1.1*exp(-(\x-0.2))});
\draw[thick,densely dashed,gray!50!black] plot[domain=0:0.2,samples=30] (\x,{1.4+1.1*(1-exp(-\x/0.05))/(1-exp(-4))})
  plot[domain=0.2:1.4,samples=80] (\x,{1.4+1.1*exp(-(\x-0.2)/0.05)});
\node[right,font=\scriptsize,blue!60!black] at (1.0,2.15) {$\tau_e=1\,\text{s}$};
\node[right,font=\scriptsize,gray!40!black] at (0.3,1.65) {$\tau_e=50\ms$};
% weights
\draw[very thick,blue!60!black] (0,0.25) -- (0.7,0.25) -- (0.8,0.8) -- (1.4,0.8) node[right,font=\scriptsize] {$\tau_e=1\,\text{s}$: weight grew};
\draw[thick,densely dashed,gray!50!black] (0,0.2) -- (1.4,0.2) node[right,font=\scriptsize,gray!40!black] {$\tau_e=50\ms$: no change};
% time axis marks
\foreach \x/\l in {0/0,0.2/0.2,0.7/0.7,1.4/1.4}{\draw[gray!60!black] (\x,-0.05) -- (\x,0.05) node[below=3pt,font=\scriptsize] {\l\,s};}
\draw[<->,thick,gray!50!black] (0.2,5.35) -- node[above,font=\scriptsize,black] {reward delay $0.5\,\text{s}$} (0.7,5.35);
\end{tikzpicture}
\caption{The eligibility trace under rule~\eqref{eq:threefactor}. Joint activity lasts $0.2\,\text{s}$; dopamine arrives half a second after it ends. The traces are shown as fractions of their own maximum. At that moment the slow trace ($\tau_e=1\,\text{s}$) is still alive, while the fast one ($\tau_e=50\ms$) has long faded.}
\label{fig:trace}
\end{figure}

\begin{keyblock}
\begin{proposition}[Learning from a broadcast signal]
\label{prop:threefactor}
Rule~\eqref{eq:threefactor} changes only those synapses that took part in the network's operation during the last $\tau_e$, in the direction set by the shared signal $\delta$. A broadcast signal together with a local trace yields an addressed correction. A delay between action and outcome is acceptable as long as it does not exceed $\tau_e$.
\end{proposition}
\end{keyblock}

\section{Why it works: noise as search}
\label{sec:dofagradient}

Why does rule~\eqref{eq:threefactor} lead to improvement? The answer comes from the noise of Chapter~\ref{sec:code}. Let the neuron's output be $y=f(u)+\xi$, where $u=\sum_jw_jx_j$ and $\xi$ is random jitter with zero mean and variance $\sigma^2$. The reward $R$ depends on $y$. Take as the trace not simply $x_jy$ but its random part $x_j\xi$, and as the third factor the error $\delta=R-\bar R$, where $\bar R$ is the expected reward. For small noise $R\approx R\bigl(f(u)\bigr)+R'\,\xi$, so
\begin{empheq}[box=\keybox]{equation}
\bigl\langle \delta\,\xi\,x_j \bigr\rangle \;=\; \sigma^2\,R'\,x_j \;=\; \frac{\sigma^2}{f'(u)}\,\frac{\partial\langle R\rangle}{\partial w_j} .
\label{eq:perturbation}
\end{empheq}
The mean step follows the gradient of expected reward~\cite{Williams1992}. Nobody computed any derivatives: the network deviated at random, dopamine reported whether things got better, and the participating synapses moved in the favorable direction (Fig.~\ref{fig:perturbation}). The noise that interfered with sending messages in Chapter~\ref{sec:code} becomes search here.

\begin{figure}[t]
\centering
\begin{tikzpicture}[>=Stealth,x=2cm,y=3cm]
\draw[->,gray!70!black] (0,0) -- (4.2,0) node[below left,font=\small,black] {output $y$};
\draw[->,gray!70!black] (0,0) -- (0,1.2) node[left,font=\small,black] {reward $R$};
% reward hill
\draw[very thick,blue!60!black] plot[domain=0:4,samples=60] (\x,{max(0,1-((\x-3)/2.2)^2)});
\node[font=\scriptsize,blue!60!black,anchor=south] at (3,1.02) {best output};
% noise around f(u)
\draw[thick,gray!60!black] plot[domain=0.6:2.4,samples=50] (\x,{0.28*exp(-((\x-1.5)/0.3)^2/2)});
\draw[densely dashed,gray!60!black] (1.5,0) -- (1.5,0.535);
\node[below,font=\small] at (1.5,0) {$f(u)$};
\node[font=\scriptsize,gray!50!black] at (1.5,-0.2) {jitter $\xi$};
% expected reward
\draw[densely dotted,gray!60!black] (0.5,0.535) node[left,font=\scriptsize,black] {$\bar R$} -- (2.1,0.535);
% two random deviations
\fill[green!45!black] (2.1,0.833) circle (0.035);
\draw[very thick,green!45!black] (2.1,0.535) -- (2.1,0.833);
\node[font=\scriptsize,green!45!black,anchor=west,align=left] at (2.18,0.6) {$\xi>0$: better,\\ $\delta>0$};
\fill[red!70!black] (0.9,0.089) circle (0.035);
\draw[very thick,red!70!black] (0.9,0.535) -- (0.9,0.089);
\node[font=\scriptsize,red!70!black,anchor=east,align=right] at (0.83,0.27) {$\xi<0$: worse,\\ $\delta<0$};
% resulting step
\draw[->,very thick,orange!85!black] (1.55,0.4) -- (1.95,0.4) node[right,font=\scriptsize,black] {step};
\end{tikzpicture}
\caption{Noise as search~\eqref{eq:perturbation}. The neuron's output jitters around $f(u)$. A random deviation to the right (green) gave more than expected, $\delta>0$; to the left (red), less, $\delta<0$. In both cases the product $\delta\,\xi$ is positive, and the weight moves to the right, where reward increases. On average the step is proportional to the slope $R'$: at the top of the hill, deviations in either direction are equally bad, and the steps cancel.}
\label{fig:perturbation}
\end{figure}

\begin{keyblock}
\begin{proposition}[Gradient descent without backward wires]
\label{prop:gradient}
If the network has random deviations of activity, and the broadcast signal equals the reward-prediction error and averages to zero in each situation, then rule~\eqref{eq:threefactor} on average changes the weights along the gradient of expected reward. No backward wires and no separate learning phase are needed.
\end{proposition}
\end{keyblock}

Now it is clear why dopamine reports not the reward but the \emph{error} in predicting it. If the third factor were the reward itself, $R\ge0$, the average in~\eqref{eq:perturbation} would not change, but two problems would appear. First, a term $\bar R\,\xi x_j$ would be added to the useful part: zero on average, but strong noise. Second, and worse, a real synapse cannot separate the random part of the trace $x_jy$ from the nonrandom part. For given inputs, $x_jf(u)$ is the same number from trial to trial; if $\delta$ averages to zero in that situation, this part of the trace changes nothing, but with $\delta\ge0$ it strengthens everything the network did, time after time, right and wrong alike. So $\bar R$ must be the expectation \emph{in the given situation}, not the average over a lifetime. Computing it is the job of the critic, discussed below.

We tested this on a model. Two groups of \nt{GLUT} neurons choose one of two actions through mutual inhibition, as in Fig.~\ref{fig:wta}; the weights from the ``go'' signal to the groups start out equal, and noise decides who wins. Action $A$ brings juice with probability $0.8$, action $B$ with probability $0.2$, and the juice arrives half a second after the choice. With $\tau_e=1\,\text{s}$ and $\delta=R-\bar R$, the fraction of $A$ choices over five hundred trials rose from $0.65$--$0.8$ in the first hundred to $0.96$--$1.0$ in the last, and the weights stayed within $0.85$--$1.35$. With $\tau_e=50\ms$, shorter than the reward delay, the network learned nothing: the fraction of $A$ choices stayed around one half. With $\delta=R$, without subtracting the expectation, the network also came to choose $A$, but the winner's weight grew a thousandfold over the same five hundred trials and kept growing; and with the same $\delta=R$ and a tenfold higher learning rate, in one of three runs the network permanently locked in its first random choice, the worse one.

\section{The price of a single wire}

There is one signal $\delta$ for everyone, and the noise it evaluates is the sum of the jitter of all $N$ neurons that influence the outcome. Each synapse sees in $\delta$ its own useful part plus the noise of its $N-1$ neighbors. To extract its own part it must average over many trials, and the learning time grows in proportion to $N$~\cite{Werfel2005}. Backpropagation does not pay this price.

\begin{keyblock}
\begin{proposition}[The price of broadcasting]
\label{prop:broadcastcost}
The learning time from a single shared signal grows in proportion to the number of neurons whose noise affects the outcome. Such learning suits small circuits choosing among a few options, but not the tuning of billions of weights.
\end{proposition}
\end{keyblock}

The brain apparently divides the work this way: the outputs of \nt{DOPA} neurons go primarily to the regions that select actions (Chapter~\ref{sec:neurontypes}), while the huge \nt{GLUT} network of the cortex, which holds the overwhelming majority of the weights, learns mostly by local rules, without an external teacher. These used to be reduced to Hebb's rules~\eqref{eq:oja}. Now there is growing evidence that the cortex has its own goal, to \emph{predict} its next input, and then the error is computed on the spot, as the difference between what was predicted and what arrived~\cite{Keller2018}: the teacher is built into the network itself. Plasticity windows lasting seconds, where the third signal is strong excitation of the recipient itself~\cite{Bittner2017}, show that a local error signal at synapses really does exist. How general a rule of the cortex this is remains a hypothesis.

\section{Where the error comes from: the critic in continuous time}

The remaining question is how \nt{DOPA} neurons themselves compute $\delta$. In reinforcement-learning programs the prediction error is computed in steps $t,t+1,\dots$ There are no steps in the organism, so we write it in continuous time~\cite{Doya2000}. Let $r(t)$ be the reward flow and $V(t)$ the expectation of all future reward, with more distant reward valued less:
\begin{equation}
V(t) \;=\; \int_t^{\infty} e^{-(s-t)/\tau_\gamma}\,r(s)\,\dd s .
\end{equation}
Differentiating gives $\dd V/\dd t=V/\tau_\gamma-r$. The residual of this equation is the prediction error:
\begin{empheq}[box=\keybox]{equation}
\delta(t) \;=\; r(t) \;+\; \frac{\dd V}{\dd t} \;-\; \frac{V}{\tau_\gamma} .
\label{eq:ctd}
\end{empheq}
The constant $\tau_\gamma$ is the planning horizon, the continuous analog of the coefficient $\gamma$; by hypothesis it is set by \nt{SERO} (Section~\ref{sec:horizon}), and then~\eqref{eq:patience} is the rule for how long it is worth waiting. Let us check three cases (Fig.~\ref{fig:ctdcases}). A lamp promising juice lights up: $V$ jumps, $\dd V/\dd t$ is large, a burst. The promised juice arrives: $r>0$, but at the same moment the expectation is used up, $\dd V/\dd t\approx-r$, and there is no burst. The juice does not come: the expectation vanishes without a reward, $\dd V/\dd t<0$, a dip.

\begin{figure}[t]
\centering
\begin{tikzpicture}[>=Stealth,x=1.15cm,y=1cm]
\foreach \col/\ttl/\lampon/\juice in {0/{unexpected juice}/0/1, 1/{promised juice}/1/1, 2/{juice omitted}/1/0}{
 \begin{scope}[xshift=\col*4.6cm]
  \node[font=\small] at (1.5,3.55) {\ttl};
  \foreach \yy in {2.4,1.2,0}{\draw[gray!60!black] (0,\yy) -- (3.1,\yy);}
  % reward r
  \ifnum\juice=1
    \fill[green!45!black] (1.95,2.4) rectangle (2.05,2.85);
    \pic[scale=0.55] at (2.0,3.1) {juice};
  \else
    \pic[scale=0.55] at (2.0,3.1) {nojuice};
  \fi
  % expectation V and error delta
  \ifnum\lampon=1
    \pic[scale=0.55] at (1.0,3.1) {lamp};
    \draw[very thick,blue!60!black] (0,1.2) -- (1,1.2) -- (1,1.45)
      plot[domain=1:2,samples=20] (\x,{1.2+0.25*exp((\x-1)/1.5)}) -- (2,{1.2+0.25*exp(1/1.5)}) -- (2,1.2) -- (3.1,1.2);
    \draw[very thick,red!70!black] (0,0) -- (0.95,0) -- (0.95,0.5) -- (1.05,0.5) -- (1.05,0) -- (3.1,0);
    \ifnum\juice=0
      \draw[very thick,red!70!black] (1.95,0) -- (1.95,-0.45) -- (2.05,-0.45) -- (2.05,0);
      \node[font=\scriptsize,red!70!black,anchor=west] at (2.1,-0.35) {dip};
    \else
      \node[font=\scriptsize,gray!50!black,anchor=north,align=center] at (2.0,-0.08) {$r$ and the fall of $V$\\ cancel};
    \fi
  \else
    \draw[very thick,blue!60!black] (0,1.2) -- (3.1,1.2);
    \draw[very thick,red!70!black] (0,0) -- (1.95,0) -- (1.95,0.5) -- (2.05,0.5) -- (2.05,0) -- (3.1,0);
  \fi
  \ifnum\col=0
    \node[font=\small,anchor=east] at (-0.1,2.55) {$r$};
    \node[font=\small,anchor=east] at (-0.1,1.35) {$V$};
    \node[font=\small,anchor=east] at (-0.1,0.1) {$\delta$};
  \fi
 \end{scope}}
\end{tikzpicture}
\caption{Three cases for the error~\eqref{eq:ctd}; from top to bottom: reward $r$, expectation $V$ and error $\delta$. Left: there is no expectation, and the juice is entirely a surprise. Middle: the lamp raises $V$ in a jump; this is a jump in $\dd V/\dd t$, and the burst of $\delta$ falls on the lamp. After that $V$ grows exactly as the horizon requires, and $\delta=0$; at the moment of the juice, the reward $r$ and the fall of $V$ cancel each other. Right: $V$ falls without a reward, a dip. These are the same burst, transfer and dip as in Fig.~\ref{fig:rpe}, but now we see where they come from.}
\label{fig:ctdcases}
\end{figure}

What does~\eqref{eq:ctd} need from the circuits we already have? Three things.

\emph{The estimate $V$ itself.} It is produced by a group of \nt{GLUT} neurons, the critic, whose weights learn by the same rule~\eqref{eq:threefactor} with the same $\delta$: if $\delta>0$, the expectation was too low, and the connections active just before are strengthened.

\emph{The derivative $\dd V/\dd t$.} Any circuit that responds to changes rather than to level computes it: direct excitation together with delayed inhibition through a \nt{GABA} intermediary, as in the direction detector of Chapter~\ref{sec:gaba}, or a depressing synapse~\eqref{eq:depression} from Chapter~\ref{sec:code}.

\emph{A sign on a single wire.} The error can be positive or negative, but a spike rate can only be positive. The solution is the same as for the \nt{SERO} neuron in Chapter~\ref{sec:serocomputer}: the \nt{DOPA} neuron is a pacemaker. Its steady few spikes per second mean $\delta=0$, a burst means $\delta>0$, and a pause means $\delta<0$. A negative error has less room: the rate cannot drop below zero. In real \nt{DOPA} neurons the dips are indeed shallower than the bursts~\cite{Bayer2005}.

That dopamine behaves like $\delta$ has been measured. Exactly how the brain computes $\dd V/\dd t$ is a hypothesis.

\section{Actor and critic}

Let us put it all together (Fig.~\ref{fig:actorcritic}). Context neurons excite the action groups, which pick a winner through \nt{GABA} (Proposition~\ref{prop:wta}); this is the \emph{actor}. They also excite the critic, which outputs $V$~\cite{SuttonBarto2018}. The \nt{DOPA} neuron receives the reward $r$ and $V$, computes~\eqref{eq:ctd} and broadcasts $\delta$ to all synapses of the actor and the critic.

\begin{figure}[t]
\centering
\begin{tikzpicture}[>=Stealth,
  nrn/.style={circle,draw,very thick,fill=blue!6,minimum size=0.9cm,inner sep=0pt,font=\small},
  gab/.style={circle,draw,very thick,fill=red!10,minimum size=0.6cm,inner sep=0pt},
  dofa/.style={circle,draw,very thick,double,double distance=1.2pt,fill=orange!15,minimum size=1.35cm,inner sep=0pt,font=\scriptsize},
  inh/.style={-{Bar[width=7pt]},thick,red!70!black},
  bc/.style={->,thick,densely dashed,orange!85!black}]
\node[nrn] (S1) at (0,2.4) {$x_1$};
\node[nrn] (S2) at (0,1.2) {$x_2$};
\node[nrn] (S3) at (0,0) {$x_3$};
\node[left=0.15cm,align=right,font=\small] at (-0.5,1.2) {context};
\node[nrn] (A) at (4,2.6) {$A$};
\node[nrn] (B) at (4,1.0) {$B$};
\node[gab] (G) at (5.2,1.8) {};
\draw[->,thick] (A) -- (G); \draw[->,thick] (B) -- (G);
\draw[inh] (G) to[bend left=35] (A);
\draw[inh] (G) to[bend right=35] (B);
\node[above,font=\small] at (4.6,3.05) {actor: action selection};
\node[nrn] (V) at (4,-1.1) {$V$};
\node[below,font=\small] at (4,-1.6) {critic};
\foreach \s in {S1,S2,S3}{\foreach \t in {A,B,V}{\draw[->,gray!60!black] (\s) -- (\t);}}
\node[dofa] (D) at (8.0,-1.1) {\nt{DOPA}};
\draw[->,thick] (V) -- node[above,font=\scriptsize] {$V,\ \dd V/\dd t$} (D);
\draw[->,thick,green!45!black] (10.0,-1.1) node[right,black,font=\small] {reward $r$} -- (D);
\pic[scale=1.1] at (10.6,-0.5) {juice};
\draw[bc] (D.north) -- (8.0,4.0) -- node[above,font=\small,black] {$\delta$ to all synapses of the actor and the critic} (2.2,4.0) -- (2.2,3.0);
\draw[bc] (D.south) -- (8.0,-2.4) -- (2.2,-2.4) -- (2.2,-0.6);
\end{tikzpicture}
\caption{A network that learns to choose actions. Gray arrows: \nt{GLUT} synapses with changing weights; red circle: a \nt{GABA} neuron; double circle: the \nt{DOPA} pacemaker computing $\delta$ by~\eqref{eq:ctd}; dashed arrows: the broadcast of $\delta$.}
\label{fig:actorcritic}
\end{figure}

The sign of a transmitter's action is chosen by the recipient (Proposition~\ref{prop:receptor}), and in the action-selection region some neurons carry dopamine receptors of one family and others of another. In slice experiments, dopamine together with joint activity strengthens synapses on the first kind and weakens them on the second~\cite{Shen2008}; in the model this means that for the second kind the sign of $\delta$ in~\eqref{eq:threefactor} is flipped. The first learn what is \emph{worth doing}, the second what is \emph{not worth doing}.

\section{Summary}

\begin{table}[t]
\centering
\footnotesize
\setlength{\tabcolsep}{4pt}
\renewcommand{\arraystretch}{1.15}
\begin{tabular}{>{\raggedright}p{2.6cm}>{\raggedright}p{2.7cm}>{\raggedright}p{3.1cm}>{\raggedright\arraybackslash}p{5.0cm}}
\toprule
Rule & What the synapse knows & What the network learns & What is known \\
\midrule
rate-based Hebb~\eqref{eq:oja} & $x$, $y$, its own weight & compression of inputs: the principal directions & strengthening under joint activity is measured; the mechanism that bounds the weights is a subject of research \\
timing-based Hebb & the order of the $x$ and $y$ spikes within $\sim20\ms$ & causal links, sequences & measured at \nt{GLUT} synapses \\
three factors~\eqref{eq:threefactor} & $x$, $y$, trace $e$, shared signal $\delta$ & actions that bring reward & the effect of dopamine within seconds after activity is measured; as the main mechanism of learning to choose, a well-founded hypothesis \\
critic~\eqref{eq:ctd} & the same & the reward expectation $V$ & the similarity of dopamine to $\delta$ is measured; the circuit computing $\dd V/\dd t$ is a hypothesis \\
backpropagation of error & a personal correction for each weight & anything, but needs backward wires and a clock & not found in the brain in this form; approximations to it are being sought \\
\bottomrule
\end{tabular}
\caption{Learning rules in a network of \nt{GLUT}, \nt{GABA} and \nt{DOPA} neurons.}
\label{tab:learning}
\end{table}

The learning rules are collected in Table~\ref{tab:learning}. \nt{GLUT} carries the data, \nt{GABA} controls the flow, and \nt{DOPA} decides which changes in the network to keep.

\section{Exercises}

\begin{exercise}[\lvlA]
\label{ex:06:1}
\emph{How long the trace lives.} Joint activity $xy=1$ lasts $T_a=0.2\,\text{s}$, starting from $e=0$, and dopamine arrives $d=0.5\,\text{s}$ after it ends. (a)~How many times larger is the trace~\eqref{eq:threefactor} at that moment for $\tau_e=1\,\text{s}$ than for $\tau_e=50\ms$? (b)~For what $\tau_e$ is it largest?
\end{exercise}

\begin{exercise}[\lvlA]
\label{ex:06:2}
\emph{Drift of Hebb's rule.} The sender and the recipient emit independent Poisson trains of $10$ spikes per second, and each pair of spikes changes the weight according to the window of Fig.~\ref{fig:stdp} with $\tau_\pm=20\ms$, $A_-=0.6A_+$. By how many $A_+$ does the weight grow in an hour of completely random activity? For what $\tau_-$ does the drift vanish?
\end{exercise}

\begin{exercise}[\lvlB]
\label{ex:06:3}
\emph{Correlations, not covariances.} The rule of Proposition~\ref{prop:oja} finds the principal eigenvector of $C=\langle\mathbf x\mathbf x^{\mathsf T}\rangle$. Rates are positive: $\mathbf x=\mathbf m+\mathbf n$, $\mathbf m=(\mu,0)$, noise covariance $\operatorname{diag}(s_1^2,s_2^2)$. Under what condition does the neuron learn $\mathbf e_1$? What does it learn for $\mu=1$, $s_1=0.1$, $s_2=0.5$?
\end{exercise}

\begin{exercise}[\lvlA]
\label{ex:06:4}
\emph{The horizon in the burst.} A lamp lights up at an unpredictable moment, and juice of size $R$ arrives $L=1\,\text{s}$ later. Show that for the learned expectation, by~\eqref{eq:ctd}, $\delta\equiv0$ between the lamp and the juice, and find the area of the burst at the lamp for $\tau_\gamma=2\,\text{s}$ and $\tau_\gamma=4\,\text{s}$.
\end{exercise}

\begin{exercise}[\lvlB]
\label{ex:06:5}
\emph{Where the price of broadcasting comes from.} $N$ neurons jitter independently, $\xi_i\sim\mathcal N(0,\sigma^2)$, the error is $\delta=a\sum_i\xi_i$, and synapse $j$ estimates the gradient as $\delta\,\xi_jx_j$. Find the signal-to-noise ratio of one trial. One neuron learns in $100$ trials of $5\,\text{s}$ each: how long does learning take for $N=10^4$ and $N=10^{10}$?
\end{exercise}

\begin{exercise}[\lvlB]
\label{ex:06:6}
\emph{Design: a circuit that computes $\delta$.} A \nt{DOPA} neuron receives $r$, $V$ with weight $k_1$, and a copy of $V$ smoothed with $\tau_d$, with weight $-k_2$. (a)~Choose $k_1$, $k_2$ so that for slow $V$ the output is~\eqref{eq:ctd}. (b)~For $V=V_0e^{t/\tau_\gamma}$ the true $\delta=0$. For what $\tau_d$ is the circuit's spurious error at most $5\,\%$ of $V/\tau_\gamma$, if $\tau_\gamma=2\,\text{s}$?
\end{exercise}

\begin{exercise}[\lvlB]
\label{ex:06:7}
\emph{Simulation: the maximum reward delay.} Add a reward delay $d$ to \texttt{run()} in a copy of \texttt{models/\allowbreak dofa\_learning.py}. Find the fraction of $A$ choices in the last hundred of $500$ trials for $d=0.5\,\text{s}$ and $\tau_e=0.05$, $0.2$, $1$, $4\,\text{s}$, and for $\tau_e=1\,\text{s}$, $d=3\,\text{s}$. At what fraction of the trace surviving until the reward (Exercise~\ref{ex:06:1}) does learning disappear?
\end{exercise}

\begin{exercise}[\lvlC]
\label{ex:06:8}
\emph{Can the price of broadcasting be avoided?} $N$ neurons: $y_i=w_i+\xi_i$, $\xi_i\sim\mathcal N(0,\sigma^2)$, reward $R=-\sum_i(y_i-w_i^*)^2$, rule $\Delta w_i=\eta\,(R-\langle R\rangle)\,\xi_i$. With the best $\eta$, in how many trials does the error $\sum_i(w_i-w_i^*)^2$ fall by a factor of $e$? And if only $m$ random neurons out of $N$ jitter? Does sparse search help?
\end{exercise}

\chapter{Alternative Modern Theories of Dopamine-Based Learning}
\chaptermark{Other Theories of \nt{DOPA}}
\label{sec:dofaalt}

\section{What actually travels on the wire}

In Chapter~\ref{sec:dofa} the \nt{DOPA} system was a single wire carrying a single number, the reward-prediction error $\delta$~\eqref{eq:ctd}. This picture explains a great deal, but the more precisely dopamine is measured, the more is found that does not fit. Some neurons respond to unpleasant events as if to reward; the dopamine concentration rises while an animal runs toward a goal, even though nothing unexpected happens; different \nt{DOPA} neurons behave differently.

For an engineer, all these findings are one question: \emph{what is the signal format on the broadcast bus?} One number or a vector? One signal on the wire, or several separated by frequency? Does it speak of the future, like $\delta$, or of the past? Below are six modern answers; for each one we separate what is measured from what is assumed.

\section{Not one number but a distribution}

The error~\eqref{eq:ctd} teaches the critic the \emph{mean} reward expectation. But a sure half-drop of juice and a whole drop with probability one half have the same mean. To tell them apart, one needs the entire reward distribution.

Take the simplest task: after a cue, a reward of random size $r$ arrives. Let there be many \nt{DOPA} neurons, each of which, the $i$th, maintains its own estimate $V_i$, so that at the moment of reward its error is $\delta_i=r-V_i$. And let each weight a positive error by $\alpha_i^+$ and a negative one by $\alpha_i^-$. After each reward the estimate shifts by
\begin{equation}
\Delta V_i \;=\; \eta\,\bigl(\alpha_i^+\,[\delta_i]_+ \;-\; \alpha_i^-\,[-\delta_i]_+\bigr).
\label{eq:asymmetric}
\end{equation}
The estimate stops changing when, on average, positive corrections balance negative ones:
\begin{empheq}[box=\keybox]{equation}
\alpha_i^+\,\bigl\langle[r-V_i]_+\bigr\rangle \;=\; \alpha_i^-\,\bigl\langle[V_i-r]_+\bigr\rangle ,
\qquad
\kappa_i \;=\; \frac{\alpha_i^+}{\alpha_i^++\alpha_i^-} .
\label{eq:expectile}
\end{empheq}
For $\kappa_i=1/2$ this is the ordinary mean. For $\kappa_i>1/2$ the neuron is an ``optimist'': its estimate is shifted toward the upper end of the distribution; for $\kappa_i<1/2$ it is a ``pessimist.''

Example: juice arrives with probability $1/2$, and the drop has size $1$. Then~\eqref{eq:expectile} gives $\kappa_i\cdot\tfrac12(1-V_i)=(1-\kappa_i)\cdot\tfrac12V_i$, that is, $V_i=\kappa_i$ (Fig.~\ref{fig:expectiles}). A set of neurons with different $\kappa_i$ records the reward distribution the same way a set of quantiles records a distribution in statistics.

\begin{figure}[t]
\centering
\begin{tikzpicture}[>=Stealth,x=7cm,y=3cm]
\draw[->,gray!70!black] (-0.08,0) -- (1.15,0) node[right,black] {\small reward};
\draw[->,gray!70!black] (-0.08,0) -- (-0.08,0.75) node[left,black] {\small probability};
\fill[blue!25] (-0.03,0) rectangle (0.03,0.5);
\fill[blue!25] (0.97,0) rectangle (1.03,0.5);
\node[below,font=\small] at (0,0) {$0$};
\node[below,font=\small] at (1,0) {$1$};
\node[left,font=\scriptsize] at (-0.08,0.5) {$1/2$};
\foreach \k/\c/\lab/\h in {0.2/blue!60!black/{pessimist, $\kappa=0.2$}/0.62, 0.5/gray!50!black/{mean, $\kappa=0.5$}/0.42, 0.8/red!70!black/{optimist, $\kappa=0.8$}/0.22}{
  \draw[very thick,\c] (\k,0) -- (\k,\h);
  \node[above,font=\scriptsize,\c] at (\k,\h) {\lab};}
\end{tikzpicture}
\caption{The reward distribution recorded by a set of \nt{DOPA} neurons. The reward is $0$ or $1$ with probability $1/2$ (blue bars); a neuron with fraction $\kappa$ arrives at the estimate $V=\kappa$~\eqref{eq:expectile}.}
\label{fig:expectiles}
\end{figure}

What is measured: different \nt{DOPA} neurons have different ``reversal points'' (the reward size at which the response switches from a burst to a dip), and neurons with a larger asymmetry between responses to positive and negative errors reverse at larger rewards, as~\eqref{eq:expectile} requires~\cite{Dabney2020}. The shape of the distribution can be reconstructed from the responses of a group of neurons. That this is the main function of the diversity rather than a side effect is a hypothesis.

\begin{keyblock}
\begin{proposition}[A distribution from asymmetry]
\label{prop:expectile}
If \nt{DOPA} neurons differ in the fraction $\kappa_i$ with which they weigh positive errors, their estimates converge to different points~\eqref{eq:expectile} of the reward distribution. The group of neurons conveys not the mean but the distribution, and with it the risk.
\end{proposition}
\end{keyblock}

\section{Not only error but also salience}

By the error formula~\eqref{eq:ctd}, an unpleasant event (a shock, a bitter taste) should cause a dip: things turned out worse than expected. Some \nt{DOPA} neurons do exactly that, but others respond to unpleasant events with a \emph{burst}, as if to a reward~\cite{Matsumoto2009}. These neurons report not the sign of the error but that something \emph{important} has happened: unexpected, strong, demanding attention~\cite{BrombergMartin2010}.

An engineer can think of this as two signals. The first is the signed error $\delta$: in which direction to change the weights. The second is its magnitude $|\delta|$ or the novelty of the event: \emph{how much} to change them, that is, the learning rate; after an unexpected event, one should learn faster. In meaning it is close to the noradrenaline gain of Chapter~\ref{sec:neurontypes}.

There is also a third option: a separate wire for a separate axis. \nt{DOPA} neurons that project to one of the action-selection regions respond not to reward but to the novelty and intensity of a stimulus, and they are needed for the animal to learn to avoid danger~\cite{Menegas2018}: the wire carries not ``better or worse'' but ``dangerous or not.''

What is measured: different groups of \nt{DOPA} neurons respond differently to unpleasant and to novel events, and different outputs teach different things. How exactly the brain uses the salience signal (as a learning rate, as an attention signal, or as a separate error along a danger axis) is debated.

\section{Not only teaching but also urging on}

Dopamine also has an immediate effect: the more there is, the more readily and quickly the animal acts. How does this fit with learning?

The engineering answer is \emph{frequency separation}. Fast bursts and dips lasting hundreds of milliseconds carry $\delta$ and teach. A slow level, changing over minutes, carries another quantity: the average reward per unit time, $\bar R$~\cite{Niv2007}. Suppose an action performed in time $\tau_a$ costs an effort $C/\tau_a$: the faster, the more expensive. But each second of delay costs $\bar R$, the reward one could have earned in it. The total cost $C/\tau_a+\bar R\,\tau_a$ is minimal at
\begin{empheq}[box=\keybox]{equation}
\tau_a^{*} \;=\; \sqrt{\frac{C}{\bar R}} .
\label{eq:vigor}
\end{empheq}
The richer the environment, the more expensive time is and the faster it pays to act: if $\bar R$ doubles, the action time falls by a factor of $\sqrt2\approx1.4$. Note that $\bar R$ is the same expected reward that was subtracted in Chapter~\ref{sec:dofa}.

Dopamine release at the target can change without any change in the spike rate: it is adjusted by local signals at the outputs themselves~\cite{Mohebi2019}. But slow components are present in the spikes too: in the experiments of the next section, a smooth rise on approaching reward is also seen in the firing rate of \nt{DOPA} neurons~\cite{Kim2020}. So the slow level combines two sources, the spike rate and local regulation of release, and which one dominates apparently depends on the output and on the task.

\begin{keyblock}
\begin{proposition}[Two signals on one wire]
\label{prop:multiplex}
The \nt{DOPA} system carries at least two quantities separated in time: a fast error $\delta$, which teaches, and a slow level, which sets the cost of time~\eqref{eq:vigor} and thereby the speed of actions. That the fast and slow components behave differently is measured; that the slow one equals precisely $\bar R$ is a hypothesis.
\end{proposition}
\end{keyblock}

\section{The ramp on approach}

When an animal runs down a corridor toward a distant reward, the dopamine concentration in the action-selection region rises smoothly over seconds as the goal approaches~\cite{Hamid2016}. According to Chapter~\ref{sec:dofa}, this should not happen: everything is going according to plan, and $\delta$ should be zero.

The first explanation: this ramp is not $\delta$ but the expectation $V$ itself, the ``goal is near, worth the effort'' signal of the previous section~\cite{Hamid2016}. The second explanation keeps $\delta$~\cite{Kim2020}. If the expectation is correct, then with horizon $\tau_\gamma$, on the way to a reward $R$, it grows as $V=Re^{-(T-t)/\tau_\gamma}$, and by the error formula~\eqref{eq:ctd} $\dd V/\dd t-V/\tau_\gamma=0$. But suppose that from far away the expectation is too low, and the estimate is refined on approach. Then the expectation grows more steeply, say as $V=Re^{-(T-t)/\tau_v}$ with $\tau_v<\tau_\gamma$, and
\begin{empheq}[box=\keybox]{equation}
\delta(t) \;=\; \frac{\dd V}{\dd t}-\frac{V}{\tau_\gamma} \;=\; V\Bigl(\frac{1}{\tau_v}-\frac{1}{\tau_\gamma}\Bigr) \;>\; 0 .
\label{eq:ramp}
\end{empheq}
The error is positive and grows together with $V$: exactly the smooth ramp (Fig.~\ref{fig:ramp}). With $\tau_\gamma=4\,\text{s}$ and $\tau_v=1\,\text{s}$ we get $\delta=0.75\,V$ per second. This version was tested in a virtual corridor by instantly teleporting the animal closer to the goal: dopamine responded with a burst, as a derivative should, rather than with a smooth level~\cite{Kim2020}.

\begin{figure}[t]
\centering
\begin{tikzpicture}[>=Stealth,x=1.6cm,y=2.6cm]
\draw[->,gray!70!black] (-4.2,0) -- (0.5,0) node[below left,black] {\small $t$};
\draw[->,gray!70!black] (-4.2,0) -- (-4.2,1.2);
\foreach \x/\l in {-4/{$T-4$},-2/{$T-2$},0/{$T$}}{\draw[gray!60!black] (\x,-0.03) -- (\x,0.03); \node[below,font=\scriptsize] at (\x,0) {\l\,s};}
\draw[thick,densely dashed,gray!60!black] plot[domain=-4:0,samples=40] (\x,{exp(\x/4)});
\draw[very thick,blue!60!black] plot[domain=-4:0,samples=60] (\x,{exp(\x)});
\draw[very thick,red!70!black] plot[domain=-4:0,samples=60] (\x,{0.75*exp(\x)});
\node[anchor=south east,font=\small,gray!50!black] at (-1.3,0.77) {$V$ with $\tau_v=\tau_\gamma$: $\delta=0$};
\node[right,font=\small,blue!60!black] at (0.05,1.0) {$V$ with $\tau_v=1\,\text{s}$};
\node[right,font=\small,red!70!black] at (0.05,0.72) {$\delta$ by~\eqref{eq:ramp}};
\end{tikzpicture}
\caption{The dopamine ramp on approaching reward as a prediction error, $\tau_\gamma=4\,\text{s}$. Dashed line: an expectation growing exactly as the horizon requires ($\delta=0$); blue: an expectation growing more steeply; red: the error~\eqref{eq:ramp}.}
\label{fig:ramp}
\end{figure}

What is measured: the smooth ramp exists, both in dopamine concentration and in the firing rate of \nt{DOPA} neurons~\cite{Kim2020}, and instantaneous jumps of the environment cause jumps in dopamine. Which of the two explanations is correct, or whether both are correct at different outputs, is debated.

\section{Looking back}

All the theories above look \emph{forward}. One modern theory reverses the direction~\cite{Jeong2022}. Suppose the brain learns not to predict reward from a cue but, having received a reward, to \emph{look back}: which event preceded it more often than others? The measure of such a link is the difference of two probabilities:
\begin{empheq}[box=\keybox]{equation}
M \;=\; P(\text{cue before reward}) \;-\; P(\text{cue in a random window}) .
\label{eq:retro}
\end{empheq}
If the cue often occurs without reward, the second term is large, and the link is weak. In this theory dopamine reports not a prediction error but how likely a given event is to be a \emph{cause} of reward.

The look-back mechanism needs the same thing as the three-factor rule of Chapter~\ref{sec:dofa}: every event must leave a fading trace, so that at the moment of reward one can read what came before it. The difference lies not in the synapse but in what \nt{DOPA} transmits.

In experiments where theory~\eqref{eq:retro} and the error~\eqref{eq:ctd} make different predictions, dopamine release in several cases followed the former~\cite{Jeong2022}. But in other experiments the dopamine response during learning gradually shifted from the moment of reward to the moment of the cue, as learning by the error~\eqref{eq:ctd} requires~\cite{Amo2022}. Which theory describes the brain is not yet known.

\section{Error not only in reward}

The last extension concerns \emph{what} the error is about. If the expected juice is replaced by another, equally valuable but with a different flavor, the value has not changed, and the error~\eqref{eq:ctd} is zero. Yet \nt{DOPA} neurons respond to such a substitution~\cite{Takahashi2017}. And artificially evoked dopamine bursts can teach an animal an association between two neutral stimuli, with no reward at all~\cite{Sharpe2017}. It seems that \nt{DOPA} reports a prediction error not only about reward but also about \emph{what} will happen.

Formally, this is a simple generalization of~\eqref{eq:ctd}: instead of a single reward flow $r$ we take a set of features of what is happening, $\varphi_k$ (taste, place, sound), and for each its own expectation $M_k$ and its own error~\cite{Gardner2018}:
\begin{empheq}[box=\keybox]{equation}
\delta_k(t) \;=\; \varphi_k(t) \;+\; \frac{\dd M_k}{\dd t} \;-\; \frac{M_k}{\tau_\gamma} ,
\qquad k=1,\dots,K .
\label{eq:vectortd}
\end{empheq}
Reward is just one of the features, and the vector of errors teaches not only what is good but also what follows what: a model of the world.

The diversity of \nt{DOPA} neurons agrees with this: in the same task, different neurons carry different quantities, some related to reward, others to movement, still others to where attention is directed~\cite{Engelhard2019}.

\begin{keyblock}
\begin{proposition}[A bundle instead of a wire]
\label{prop:bundle}
The signal of the \nt{DOPA} system is not a single number. It is spread across neurons (the distribution~\eqref{eq:expectile}, different features~\eqref{eq:vectortd}, a separate danger axis) and across time (a fast error and a slow level~\eqref{eq:vigor}). The scalar error~\eqref{eq:ctd} is a first approximation to this signal, just as a summing unit with a threshold is a first approximation to a neuron.
\end{proposition}
\end{keyblock}

\section{Summary}

\begin{table}[t]
\centering
\footnotesize
\setlength{\tabcolsep}{4pt}
\renewcommand{\arraystretch}{1.15}
\begin{tabular}{>{\raggedright}p{2.5cm}>{\raggedright}p{3.2cm}>{\raggedright}p{3.6cm}>{\raggedright\arraybackslash}p{4.2cm}}
\toprule
Theory & What travels on the wire & Key measurement & What is known \\
\midrule
prediction error (Ch.~\ref{sec:dofa}) & scalar $\delta$~\eqref{eq:ctd} & burst, transfer to the cue, dip & measured; a first approximation \\
distribution & a set of $\delta_i$ with different asymmetry & different reversal points in different neurons & agrees with experiment; the role of the diversity is a hypothesis \\
salience & $|\delta|$, novelty; a danger axis & bursts to unpleasant events in some neurons & measured; how it is used is debated \\
cost of time & slow level $\bar R$ & dopamine speeds up actions; a slow component exists both in release at the outputs and in the spike rate & separation of fast and slow is measured; $\bar R$ is a hypothesis \\
ramp on approach & $V$, or $\delta$ as the estimate is refined~\eqref{eq:ramp} & smooth ramp, jumps on teleportation in the corridor & the ramp is measured; the explanation is debated \\
looking back & a measure of causal link~\eqref{eq:retro} & experiments where the two theories' predictions diverge & the results contradict one another \\
vector of errors & $\delta_k$ per feature~\eqref{eq:vectortd} & response to a flavor switch; learning without reward & measured; the vector form is a hypothesis \\
\bottomrule
\end{tabular}
\caption{What the \nt{DOPA} system transmits: the standard picture of Chapter~\ref{sec:dofa} and its modern extensions.}
\label{tab:dofaalt}
\end{table}

The theories of Table~\ref{tab:dofaalt} do not exclude one another: almost all of them keep the prediction error as the foundation and extend its format. Only looking back seriously competes with it, and that dispute will be settled by experiment. For an engineer the conclusion is simple: the price of broadcasting from Proposition~\ref{prop:broadcastcost} falls if the wire is really many wires with different destinations.

\section{Exercises}

\begin{exercise}[\lvlA]
\label{ex:07:1}
\emph{Distribution points for a single drop.} Juice of size $1$ arrives with probability $p=0.2$; otherwise the reward is $0$. Using~\eqref{eq:expectile}, find $V_i$ for $\kappa_i=0.2$, $0.5$, $0.8$. What is the median of this reward, and how does the point~\eqref{eq:expectile} differ from a quantile?
\end{exercise}

\begin{exercise}[\lvlB]
\label{ex:07:2}
\emph{A distribution from two neurons.} The reward is $0$ or $x$, and the probability of $x$ is $p$. Two neurons with rule~\eqref{eq:asymmetric} report $V(1/2)=0.25$ and $V(0.8)=0.5$. Reconstruct $p$, $x$ and the standard deviation of the reward. What will a neuron with $\kappa=0.2$ report?
\end{exercise}

\begin{exercise}[\lvlB]
\label{ex:07:3}
\emph{Simulation: a distribution from asymmetry.} Simulate nine \nt{DOPA} neurons with $\kappa_i=0.1,\dots,0.9$ and rule~\eqref{eq:asymmetric}, $\alpha_i^+=2\kappa_i$, $\alpha_i^-=2(1-\kappa_i)$, for a reward uniform on $[0,1]$. How does the scatter of $V_i$ depend on $\eta$, and what $\eta$ is needed to distinguish this distribution from two equally likely drops of $0.25$ and $0.75$?
\end{exercise}

\begin{exercise}[\lvlA]
\label{ex:07:4}
\emph{The cost of time.} (a)~Derive~\eqref{eq:vigor} and find the total cost at the optimum. (b)~The brain has misestimated $\bar R$ by a factor of $k$. Find the relative overpayment for $k=2$ and $k=4$. How precisely must the slow dopamine level report $\bar R$?
\end{exercise}

\begin{exercise}[\lvlB]
\label{ex:07:5}
\emph{A burst on teleportation.} The expectation grows by~\eqref{eq:ramp} with $\tau_v=1\,\text{s}$, $\tau_\gamma=4\,\text{s}$; the animal is instantly moved from point $T-2\,\text{s}$ to $T-1\,\text{s}$. Find the area of the $\delta$ burst as a fraction of $R$, and how many seconds of the ramp before teleportation it replaces. What will teleportation show if dopamine reports $V$ itself?
\end{exercise}

\begin{exercise}[\lvlB]
\label{ex:07:6}
\emph{Design: two signals on a wire.} The \nt{DOPA} wire carries a level rising at $s=1/60$ spikes per second per second, and events of $A=3$ spikes each. A synapse with trace $\tau_e=1\,\text{s}$ subtracts from the rate a copy of it smoothed with $\tau_f$. For which $\tau_f$ is at most $10\,\%$ of an event lost, while the level's contribution over the trace time is below $10\,\%$ of the event's?
\end{exercise}

\begin{exercise}[\lvlB]
\label{ex:07:7}
\emph{Designing an experiment: looking back versus error.} In a window the cue occurs with probability $0.1$, and after it a reward with probability $0.8$. Compare $\Delta P=P(\text{rew.}\mid\text{cue})-P(\text{rew.}\mid\text{none})$ and $M$ from~\eqref{eq:retro} if (a)~rewards without a cue are added, $0.08$ per window; (b)~the cue rate is doubled with no rewards.
\end{exercise}

\begin{exercise}[\lvlC]
\label{ex:07:8}
\emph{The bundle and the price of broadcasting.} In the model of Exercise~\ref{ex:06:8}, $N$ neurons are split into $K$ groups with their own wires, and a fraction $\rho$ of other groups' errors leaks into each wire. Show that the learning constant equals $N/K+2+\rho^2N(1-1/K)$. How many trials does it give for $K\to\infty$, $N=10^{10}$ and $\rho=0.01$?
\end{exercise}

\chapter{When to Explore and When to Decide: Adding \nt{NORA}}
\chaptermark{Adding \nt{NORA}}
\label{sec:nora}

\section{What is missing}

In Chapter~\ref{sec:dofa} the network learned thanks to noise: random deviations of activity were the search, and \nt{DOPA} reported whether they made things better. But one parameter was left unset there: \emph{how much} to search. If the noise is too strong, the network thrashes about and does not use what it already knows. If it is too weak, the network picks the same thing again and again and fails to notice that something nearby has become better. In reinforcement learning this is called the exploration--exploitation trade-off, and every algorithm has a knob for it. In Chapter~\ref{sec:neurontypes} we proposed that in the brain this knob is turned by \nt{NORA}.

A human has a few tens of thousands of noradrenaline neurons (Table~\ref{tab:types}), almost all of them gathered in one small group, and their outputs spread through practically the whole brain, although different parts of this group, it has turned out, partly serve different regions~\cite{Poe2020}. This is a broadcast channel even narrower than \nt{DOPA}. It works in two modes~\cite{AstonJones2005}. In the \emph{tonic} mode the neurons fire steadily, at rates from a fraction of a spike to a few spikes per second, and this rate changes slowly with wakefulness. In the \emph{phasic} mode they respond with a short burst to an event that matters for the current task, at roughly the moment of decision. A convenient but imprecise test point is pupil diameter: it changes with the activity of these neurons, but also with the activity of other regions~\cite{Joshi2016} and of cholinergic outputs in the cortex~\cite{Reimer2016}. The pupil shows the overall level of arousal, not \nt{NORA} alone.

\section{Gain}

Here is what is measured: noradrenaline suppresses the background activity of neurons more strongly than their response to a stimulus, and the signal-to-background ratio rises~\cite{Foote1975NE}. That the slope of an individual neuron's transfer function is literally multiplied by a factor does not follow from the experiment. As in Chapter~\ref{sec:neurontypes}, we take a simple model that captures this effect: noradrenaline changes the slope of the recipient's transfer function~\cite{ServanSchreiber1990}. Weak, subthreshold inputs (the background) then give even less, and strong ones give even more. Let the neuron respond with the rate
\begin{empheq}[box=\keybox]{equation}
r \;=\; F\bigl(g\,(u-\theta)\bigr), \qquad F(z)=\frac{r_{\max}}{1+e^{-z}} ,
\label{eq:gain}
\end{empheq}
where $u$ is the input current, $\theta$ is the threshold, and $g$ is the gain, which in the model is controlled by \nt{NORA}. At small $g$ the rate follows the input smoothly over a wide range: the neuron is an \emph{analog amplifier}. At large $g$ almost any input above threshold gives the full rate, and below threshold zero: the neuron is a \emph{comparator}, an almost digital element (Fig.~\ref{fig:gaincurves}). A single knob switches the same neuron between two modes that in electronics are served by different circuits.

\begin{figure}[t]
\centering
\begin{tikzpicture}[>=Stealth,x=1.1cm,y=2.4cm]
\draw[->,gray!70!black] (-3.3,0) -- (3.4,0) node[below left,black] {\small $u$};
\draw[->,gray!70!black] (-3.3,0) -- (-3.3,1.2) node[left,black] {\small $r$};
\draw[densely dashed,gray!60!black] (0,0) -- (0,1.1);
\node[below,font=\small] at (0,0) {$\theta$};
\draw[densely dotted,gray!60!black] (-3.3,1) -- (3.2,1) node[right,black,font=\small] {$r_{\max}$};
\draw[very thick,blue!60!black] plot[domain=-3.2:3.2,samples=60] (\x,{1/(1+exp(-0.8*\x))});
\draw[very thick,orange!85!black] plot[domain=-3.2:3.2,samples=60] (\x,{1/(1+exp(-2*\x))});
\draw[very thick,red!70!black] plot[domain=-3.2:3.2,samples=120] (\x,{1/(1+exp(min(-8*\x,9)))});
\node[blue!60!black,font=\small,anchor=west] at (0.5,0.12) {small $g$: amplifier};
\node[red!70!black,font=\small,anchor=east] at (-0.3,0.85) {large $g$: comparator};
\end{tikzpicture}
\caption{The transfer function~\eqref{eq:gain} for three values of the gain $g$.}
\label{fig:gaincurves}
\end{figure}

Does a large gain help against noise? Not always. Suppose two inputs differ by $\Delta u$, and we must say which is larger. If noise $\sigma_{\text{pre}}$ is added to the input \emph{before} the amplifier, both signal and noise are amplified, and their ratio does not change. But if noise $\sigma_{\text{post}}$ is added \emph{after} the amplifier (from unreliable synapses and random spikes of the network itself, as in Chapter~\ref{sec:code}), then the signal-to-noise ratio
\begin{empheq}[box=\keybox]{equation}
d' \;=\; \frac{g\,\Delta u}{\sqrt{g^2\sigma_{\text{pre}}^2+\sigma_{\text{post}}^2}}
\label{eq:gainsnr}
\end{empheq}
grows with $g$ until $g\sigma_{\text{pre}}$ becomes comparable to $\sigma_{\text{post}}$~\cite{ServanSchreiber1990}.

\begin{keyblock}
\begin{proposition}[When gain helps]
\label{prop:gainnoise}
By raising the gain, \nt{NORA} improves the discrimination of inputs only against noise that arises \emph{after} the amplifier, in the network itself. Gain cannot remove noise that arrives together with the input.
\end{proposition}
\end{keyblock}

\section{Gain in a two-way choice}

Let us return to the choice of Chapter~\ref{sec:gaba}: two groups of \nt{GLUT} neurons with inputs $I_1$, $I_2$ inhibit each other with strength $\beta$. Now each group has a gain $g$: in the choice equations~\eqref{eq:wta}, the expression in brackets is multiplied by $g$. While both groups are active, the steady-state rates are found from $r_1=g(I_1-\beta r_2)$, $r_2=g(I_2-\beta r_1)$. Adding and subtracting these equations gives the sum $r_1+r_2=g(I_1+I_2)/(1+g\beta)$ and the difference $r_1-r_2=g(I_1-I_2)/(1-g\beta)$. Their ratio tells how sharply the network prefers one input to the other:
\begin{empheq}[box=\keybox]{equation}
\frac{r_1-r_2}{r_1+r_2} \;=\; \varepsilon\,\frac{1+g\beta}{1-g\beta}, \qquad \varepsilon=\frac{I_1-I_2}{I_1+I_2}, \quad g\beta<1 .
\label{eq:wtagain}
\end{empheq}
A small input difference $\varepsilon$ is amplified by a factor of $(1+g\beta)/(1-g\beta)$, and this factor grows without bound as $g\beta$ approaches one. For $g\beta>1$ the state with both groups active is unstable, as in Proposition~\ref{prop:wta}: one group wins, the other falls silent (Fig.~\ref{fig:pitchfork}).

\begin{figure}[t]
\centering
\begin{tikzpicture}[>=Stealth,x=3.2cm,y=1.5cm]
\draw[->,gray!70!black] (0,0) -- (2.15,0) node[below left,black] {\small $g\beta$};
\draw[->,gray!70!black] (0,-1.25) -- (0,1.35) node[above,black] {\small $(r_1-r_2)/(r_1+r_2)$};
\draw[gray!60!black] (1,-0.04) -- (1,0.04); \node[below right,font=\small] at (1,0) {$1$};
\node[left,font=\scriptsize] at (0,1) {$1$}; \node[left,font=\scriptsize] at (0,-1) {$-1$};
\draw[very thick,blue!60!black] plot[domain=0:0.905,samples=60] (\x,{0.05*(1+\x)/(1-\x)}) -- (1,1) -- (2.05,1);
\draw[very thick,blue!60!black] (1.105,-1) -- (2.05,-1);
\draw[thick,densely dashed,gray!60!black] (1,0) -- (2.05,0);
\node[font=\small,align=center] at (0.45,-0.62) {deliberating:\\both groups\\active};
\node[font=\small,align=center] at (1.55,0.45) {decided:\\one group active};
\node[font=\small,blue!60!black,anchor=south west] at (1.05,-0.97) {weak input won earlier and holds};
\end{tikzpicture}
\caption{A two-way choice at different gains, with inputs differing by $\varepsilon=0.05$. For $g\beta>1$ both decisions are stable (solid lines; the weak input's win is stable for $g\beta>(1+\varepsilon)/(1-\varepsilon)\approx1.1$), and the network holds whichever it has already made; the symmetric state (dashed line) is unstable.}
\label{fig:pitchfork}
\end{figure}

\begin{keyblock}
\begin{proposition}[Gain switches the choice mode]
\label{prop:gainwta}
In a choice network with mutual inhibition $\beta$, the gain $g$ carries the network across the boundary $g\beta=1$. Below it the network ``deliberates'': both groups are active, and the input difference is amplified by~\eqref{eq:wtagain}. Above it the network ``has decided'': one group is active, and the decision holds. By changing $g$, \nt{NORA} can switch the network between these modes without touching a single weight.
\end{proposition}
\end{keyblock}

\section{Gain as temperature}

Now let us add to the choice noise that arises in the network itself, after the amplifier. Which group wins is decided by the sign of $g(u_A-u_B)+\eta$, where $u_A-u_B$ is the input difference, that is, the difference of the learned weights from Chapter~\ref{sec:dofa}, and $\eta$ is noise of scale $\sigma$. If the noise has a logistic distribution (for a Gaussian the curve is almost the same), the probability of choosing $A$ is
\begin{empheq}[box=\keybox]{equation}
P(A) \;=\; \frac{1}{1+e^{-g\,(u_A-u_B)/\sigma}} .
\label{eq:softmax}
\end{empheq}
A programmer will recognize softmax choice with temperature $T=\sigma/g$. The gain is the inverse temperature. At small $g$ even a clearly better option is chosen only slightly more often than a worse one: the network explores a lot. At large $g$ the best known option is chosen almost always: the network exploits what it knows.

Let us be precise about what $g$ in~\eqref{eq:wtagain} and~\eqref{eq:softmax} is: it is the \emph{effective} gain on the input difference of the current task at the moment of choice. The two modes of \nt{NORA} act on it in opposite ways. A phasic burst at the moment of decision raises it, and the network locks in its choice. High tonic activity, by contrast, goes together with exploration: in humans the baseline pupil is wider before a random choice~\cite{JepmaNieuwenhuis2011}, and in rats \nt{NORA} input to the anterior cingulate cortex switches choice into a random mode~\cite{Tervo2014stochastic}. So high tonic \nt{NORA} corresponds to a \emph{small} effective $g$, and a burst to a large one.

\section{How much to explore}

Which gain is best? We tested this on a model. The network chooses one of two actions by~\eqref{eq:softmax}; each brings reward with some probability, and the network estimates it from the rewards of the chosen action, as in Chapter~\ref{sec:dofa}. From time to time the world changes: both probabilities are drawn anew. The change can affect the chosen action or the one the network has not tried for a long time; about the latter the network can learn only by exploring. Figure~\ref{fig:invertedU} shows the fraction of the best possible reward that the network collects at different constant $g$.

\begin{figure}[t]
\centering
\begin{tikzpicture}[>=Stealth,x=1.2cm,y=1cm]
\draw[->,gray!70!black] (0,0) -- (7.6,0) node[below left,black] {\small $g$};
\draw[->,gray!70!black] (0,0) -- (0,4.4) node[above,black,font=\small] {fraction of best reward};
\foreach \x/\l in {0/0.5,1/1,2/2,3/4,4/8,5/16,6/32,7/64}{\draw[gray!60!black] (\x,-0.06) -- (\x,0.06); \node[below,font=\scriptsize] at (\x,0) {\l};}
\foreach \v/\y in {0.8/0.8,0.9/2.4,1.0/4.0}{\draw[gray!60!black] (-0.06,\y) -- (0.06,\y); \node[left,font=\scriptsize] at (0,\y) {\v};}
\draw[very thick,blue!60!black] plot[mark=*,mark size=1.3pt] coordinates {(0,0.368) (1,0.880) (2,1.728) (3,2.784) (4,3.456) (5,3.616) (6,3.488) (7,3.264)};
\draw[very thick,red!70!black] plot[mark=*,mark size=1.3pt] coordinates {(0,0.368) (1,0.672) (2,1.184) (3,1.840) (4,2.176) (5,1.920) (6,1.456) (7,1.104)};
\node[blue!60!black,font=\small,anchor=south] at (5,3.7) {calm world};
\node[red!70!black,font=\small,anchor=north] at (5.1,1.25) {fast-changing world};
\draw[->,thick,blue!60!black] (5,3.25) -- (5,3.55);
\draw[->,thick,red!70!black] (4,1.8) -- (4,2.1);
\node[font=\small\itshape,gray!35!black,anchor=west] at (0.15,2.3) {searches blindly};
\node[font=\small\itshape,gray!35!black] at (6.45,2.55) {misses changes};
\end{tikzpicture}
\caption{Fraction of the best possible reward at a constant gain $g$ (logarithmic $g$ scale). Blue curve: the world changes on average once per thousand trials; red: once per twenty. Arrows mark the best $g$: in a fast-changing world it is half as large. Average over 60 runs of 4000 trials.}
\label{fig:invertedU}
\end{figure}

At $g=0.5$ the network hardly uses its knowledge and collects about three quarters of what is possible; at very large $g$ it misses changes. The best value is $g=16$ when the world changes once per thousand trials (fraction $0.976$), and $g=8$ when it changes once per twenty ($0.886$).

\begin{keyblock}
\begin{proposition}[The inverted U]
\label{prop:explore}
The reward the network collects depends on the gain as an inverted U: too small a $g$ wastes trials on exploration, too large a $g$ misses changes. The faster the world changes, the smaller the best $g$.
\end{proposition}
\end{keyblock}

The inverted U is also observed in experiments: animals perform a task best at moderate tonic activity of \nt{NORA} neurons with sharp phasic responses, while at high tonic activity they become distractible and switch between tasks~\cite{AstonJones2005}. This agrees with the distinction between the modes in the previous section: a phasic burst at the moment of decision gives a large effective $g$ exactly when a choice must be made, while high tonic activity without bursts amplifies everything indiscriminately, including irrelevant inputs, so the effective $g$ for the current task falls: this is exploration.

\section{Who turns the knob}

Since the best gain depends on how fast the world changes, it would be good to adjust it. But how can \nt{NORA} neurons know that the world has changed? A natural sign is \emph{surprise}: prediction errors have become larger than usual. It is important to distinguish two kinds of uncertainty. If the reward is simply random, the errors are always large, and that is expected. But if the errors suddenly grow relative to their usual size, something has changed. According to one hypothesis, it is this unexpected uncertainty that \nt{NORA} reports, while \nt{ACET} reports the expected one~\cite{YuDayan2005}.

What should be done with this signal? One can lower the gain and explore more. One can raise the learning rate so as to forget outdated estimates faster: experiments show that after an abrupt change people learn faster, and their pupils dilate~\cite{Nassar2012}; recall that the pupil is an indicator not only of \nt{NORA} but also of \nt{ACET}~\cite{Reimer2016}. Or one can do both at once: according to yet another hypothesis, the phasic \nt{NORA} burst works as an \emph{interrupt}, a signal on which the network resets its current state and reconfigures for the new situation~\cite{BouretSara2005}, like a processor's global interrupt line.

To be honest, here is what we did not find. In the model we tried both simple adjustment schemes: lowering $g$ or raising the learning rate when the smoothed error exceeds its long-term average. In a world that alternately stays put for long stretches and changes often, the best settings of both schemes gave $0.926$--$0.927$ of the best reward, almost the same as the best constant $g$ ($0.925$ at $g=8$) or a constant but sufficiently large learning rate ($0.928$), and poor settings gave less. The curves in Fig.~\ref{fig:invertedU} are flat near the top, and in such a world there is almost nothing to adjust. Adjustment should pay off where different regimes require very different settings. Exactly which control law \nt{NORA} implements has not yet been established.

\section{Summary}

\begin{table}[t]
\centering
\footnotesize
\setlength{\tabcolsep}{4pt}
\renewcommand{\arraystretch}{1.15}
\begin{tabular}{>{\raggedright}p{3.0cm}>{\raggedright}p{4.4cm}>{\raggedright\arraybackslash}p{6.0cm}}
\toprule
What \nt{NORA} does & How & What is known \\
\midrule
changes the slope of neuronal responses & the gain $g$ in~\eqref{eq:gain} & the rise in signal-to-noise ratio is measured; the multiplicative model is a simplification \\
switches the choice & carries the network across $g\beta=1$ & follows from the model~\eqref{eq:wtagain}; no direct measurements of the threshold \\
sets how much to explore & $g$ is the inverse temperature~\eqref{eq:softmax}; tonic: small effective $g$ (exploration), burst: large (decision) & the inverted U and the two modes of operation are measured; the link to exploration is a well-founded hypothesis \\
reports changes & unexpected uncertainty & pupil size and learning rate rise after changes: measured; that this is a \nt{NORA} signal is a hypothesis \\
interrupts and resets & a phasic burst to an unexpected event & bursts to important events are measured; the ``interrupt'' is a hypothesis \\
\bottomrule
\end{tabular}
\caption{What \nt{NORA} adds to a network of \nt{GLUT}, \nt{GABA} and \nt{DOPA}.}
\label{tab:nora}
\end{table}

\nt{DOPA} tells the network \emph{where} to change the weights. \nt{NORA} changes no weights at all, but decides \emph{how} the network uses what it already knows (Table~\ref{tab:nora}): a single knob, the gain, turns a neuron from an amplifier into a comparator, a choice network from ``deliberating'' into ``decided,'' and choice from exploration into exploitation. How \nt{NORA} finds out how fast the world is changing is an open question.

\section{Exercises}

\begin{exercise}[\lvlA]
\label{ex:08:1}
\emph{One knob for the whole brain.} (a)~Using Table~\ref{tab:types}, estimate how many brain neurons there are per \nt{NORA} neuron. (b)~The noise before the amplifier is $\sigma_{\text{pre}}=0.5$, after it $\sigma_{\text{post}}=4$. At what $g$ does $d'$ from~\eqref{eq:gainsnr} reach $90\,\%$ of its limit?
\end{exercise}

\begin{exercise}[\lvlB]
\label{ex:08:2}
\emph{Choice with a gain.} $\tau\,\dd r_1/\dd t=-r_1+g[I_1-\beta r_2]_+$, $\tau\,\dd r_2/\dd t=-r_2+g[I_2-\beta r_1]_+$, $\tau=20\ms$. (a)~How fast does $r_1-r_2$ decay for $g\beta=0.5$ and $0.9$? (b)~For $\varepsilon=0.05$, find the $g\beta$ below which both groups are active and above which the weak input's win is stable.
\end{exercise}

\begin{exercise}[\lvlA]
\label{ex:08:3}
\emph{Softmax from noise.} Each of $K$ groups receives its own Gumbel noise, $P(\eta_k<z)=\exp(-e^{-z/\sigma})$, and the largest $gu_k+\eta_k$ wins. Find $P(k)$. At what $g$ is the better of two options with $u_A-u_B=0.1$, $\sigma=1$, chosen $90\,\%$ of the time?
\end{exercise}

\begin{exercise}[\lvlB]
\label{ex:08:4}
\emph{Estimating the best gain.} In the model of Fig.~\ref{fig:invertedU}, reward probabilities after a change are uniform on $[0,1]$. The network tries the worse action once every $e^{g\Delta}$ trials; setting this equal to $1/h$ gives $g^*\approx\ln(1/h)/\bar\Delta$. Find $\bar\Delta=\langle|p_A-p_B|\rangle$ and $g^*$ for $h=0.001$, $0.01$, $0.05$.
\end{exercise}

\begin{exercise}[\lvlB]
\label{ex:08:5}
\emph{Simulation: gain and the rate of change.} Find $g^*(h)$ for $h$ from $0.001$ to $0.2$ with a copy of \texttt{models/\allowbreak nora\_explore.py} (it writes files to the current folder). Where is the estimate of Exercise~\ref{ex:08:4} correct, and why does $g^*$ stop decreasing when changes are fast?
\end{exercise}

\begin{exercise}[\lvlB]
\label{ex:08:6}
\emph{Design: an interrupt for a choice network.} The network of Exercise~\ref{ex:08:2} with $\beta=2$, $\tau=20\ms$, $g=1$ holds $r_1=0.9$, $r_2=0$, although the inputs are now $I_1=0.9$, $I_2=1.0$. \nt{NORA} lowers $g$ to $g_{\text{lo}}$ for a time $T_p$. Find the smallest $T_p$ analytically for $g_{\text{lo}}=0$ and by simulation for $g_{\text{lo}}=0.25$.
\end{exercise}

\begin{exercise}[\lvlC]
\label{ex:08:7}
\emph{When adjustment pays off.} An oracle knows whether the epoch is calm or volatile and sets its own $g$ in each. (a)~By simulation, find its advantage over the best constant $g$ in the chapter's world (epochs of $1000$ trials, $h=0.001$ and $0.05$). (b)~Construct a world where the advantage is larger, and find what part of it surprise-driven adjustment captures.
\end{exercise}

\chapter{When to Write and When to Read: Adding \nt{ACET}}
\chaptermark{Adding \nt{ACET}}
\label{sec:acet}
\begin{chaptergoals}
\item why a memory that stores and learns in the same connections needs a write-enable line;
\item how \nt{ACET} strengthens the input, weakens internal connections and eases weight changes;
\item why no single \nt{ACET} level serves both writing and reading, so a switch is needed;
\item why the optimal filter sets input weight and learning rate with one number, $P/R$.
\end{chaptergoals}

\section{What is missing}

Besides its address and data lines, a memory chip has one more line: write enable. While it is off, the memory is only read; when it is on, whatever is on the data lines is written. Without such a line memory does not work: if every read also writes something, whatever was read, including anything read in error, is immediately written back.

In the brain this problem is more acute than in a chip. In Chapter~\ref{sec:glutcomputer}, memory was a group of \nt{GLUT} neurons kept switched on by its own connections (Fig.~\ref{fig:bistable}). In Chapter~\ref{sec:dofa} we saw that these connections learn by Hebb's rule right in the course of operation. So the same synapses both store the old and take in the new, and there is no separate write phase, just as there is no clock. How does the network tell the moment when it should memorize what it sees from the moment when it should recall what it knows? In Chapter~\ref{sec:neurontypes} we proposed that this line is held by \nt{ACET}. In this chapter we examine how such a line should work and why it cannot be done without.

\section{Three actions of one wire}

The acetylcholine neurons that work inside the brain are few, gathered in a few small groups, and their outputs spread throughout the cortex. This is a broadcast channel, just like \nt{DOPA} and \nt{NORA}. The acetylcholine level in the cortex is high during attentive waking and low in slow-wave sleep~\cite{Hasselmo1999}. As with dopamine, the meaning of the signal is set by the recipient's receptor (Proposition~\ref{prop:receptor}): acetylcholine has both fast receptors, which open a channel, and slow ones, which trigger reactions inside the cell. Three actions matter to us.

\emph{First: weaken the internal connections.} Slice experiments on the olfactory cortex showed that acetylcholine strongly suppresses transmission along connections between neurons of the same region and hardly touches the inputs arriving from outside~\cite{HasselmoBower1992}.

\emph{Second: strengthen the inputs.} Acetylcholine raises the excitability of neurons and facilitates transmission along some input pathways; the signal from the senses becomes louder than the network's own activity~\cite{Hasselmo2006}.

\emph{Third: make weight changes easier.} At a high acetylcholine level, connections change more easily~\cite{Hasselmo2006}.

For an engineer, all three actions are a single operation: switching from ``read'' mode to ``write'' mode. The network stops listening to itself, starts listening to the input, and memorizes what it hears.

\section{A network that completes}

Take $N$ \nt{GLUT} neurons connected to one another. Store several patterns in their connections by Hebb's rule~\cite{Hebb1949}, each pattern being a set of $pN$ simultaneously active neurons. If half a pattern is presented to such a network, the connections excite the missing half: the network \emph{completes} the pattern from a part, just as the group in Fig.~\ref{fig:bistable} sustained itself. This is associative memory: part of the content serves as the address. \nt{GABA}, as in Chapter~\ref{sec:gaba}, keeps the total number of active neurons near $pN$, so that two patterns cannot light up together.

Denote the acetylcholine level by $a$, from $0$ to $1$, and write its three actions directly into the model:
\begin{empheq}[box=\keybox]{equation}
\tau\frac{\dd r_i}{\dd t} = -r_i + F\Bigl(\tfrac{1+a}{2}\,x_i + (1-a)\sum_j W_{ij}\,r_j - g\Bigr),
\qquad
\frac{\dd W_{ij}}{\dd t} = \eta\,a\,(r_i-p)(r_j-p) .
\label{eq:acetnet}
\end{empheq}
Here $r_i$ is the rate of neuron $i$, $x_i$ is its input from the senses, $F$ is a threshold transfer function, and $g$ is the global inhibition from \nt{GABA}, which grows when more than $pN$ neurons are active. The factor $\frac{1+a}{2}$ is the input gain, $1-a$ the weakening of internal connections, and $\eta a$ the learning rate. The second equation is Hebb's rule~\eqref{eq:hebb} in a form suited to sparse patterns~\cite{Tsodyks1988}: the weight grows when both neurons are more active than usual and decreases when only one is active. The negative part of the weights, as always, is implemented by \nt{GABA} intermediaries. There is no clock: both neurons and weights change continuously.

\begin{figure}[t]
\centering
\begin{tikzpicture}[>=Stealth,
  nrn/.style={circle,draw,very thick,fill=blue!6,minimum size=0.8cm,inner sep=0pt,font=\small},
  acet/.style={circle,draw,very thick,double,double distance=1.2pt,fill=orange!15,minimum size=1.35cm,inner sep=0pt,font=\scriptsize},
  bc/.style={->,thick,densely dashed,orange!85!black}]
\foreach \i/\y in {1/2.4,2/1.2,3/0}{
  \node[nrn] (N\i) at (3,\y) {$r_\i$};
  \draw[->,thick,blue!60!black] (0,\y) node[left,black] {\small $x_\i$} -- (N\i);}
\node[left,align=right,font=\small] at (-0.5,1.2) {from the\\senses};
\draw[->,thick,gray!60!black] (N1) to[bend left=30] (N2);
\draw[->,thick,gray!60!black] (N2) to[bend left=30] (N1);
\draw[->,thick,gray!60!black] (N2) to[bend left=30] (N3);
\draw[->,thick,gray!60!black] (N3) to[bend left=30] (N2);
\draw[->,thick,gray!60!black] (N1.east) .. controls (4.6,2.0) and (4.6,0.4) .. (N3.east);
\node[right,font=\small,gray!40!black,align=left] at (4.7,1.2) {internal\\connections $W$};
\node[acet] (A) at (1.2,-1.9) {\nt{ACET}};
\draw[bc] (A) -- (1.6,-0.15);
\node[font=\small,orange!60!black,anchor=east] at (0.95,-0.9) {$+$ input};
\draw[bc] (A) -- (4.3,0.6);
\node[font=\small,orange!60!black,anchor=west] at (4.3,0.1) {$-$ internal connections, $+$ learning rate};
\node[right,font=\small,align=left] at (2.2,-1.9) {high level: ``write''\\low level: ``read''};
\end{tikzpicture}
\caption{\nt{ACET} as a write-enable line. Neurons $r_i$ receive input $x_i$ from the senses (blue arrows) and excite one another through internal connections $W$ (gray), in which the patterns are stored. Acetylcholine (dashed arrows) strengthens the input, weakens the internal connections and makes weight changes easier, as in~\eqref{eq:acetnet}.}
\label{fig:acetnet}
\end{figure}

\section{Writing and reading interfere}

Let us test model~\eqref{eq:acetnet} on a task in which writing and reading really do conflict. A network of $200$ neurons already stores five patterns of $20$ active neurons each. Now it must memorize a sixth, which half overlaps one of the old ones: they share ten neurons. This happens all the time: a new room resembles a familiar one, a new face an old one.

\emph{Writing at low $a$.} First let us separate the first two actions of acetylcholine from the third: we keep the learning rate at full value and vary only the input gain and the weakening of internal connections. At $a=0$ the input lights up the ten shared neurons, and they, through the internal connections, light up the remaining half of the \emph{old} pattern. \nt{GABA} does not let both burn at once, and the old pattern, supported by its own connections, displaces the new one, which is held only by the input. The network sees not what is in front of it but what it remembers, and Hebb's rule diligently writes down what it sees.

As a result, the new pattern is not memorized at all: from its distinctive half, the network recalls nothing. And the old one is corrupted: when evoked by its own distinctive half, it drags along on average six foreign neurons out of ten. At $a=0.1$ the new pattern is memorized, but it merges with the old one, and the corruption is even worse: each of the two, evoked by its own half, drags along about eight foreign neurons; from $a=0.2$ on, writing is clean: fewer than one extra neuron on recall (averaged over ten runs).

\emph{Writing with the actual learning rate.} Now let acetylcholine, as in~\eqref{eq:acetnet}, also set the learning rate. After a single presentation the new pattern is memorized in full only for $a\ge0.9$; at $a=0.75$, to eight tenths; for $a\le0.5$, not at all.

\emph{Reading.} Present half of a pattern to the network with five cleanly written patterns, and then remove it too. For $a\le0.2$ the network completes the pattern in full and holds it on its own; at $a=0.3$, almost in full; for $a\ge0.4$ the internal connections are weakened too much, and once the cue is gone the activity dies out (Fig.~\ref{fig:acetsim}).

\begin{figure}[t]
\centering
\begin{tikzpicture}[>=Stealth,x=9cm,y=3.2cm]
\draw[->,gray!70!black] (0,0) -- (1.1,0) node[right,black] {\small $a$};
\draw[->,gray!70!black] (0,0) -- (0,1.2);
\foreach \x in {0,0.2,0.4,0.6,0.8,1.0}{\draw[gray!60!black] (\x,-0.02) -- (\x,0.02); \node[below,font=\scriptsize] at (\x,0) {$\x$};}
\foreach \y in {0,0.5,1}{\draw[gray!60!black] (-0.01,\y) -- (0.01,\y); \node[left,font=\scriptsize] at (0,\y) {$\y$};}
\node[rotate=90,font=\small] at (-0.1,0.6) {fraction of pattern recovered};
\draw[very thick,blue!60!black] plot[mark=*,mark size=1.6pt] coordinates {(0,1.00) (0.1,1.00) (0.2,1.00) (0.3,0.96) (0.4,0.04) (0.5,0.00) (0.75,0.00) (1.0,0.00)};
\draw[very thick,red!70!black] plot[mark=*,mark size=1.6pt] coordinates {(0.1,0.00) (0.2,0.00) (0.3,0.00) (0.5,0.00) (0.75,0.80) (0.9,1.00) (1.0,1.00)};
\node[font=\small,blue!60!black,anchor=west,align=left] at (0.03,0.5) {reading:\\pattern\\from half};
\node[font=\small,red!70!black,anchor=east,align=right] at (1.0,0.35) {writing:\\new pattern\\in one go};
\end{tikzpicture}
\caption{Model~\eqref{eq:acetnet}: $200$ neurons, patterns of $20$ active neurons, averaged over ten runs. Blue points, reading: the fraction of a pattern recovered from its half after the cue is removed. Red points, writing: the fraction of a new pattern, half similar to an old one, recalled after a single presentation, when acetylcholine also sets the learning rate. Reading works for $a\le0.3$, writing for $a\ge0.9$.}
\label{fig:acetsim}
\end{figure}

\begin{keyblock}
\begin{proposition}[Writing and reading require different modes]
\label{prop:writeread}
In model~\eqref{eq:acetnet}, where the same connections store the old and learn the new, and acetylcholine also sets the learning rate, there is no level $a$ at which writing and reading both work equally well. For writing, the internal connections must be weakened, or else the network writes down not the input but what it recalled; for reading, they must be strengthened, or else it cannot complete a pattern from a part. A switch is needed, and a single broadcast wire can serve as one.
\end{proposition}
\end{keyblock}

If acetylcholine did not change the learning rate, a narrow band $0.2\le a\le0.3$ would suit both tasks. But then the network would also learn during every read, that is, it would rewrite what it read, errors included. The idea itself, suppressing internal connections during learning, comes from models built on slice measurements~\cite{HasselmoAndersonBower1992}. The numbers above refer to our simplified model; exactly where the mode boundaries lie in the brain is unknown.

\section{How much to trust your eyes}

The write/read switch is an extreme case of a more general question: how much to trust the input and how much to trust memory? Engineers have an exact answer to this question, and it sheds light on why one wire controls both the mode and the learning rate.

Let the network track a quantity $s(t)$ that drifts slowly: its variance grows by $q$ per second. The senses report $y(t)=s(t)+{}$noise of intensity $R$. The best estimate $\hat s$ in continuous time is given by the Kalman--Bucy filter~\cite{KalmanBucy1961}:
\begin{empheq}[box=\keybox]{equation}
\frac{\dd\hat s}{\dd t} \;=\; \frac{P}{R}\,\bigl(y-\hat s\bigr),
\qquad
\frac{\dd P}{\dd t} \;=\; q-\frac{P^2}{R} ,
\label{eq:kalman}
\end{empheq}
where $P$ is the variance of the estimation error, that is, how unsure the network itself is of what it knows. The first equation is the same RC circuit as the membrane in the neuron model~\eqref{eq:glutneuron}, but with a time constant $R/P$ that varies. When the network is unsure, $P$ is large, the time constant is small, and the estimate quickly follows the input: this is ``write.'' When it is sure, $P$ is small, and the estimate hardly listens to the input, holding on to memory: this is ``read.''

In steady state $P=\sqrt{qR}$, and the memory time constant is $\sqrt{R/q}$: if the world drifts by one unit in a hundred seconds, $q=0.01$, and the sensory noise is $R=1$, then memory should last about ten seconds.

The main point here is that one number, $P/R$, sets two things at once: the weight of the input relative to memory, and the rate at which the stored estimate changes. This is exactly what acetylcholine does in~\eqref{eq:acetnet}. According to a hypothesis~\cite{YuDayan2005}, acetylcholine tells the network a quantity like $P$: the \emph{expected} uncertainty, the kind the network knows about in advance. This channel is not always slow: some \nt{ACET} neurons respond to reward and punishment within twenty milliseconds or so, and the more unexpected the reinforcement, the stronger the response~\cite{Hangya2015}.

\begin{keyblock}
\begin{proposition}[One knob for input weight and learning rate]
\label{prop:kalman}
In the optimal filter~\eqref{eq:kalman}, the weight with which the input is mixed with memory and the learning rate are one and the same quantity, $P/R$. It therefore makes sense to control them with one signal. When the properties of the world do not change, this signal settles at a constant level $\sqrt{q/R}$, and it needs adjusting only when the properties of the world change.
\end{proposition}
\end{keyblock}

The second sentence of the proposition explains the result of Chapter~\ref{sec:nora}, where surprise-driven adjustment of the learning rate did not beat the best constant rate: in a world whose rate of change is fixed, the best learning rate is a constant. A gain appears when the world suddenly becomes different. Then the estimate is suddenly far from the input, and $P$ knows nothing about it. The right action is to raise $P$ sharply and thereby switch to write mode. According to the hypothesis we discussed in Chapter~\ref{sec:nora}, such an \emph{unexpected} change is reported by \nt{NORA}~\cite{YuDayan2005}. The division of labor is natural: \nt{NORA} notices that the world has changed, and \nt{ACET} keeps the network in write mode until the new situation has been learned.

\section{Sleep as read mode}

If a high acetylcholine level is write mode, what does the network do at a low one? In slow-wave sleep the acetylcholine level falls, the internal connections are released from suppression, and there is almost no input from the senses. The network, in read mode with no cue, replays what is stored in it. Recordings show that neurons that worked together during the day fire together again in slow-wave sleep~\cite{WilsonMcNaughton1994}, and even in the same order~\cite{Skaggs1996}. According to a hypothesis~\cite{Hasselmo1999}, these replays copy what was quickly learned during the day into long-term memory (artificially lengthening the brief bursts of replay improves memory~\cite{FernandezRuiz2019}): by day, writing from the input; by night, rewriting from within. For an engineer this resembles writing to a fast buffer during the day and flushing it to slow permanent storage at night.

\section{Summary}

\begin{table}[t]
\centering
\footnotesize
\setlength{\tabcolsep}{4pt}
\renewcommand{\arraystretch}{1.15}
\begin{tabular}{>{\raggedright}p{3.2cm}>{\raggedright}p{4.2cm}>{\raggedright\arraybackslash}p{6.0cm}}
\toprule
What \nt{ACET} does & How & What is known \\
\midrule
weakens internal connections & factor $1-a$ in~\eqref{eq:acetnet} & measured in slices \\
strengthens the input & factor $\frac{1+a}{2}$ & increased excitability and transmission along input pathways are measured \\
facilitates learning & rate $\eta a$ & measured \\
switches ``write/read'' & Proposition~\ref{prop:writeread} & follows from models; no direct measurement of the modes \\
reports expected uncertainty & $P$ in~\eqref{eq:kalman} & hypothesis \\
frees the network for replay in sleep & low $a$ in slow-wave sleep & the low level in sleep and the replays are measured; rewriting into long-term memory is a hypothesis \\
\bottomrule
\end{tabular}
\caption{What \nt{ACET} adds to a network of \nt{GLUT}, \nt{GABA}, \nt{DOPA} and \nt{NORA}.}
\label{tab:acet}
\end{table}

\nt{DOPA} tells the network where to change the weights, \nt{NORA} how decisively to use what it knows, and \nt{ACET} whether to listen to the world or to itself (Table~\ref{tab:acet}). This is the last broadcast control line of Table~\ref{tab:types}: write enable, without which a memory that stores patterns in the same connections it reads them with would corrupt itself on every read.

\section{Exercises}

\begin{exercise}[\lvlA]
\label{ex:09:1}
\emph{How long to remember.} The filter~\eqref{eq:kalman} with $q=0.01$, $R=1$ remembers for about $10$~s. Find $P_\infty$ and the memory $\sqrt{R/q}$ if (a)~the sensor noise doubles in amplitude; (b)~the world starts drifting a hundred times faster. (c)~By what factor must the noise grow for the network to remember for a minute?
\end{exercise}

\begin{exercise}[\lvlB]
\label{ex:09:2}
\emph{Write mode after a change.} The world has changed, and the uncertainty of the filter~\eqref{eq:kalman} with $q=0.01$, $R=1$ has jumped to $P_0\to\infty$. (a)~With what time constant does $P$ return to $P_\infty$? (b)~After how many seconds does the learning rate exceed its steady-state value by no more than $10\,\%$?
\end{exercise}

\begin{exercise}[\lvlB]
\label{ex:09:3}
\emph{A constant learning rate.} The network learns as $\dd\hat s/\dd t=(y-\hat s)/T$; the world drifts, $\dd s/\dd t=w$, $y=s+n$, where $w$ and $n$ are white noises of intensities $q$ and $R$. Find the best $T$ and the error variance at it. By what factor does the variance grow if $T$ is off by a factor of $2$ and of $10$?
\end{exercise}

\begin{exercise}[\lvlB]
\label{ex:09:4}
\emph{Where the write boundary lies.} In \texttt{models/\allowbreak acet\_memory.py} the threshold is $\theta=0.35$ and the weight within a pattern is $w=0.041$ (ideally $0.045$). A new pattern shares $m=10$ neurons with an old one. For what $a$ in~\eqref{eq:acetnet} do the remaining neurons of the old pattern not light up from these $m$ (sharp threshold)?
\end{exercise}

\begin{exercise}[\lvlB]
\label{ex:09:5}
\emph{Simulation: memory capacity.} Modify the function \texttt{read\_sweep} in \texttt{models/\allowbreak acet\_memory.py}: at $a=1$ write $M$ patterns ($M$ from $5$ to $100$), then at $a=0$ read the first ten from halves. At what $M$ do errors appear, and how does the limit $M_{\max}$ depend on $N$ and $p$?
\end{exercise}

\begin{exercise}[\lvlB]
\label{ex:09:6}
\emph{Design: writing without leakage.} With learning rate $\eta a$, the network also learns while reading. Propose a monotonic $\eta(a)\le\eta$ that is zero for $a\le0.3$ and no worse than $\eta a$ for $a\ge0.9$. Check: after $20$ reads at $a=0.2$ with a cue containing three foreign neurons, how many of them have entered the pattern?
\end{exercise}

\begin{exercise}[\lvlC]
\label{ex:09:7}
\emph{How to rewrite the buffer at night.} (a)~In sleep $a=0.05$ and a replay lasts $0.1\,\text{s}$, versus $0.2\,\text{s}$ by day at $a=1$. How many replays give the weight change of one daytime presentation at rate $\eta a$ and at $\eta(a)$ of Exercise~\ref{ex:09:6}? (b)~The store learns $100$ times slower than the buffer and hears the pattern $20$ times a night, $0.1\,\text{s}$ each. How many nights are needed?
\end{exercise}

\chapter{The Whole Machine}
\chaptermark{The Whole Machine}
\label{sec:synthesis}

\section{What we have assembled}

We took the neuron types of Table~\ref{tab:types} one at a time and asked what each adds to a computing machine. The rules were the same every time: only what exists in a living organism, and everything runs in continuous time. Now we put the parts together into one machine and watch how they work jointly.

The picture from Chapter~\ref{sec:neurontypes} has held up and become sharper. A huge addressed network of \nt{GLUT} neurons carries the data. \nt{GABA} neurons control the flow of that data. Above the network sit four broadcast channels, each with its own knob: \nt{DOPA} says which weight changes to keep, \nt{SERO} holds the overall level and, by hypothesis, the planning horizon, \nt{NORA} chooses between searching and deciding, and \nt{ACET} chooses between writing and reading (Fig.~\ref{fig:machine}).

\begin{figure}[t]
\centering
\begin{tikzpicture}[>=Stealth,
  blk/.style={draw,very thick,rounded corners=3pt,align=center,font=\small},
  reg/.style={draw,very thick,double,double distance=1.2pt,circle,minimum size=1.25cm,inner sep=0pt,font=\scriptsize,fill=orange!15},
  bc/.style={->,thick,densely dashed,orange!85!black}]
% sensors, bus, actions
\node[blk,fill=green!8,text width=1.7cm] (S) at (0.9,0) {sensors\\\scriptsize on/off};
\draw[very thick,rounded corners=4pt,fill=blue!5] (2.6,-1.35) rectangle (9.0,1.35);
\node[font=\small] at (5.8,0.95) {\nt{GLUT} bus: data};
\node[font=\scriptsize,align=center] at (4.3,0.25) {coincidences, delays,\\logic (Ch.~\ref{sec:glutcomputer})};
\node[font=\scriptsize,align=center] at (7.4,0.25) {memory, code\\(Ch.~\ref{sec:code}, \ref{sec:acet})};
\draw[thick,red!70!black,fill=red!8,rounded corners=2pt] (2.8,-1.15) rectangle (8.8,-0.55);
\node[font=\scriptsize,red!60!black] at (5.8,-0.85) {\nt{GABA} (Ch.~\ref{sec:gaba}): veto, choice, division, rhythm};
\node[blk,fill=blue!5,text width=2.0cm] (A) at (10.9,0) {action\\selection};
\draw[->,very thick] (S) -- (2.6,0);
\draw[->,very thick] (9.0,0) -- (A);
\draw[->,very thick] (A) -- (12.6,0) node[right,font=\small] {muscles};
\pic[scale=1.1] at (13.2,-0.5) {spring};
\pic[scale=1.15] at (0.4,0.85) {eyepic};
\pic[scale=1.05] at (1.35,0.85) {speaker};
% registers
\node[reg] (D) at (3.4,3.2) {\nt{DOPA}};
\node[reg] (C) at (5.6,3.2) {\nt{SERO}};
\node[reg] (N) at (7.8,3.2) {\nt{NORA}};
\node[reg] (H) at (10.0,3.2) {\nt{ACET}};
\node[font=\scriptsize,align=center,above=1pt] at (D.north) {error $\delta$};
\node[font=\scriptsize,align=center,above=1pt] at (C.north) {level, $\tau_\gamma$};
\node[font=\scriptsize,align=center,above=1pt] at (N.north) {gain $g$};
\node[font=\scriptsize,align=center,above=1pt] at (H.north) {write $a$};
\draw[bc] (D.south) -- (3.4,1.35);
\draw[bc] (C.south) -- (5.6,1.35);
\draw[bc] (N.south) -- (7.8,1.35);
\draw[bc] (H.south) -- (10.0,2.1) -- (8.8,1.35);
\draw[bc] (H.south east) to[bend left=20] (A.north east);
\draw[->,thick] (2.85,1.35) -- (2.85,2.75) -- (D.west) node[pos=0.25,left,font=\scriptsize] {$V$};
\draw[->,thick,green!45!black] (0.4,3.2) node[left,black,font=\small] {reward $r$} -- (2.2,3.2) -- (D.west);
\pic[scale=0.9] at (-0.45,3.7) {juice};
\draw[->,thick,densely dotted,gray!50!black] ([yshift=0.45cm]D.north) to[bend left=18] node[above,font=\scriptsize,black] {surprise} ([yshift=0.45cm]N.north);
\end{tikzpicture}
\caption{The whole machine. The \nt{GLUT} bus carries data from the sensors to action selection; \nt{GABA} controls the flow inside it. Double circles are broadcast channels; dashed arrows send their signals to the whole network, each with its own knob. \nt{DOPA} receives the reward $r$ and the expectation $V$ from a critic inside the bus (Ch.~\ref{sec:dofa}); the dotted arrow is the hypothesis that \nt{NORA} learns of surprise from unusually large errors (Ch.~\ref{sec:nora}).}
\label{fig:machine}
\end{figure}

\section{One formula for all the knobs}

The knobs can be gathered into one schematic formula. This is not a new model but a summary of the models of Chapters~\ref{sec:glutcomputer}--\ref{sec:acet}: each symbol comes from its own chapter.
\begin{empheq}[box=\keybox]{equation}
\begin{aligned}
\tau\,\frac{\dd r_i}{\dd t} \;&=\; -r_i + F\Bigl(g\,\Bigl[\tfrac{1+a}{2}\,x_i + (1-a)\sum_j W_{ij}\,r_j - \sum_k M_{ik}\,q_k - \theta\Bigr]\Bigr),\\
\frac{\dd W_{ij}}{\dd t} \;&=\; \eta\,a\,\delta(t)\,e_{ij}(t),
\qquad \tau_e\,\frac{\dd e_{ij}}{\dd t} = -e_{ij} + r_i\,r_j,
\qquad \delta = r + \frac{\dd V}{\dd t} - \frac{V}{\tau_\gamma} .
\end{aligned}
\label{eq:machine}
\end{empheq}
Here $r_i$ are the rates of the \nt{GLUT} neurons, $x_i$ is the input from the sensors, and $W_{ij}\ge0$ are their mutual weights (Chapter~\ref{sec:glutcomputer}); $q_k$ are the rates of the \nt{GABA} neurons and $M_{ik}\ge0$ the strength of their inhibition (Chapter~\ref{sec:gaba}); $e_{ij}$ is the eligibility trace and $\delta$ the \nt{DOPA} error (Chapter~\ref{sec:dofa}); $\tau_\gamma$ is the horizon, set by \nt{SERO} according to the hypothesis; $g$ is the \nt{NORA} gain (Chapter~\ref{sec:nora}); $a$ is the \nt{ACET} level (Chapter~\ref{sec:acet}).

Look at what is missing from~\eqref{eq:machine}. There is no clock tick: everything is written as differential equations, and each neuron changes on its own. There is no separate memory: the same $W_{ij}$ that carry the computation do the remembering, and so does the same activity $r_i$. There are no individual corrections for the vast majority of synapses: their learning is the product of a local trace and one shared number; a separate error wire to each output exists only in the circuit that learns movements (Chapter~\ref{sec:motorlearning}). And there are only four kinds of control signal, $\delta$, $\tau_\gamma$, $g$, and $a$, for eighty billion neurons. Each of them is really not a single number but a small bundle of parallel lines (Proposition~\ref{prop:bundle}), yet a bundle holds only a handful of lines.

\begin{keyblock}
\begin{proposition}[Architecture]
\label{prop:architecture}
The brain's machine, as far as we understand it today, can be described in three layers. Data travel along the \nt{GLUT} address bus. Addressed \nt{GABA} neurons steer, limit, and normalize the data flow. The operating mode is set by a few broadcast channels, each a small bundle of parallel lines shared by large parts of the network. The first two layers are firmly established; the division of roles among the broadcast channels is mostly hypothesis.
\end{proposition}
\end{keyblock}

\section{All the types in one table}

Table~\ref{tab:synthesis} collects all the neuron types of the book. The four types that did not get a chapter of their own each take one row in it; two of them found a role in the chapters on the machine's outputs: \nt{GLYC} in the loops of the spinal cord (Chapter~\ref{sec:muscles}) and \nt{ADRE} in the chemical output (Chapter~\ref{sec:chemoutput}). The role of the other two in computation is either simple or unknown. Substances that do not define a neuron type are collected in Appendix~\ref{sec:othermediators}.

\begin{table}[t]
\centering
\footnotesize
\setlength{\tabcolsep}{3pt}
\renewcommand{\arraystretch}{1.15}
\begin{tabular}{>{\raggedright}p{1.3cm}>{\raggedright}p{1.0cm}>{\raggedright}p{3.9cm}>{\raggedright}p{1.5cm}>{\raggedright}p{1.8cm}>{\raggedright\arraybackslash}p{3.8cm}}
\toprule
Type & Ch. & What it computes & Time & Bus & What is known \\
\midrule
\nt{GLUT} & \ref{sec:glutcomputer}, \ref{sec:code}, \ref{sec:sensors}, \ref{sec:motorlearning}, \ref{sec:dendrites} & data: coincidences, delays, logic, memory, sequences & ms & address & established \\
\nt{GABA} & \ref{sec:gaba}, \ref{sec:code}, \ref{sec:pain}, \ref{sec:motorlearning} & flow control: veto, selection, division, stability, rhythm; the pain gate & ms & address & established; division of rate only together with background noise \\
\nt{SERO} & \ref{sec:serocomputer}, \ref{sec:pain}, \ref{sec:sleep} & level regulator, smoothing; horizon $\tau_\gamma$; pain volume; sleep modes & seconds & volume & self-inhibition measured; horizon is a hypothesis \\
\nt{DOPA} & \ref{sec:dofa}, \ref{sec:dofaalt} & prediction error $\delta$; slow level is the price of time & $0.1$\,s and minutes & broadcast & $\delta$ measured; extensions in Ch.~\ref{sec:dofaalt} \\
\nt{NORA} & \ref{sec:nora}, \ref{sec:pain}, \ref{sec:chemoutput}, \ref{sec:sleep} & gain $g$: search or decide; surprise; ``accelerator'' for the organs & seconds & broadcast & rise in signal-to-noise ratio measured; role in search is a hypothesis \\
\nt{ACET} & \ref{sec:acet}, \ref{sec:muscles}, \ref{sec:chemoutput}, \ref{sec:sleep} & write or read; learning rate; output to muscle; ``brake'' for the organs & seconds & both & three effects measured; mode switching is a hypothesis \\
\nt{GLYC} & \ref{sec:muscles}, \ref{sec:pain} & negative weight in the spinal cord and brainstem: veto on the opposing muscle & ms & address & established \\
\nt{HIST} & --- & global ``stay awake'' signal & slow & broadcast & role in computation not studied \\
\nt{ADRE} & \ref{sec:chemoutput} & ``accelerator'' for the whole body through the blood; unknown in the brain & slow & broadcast & established outside the brain; role in the brain unclear \\
\nt{ASPA} & --- & weak positive weight & ms & address & role disputed \\
\bottomrule
\end{tabular}
\caption{All the neuron types of the book: what each computes and how well that is known.}
\label{tab:synthesis}
\end{table}

\section{How the knobs work together}

So far each chapter turned one knob. Let us now follow one event through the whole machine (Fig.~\ref{fig:timeline}). The animal learned long ago that action $A$ brings juice. At some moment the world changes: $A$ no longer brings juice, and it is action $B$, which the animal hardly ever tries, that does.

\emph{Error.} No juice arrives after $A$, and the \nt{DOPA} neurons respond with a dip, $\delta<0$ by the error formula~\eqref{eq:ctd}. The synapses that just chose $A$ carry a trace and weaken.

\emph{Surprise.} A single dip is ordinary chance. But the dips repeat, and the errors grow larger than usual. By the hypothesis of Chapter~\ref{sec:nora}, this is exactly the \nt{NORA} signal: the world has changed. The gain $g$ drops, the choice becomes soft, and noise leads to a trial of $B$ more often.

\emph{Writing.} By the hypothesis of Chapter~\ref{sec:acet}, after the change \nt{ACET} rises and puts the network into write mode: the internal connections that store the old habit are weakened, the input from the sensors is amplified, and the weights change easily.

\emph{New habit.} Trying $B$ brings juice that nobody expected: a burst, $\delta>0$, and the synapses that chose $B$ strengthen. A few such trials, and the expectation $V$ catches up with the new world; the errors are small again.

\emph{Return.} The surprise is over, \nt{NORA} restores high gain, and the choice is hard again; \nt{ACET} falls, and the network, in read mode, uses what it has learned. All this time \nt{SERO}, by hypothesis, sets the horizon $\tau_\gamma$: how distant a juice is still worth waiting for.

\begin{figure}[t]
\centering
\begin{tikzpicture}[>=Stealth,x=1.1cm,y=0.95cm]
\foreach \y/\lab in {4/{$\delta$ (\nt{DOPA})},3/{$g$ (\nt{NORA})},2/{$a$ (\nt{ACET})},1/{choice of $B$}}{
  \draw[gray!60!black] (0,\y-0.35) -- (10,\y-0.35);
  \node[left,font=\scriptsize] at (0,\y) {\lab};}
\draw[densely dashed,gray!60!black] (2,0.4) -- (2,4.55) node[above,font=\scriptsize,black] {the world changed};
\draw[densely dashed,gray!60!black] (5,0.4) -- (5,4.55) node[above,font=\scriptsize,black] {first juice for $B$};
\pic[scale=0.85] at (2,5.25) {nojuice};
\pic[scale=0.85] at (5,5.25) {juice};
% delta: dips after the change, burst at the first juice for B, then back to zero
\draw[thick,red!70!black] (0,3.65) -- (2.3,3.65) -- (2.3,3.3) -- (2.5,3.3) -- (2.5,3.65) -- (3.2,3.65) -- (3.2,3.3) -- (3.4,3.3) -- (3.4,3.65) -- (4.1,3.65) -- (4.1,3.35) -- (4.3,3.35) -- (4.3,3.65) -- (5.0,3.65) -- (5.0,4.35) -- (5.2,4.35) -- (5.2,3.65) -- (5.8,3.65) -- (5.8,4.1) -- (6.0,4.1) -- (6.0,3.65) -- (6.6,3.65) -- (6.6,3.85) -- (6.8,3.85) -- (6.8,3.65) -- (10,3.65);
% gain: high, drops after surprise, recovers
\draw[thick,blue!60!black] (0,3.2) -- (3.0,3.2) -- (3.4,2.75) -- (5.8,2.75) -- (7.2,3.2) -- (10,3.2);
% ACh: low, rises after the change, falls after learning
\draw[thick,green!45!black] (0,1.75) -- (2.8,1.75) -- (3.3,2.25) -- (6.8,2.25) -- (7.8,1.75) -- (10,1.75);
% choice of B
\draw[thick,black] (0,0.7) -- (3.0,0.7) plot[domain=3:10,samples=40] (\x,{0.7+0.6/(1+exp(-(\x-5.3)*1.8))});
\draw[->,gray!70!black] (0,0.2) -- (10.2,0.2) node[below left,font=\scriptsize,black] {time};
\end{tikzpicture}
\caption{One event followed through the whole machine. This is a qualitative sketch based on the models and hypotheses of Chapters~\ref{sec:dofa}--\ref{sec:acet}, not the result of a calculation. After the change \nt{DOPA} responds with dips; repeated dips, by hypothesis, raise the surprise signal, \nt{NORA} lowers the gain, and \nt{ACET} turns on writing; the first juice for $B$ gives a burst, the choice moves to $B$, and the knobs return.}
\label{fig:timeline}
\end{figure}

Notice that no knob knows what the others are doing. Each responds to its own cues (the error, its unusual size, uncertainty), and the order of operations arises by itself, as in an asynchronous circuit where each block fires when its inputs are ready. As far as we know, this coordinated operation as a whole has not yet been tested: the knobs have been measured one by one, but their joint operation is a hypothesis.

\section{How this machine differs from a computer}

\begin{table}[t]
\centering
\footnotesize
\setlength{\tabcolsep}{4pt}
\renewcommand{\arraystretch}{1.15}
\begin{tabular}{>{\raggedright}p{2.2cm}>{\raggedright}p{4.2cm}>{\raggedright\arraybackslash}p{6.6cm}}
\toprule
& Ordinary computer & Brain \\
\midrule
time & a shared clock of about a gigahertz & no clock; each neuron changes on its own, windows are set by events and local rhythms (Ch.~\ref{sec:glutcomputer}, \ref{sec:gaba}, \ref{sec:code}) \\
elements & reliable binary gates & analog threshold elements; a synapse loses most spikes (Ch.~\ref{sec:code}) \\
sign of a connection & any & set by the sending neuron (Ch.~\ref{sec:neurontypes}) \\
memory & separate from the processor & in the same connections as the computation, and in self-sustaining activity (Ch.~\ref{sec:glutcomputer}, \ref{sec:acet}) \\
learning & error backpropagation & local rules, one broadcast number $\delta$ (Ch.~\ref{sec:dofa}), and error wires to the outputs of the movement circuit (Ch.~\ref{sec:motorlearning}) \\
control & registers and a command bus & four broadcast channels, each a bundle of a few lines (Ch.~\ref{sec:serocomputer}, \ref{sec:dofa}--\ref{sec:acet}) \\
energy & tens to hundreds of watts per processor & about $20$ watts for the whole brain, on average about one spike per second per neuron (Ch.~\ref{sec:code}) \\
\bottomrule
\end{tabular}
\caption{The brain and an ordinary computer.}
\label{tab:computer}
\end{table}

Table~\ref{tab:computer} summarizes the main differences. Each of them is not a flaw that nature did not get around to fixing but part of a single design. Without a shared clock you need on/off sensors and coincidence detectors. Unreliable elements need the redundancy of many neurons. Memory combined with computation needs \nt{ACET} so that writing does not spoil reading. Learning without backward wires needs \nt{DOPA} and noise. And the energy limit of one spike per second forces the whole language to be frugal and to spend spikes on news.

\section{What we do not know}

The most honest way to end this summary is with a list of what we do not know. Each item is a problem for an engineer.

\emph{The code.} We know the capacity of the spike channel, and we know that the recipient chooses which part of the message to read (Chapter~\ref{sec:code}). But how much the cortex uses the precise timing of spikes, and whether neurons have ``words,'' is unknown.

\emph{Learning through many layers.} A broadcast signal teaches slowly, and more slowly with each neuron that affects the result (Proposition~\ref{prop:broadcastcost}). How the cortex then tunes billions of weights in multilayer circuits is unknown; there are models in which bursts of spikes carry the error between layers~\cite{Payeur2021}. Part of the answer may lie in computation inside the branches of the neuron itself, where a single neuron behaves like a small multilayer network: to reproduce the response of a model of one cortical neuron, an artificial network needs five to eight layers~\cite{Beniaguev2021}, and the branches of human neurons can even compute exclusive OR~\cite{Gidon2020}. Another candidate is self-supervised learning: if the cortex learns to predict its own inputs, the error arises on the spot, in every layer, and no shared signal is needed for it (Chapter~\ref{sec:dofa})~\cite{Keller2018}.

\emph{What the broadcast channels really carry.} For \nt{DOPA} there is a standard picture and six extensions of it (Chapter~\ref{sec:dofaalt}); for \nt{NORA} and \nt{ACET} there are plausible hypotheses (Chapters~\ref{sec:nora}, \ref{sec:acet}); for \nt{SERO} there are mostly guesses.

\emph{Coordination in time.} We have seen how to compute a function, measure time from an event, and slice the stream into windows with a local rhythm. But how a huge network with no shared clock coordinates the work of distant areas is understood only in part.

\emph{Energy.} Twenty watts for everything. We understand why the code must be sparse (Proposition~\ref{prop:economy}), but nobody yet knows how to build a machine that would do the same thing at the same power~\cite{MehonicKenyon2022}.

The brain is a computing machine built not by an engineer, yet by laws any engineer will recognize: RC circuits, thresholds, feedback, filters, buses, and broadcast registers. We have read its schematic only in part. It will be read to the end by those who know how to read schematics.

\section{What comes next in the book}

This chapter assembled the computational core of the machine. The remaining chapters go beyond it. Chapters~\ref{sec:sensors} and~\ref{sec:recognition} are about the inputs: how sensors turn light, sound, and molecules into spikes, and how the machine recognizes objects in them. Chapters~\ref{sec:muscles}--\ref{sec:chemoutput} are about the outputs and what serves them: muscles, pain as an alarm signal, learning movements through dedicated error wires, and the slow chemical output to the body. Chapter~\ref{sec:debugging} is about how the machine is debugged from outside, including through brain--computer interfaces. Chapters~\ref{sec:eetypes} and~\ref{sec:dendrites} sort the neurons once more, now by electrical function and by the number of inputs. Chapter~\ref{sec:sleep} is about what the machine does when it is disconnected from the world, and Chapters~\ref{sec:livingmachine} and~\ref{sec:dishbrain} are about machines built from living neurons on a chip.

\section{Exercises}

The exercises of this chapter connect results from different chapters: each relies on at least two of them.

\begin{exercise}[\lvlA]
\label{ex:10:1}
\emph{Bus and registers.} The four broadcast channels are bundles of about five lines each (Proposition~\ref{prop:bundle}); each line changes once per second and has four distinguishable levels. How many years would the control channels need to transmit what the \nt{GLUT} bus ($8.6\cdot10^{10}$ neurons at the rate~\eqref{eq:energy}, $1\ms$ precision by~\eqref{eq:spikeentropy}) transmits in one second?
\end{exercise}

\begin{exercise}[\lvlB]
\label{ex:10:2}
\emph{Two knobs on one loop.} In the loop~\eqref{eq:ei} the knobs of~\eqref{eq:machine} act on excitation: $\tau_E\dot E=-E+g[(1-a)w_{EE}E-w_{EI}I+h]$, and they do not control \nt{GABA}. With the numbers of Fig.~\ref{fig:ei}, a \nt{NORA} burst raises $g$ to $6$. What is the smallest $a$ for which the loop is stable? Can the \nt{NORA} and \nt{ACET} knobs be turned independently?
\end{exercise}

\begin{exercise}[\lvlB]
\label{ex:10:3}
\emph{Where the cost of broadcasting comes from.} The outputs of $N$ neurons are $y_k=f(u_k)+\xi_k$ with independent Gaussian noise of variance $\sigma^2$; $R\approx R_0+\sum_kR'_k\xi_k$ with equal $R'_k$, and \nt{DOPA} broadcasts $\delta=R-R_0$. Find the signal-to-noise ratio of the correction $\delta\,\xi_jx_j$ in one trial. How does the number of trials grow with $N$?
\end{exercise}

\begin{exercise}[\lvlB]
\label{ex:10:4}
\emph{Binding a sound and a flash.} A sound reaches the bus after $5\ms$, a flash after $30\ms$ plus the first-spike latency~\eqref{eq:latency} ($\tau=10\ms$, $U$ from $1.2\theta$ to $4\theta$); the background on each input is $5$ spikes per second. Choose the delay of the sound line and the smallest coincidence window for all contrasts. How many false bindings per second will the background produce?
\end{exercise}

\begin{exercise}[\lvlB]
\label{ex:10:5}
\emph{Simulation: who notices the change.} A critic sees $y=s+\text{noise}$ ($R=1$); with probability $1/200$, $s$ jumps to a new value with variance $J^2=4$. Simulate a Kalman filter with $q=0$ that raises $P$ to $J^2$ when $E\leftarrow0.9E+\delta^2/(P+R)$ exceeds $20$. How much smaller is its error than with the best constant $\alpha$?
\end{exercise}

\begin{exercise}[\lvlA]
\label{ex:10:6}
\emph{How many trials to change a habit.} A network has learned $Q_A=0.8$, $Q_B=0.2$ and chooses by~\eqref{eq:softmax}, $\sigma=1$, $g=8$. Now $A$ gives juice with probability $0.2$; $Q_A\leftarrow Q_A+\alpha(r-Q_A)$, and $Q_B$ does not change. After how many choices of $A$ does the probability of choosing $A$ fall to $0.6$ for $\alpha=0.1$ and $0.5$? And if \nt{NORA} lowers $g$ to $2$?
\end{exercise}

\begin{exercise}[\lvlC]
\label{ex:10:7}
\emph{A bundle instead of a wire.} Outputs $y_k=w_k+\xi_k$, reward $R=-\sum_ky_k^2$; a line of the bundle sends a group of $G$ neurons the error of its outputs, $w_k\leftarrow w_k+\eta\,\delta_l\xi_k$. Show that with the best $\eta$, reaching $\sum w_k^2=\varepsilon\sum w_k^2(0)$ takes $\approx(G+2)\ln(1/\varepsilon)$ trials. For $\varepsilon=0.01$, how long does a circuit of $10^6$ neurons on five lines take to learn, with one trial every three seconds?
\end{exercise}

\chapter{The Machine's Inputs: How the Eyes, Ears, and Other Sensors Connect to It}
\chaptermark{The Machine's Inputs}
\label{sec:sensors}

\section{A sensor as a transducer}

In Fig.~\ref{fig:machine} data entered the machine from the left, from the ``sensors'' block. Now we open that block. To an engineer every sense organ is a \emph{transducer} with an analog front end and an analog-to-digital converter at the end, except that the digital output is not numbers but spikes.

All sense organs share the same scheme (Fig.~\ref{fig:transducer}). A physical quantity (light, pressure, vibration, a molecule) acts on a receptor protein in the membrane of a sensory cell. This protein works as a gate: by itself or through a short chain of reactions, it opens an ion channel. A \emph{receptor current} arises, and with it an analog voltage across the membrane. The voltage then passes through gain control and enters a spike encoder, the familiar integrating neuron~\eqref{eq:glutneuron}, which turns a level into a rate and spike times. The output is almost always the same: the sensory neurons of the eye, the ear, smell, and skin release glutamate, so the inputs connect directly to the \nt{GLUT} bus.

\begin{figure}[t]
\centering
\begin{tikzpicture}[>=Stealth,
  blk/.style={draw,very thick,rounded corners=2pt,align=center,font=\footnotesize,minimum height=1.0cm,text width=1.9cm,fill=blue!5}]
\node[align=center,font=\small] (P) at (0.1,0) {light, sound,\\pressure,\\molecule};
\node[blk] (R) at (3.0,0) {receptor\\gate};
\node[blk] (G) at (6.5,0) {gain\\control};
\node[blk] (E) at (10.0,0) {spike\\encoder};
\node[align=center,font=\small] (B) at (13.4,0) {\nt{GLUT}\\bus};
\draw[->,very thick] (P) -- (R);
\foreach \ic/\xx in {sun/-1.1,speaker/-0.3,press/0.5,molecule/1.3}{\pic[scale=0.85] at (\xx,1.05) {\ic};}
\draw[->,very thick] (R) -- node[above,font=\scriptsize] {current} (G);
\draw[->,very thick] (G) -- node[above,font=\scriptsize] {level} (E);
\draw[->,very thick,red!70!black] (E) -- node[above,font=\scriptsize,black] {spikes} (B);
\draw[decorate,decoration={brace,amplitude=5pt,mirror},gray!60!black] (1.8,-0.8) -- (7.7,-0.8) node[midway,below=5pt,font=\scriptsize,black] {analog part};
\draw[decorate,decoration={brace,amplitude=5pt,mirror},gray!60!black] (8.8,-0.8) -- (11.2,-0.8) node[midway,below=5pt,font=\scriptsize,black] {spikes};
\end{tikzpicture}
\caption{Any sense organ as a chain of conversions. The receptor protein opens a channel and produces a current; gain control adapts it to the conditions; the encoder~\eqref{eq:glutneuron} turns the level into spikes, which go out onto the \nt{GLUT} bus. In the eye and the ear the first stages produce no spikes at all: the signal stays analog until it has to be sent far.}
\label{fig:transducer}
\end{figure}

One detail is very familiar to an electronics engineer. In the eye and the ear the very first cells produce no spikes at all: they pass a smooth voltage to their neighbors, like an analog stage on a circuit board. Spikes appear only where the signal must travel far. Over millimeters an analog signal decays and picks up noise, while a spike, as we saw in Chapter~\ref{sec:neurontypes}, reaches the end of the wire at full height. Digitization happens where the long line begins.

\section{At the limit of physics}

The brain's sensors work at the physical limits of sensitivity. A light sensor in the dark responds to a \emph{single photon}: absorbing one particle of light gives a current of about a picoampere, visible above the sensor's own noise~\cite{Baylor1979}. A sound sensor detects a displacement of its sensitive bundle of less than a nanometer~\cite{Hudspeth1989}.

When light is scarce, precision is limited not by the sensor but by the light itself: photons arrive at random, as a Poisson stream. If on average $N=\Phi T$ photons reach the sensor during the integration time $T$, their count fluctuates by $\sqrt N$. For an increment $\Delta N$ to be noticeable, it must be several times larger than this fluctuation, and the detectable fractional change in brightness is
\begin{empheq}[box=\keybox]{equation}
\frac{\Delta I}{I} \;\approx\; \frac{k}{\sqrt{\Phi\,T}} .
\label{eq:photonnoise}
\end{empheq}
This is the same formula as for the spike count~\eqref{eq:rateerror}: the shot noise of photons and the shot noise of spikes obey one law. In the dark the eye has only one way to improve precision, which is to collect photons longer, and the integration time does indeed grow to tenths of a second. That is why we see more slowly at dusk.

\section{Dynamic range: gain control}

The hardest engineering problem for sensors is range. Illuminance varies by about ten orders of magnitude from a starlit night to sunlit snow, and sound intensity by twelve from the threshold of hearing to the threshold of pain. Yet the spike rate varies from zero to a few hundred per second, less than three orders of magnitude. Such an input cannot be mapped linearly onto such an output.

The solution is the same as in any receiver: \emph{automatic gain control}. The responses of the sensory cells of the eye are well described by the formula
\begin{empheq}[box=\keybox]{equation}
r \;=\; r_{\max}\,\frac{I}{I+\sigma} ,
\label{eq:nakarushton}
\end{empheq}
where $\sigma$ is the brightness at which the response reaches half its maximum~\cite{NakaRushton1966}. By itself,~\eqref{eq:nakarushton} covers only two or three orders of magnitude. But $\sigma$ is not constant: during prolonged work against a bright background it grows along with the mean brightness. The response curve shifts along the logarithmic brightness axis and each time places its steep part where the input currently is (Fig.~\ref{fig:adaptation}). This is the same division by total activity as the shunting inhibition of Chapter~\ref{sec:gaba}~\cite{Carandini2012}, except that the denominator here is the mean brightness over the last few seconds.

\begin{figure}[t]
\centering
\begin{tikzpicture}[>=Stealth,x=1.25cm,y=2.8cm]
\draw[->,gray!70!black] (-2,0) -- (6.4,0) node[below left,font=\small,black] {$\log_{10} I$};
\draw[->,gray!70!black] (-2,0) -- (-2,1.15) node[left,font=\small,black] {$r/r_{\max}$};
\foreach \u in {-2,0,2,4,6}{\draw[gray!60!black] (\u,-0.02) -- (\u,0.02); \node[below,font=\scriptsize] at (\u,0) {$\u$};}
\draw[densely dashed,gray!60!black] (-2,0.5) -- (6.2,0.5);
\foreach \s/\c/\lab in {0/blue!60!black/{dark},2/green!45!black/{medium},4/red!70!black/{bright}}{
  \pgfmathsetmacro{\lo}{max(-2,\s-3)}
  \draw[very thick,\c] (-2,0) -- (\lo,0) plot[domain=\lo:6.2,samples=80] (\x,{1/(1+10^(\s-\x))});
  \node[font=\scriptsize,\c,anchor=west] at (\s+0.15,0.42) {\lab};}
\foreach \s/\ic in {0/moon,2/cloud,4/sun}{\pic at (\s+0.6,0.64) {\ic};}
\end{tikzpicture}
\caption{Gain control of a light sensor by~\eqref{eq:nakarushton}. Each curve covers two or three orders of magnitude of brightness; on moving to a brighter background $\sigma$ grows, and the curve shifts along the logarithmic axis. Together the shifting curves cover the whole range.}
\label{fig:adaptation}
\end{figure}

Gain control has a price, familiar from Chapter~\ref{sec:code}: the sensor reports not brightness but brightness \emph{relative to the background}. A constant input gradually stops evoking a response, while every change evokes one. From the very start, the machine's inputs speak of changes, not of levels, and this is the same economy of spikes as in Proposition~\ref{prop:economy}~\cite{Barlow1961}. Hence also the on and off sensors of Chapter~\ref{sec:glutcomputer}~\cite{Schiller1992}: some report that it got brighter, others that it got darker.

Engineers have copied this trick in \emph{event cameras}. Such a camera does not take frames on a timer: each of its pixels, on its own and with no shared clock, emits a ``brighter'' or ``darker'' event when the logarithm of brightness has changed by a set fraction~\cite{Lichtsteiner2008}. This is exactly the two-wire on/off code of Chapter~\ref{sec:glutcomputer}, implemented in silicon, and it gives the camera a range and speed out of reach of an ordinary sensor array~\cite{Gallego2022}.

\section{The eye: compression down to a cable}

The eye has about $10^8$ light sensors~\cite{Curcio1990}, while the output neurons that carry the image to the brain number about a million. Before the signal leaves the eye, it is compressed about a hundredfold. The main compression tricks are familiar. Each output neuron responds to the difference between the brightness in a small spot and the brightness around it; this is subtraction of neighbors through inhibition, as in Chapter~\ref{sec:gaba}. Uniform regions of the picture cost almost nothing, while edges and changes are transmitted. Different output neurons are specialized: some report brightening, others darkening, still others motion in a particular direction (Fig.~\ref{fig:direction}); in the mouse more than thirty such output channels have been counted~\cite{Baden2016}.

What is the capacity of the resulting cable? In the guinea pig, recordings from output neurons gave about $875$ thousand bits per second for $\sim10^5$ such neurons; scaling to the human million output neurons gives on the order of $10^7$ bits per second~\cite{Koch2006}, as much as an old ten-megabit network cable. At a mean rate of a few spikes per second per neuron this is a few bits per spike, as we found in Chapter~\ref{sec:code}.

\section{The ear: a filter bank}

The ear solves a different problem: sound must be decomposed by frequency. An engineer would install a bank of bandpass filters, and the ear does exactly that, mechanically. The sound vibration travels along an elastic partition whose properties change smoothly along its length, and each section resonates at its own frequency. The sensors sitting on a section respond only to their own band (Fig.~\ref{fig:filterbank}). Frequencies are laid out by place roughly logarithmically: each octave gets a similar share of the sensors. The index of the sensor that fired is the frequency, the place code of Chapter~\ref{sec:code}.

\begin{figure}[t]
\centering
\begin{tikzpicture}[>=Stealth,x=1.35cm,y=2.2cm]
\draw[->,gray!70!black] (-0.3,0) -- (7.6,0) node[above left,font=\small,black] {frequency, Hz};
\draw[->,gray!70!black] (-0.3,0) -- (-0.3,1.15) node[left,font=\small,black] {response};
\foreach \k/\f in {0/125,1/250,2/500,3/1000,4/2000,5/4000,6/8000}{
  \draw[gray!60!black] (\k,-0.02) -- (\k,0.02); \node[below,font=\scriptsize] at (\k,0) {\f};
  \draw[thick,blue!60!black] plot[domain={max(\k-1.2,-0.3)}:{\k+0.7},samples=60] (\x,{1/(1+((\x-\k)/(\x<\k?0.45:0.2))^4)});}
% a piano keyboard: seven white keys per octave, like the octaves on the axis
\foreach \i in {0,...,48}{\draw[gray!50!black,fill=white,line width=0.3pt] ({\i/7},-0.44) rectangle ({(\i+1)/7},-0.26);}
\foreach \o in {0,...,6}{\foreach \j in {0,1,3,4,5}{\fill[black] ({\o+(\j+1)/7-0.035},-0.35) rectangle ({\o+(\j+1)/7+0.035},-0.26);}}
\end{tikzpicture}
\caption{The ear as a bank of bandpass filters, schematic. Each curve is the response of the sensors of one section to sounds of different frequency; the center frequencies are an octave apart, like keys on a logarithmic scale. Real filters have a steep high-frequency slope and a gentle low-frequency one. Below is a keyboard: on it, as in the ear, each octave gets the same amount of space.}
\label{fig:filterbank}
\end{figure}

The resonators of the ear are not passive. Some of the sensors themselves drive the partition in time with the sound and amplify weak vibrations thousands of times in amplitude (by $66$--$76$\,dB), and the gain falls as loudness grows~\cite{Ruggero1997,Robles2001}. To an engineer this is an amplifier with positive feedback, set right at the edge of self-oscillation: the same life on the edge as in the network of Chapter~\ref{sec:gaba}, where near the boundary the system is especially sensitive to small signals. Loudness compression is done by the same amplifier: quiet sounds are amplified strongly, loud ones weakly.

The ear also has a second code. At low frequencies neurons fire in phase with the sound: a neuron fires at a particular part of each wave, though not on every wave. In the guinea pig, phase locking begins to weaken at about $600$\,Hz and is still noticeable up to about $3.5$\,kHz~\cite{PalmerRussell1986}. In this way spike times preserve the precise timing of the sound, and from it the difference in arrival time at the two ears, from which the network of Chapter~\ref{sec:glutcomputer} finds the direction~\cite{Grothe2010}. At high frequencies the spikes cannot keep up with the wave, and only the place code remains. A single ear uses both codes of Chapter~\ref{sec:code}: timing where the neurons can keep up, and place where they cannot.

\section{The other sensors}

\begin{table}[t]
\centering
\footnotesize
\setlength{\tabcolsep}{3pt}
\renewcommand{\arraystretch}{1.15}
\begin{tabular}{>{\raggedright}p{1.9cm}>{\raggedright}p{2.6cm}>{\raggedright}p{4.0cm}>{\raggedright\arraybackslash}p{4.6cm}}
\toprule
Sense & What it measures & How the input is built & Trick from the book \\
\midrule
vision & light & $\sim10^8$ sensors, compression by $\sim100$, $\sim10^7$ bits/s & gain control, subtraction of neighbors, two wires, direction of motion \\
hearing & air pressure & mechanical filter bank, active amplifier & place code and timing code, delays \\
touch & pressure, vibration & fast- and slow-adapting sensors & a ``high-pass filter --- low-pass filter'' pair \\
temperature & heat & separate sensors for warm and cold & two-wire code \\
smell & molecules & $\sim400$ receptor types, one type per sensor & combinatorial code: an odor is the set of types that fired \\
taste & substances in food & five qualities: sweet, bitter, salty, sour, umami & labeled lines; some cells release ATP or serotonin rather than glutamate \\
body position & muscle stretch & sensors of stretch length and velocity & feedback for the movement controller \\
\bottomrule
\end{tabular}
\caption{How the sensors are connected. All the devices in the third column have been measured; the right column links them to tricks from earlier chapters.}
\label{tab:sensors}
\end{table}

The other senses are collected in Table~\ref{tab:sensors}, and almost every row shows a trick from earlier chapters. Touch uses a pair of filters: some pressure sensors respond to contact and fall silent at once, while others keep responding as long as the pressure lasts. Temperature is carried by two wires, separate for warm and cold. The most unusual input is smell. Humans have about four hundred types of olfactory receptor, and each olfactory neuron carries only one type~\cite{Mombaerts2004}. An odor activates not one type but a set, and the message is \emph{which} types fired. To an engineer this is a binary vector of length about four hundred, with an astronomical number of possible values.

\begin{keyblock}
\begin{proposition}[The machine's inputs]
\label{prop:sensors}
Each sense is a transducer that works at the physical limit of sensitivity, compresses a huge dynamic range by automatic gain control, and sends mainly \emph{changes} over the \nt{GLUT} bus. To do this, the sensors use the same tricks as the machine itself: the two-wire code, the place code, the timing code, and division by the background. The structure of the sensors has been measured; that their codes are tuned for the greatest information per spike is a well-founded hypothesis.
\end{proposition}
\end{keyblock}

The inputs, like everything else, work asynchronously. Each sensor emits a spike when it has something to report, and the different senses arrive with different delays: sound after milliseconds, light after tens of milliseconds. How the machine combines them into one picture of the world with no shared clock is one of the problems of coordination in time from Chapter~\ref{sec:synthesis}.

\section{Exercises}

\begin{exercise}[\lvlA]
\label{ex:11:1}
\emph{How long to collect photons.} At dusk $100$ photons per second reach a sensor; an increment is noticeable if it is $3$ times the fluctuation~\eqref{eq:photonnoise}. (a)~How long does one sensor need to notice a $10\,\%$ change in brightness? (b)~How many sensors must be summed to manage it within $0.2\,\text{s}$, and by what factor does the spatial resolution along one axis drop?
\end{exercise}

\begin{exercise}[\lvlA]
\label{ex:11:2}
\emph{How many bits per spike.} (a)~The eye transmits $\sim10^7$ bits/s over $10^6$ output neurons firing $3$--$5$ spikes per second. What fraction of the capacity~\eqref{eq:spikeentropy} at $\Delta t=1\ms$ does it use? (b)~An odor activates $20$ of $\approx400$ receptor types. How many bits does such a set carry, and how many types do two random odors share on average?
\end{exercise}

\begin{exercise}[\lvlB]
\label{ex:11:3}
\emph{What one curve~\eqref{eq:nakarushton} can do.} (a)~Over what range of brightness does the response rise from $10\,\%$ to $90\,\%$ of $r_{\max}$? How many orders of magnitude is that? (b)~How many positions of the shifting curve are needed to cover ten orders of magnitude of brightness? Twelve orders of magnitude of loudness?
\end{exercise}

\begin{exercise}[\lvlB]
\label{ex:11:4}
\emph{The limit of the timing code in the ear.} The difference in arrival time at the two ears is up to $0.22\,\text{m}/343\,\text{m/s}$. (a)~Spikes are phase-locked to a tone: up to what frequency do they determine the direction unambiguously? (b)~A spike jitters by $0.5\ms$, and $20$~\textmu s must be resolved. How many neurons at $100$ spikes per second must be averaged over $50\ms$?
\end{exercise}

\begin{exercise}[\lvlB]
\label{ex:11:5}
\emph{An event camera on the eye's cable.} A camera of $10^6$ pixels with an output of $10^7$ bits/s sends an event (address and sign) when $\ln I$ in a pixel has changed by $\ln1.2$. What is the highest speed at which an edge $500$ pixels long with a brightness step of $4$ can move? How many frames per second would an ordinary camera at $8$ bits per pixel get through the same output?
\end{exercise}

\begin{exercise}[\lvlB]
\label{ex:11:6}
\emph{Simulation: gain control.} The sensor~\eqref{eq:nakarushton} adjusts $\sigma$ with $\tau=1\,\text{s}$ either logarithmically, $\tau\,\dd\ln\sigma/\dd t=\ln I-\ln\sigma$, or linearly, $\tau\,\dd\sigma/\dd t=I-\sigma$; the background jumps $1\to100\to1$. For each law, how many seconds after the upward and the downward jump does the response to a $+10\,\%$ probe recover to half its adapted value?
\end{exercise}

\begin{exercise}[\lvlC]
\label{ex:11:7}
\emph{What world the curve is tuned for.} With small constant output noise, the information about brightness is greatest when $r(I)=r_{\max}\Phi_I(I)$, where $\Phi_I$ is the distribution function of brightness. (a)~For which distribution is the curve~\eqref{eq:nakarushton} optimal, and what is its spread in $\log_{10}I$? (b)~Which curve is optimal with Poisson output noise?
\end{exercise}

\chapter{Recognizing Images and Video}
\chaptermark{Recognition}
\label{sec:recognition}
\begin{chaptergoals}
\item why recognition must ignore shift, scale and lighting, and why template matching fails;
\item how a pipeline alternating AND over parts and OR over positions recognizes in $150\ms$;
\item how a Hebb rule with a trace learns invariance from unlabeled video;
\item how the brain differs from a convolutional network, and what illusions reveal about it.
\end{chaptergoals}

\section{The problem: the same thing, different pixels}

Chapter~\ref{sec:sensors} brought the image to the input of the machine: about a million output neurons of the eye, a compressed stream that carries mostly edges and changes. But between the sensor and the action there is a step we have not discussed yet. A cat in a picture is not a set of pixels. Shift it by half a frame, rotate it, move it farther away, turn off half the light: almost every pixel changes, yet the answer ``cat'' must stay the same. An engineer calls this \emph{invariance}: the output of a classifier must not change with shift, scale, rotation, or lighting.

A simple experiment shows how hard this is. Take a one-dimensional ``retina'' of $24$ points and two shapes of five points each that can sit anywhere. A classifier that compares the input with a template stored at one position is right $49$--$52\%$ of the time, as often as a coin. A shift breaks pixel-by-pixel comparison completely. We need a different circuit.

\section{A pipeline of stages}

We already know from Chapter~\ref{sec:code} how fast the brain does this: an animal in a briefly flashed picture is recognized in about $150\ms$, and the signal passes through about ten stages, $10$--$20\ms$ each~\cite{Thorpe1996}. At each stage a neuron has time to fire zero or one spike. So fast recognition is a single pass through a pipeline, with no long deliberation and no repeats. The question is what each stage does.

The answer suggested by measurements at the first stages of vision~\cite{HubelWiesel1962} consists of two alternating operations (Fig.~\ref{fig:conveyor}).

\emph{Selectivity.} A neuron of stage $S$ looks at a small patch of the input and fires only if a particular combination of parts is there: this edge, and that one next to it. This is AND over parts: the threshold neuron of Chapter~\ref{sec:glutcomputer}, whose threshold is higher than any single part can supply.

\emph{Invariance.} A neuron of stage $C$ collects the outputs of many $S$ neurons that look for the same combination at different places, and fires if it was found anywhere. This is OR over positions, or more precisely, taking the maximum.

For the one-dimensional model the two operations read
\begin{empheq}[box=\keybox]{equation}
s_{f,p} \;=\; \Bigl[\sum_i t_{f,i}\,x_{p+i} - \theta\Bigr]_+ ,
\qquad
c_f \;=\; \max_p\, s_{f,p} ,
\label{eq:scstage}
\end{empheq}
where $x$ is the input, $t_{f,i}$ is the template of feature $f$, $p$ is the position, and $\theta$ is the threshold. The first expression makes the response selective, the second makes it independent of position. The network gets the maximum of many inputs with the tools of Chapter~\ref{sec:gaba}: mutual inhibition keeps the strongest input (Proposition~\ref{prop:wta}), and division by the total activity equalizes responses across brightness levels~\cite{Carandini2012}.

\begin{figure}[t]
\centering
\begin{tikzpicture}[>=Stealth,
  px/.style={draw,minimum size=0.36cm,inner sep=0pt},
  on/.style={px,fill=blue!55!black},
  snode/.style={circle,draw,very thick,minimum size=0.62cm,inner sep=0pt,font=\scriptsize,fill=blue!6},
  cnode/.style={circle,draw,very thick,minimum size=0.8cm,inner sep=0pt,font=\scriptsize,fill=red!10}]
% two inputs: the same shape at two positions
\foreach \row/\sh/\lab in {0/1/{frame 1},1/7/{frame 2}}{
  \begin{scope}[yshift=-\row*1.0cm]
  \node[left,font=\scriptsize] at (-0.25,0) {\lab};
  \foreach \i in {0,...,13}{\node[px] at (\i*0.4,0) {};}
  \foreach \d in {0,1,3,4}{\node[on] at ({(\sh+\d)*0.4},0) {};}
  \end{scope}}
% S stage: detectors of the shape at several positions
\foreach \k in {0,...,5}{
  \node[snode] (S\k) at ({0.8+0.8*\k},1.7) {$S$};}
\node[font=\scriptsize,left] at (-0.25,1.7) {stage $S$: AND over parts};
\draw[->,thick,blue!60!black] (1.2,0.25) -- (S1.south);
\draw[->,thick,orange!85!black] (3.6,0.25) -- (S4.south);
\node[font=\scriptsize,blue!60!black,anchor=east] at (1.25,0.85) {frame 1};
\node[font=\scriptsize,orange!85!black,anchor=west] at (3.9,0.85) {frame 2};
% C stage
\node[cnode] (C) at (2.8,3.25) {$C$};
\node[font=\scriptsize,left] at (-0.25,3.25) {stage $C$: OR over positions};
\foreach \k in {0,...,5}{\draw[->,gray!60!black] (S\k.north) -- (C);}
\draw[->,very thick,red!70!black] (C.north) -- ++(0,0.6) node[above,font=\scriptsize,black] {``shape present'' in both frames};
% next stages
\draw[decorate,decoration={brace,amplitude=5pt},gray!60!black] (6.3,1.4) -- (6.3,-1.3);
\node[right,font=\scriptsize,align=left] at (6.5,0.05) {the same\\object, different\\pixels};
\draw[->,thick,densely dashed,gray!60!black] (3.6,3.25) -- (5.9,3.25) node[right,font=\scriptsize,black,align=left] {next pairs\\$S$, $C$: larger,\\more complex, more invariant};
\end{tikzpicture}
\caption{One pair of pipeline stages. The $S$ neurons look for the same combination of parts at different places; the $C$ neuron responds if it was found anywhere~\eqref{eq:scstage}. The shape in frame 1 and in frame 2 excites different $S$ neurons but the same $C$. The next pairs of stages repeat this on ever larger and more complex combinations.}
\label{fig:conveyor}
\end{figure}

In the same one-dimensional model, a pair of stages~\eqref{eq:scstage} recognizes the shape at any position without errors; if each input point is randomly flipped with probability $0.1$, it is right $86\%$ of the time, and with probability $0.2$, $70\%$. The coin still gives one half. Stack several such pairs on top of one another: the first looks for edges, the next for combinations of edges, higher up for parts of objects, and at the top for whole objects. With each pair the combinations get larger, and the answer depends less and less on where and how the object lies~\cite{Riesenhuber1999}.

\begin{keyblock}
\begin{proposition}[A pipeline of selectivity and invariance]
\label{prop:conveyor}
Fast recognition is a single pass through about ten stages. At each pair of stages, AND over parts~\eqref{eq:scstage} makes the response selective, and OR over positions makes it independent of shift. Alternating these two operations gives a response that tells objects apart and barely depends on where and how they are seen. The structure of the first stages and the speed of the pass are measured; that the whole chain reduces to these two operations is a simplification.
\end{proposition}
\end{keyblock}

If this design looks familiar, it should. Exactly this alternation, searching for a feature at every position and then pooling over neighboring positions, is what engineers carried over into artificial \emph{convolutional} networks~\cite{Fukushima1980}. Recently the traffic went the other way: convolutional networks trained to recognize objects turned out to be the best known model of neuron responses at the upper stages of vision~\cite{Yamins2014}. The result at the top is easy to read out: from the summed rates of about a hundred neurons of the last stage, a linear decision unit recognizes the object a tenth of a second after it is shown~\cite{Hung2005}. This is the population rate code of Chapter~\ref{sec:code}.

\section{The missing teacher: learning from video}

A convolutional network is trained on millions of labeled pictures. Almost nobody gives the brain labels: a child looks at objects for hours but hears the word ``cup'' only now and then. How, then, do the $C$ stages know which $S$ neurons to pool? To pool ``this shape here'' with ``this shape there,'' you must already know that it is the same shape.

The video itself gives the hint. The world changes continuously: an object in view remains the same object while it moves, turns, and changes lighting, and only occasionally does the gaze move to another one. Whatever changes \emph{slowly} most likely belongs to the object, and whatever changes quickly belongs to its position and view~\cite{WiskottSejnowski2002}. So the $C$ neuron needs only the Hebb rule of Chapter~\ref{sec:dofa} with an eligibility trace: link what came in sequence, not just what came at the same time~\cite{Foldiak1991}:
\begin{empheq}[box=\keybox]{equation}
\tau_{\text{tr}}\,\frac{\dd\bar y_k}{\dd t} \;=\; -\bar y_k + y_k ,
\qquad
\frac{\dd\mathbf w_k}{\dd t} \;=\; \eta\,\bar y_k\,\bigl(\mathbf x - \mathbf w_k\bigr) .
\label{eq:traceinvariance}
\end{empheq}
Here $y_k$ is the activity of $C$ neuron $k$, which won the competition through inhibition, $\bar y_k$ is its trace (the same RC circuit as in rule~\eqref{eq:threefactor}), and $\mathbf x$ is the output of the $S$ neurons. The neuron that won on the first view of an object still remembers it when the second view arrives, and pulls that view in too.

\begin{figure}[t]
\centering
\begin{tikzpicture}[>=Stealth,x=1.0cm,y=1.0cm]
\foreach \i/\lab in {0/1,1/2,2/3,3/4}{
  \node[draw,thick,fill=blue!8,minimum width=0.9cm,minimum height=0.55cm,font=\scriptsize] at (\i+0.5,2.2) {$A$ at \lab};}
\foreach \i/\lab in {0/4,1/3,2/2,3/1}{
  \node[draw,thick,fill=orange!12,minimum width=0.9cm,minimum height=0.55cm,font=\scriptsize] at (\i+6.5,2.2) {$B$ at \lab};}
\node[font=\scriptsize,gray!40!black] at (5.0,2.2) {pause};
\draw[->,gray!70!black] (0,0) -- (10.4,0) node[below left,font=\scriptsize,black] {time};
% traces
\draw[thick,blue!60!black] (0,0.05) .. controls (0.4,1.2) and (0.8,1.3) .. (1.2,1.35) -- (4,1.4) .. controls (4.6,0.5) and (5.2,0.15) .. (6,0.08) -- (10,0.05);
\draw[thick,orange!85!black] (0,0.05) -- (6,0.05) .. controls (6.4,1.2) and (6.8,1.3) .. (7.2,1.35) -- (10,1.4);
\node[font=\scriptsize,blue!60!black,anchor=west] at (1.4,1.65) {trace $\bar y_1$ lasts the whole pass of $A$};
\node[font=\scriptsize,orange!85!black,anchor=west] at (7.2,1.65) {trace $\bar y_2$};
\draw[decorate,decoration={brace,amplitude=4pt},gray!60!black] (4.0,2.6) -- (0,2.6);
\draw[decorate,decoration={brace,amplitude=4pt,mirror},gray!60!black] (0,-0.35) -- (4,-0.35) node[midway,below=4pt,font=\scriptsize,black] {all views of $A$ are linked to neuron 1};
\end{tikzpicture}
\caption{How video teaches invariance, schematic. Object $A$ passes through four positions, then a pause, then object $B$. The trace~\eqref{eq:traceinvariance} of the neuron that won on the first view of $A$ lasts the whole pass, and all views of $A$ end up linked to one neuron. The pause erases the trace, and $B$ goes to another neuron.}
\label{fig:tracevideo}
\end{figure}

We tested this in a model (Fig.~\ref{fig:tracevideo}). Two $C$ neurons compete for inputs from $16$ $S$ neurons: two shapes at eight positions. A shape sweeps through all positions, $50\ms$ at each, then comes a $0.3\,\text{s}$ pause and the next random shape. No labels at all. If the trace is short, $\tau_{\text{tr}}=5\ms$, the neurons split the input by \emph{position}, and their answer matches the shape no better than a coin: $52\%$ on average over ten runs. If the trace is comparable to the time one position is shown, from $\tau_{\text{tr}}\approx20\ms$ up, each neuron collects \emph{all} positions of its shape (at $20$, $50$, and $200\ms$, in all ten runs): invariance is learned from video alone. The pause between objects matters: without it the trace spills over onto the next object and confuses the two.

The same thing shows up in experiment. If an object is covertly swapped while the eye jumps from one place to another, neurons of the upper stages begin, within an hour or two, to respond to the swapped object as if it were the same one: they link what came in sequence~\cite{LiDiCarlo2008}. Invariance in the brain really is learned from the continuity of time. How much of visual learning this rule explains is a hypothesis.

\section{Video is motion}

Video is not only a stream of frames to learn from but also a task in itself: what is moving, where, and how fast. The first step is familiar: the direction detector of Chapter~\ref{sec:gaba} (Fig.~\ref{fig:direction}), in which delayed inhibition cancels motion in one direction. Such detectors cover the whole visual field, and from there on they are treated just like shape features: stages that collect many detectors of the same direction at different places respond to the motion of a large object regardless of where it is. This is the same AND and OR pair~\eqref{eq:scstage}, except that the feature is not ``an edge'' but ``an edge that moved within time $d$.''

Video is also predictable: the next frame is almost the same as the previous one. According to the \emph{predictive coding} hypothesis, the upper stages send a prediction down to the lower ones, and what goes up is mainly the error, whatever the prediction missed~\cite{RaoBallard1999}. This is the same spike economy as in Proposition~\ref{prop:economy}: pay for the news, not for what is already known. Some observations agree with this hypothesis, but it has not been proven as the general design of the pipeline.

\section{The brain and the convolutional network}

How far does the similarity go? For fast recognition of a single object it runs deep~\cite{Yamins2014}. But the brain also has things a plain convolutional network lacks.

\emph{Feedback.} The brain's pipeline is not only feedforward: every stage has connections backward and sideways. For hard pictures (occluded, poorly lit) the response of the upper stages forms later, and these late responses are better described by networks with feedback than by feedforward ones~\cite{Kar2019}. A single pass solves the easy cases; the hard ones need several rounds.

\emph{Robustness.} An artificial network can be fooled by a pixel forgery invisible to the eye~\cite{Szegedy2014}: a change a human cannot see turns a ``panda'' into a ``gibbon''~\cite{Goodfellow2015}. Human vision is far more robust to such forgeries, but not immune: images that fool many networks at once also shift human choices when shown briefly~\cite{Elsayed2018}. Why robustness differs so much is an open question; the candidates are noise, division by the total activity, and feedback.

\emph{The teacher.} A network needs millions of labeled examples; the brain needs mostly unlabeled video~\eqref{eq:traceinvariance} and a few hints from \nt{DOPA} about what matters.

\emph{Energy.} The whole brain uses about $20$ watts (Proposition~\ref{prop:economy}), and recognition takes only part of that; a trained convolutional network on a graphics accelerator uses hundreds of watts.

\section{Illusions: where the machine errs and why}
\label{sec:illusions}

An engineer who wants to understand someone else's circuit feeds it unusual signals and watches where it fails. For vision such signals were invented long ago: they are visual illusions. An illusion is not a malfunction. Each one points to a particular mechanism of the machine that works correctly in ordinary conditions and gives itself away on a specially chosen picture.

\emph{Reasonable expectations.} The input is noisy, and the machine fills it in with what is usual in the world. Slow motions are more common than fast ones; when the input is unreliable, for example when the picture has little contrast, the speed estimate shifts toward slow, and a faint object seems to move more slowly than a bright one~\cite{Weiss2002}. Many brightness illusions are explained the same way: we see not the brightness of a pixel but the surface that such brightness most often corresponded to in the past~\cite{Purves2011}. The photograph of a dress that some see as white and gold and others as blue and black is the same mechanism: the picture is ambiguous, and people differ in what they unconsciously assume the lighting to be~\cite{LaferSousa2015,Wallisch2017}. The differences between people are measured; that the brain here combines input with expectation by the rules of probability theory is a well-founded hypothesis.

\emph{Gain control.} Stare at a waterfall for a minute, then at motionless rocks: the rocks will drift upward. The ``down'' and ``up'' channels are a pair of wires, as in~\eqref{eq:dualrail}; gain control~\eqref{eq:nakarushton} has turned down the fatigued channel, and the balance of the pair has shifted. Tilt and contrast illusions, in which the surround changes how the center looks, are explained by division by the activity of the neighbors~\cite{Schwartz2009}, as in Chapter~\ref{sec:gaba}. There is also an instructive case for caution: the dark spots at the intersections of the white bars of a grid were long explained by simple inhibition from neighbors, but experiments with a modified grid showed that this explanation does not work~\cite{SchillerCarvey2005}.

\emph{Delay compensation.} The path from the eye to the upper stages takes tens of milliseconds, and a moving object shifts in that time. If a stationary dot flashes next to it, the object appears ahead of the flash, although the two were aligned~\cite{Nijhawan1994}. According to one explanation, the machine extrapolates the motion forward to compensate for the delay, as in Chapter~\ref{sec:muscles}: with delayed feedback one has to predict. This illusion has several explanations, and the debate is not settled.

\emph{An error in the timing code.} A still image of rings with sharp color transitions appears to rotate. High-contrast regions are processed faster than faint ones, and the difference in latencies~\eqref{eq:latency} is read by direction detectors as motion~\cite{Conway2005}. The illusion is triggered by small involuntary jumps of the eye and by blinks: each jump refreshes the input, and the detectors see ``motion'' again~\cite{OteroMillan2012}. This is the timing code of Chapter~\ref{sec:code}, read incorrectly.

\emph{Two pictures in one.} A wire-frame cube that faces us first with one side and then with another, or different pictures in the two eyes: perception flips every few seconds. The best models are mutual inhibition between the two interpretations, slow fatigue, and noise~\cite{MorenoBote2007}, that is, the very circuit~\eqref{eq:halfcentre} that in Chapter~\ref{sec:muscles} generated the stepping rhythm. The switching time is set by noise and fatigue; in the model of~\cite{MorenoBote2007} noise plays the main role, and fatigue only helps.

And artificial networks? A network trained to predict the next video frame sees rotation in the same still rings and does not see it in control pictures~\cite{Watanabe2018}. Networks trained to denoise and sharpen images fall for the same color and brightness illusions as people~\cite{GomezVilla2020}. So some illusions follow from the task the machine learns, not from peculiarities of neural tissue. Which illusions belong to the brain alone is a good test for any model of vision.

\begin{keyblock}
\begin{proposition}[Illusions as fingerprints of mechanisms]
\label{prop:illusions}
An illusion arises when a mechanism that is useful in the ordinary world receives an unusual input: expectation overrides an unreliable input, gain control shifts a pair of channels, delay compensation carries an object forward, a difference in latencies is read as motion, mutual inhibition with noise and fatigue flips. Each illusion is a test signal pointing to its own mechanism. The illusions themselves are measured; their explanations range from well founded to disputed.
\end{proposition}
\end{keyblock}

\section{Summary}

\begin{table}[t]
\centering
\footnotesize
\setlength{\tabcolsep}{4pt}
\renewcommand{\arraystretch}{1.15}
\begin{tabular}{>{\raggedright}p{3.0cm}>{\raggedright}p{4.6cm}>{\raggedright\arraybackslash}p{5.2cm}}
\toprule
What the machine does & How & What is known \\
\midrule
recognizes in $150\ms$ & a single pass through about ten stages & speed measured \\
makes the response selective & AND over parts, stage $S$~\eqref{eq:scstage} & measured at the first stages \\
makes the response invariant & OR over positions, stage $C$; inhibition and division & measured at the first stages; for the upper ones, a simplification \\
gives the answer & rates of about a hundred neurons of the last stage, linear readout & measured \\
learns without labels & Hebb rule with a trace~\eqref{eq:traceinvariance} on continuous video & object swap changes responses: measured; as the main rule: hypothesis \\
sees motion & direction detectors, then the same AND/OR pair & measured \\
saves on the predictable & the prediction error goes up & hypothesis \\
solves hard cases & feedback, several rounds & late responses and the role of feedback measured \\
errs predictably & illusions: expectations, gain control, delays, mutual inhibition with noise and fatigue & illusions measured; explanations from well founded to disputed \\
\bottomrule
\end{tabular}
\caption{How the machine recognizes images and video.}
\label{tab:recognition}
\end{table}

Recognition is built from parts we already had (Table~\ref{tab:recognition}): the threshold gives AND over parts, mutual inhibition gives OR over positions, division gives independence from brightness, and the Hebb rule with a trace stands in for the missing teacher. Engineers took from vision the alternation of selectivity and invariance and built convolutional networks on it; the brain has not yet given away its main secret: learning from unlabeled video on almost no energy.

\section{Exercises}

\begin{exercise}[\lvlA]
\label{ex:12:1}
\emph{Rate at a single stage.} A pipeline stage lasts $15\ms$, neurons fire at $20$ spikes/s, and the noise is correlated with $c=0.1$~\eqref{eq:correlated}. What is the relative error of the rate of a group of a hundred neurons, and its limit for an arbitrarily large group? How many rate levels can a stage distinguish?
\end{exercise}

\begin{exercise}[\lvlA]
\label{ex:12:2}
\emph{The energy of one recognition.} In a single $150\ms$ pass, ten stages of $10^7$ neurons each fire at most one spike each, at $10^{-10}$~J per spike. (a)~What fraction of the whole brain's energy over those $150\ms$ does recognition cost? (b)~How many times more does a $300$~W accelerator spend if it recognizes a picture in $10\ms$?
\end{exercise}

\begin{exercise}[\lvlB]
\label{ex:12:3}
\emph{Thresholds of stage $S$.} A ``retina'' of $24$ points, shapes $A=11011$ and $B=10101$; the template~\eqref{eq:scstage} is $+1$ on the ones of the shape and $-1$ on the zeros. (a)~For which thresholds does an $S$ neuron forgive one flipped point but stay silent on the other shape? (b)~How many $S$ neurons are needed for ten $5\times5$ features on a $100\times100$ picture?
\end{exercise}

\begin{exercise}[\lvlB]
\label{ex:12:4}
\emph{How long one of two pictures lasts.} In the flipping circuit~\eqref{eq:halfcentre} with $\tau\ll\tau_a$, while group~1 wins, $r_2=0$ and $r_1=I-a_1$. (a)~Find the duration of a win for $I=1$, $\beta=2$, $b=2$, $\tau_a=0.4\,\text{s}$. (b)~What $\tau_a$ gives a switch every three seconds?
\end{exercise}

\begin{exercise}[\lvlB]
\label{ex:12:5}
\emph{Simulation: the cost of invariance.} Using the function \texttt{two\_stage} from \texttt{models/\allowbreak recognition\_models.py} ($4000$ trials), find the point-flip probability $p$ at which the $S$, $C$ pair with $N=24$ errs in one case out of four. How does the fraction of correct answers at $p=0.1$ change from $N=8$ to $N=192$?
\end{exercise}

\begin{exercise}[\lvlB]
\label{ex:12:6}
\emph{Simulation: the trace window.} With ten runs of \texttt{trace\_learning} from \texttt{models/\allowbreak recognition\_models.py} and a $0.3\,\text{s}$ pause, find for which $\tau_{\text{tr}}$ between $5\ms$ and $2\,\text{s}$ all runs learn invariance without errors. Where does the upper edge of this window come from?
\end{exercise}

\begin{exercise}[\lvlB]
\label{ex:12:7}
\emph{Motion anywhere.} Direction detectors of Chapter~\ref{sec:gaba} sit at neighboring points of a row spaced $0.1^\circ$ apart: inhibition from $A$ with delay $d$ cancels excitation from $B$ if the two are at most $5\ms$ apart; a stage $C$ sits above them. What delays keep $C$ silent for reverse motion at speeds from $5$ to $20^\circ$ per second?
\end{exercise}

\begin{exercise}[\lvlC]
\label{ex:12:8}
\emph{Pixel forgery.} How many flipped points does it take to change the answer of the \texttt{two\_stage} model ($N=24$) on a clean shape? And for the classifier ``more than $3.5$ ones in the input means $A$,'' which is also shift invariant? What is its fraction of correct answers at $p=0.1$, against $0.86$ for the pair of stages?
\end{exercise}

\chapter{The Machine's Outputs: How the Brain Controls Muscles}
\chaptermark{The Machine's Outputs}
\label{sec:muscles}
\begin{chaptergoals}
\item how a muscle turns spikes into force, acting as a low-pass digital-to-analog converter;
\item why recruiting units by size sets force in floating point, with steps proportional to force;
\item how a pair of pulling muscles sets torque by the difference and stiffness by the sum;
\item why delayed feedback limits correction speed, and how inhibition with fatigue makes a rhythm.
\end{chaptergoals}

\section{The fast output}

Everything the machine does quickly in the world, it does with muscles. A word, a glance, a step, a smile: these are muscle contractions. The second, slow output, the chemistry of the body, is covered in Chapter~\ref{sec:chemoutput}. In Fig.~\ref{fig:machine} the output was an arrow labeled ``muscles''; now let us look at what lies behind that arrow.

The last link of the chain is the \nt{ACET} neurons, the very ones listed in Table~\ref{tab:types} as ``output to muscle.'' Each muscle fiber receives input from exactly one such neuron, and one neuron controls a group of fibers, from several hundred to a couple of thousand~\cite{Feinstein1955mu}. In muscles that need fine adjustment the units are smaller; in the large leg muscles they are larger. A neuron together with its fibers is called a \emph{motor unit}: it is the smallest piece of muscle the brain can switch on separately.

The synapse of a neuron on a muscle is built quite differently from the cortical synapses of Chapter~\ref{sec:code}. There, most spikes were lost. Here, each spike releases several times more acetylcholine than is needed to make the fiber fire~\cite{WoodSlater2001}. This is a synapse with a \emph{safety factor}: one spike of the neuron means exactly one contraction of its fibers. Inside the machine, unreliable connections are affordable and errors can be corrected by redundancy; at the output, an error costs a movement, and reliability is bought directly in hardware.

\section{Force: rate and number of units}

One spike produces a single contraction, a twitch: the force rises over a few tens of milliseconds and then decays. A good approximation for it is~\cite{Fuglevand1993}
\begin{equation}
h(t) \;=\; F_1\,\frac{t}{T_c}\,e^{\,1-t/T_c},
\label{eq:twitch}
\end{equation}
where $T_c$ is the rise time, of order $50\ms$, and $F_1$ is the force of one twitch. If spikes come more often, the twitches add up:
\begin{empheq}[box=\keybox]{equation}
F(t) \;=\; \sum_k h\bigl(t-t_k\bigr) .
\label{eq:forcesum}
\end{empheq}
This is the same low-pass filter that turned spikes into concentration in Chapter~\ref{sec:serocomputer}, except that now the output is force rather than concentration. A muscle is a digital-to-analog converter from spikes to force (Fig.~\ref{fig:tetanus}). In our model, at five spikes per second the twitches barely overlap and the force jumps between $0.2F_1$ and $1.1F_1$; at fifteen it already stays between $1.8F_1$ and $2.2F_1$; at forty it becomes almost flat, about $5.4F_1$. A real muscle saturates at high spike rates, so the linear sum~\eqref{eq:forcesum} holds only up to a few tens of spikes per second.

\begin{figure}[t]
\centering
\begin{tikzpicture}[>=Stealth,x=20cm,y=0.55cm]
\draw[->,gray!70!black] (0,0) -- (0.66,0) node[right,font=\small,black] {$t$, s};
\draw[->,gray!70!black] (0,0) -- (0,6.4) node[left,font=\small,black] {$F/F_1$};
\foreach \t in {0,0.2,0.4,0.6}{\draw[gray!60!black] (\t,-0.1) -- (\t,0.1); \node[below,font=\scriptsize] at (\t,0) {\t};}
\foreach \f in {1,3,5}{\draw[gray!60!black] (-0.004,\f) -- (0.004,\f); \node[left,font=\scriptsize] at (0,\f) {\f};}
\draw[thick,blue!60!black] plot file {figures/muscle_twitch_5.dat};
\draw[thick,green!45!black] plot file {figures/muscle_twitch_15.dat};
\draw[thick,red!70!black] plot file {figures/muscle_twitch_40.dat};
\node[font=\scriptsize,blue!60!black,anchor=west] at (0.605,0.9) {5 spikes/s};
\node[font=\scriptsize,green!45!black,anchor=west] at (0.605,2.1) {15 spikes/s};
\node[font=\scriptsize,red!70!black,anchor=west] at (0.605,5.4) {40 spikes/s};
\end{tikzpicture}
\caption{A muscle as a digital-to-analog converter: the force~\eqref{eq:forcesum} for regular spikes of a single motor unit, $T_c=50\ms$. Sparse spikes give separate twitches; frequent ones fuse into a smooth force, just as frequent spikes fuse into a smooth concentration in Fig.~\ref{fig:lowpass}.}
\label{fig:tetanus}
\end{figure}

The brain has two force knobs: the spike rate of each unit and the number of units switched on. The second knob is very ingeniously designed. The units in a muscle differ: the weakest are a hundred times weaker than the strongest. As more and more force is needed, units are switched on \emph{in order of size}, from small to large~\cite{Henneman1957}. Nobody computes this order: small neurons are simply easier to excite by the same common input, so the recruitment order is set by the hardware itself.

What does this buy? Let each next unit in the order be $q$ times stronger than the previous one, where $q$ is slightly larger than one. The force of all recruited units is the sum of a geometric series, and almost all of it comes from the last units. Then each newly recruited unit adds
\begin{empheq}[box=\keybox]{equation}
\frac{\Delta F}{F} \;\approx\; q-1 \;=\; \text{const} ,
\label{eq:sizeprinciple}
\end{empheq}
that is, the force step is proportional to the force itself. For a weak effort the brain doles it out in small portions, for a strong one in large portions. This is a \emph{floating-point number}: the exponent is set by how many units are recruited, and the mantissa by their spike rate. It is the same logarithmic scale as in the sensors of Chapter~\ref{sec:sensors}: the input and the output of the machine measure the world in relative terms.

Floating point has a downside. The scatter of the force also grows in proportion to the force itself: a strong movement is inevitably less precise than a weak one. According to one influential hypothesis, this signal-dependent noise explains why our movements are smooth: a smooth command gives the smallest scatter at the end point~\cite{HarrisWolpert1998}.

\section{Pull, never push: two wires per joint}

A muscle can only pull. It cannot push, just as a \nt{GLUT} neuron cannot output a negative signal. The solution is familiar from Chapter~\ref{sec:glutcomputer}: \emph{two wires}. Each joint is served by a pair of muscles pulling in opposite directions (Fig.~\ref{fig:antagonists}). If their forces are $F_1$ and $F_2$ and the lever arm is $\ell$, the torque and stiffness of the joint are
\begin{empheq}[box=\keybox]{equation}
M \;=\; \ell\,(F_1-F_2), \qquad K \;\approx\; \kappa_m\,(F_1+F_2),
\label{eq:antagonists}
\end{empheq}
where $\kappa_m$ is the coefficient relating muscle force to muscle stiffness: a tensed muscle is springier than a relaxed one. The difference of the forces sets the torque, the sum sets the stiffness~\cite{Hogan1984}. With two wires the brain controls two independent quantities: which way to turn the joint and how firmly to hold it. When we need to hold a cup without spilling, we tense both muscles at once: the torque is the same, the stiffness is higher, and a random bump knocks the arm off less.

\begin{figure}[t]
\centering
\begin{tikzpicture}[>=Stealth,
  nrn/.style={circle,draw,very thick,minimum size=0.95cm,inner sep=0pt,font=\scriptsize},
  inh/.style={-{Bar[width=7pt]},thick,red!70!black}]
% the joint: a bar on a pivot, two springs pulling down on either side
\fill[gray!30] (4.6,-0.25) -- (5.4,-0.25) -- (5,0.15) -- cycle;
\draw[line width=3pt,gray!50!black] (2.4,0.2) -- (7.6,0.2);
\fill[black] (5,0.2) circle (0.08);
\draw[decorate,decoration={coil,segment length=5pt,amplitude=4pt},very thick,blue!60!black] (3.0,0.2) -- (3.0,-1.6);
\draw[decorate,decoration={coil,segment length=5pt,amplitude=4pt},very thick,orange!85!black] (7.0,0.2) -- (7.0,-1.6);
\draw[very thick] (2.5,-1.6) -- (3.5,-1.6); \draw[very thick] (6.5,-1.6) -- (7.5,-1.6);
\node[font=\small,blue!60!black] at (2.35,-0.75) {$F_1$};
\node[font=\small,orange!85!black] at (7.65,-0.75) {$F_2$};
\draw[->,thick] (5.9,0.75) arc (60:120:1.8) node[above,font=\scriptsize,pos=0.5] {$M=\ell(F_1-F_2)$};
% motor neurons and reflex circuit
\node[nrn,fill=green!12] (M1) at (1.6,-3.0) {\nt{ACET}};
\node[nrn,fill=green!12] (M2) at (8.4,-3.0) {\nt{ACET}};
\node[nrn,fill=red!10] (G) at (5.0,-3.0) {\nt{GLYC}};
\node[nrn,fill=blue!6] (S) at (1.6,-5.0) {\nt{GLUT}};
\draw[->,very thick] (M1) -- (3.0,-1.7);
\draw[->,very thick] (M2) -- (7.0,-1.7);
\draw[->,thick] (S) -- (M1);
\draw[->,thick] (S) -- (G);
\draw[inh] (G) -- (M2);
\draw[->,thick,densely dashed,gray!60!black] (3.15,-1.2) to[bend left=25] node[pos=0.35,right,font=\scriptsize,black] {stretch} (S.north east);
\draw[->,thick,blue!60!black] (-0.3,-3.0) node[left,font=\scriptsize,black,align=right] {command\\from brain} -- (M1);
\draw[->,thick,blue!60!black] (10.3,-3.0) node[right,font=\scriptsize,black,align=left] {command\\from brain} -- (M2);
\end{tikzpicture}
\caption{Two muscles on one joint and the fastest feedback loop. \nt{ACET} neurons receive commands from the brain and pull the joint in opposite directions; the difference of the forces turns it, the sum makes it stiffer. The stretch sensor (\nt{GLUT}) of the left muscle excites that muscle's own \nt{ACET} neuron and, through a \nt{GLYC} intermediary, inhibits the neuron of the opposing muscle: the veto of Chapter~\ref{sec:gaba}.}
\label{fig:antagonists}
\end{figure}

\section{Feedback with delay}

Muscles do not work blind. Each one contains the stretch sensors mentioned in Table~\ref{tab:sensors}, and the shortest loop closes right in the spinal cord (Fig.~\ref{fig:antagonists}). When a muscle is unexpectedly stretched, the stretch sensor excites the muscle's own \nt{ACET} neuron, the muscle contracts and returns the joint. At the same time, an inhibitory intermediary neuron switches off the opposing muscle so that it does not interfere. This intermediary is \nt{GLYC}, the very type from Table~\ref{tab:types} that did not get a chapter of its own: its place is right here, in the fast loops of the spinal cord.

Every loop has a delay $d$: the spinal loop about $25$--$30\ms$ for the arm, a loop through the cortex one and a half to three times longer, a loop through the eye $100$--$200\ms$. And delayed feedback, as we saw in Proposition~\ref{prop:ei}, makes a system oscillate. For the simplest controller, which corrects a deviation $x$ at rate $G$ using information $d$ seconds old, $\dd x/\dd t=-G\,x(t-d)$, the stability condition is~\cite{Hayes1950}
\begin{empheq}[box=\keybox]{equation}
G\,d \;<\; \frac{\pi}{2} ,
\label{eq:delaystable}
\end{empheq}
and at the boundary the system oscillates at frequency $1/(4d)$. We checked this in a model: at $d=30\ms$ the deviation decays for $G=40$ per second, and for $G=55$ per second, just above the boundary of $52$, it grows into oscillation.

Two consequences follow. First, the longer the loop, the more slowly it must correct. With a delay of a hundred and fifty milliseconds through the eye, the arm cannot be corrected faster than in about a tenth of a second, or it will tremble. Second, the stability boundary of the spinal loop, $1/(4\cdot30\ms)\approx8$ oscillations per second, is close to the frequency of the fine tremor present in every healthy person, eight to twelve oscillations per second. The stretch loop is one of the likely sources of this tremor~\cite{McAuleyMarsden2000}.

How, then, does the brain move quickly and accurately if its feedback is late? According to a widely held hypothesis, it \emph{predicts}: it keeps an internal model of the body and computes the result of a command without waiting for the sensors~\cite{WolpertGhahramani2000}. The sensors are needed to correct the model, not every movement. An engineer would call this feedforward control with a state observer.

\section{Rhythm generators for movement}

Walking, breathing, chewing are rhythmic movements. Do they need a clock? No. The spinal cord and the brainstem contain small networks that generate a rhythm by themselves, even when nothing rhythmic comes into them~\cite{MarderBucher2001}. The simplest such network is built from parts we already have: two groups of \nt{GLUT} neurons that inhibit each other through \nt{GABA} or \nt{GLYC} intermediaries, as in Fig.~\ref{fig:wta}, and slowly fatigue, like the memory of Chapter~\ref{sec:glutcomputer}:
\begin{empheq}[box=\keybox]{equation}
\tau\,\frac{\dd r_i}{\dd t} = -r_i + \bigl[I - \beta\,r_j - a_i\bigr]_+ ,
\qquad
\tau_a\,\frac{\dd a_i}{\dd t} = -a_i + b\,r_i ,
\label{eq:halfcentre}
\end{empheq}
where $a_i$ is the fatigue of group $i$ and $j$ is the other group. Mutual inhibition with $\beta>1$ picks a winner (Proposition~\ref{prop:wta}). The winner works and tires; when fatigue eats up its input, the loser, already rested, pulls ahead and starts to tire in turn. The groups work in alternation, and each drives its own muscle of the pair (Fig.~\ref{fig:halfcentre}).

\begin{figure}[t]
\centering
\begin{tikzpicture}[>=Stealth,x=3.2cm,y=2.4cm]
\draw[->,gray!70!black] (0,0) -- (4.2,0) node[right,font=\small,black] {$t$, s};
\draw[->,gray!70!black] (0,0) -- (0,1.2) node[left,font=\small,black] {$r$};
\foreach \t in {0,1,2,3,4}{\draw[gray!60!black] (\t,-0.03) -- (\t,0.03); \node[below,font=\scriptsize] at (\t,0) {\t};}
\draw[thick,blue!60!black] plot file {figures/halfcentre1.dat};
\draw[thick,orange!85!black] plot file {figures/halfcentre2.dat};
\node[font=\scriptsize,blue!60!black,anchor=west] at (1.6,1.28) {group 1: ``forward''};
\node[font=\scriptsize,orange!85!black,anchor=west] at (2.7,1.28) {group 2: ``back''};
\end{tikzpicture}
\caption{A rhythm generator from model~\eqref{eq:halfcentre}: $I=1$, $\beta=2$, $b=2$, $\tau=20\ms$, $\tau_a=0.4\,\text{s}$. Two groups with mutual inhibition and fatigue work in alternation with a period of $0.84\,\text{s}$, although the input is a constant signal.}
\label{fig:halfcentre}
\end{figure}

In our model, with a constant input the groups take turns with a period of $0.84\,\text{s}$, about twice the fatigue time $\tau_a$. A constant command from above, ``walk,'' turns into a rhythm on the spot. This is one more local rhythm, like the rhythm of the \nt{GLUT}--\nt{GABA} loop in Chapter~\ref{sec:gaba}: it arises only while the network is working, and it paces only the muscles it controls.

\begin{keyblock}
\begin{proposition}[The machine's output]
\label{prop:muscles}
The brain controls muscles as a hierarchy of controllers. At the top an action is selected (Chapter~\ref{sec:dofa}); below, local generators turn constant commands into rhythm, and fast spinal loops hold the joints. Force is set in floating point: the exponent by the number of units, recruited by size, the mantissa by the spike rate. Each joint is controlled by a pair of pulling muscles, the difference of whose forces sets the torque and the sum the stiffness. These mechanisms are measured; that the brain relies on an internal model of the body is a well-founded hypothesis.
\end{proposition}
\end{keyblock}

\section{Summary}

\begin{table}[t]
\centering
\footnotesize
\setlength{\tabcolsep}{4pt}
\renewcommand{\arraystretch}{1.15}
\begin{tabular}{>{\raggedright}p{3.0cm}>{\raggedright}p{4.6cm}>{\raggedright\arraybackslash}p{5.2cm}}
\toprule
What the output does & How & What is known \\
\midrule
turns spikes into force & sum of twitches~\eqref{eq:forcesum}, a low-pass filter & measured; the linear sum is a simplification \\
doses force & recruiting units by size, floating point~\eqref{eq:sizeprinciple} & recruitment order measured \\
pushes while only able to pull & a pair of muscles: torque is the difference, stiffness the sum~\eqref{eq:antagonists} & measured \\
holds a joint & stretch loop, veto of the opponent via \nt{GLYC} & measured \\
does not tremble & $Gd<\pi/2$~\eqref{eq:delaystable} & follows from control theory; contribution to tremor is a likely cause \\
moves quickly & prediction from an internal model & well-founded hypothesis \\
walks and breathes rhythmically & rhythm generator~\eqref{eq:halfcentre} & generators measured; the simplest circuit is a model \\
\bottomrule
\end{tabular}
\caption{How the machine controls muscles.}
\label{tab:muscles}
\end{table}

The machine's output is built with the same tricks as everything else (Table~\ref{tab:muscles}): two wires instead of a sign, a low-pass filter instead of a DAC, a logarithmic scale instead of a linear one, mutual inhibition and fatigue instead of a clock generator. The input (Chapter~\ref{sec:sensors}) and the output measure the world in relative terms, and between them works the machine this book is about.

\section{Exercises}

\begin{exercise}[\lvlA]
\label{ex:13:1}
\emph{Floating point in numbers.} A muscle has $N=100$ motor units with forces $f_n=f_1q^{\,n-1}$; the strongest is $100$ times stronger than the weakest. (a)~Find $q$, the relative step~\eqref{eq:sizeprinciple}, and the dynamic range in decibels. (b)~What is the step at force $10f_1$ for a ``linear DAC'' of $100$ identical units with the same total force?
\end{exercise}

\begin{exercise}[\lvlA]
\label{ex:13:2}
\emph{The price of the safety factor.} On each spike the synapse on a muscle fiber releases a Poisson number of quanta with mean $m$, and the fiber fires at $n_{\text{th}}=20$ quanta or more; the safety factor is $S=m/n_{\text{th}}$. How many failures per day does a unit firing at $20$ spikes/s produce at $S=2$ and at $S=4$?
\end{exercise}

\begin{exercise}[\lvlB]
\label{ex:13:3}
\emph{The muscle as a filter.} The twitch~\eqref{eq:twitch} with $T_c=50\ms$ is the impulse response of a linear system. (a)~Find its transfer function $H(s)$ and cutoff frequency. (b)~At what rate of regular spikes is the peak-to-peak force ripple $10\%$ of the mean force?
\end{exercise}

\begin{exercise}[\lvlB]
\label{ex:13:4}
\emph{No correction faster than the delay.} The controller is $\dd x/\dd t=-G\,x(t-d)$. (a)~Show that the deviation decays fastest without oscillation at $Gd=1/e$, at rate $1/d$. (b)~Find this $G$ and the decay time constant for the spinal loop, $d=30\ms$, and for the loop through the eye, $d=150\ms$.
\end{exercise}

\begin{exercise}[\lvlB]
\label{ex:13:5}
\emph{The period of the rhythm generator.} For $\tau\ll\tau_a$ the period $T$ of model~\eqref{eq:halfcentre} is given by $(\beta-1)(1-E^{2+b})/(\beta-E)=b\,(1-E^{1+b})/(1+b)$, where $E=e^{-T/(2\tau_a)}$. (a)~Find $T$ for $\beta=b=2$, $\tau_a=0.4$~s. (b)~Show that $T$ does not depend on the input $I$, and that a rhythm is possible only for $b>\beta-1$.
\end{exercise}

\begin{exercise}[\lvlB]
\label{ex:13:6}
\emph{Simulating the generator.} Import the function \texttt{halfcentre} from \texttt{models/\allowbreak muscle\_models.py} (do not run the file itself) and measure the period, discarding the first two cycles, for $b=0.8$, $1.05$, $1.2$, $1.5$, $2$, $4$ and for $I=0.5$, $1$, $2$. Can the command ``walk faster'' change the tempo through $I$?
\end{exercise}

\begin{exercise}[\lvlB]
\label{ex:13:7}
\emph{Stiffness versus reflex.} A joint is enclosed in a stretch loop: $\dd\theta/\dd t=-a\theta(t)-G\theta(t-d)$, $a=K/B$, $d=30\ms$. Reduce it to Exercise~\ref{ex:13:4} and find the smallest $a$ at which the deviation decays without oscillation with a time constant of $15\ms$, and the required $G$. How should $G$ change as muscle co-contraction increases?
\end{exercise}

\begin{exercise}[\lvlC]
\label{ex:13:8}
\emph{A tempo knob.} By Exercises~\ref{ex:13:5} and~\ref{ex:13:6}, the input $I$ of generator~\eqref{eq:halfcentre} changes only the amplitude, not the tempo. Propose a minimal change to the model, built from parts in the book, such that below some $I_{\min}$ there is no rhythm, and above it the period decreases with $I$, at least halving when $I$ triples. Find $I_{\min}$ for $\tau\ll\tau_a$ and check by simulation.
\end{exercise}

\chapter{Pain: An Alarm Signal}
\chaptermark{Pain}
\label{sec:pain}

\section{An alarm, not a measurement}

All the sensors of Chapter~\ref{sec:sensors} measured the world: how much light, what pitch of sound, what pressure. Pain works differently. It reports not ``what is out there'' but ``danger, drop everything.'' To an engineer this is not a measurement channel but an \emph{alarm signal}: a top-priority line that must break through whatever the machine is currently doing.

An alarm differs from a measurement channel in three properties, and pain has all three. First, it is hard to mute by adaptation: pain sensors adapt far less than touch sensors, and after mild injury they even become more sensitive~\cite{LaMotte1982hyper}; a weakening of the response after strong heating has been measured in only one of their types~\cite{Treede1995heat}. A light sensor falls silent against a constant background (Fig.~\ref{fig:adaptation}), but an alarm that goes quiet after a minute is useless. Second, it has a high threshold: pain sensors fire only when a stimulus is close to dangerous; for heating, for example, the thresholds of different sensors range from $41^\circ$ to $46$--$48^\circ$, depending on the type~\cite{Treede1995heat, Weidner1999cfib}. Third, and this is the main point of the chapter, the volume of the alarm is decided not only by the sensor but by the machine itself. The same wound hurts differently in different circumstances. Let us see which familiar parts this is built from.

Pain sensors are \nt{GLUT} neurons, like the other sensory neurons of Chapter~\ref{sec:sensors}. Some of them release, along with glutamate, substance~P from Appendix~\ref{sec:othermediators}, which slowly enhances transmission.

\section{Two wires: fast and slow pain}

Hit your finger and you feel pain twice: first short and sharp, then, a second later, dull and lasting. These are two different wires. The fast wire is insulated and conducts spikes at a speed $v$ of $5$ to $30$ meters per second; the slow one is thin and bare, $0.5$--$2$ meters per second. A signal from the foot travels about a meter, and the arrival time
\begin{empheq}[box=\keybox]{equation}
t \;=\; \frac{L}{v}
\label{eq:paintime}
\end{empheq}
is tens of milliseconds for the fast wire and about a second for the slow one. The two wires carry two messages. The fast one says \emph{where} and \emph{when}: its spikes arrive earlier and more precisely, as in the timing code of Chapter~\ref{sec:code}. The slow one says \emph{how bad}, and its long dull pain keeps the machine in alarm mode until the damage has passed.

\section{A gate in the spinal cord}

The first place a pain signal arrives is the spinal cord, and there it passes through a \emph{gate}~\cite{MelzackWall1965}; the specific neurons of this gate were found relatively recently~\cite{Duan2014}. The pain wire and the touch wire converge on an output \nt{GLUT} neuron that passes pain on to the brain. Next to it sits an inhibitory interneuron, \nt{GABA} or \nt{GLYC}, which is excited by touch (Fig.~\ref{fig:gate}). Touch closes the gate. In the original gate theory~\cite{MelzackWall1965}, pain, conversely, switches this interneuron off, and strong pain opens the gate; this is an assumption of the theory and has not been shown directly on the gate neurons that were found.

\begin{figure}[t]
\centering
\begin{tikzpicture}[>=Stealth,
  nrn/.style={circle,draw,very thick,minimum size=1.0cm,inner sep=0pt,font=\scriptsize},
  inh/.style={-{Bar[width=7pt]},thick,red!70!black},
  bc/.style={->,thick,densely dashed,orange!85!black}]
\node[nrn,fill=blue!6] (Z) at (6,0) {\nt{GLUT}};
\node[nrn,fill=red!10] (Q) at (3.6,-1.6) {\nt{GABA}};
\draw[->,very thick,red!70!black] (0,0.6) node[left,font=\small,black,align=right] {pain $\nu$} -- (5.45,0.15);
\draw[->,very thick,blue!60!black] (0,-2.6) node[left,font=\small,black,align=right] {touch $\mu$} -- (2.9,-2.0);
\draw[->,thick,blue!60!black] (1.8,-2.35) -- (5.5,-0.35) node[pos=0.8,below right=-2pt,font=\scriptsize] {$w_{z\mu}$};
\draw[inh] (1.6,0.5) -- (3.35,-1.1) node[pos=0.5,left=1pt,font=\scriptsize] {$w_{q\nu}$};
\draw[inh] (Q) -- (Z) node[pos=0.5,above left=-1pt,font=\scriptsize] {$\beta$};
\draw[->,very thick] (Z) -- (8.2,0) node[right,font=\small,align=left] {to the brain: $z$};
\draw[bc] (6,2.1) node[above,font=\scriptsize,black,align=center] {from above: \nt{SERO}, \nt{NORA}, opioids} -- (Z.north) node[pos=0.45,right,font=\scriptsize,black] {gain $g$};
\end{tikzpicture}
\caption{The pain gate in the spinal cord according to the model~\eqref{eq:gate}. The output \nt{GLUT} neuron receives pain $\nu$ and weak excitation from touch $\mu$. The inhibitory interneuron (\nt{GABA} or \nt{GLYC}) is excited by touch and, by an assumption of gate theory, switched off by pain. The gain $g$ arrives from above over the broadcast channels.}
\label{fig:gate}
\end{figure}

Let us write this for rates, as in Chapter~\ref{sec:gaba}. The interneuron rate $q$ and the output rate $z$ are
\begin{empheq}[box=\keybox]{equation}
q \;=\; \bigl[w_{q\mu}\,\mu + q_0 - w_{q\nu}\,\nu\bigr]_+ ,
\qquad
z \;=\; \bigl[\,g\,(\nu + w_{z\mu}\,\mu) - \beta\,q\,\bigr]_+ ,
\label{eq:gate}
\end{empheq}
where $\nu$ is the pain input, $\mu$ the touch input, $w$ the connection weights (first index: to whom; second: from whom), $\beta$ the strength of inhibition, $q_0$ the interneuron's own background activity, and $g$ the gain set from above (more on it below). This is the veto~\eqref{eq:veto} of Chapter~\ref{sec:gaba}, ``pain, but not with touch,'' with one amendment taken from gate theory as an assumption: strong pain itself switches off the vetoing neuron.

We computed the model with $w_{q\mu}=1$, $q_0=0.2$, $w_{q\nu}=0.5$, $w_{z\mu}=0.3$, and $\beta=1.5$ (Fig.~\ref{fig:gatecurves}). Without touch, pain begins to get through at $\nu\approx0.18$. With touch the threshold rises to $0.86$, and at $\nu=1$ the output drops fourfold, from $1$ to $0.25$. But with strong pain, $\nu=2$, touch no longer changes anything: the gate is wide open. It is the same in life: rubbing a bruise helps, but not with a serious injury. Touch by itself, without pain, does not get through the gate: a touch does not raise the alarm.

\begin{figure}[t]
\centering
\begin{tikzpicture}[>=Stealth,x=5.2cm,y=1.35cm]
\draw[->,gray!70!black] (0,0) -- (2.12,0) node[right,font=\small,black] {pain $\nu$};
\draw[->,gray!70!black] (0,0) -- (0,4.3) node[left,font=\small,black] {$z$};
\foreach \t in {0,0.5,1,1.5,2}{\draw[gray!60!black] (\t,-0.05) -- (\t,0.05); \node[below,font=\scriptsize] at (\t,0) {\t};}
\foreach \f in {1,2,3,4}{\draw[gray!60!black] (-0.01,\f) -- (0.01,\f); \node[left,font=\scriptsize] at (0,\f) {\f};}
\draw[very thick,red!70!black] plot file {figures/pain_gate_notouch.dat};
\draw[very thick,blue!60!black] plot file {figures/pain_gate_touch.dat};
\draw[very thick,orange!85!black,densely dashed] plot file {figures/pain_gate_sens.dat};
\draw[very thick,red!70!black] (0.08,3.9) -- (0.2,3.9) node[right,font=\scriptsize,black] {without touch};
\draw[very thick,blue!60!black] (0.08,3.45) -- (0.2,3.45) node[right,font=\scriptsize,black] {with touch};
\draw[very thick,orange!85!black,densely dashed] (0.08,3.0) -- (0.2,3.0) node[right,font=\scriptsize,black] {after sensitization, $g=2$};
\end{tikzpicture}
\caption{Output of the gate~\eqref{eq:gate} as a function of pain intensity. Touch raises the threshold and dampens weak and moderate pain, but not strong pain. Sensitization (dashed line) doubles the gain: the same pain is transmitted twice as loudly, and the threshold drops.}
\label{fig:gatecurves}
\end{figure}

\section{A volume knob from above}

The gate is controlled not only from below. Descending pathways run to it from the brainstem and change the gain $g$ in~\eqref{eq:gate}: \nt{SERO}, \nt{NORA}, and the opioid peptides of Appendix~\ref{sec:othermediators}~\cite{Fields2004}. These are the same broadcast channels as in Chapters~\ref{sec:serocomputer}--\ref{sec:acet}, except that here their knob is the volume of the alarm. They can turn it both ways: the same descending circuits can both dampen pain and amplify it.

Why would the machine turn the alarm down? Because an alarm is an interrupt, and an interrupt is not always appropriate. An animal fleeing a predator should not stop because of an injured paw: run first, lick the wound later. Expecting relief also dampens pain, and part of this effect disappears when the opioid receptors are blocked~\cite{Levine1978}: an expectation computed in the brain descends through the opioid channel to the gate. This picture itself has been measured; exactly how the brain decides what the volume should be is a hypothesis.

\section{Sensitization: an alarm that learns}

After injury the output neurons of the gate become more sensitive: the same input produces a larger output, and the threshold drops~\cite{Woolf1983}. In the model~\eqref{eq:gate} this is a rise in the gain $g$, and its simplest description is a slow RC circuit charged by the pain itself:
\begin{empheq}[box=\keybox]{equation}
\tau_s\,\frac{\dd g}{\dd t} \;=\; -(g-1) + k_s\,z ,
\label{eq:sensitization}
\end{empheq}
where $\tau_s$ is the time it takes sensitivity to return to normal (it is longer than the pain itself, because sensitization persists after the damaging stimulus has stopped~\cite{Woolf1983}) and $k_s$ is the charging strength. While the tissue is damaged, the gate output holds $g$ above one, and everything that happens near the wound is transmitted more loudly (the dashed line in Fig.~\ref{fig:gatecurves}: at $g=2$ the threshold drops to $0.11$, and the output doubles). This is a sensible alarm setting: until the damage has healed, it is better to be safe than sorry.

To an engineer this is positive feedback with a slow leak: pain raises the gain, and the gain raises the pain. While the loop is weak, $g$ returns to normal once the wound has healed. If the loop is strong, however, the alarm can get stuck in the on state even after the damage is gone, like the group with strong feedback in Fig.~\ref{fig:bistable}. The model~\eqref{eq:sensitization} is a simplification; that central sensitization exists has been measured.

\section{Pain as teacher and as interrupt}

Besides raising the alarm, pain has two jobs in the machine.

\emph{Teacher.} Pain is a negative reward: in the error formula~\eqref{eq:ctd} it enters as $r<0$. Everything said in Chapter~\ref{sec:dofa} works here too: a predictor of pain starts to produce an error in advance, and the network learns to avoid what leads to pain. Experiments in humans show that brain activity during learning from pain follows exactly the temporal-difference error~\cite{Seymour2004}, and some \nt{DOPA} neurons also respond to unpleasant events (Chapter~\ref{sec:dofaalt}).

\emph{Interrupt.} Pain pulls attention away from the current task; this is its purpose, not a side effect~\cite{EcclestonCrombez1999}. In Chapter~\ref{sec:nora} the same ``interrupt'' role was discussed for the phasic \nt{NORA} burst. And the fastest response does not even reach the brain: pulling the hand away from something hot is wired directly in the spinal cord, using the same pair of muscles and the same veto on the opponent as in Fig.~\ref{fig:antagonists}.

\begin{keyblock}
\begin{proposition}[Alarm signal]
\label{prop:pain}
Pain is not a measurement but a high-priority alarm signal. It adapts far less than the measurement channels, travels over two wires (the fast one for ``where,'' the slow one for ``how bad''), passes through a gate that touch closes (and strong pain, by an assumption of gate theory, opens), and its volume is set from above by the broadcast channels. These mechanisms have been measured, except the opening of the gate by pain; the simple models~\eqref{eq:gate} and~\eqref{eq:sensitization} are simplifications.
\end{proposition}
\end{keyblock}

\section{Summary}

\begin{table}[t]
\centering
\footnotesize
\setlength{\tabcolsep}{4pt}
\renewcommand{\arraystretch}{1.15}
\begin{tabular}{>{\raggedright}p{3.0cm}>{\raggedright}p{4.9cm}>{\raggedright\arraybackslash}p{4.9cm}}
\toprule
What pain does & How & What is known \\
\midrule
raises the alarm & high threshold, weaker adaptation than touch & threshold measured; comparison of adaptation is an estimate \\
reports ``where'' and ``how bad'' & two wires of different speed~\eqref{eq:paintime} & measured \\
passes through a gate & veto through \nt{GABA}/\nt{GLYC}~\eqref{eq:gate} & gate measured; opening by pain is an assumption of the theory; the model is a simplification \\
changes the volume & descending \nt{SERO}, \nt{NORA}, opioids & measured; the rule for choosing the volume is a hypothesis \\
learns after injury & rise in gain~\eqref{eq:sensitization} & sensitization measured; the model is a simplification \\
teaches avoidance & negative reward in~\eqref{eq:ctd} & similarity to the temporal-difference error measured \\
interrupts work & capture of attention; withdrawal reflex & measured \\
\bottomrule
\end{tabular}
\caption{Pain as an alarm signal.}
\label{tab:pain}
\end{table}

Pain is assembled from parts we have already seen (Table~\ref{tab:pain}): \nt{GLUT} sensors, two wires, a veto through an inhibitory interneuron, a broadcast volume knob, a slow RC circuit, a prediction error, and an interrupt. Only one thing about it is new: it is the only input of the machine whose volume the machine sets itself, an alarm that its owner can both dampen and retune.

\section{Exercises}

\begin{exercise}[\lvlA]
\label{ex:14:1}
\emph{Two wires.} A signal comes from the foot, $L=1$~m. (a)~A volley travels along a bundle of wires of one type whose speeds fill the range~\eqref{eq:paintime}. How much does it spread out by the time it reaches the spinal cord, for the fast and the slow type? (b)~Pain waves along wires of $15$ and $1$~m/s are distinguishable if they differ by at least $0.1$~s. From what distance do they merge?
\end{exercise}

\begin{exercise}[\lvlB]
\label{ex:14:2}
\emph{When touch hurts.} The gate~\eqref{eq:gate} has the parameters of the text, and the gain $g$ grows with sensitization. Up to what $g$ does touch without pain fail to get through the gate at any strength $\mu$? At what $g$ does a touch of $\mu=1$ start to get through?
\end{exercise}

\begin{exercise}[\lvlB]
\label{ex:14:3}
\emph{Stability of sensitization.} The inputs are constant, and the gate output is linear in $g$: $z=gA-C$, $A=\nu+w_{z\mu}\mu$, $C=\beta q\ge0$. (a)~Find the equilibrium $g^*$ of the model~\eqref{eq:sensitization} and the condition for its stability. (b)~For $k_s=0.5$, find the pain intensity above which the model runs away, without touch and with touch $\mu=1$.
\end{exercise}

\begin{exercise}[\lvlA]
\label{ex:14:4}
\emph{A predictor of pain.} A cue at time $0$ predicts pain at time $T=2$~s; in~\eqref{eq:ctd} $r(t)=-P\,\delta(t-T)$ and $\tau_\gamma=2$~s. (a)~Find the area of the $\delta$ transient at the moment of the cue. (b)~An action at time $t_e=1$~s cancels the pain. What is the $\delta$ burst at that moment, and what will it teach the network?
\end{exercise}

\begin{exercise}[\lvlB]
\label{ex:14:5}
\emph{An alarm that latches.} Extend the gate of \texttt{models/\allowbreak pain\_gate.py} with the dynamics~\eqref{eq:sensitization} and saturation $z\le4$. Touch $\mu=1$ is always present, damage $\nu=2$ lasts a time $T$, then $\nu=0$. By simulation, find the smallest $T$ (in units of $\tau_s$) after which the alarm stays on, for $k_s=4$, $4.6$, $5$, $6$, and $8$.
\end{exercise}

\begin{exercise}[\lvlB]
\label{ex:14:6}
\emph{Designing the gate.} With $\beta=1.5$, $w_{z\mu}=0.3$, and $g=1$, choose $q_0$, $w_{q\nu}$, and $w_{q\mu}$ in~\eqref{eq:gate} so that the pain threshold is $0.2$ without touch and $1.0$ with touch $\mu=1$, and touch stops dampening pain at $\nu_\times=2.5$. Up to what gain $g$ does touch without pain fail to get through the gate?
\end{exercise}

\begin{exercise}[\lvlC]
\label{ex:14:7}
\emph{The volume knob.} Work yields value $V$ per unit time, working with damage $\nu$ costs $c\nu$, so interrupting pays when $\nu>V/c$; pain interrupts when $z>\theta=0.5$. (a)~What gain $g^*$ of the gate~\eqref{eq:gate} without touch gives this threshold for $V/c=0.25$, $0.5$, and $2$? (b)~From which signals of the book could the machine estimate $V$ and $c$?
\end{exercise}

\chapter{How the Machine Learns to Move}
\chaptermark{How the Machine Learns to Move}
\label{sec:motorlearning}

\section{Three teachers}

In Chapter~\ref{sec:muscles} we saw how the machine moves its muscles, and we stopped at a hypothesis: fast and precise movements require an internal model of the body that predicts the result of a command~\cite{WolpertGhahramani2000}. Where does this model come from? It has to be learned, and learned differently from the way we have learned so far.

The book has already had two teachers. The first is \emph{none at all}: the Hebb rule (Chapter~\ref{sec:dofa}) strengthens what often happens together and compresses the inputs. The second is \emph{reward}: \nt{DOPA} sends everyone a single number, ``better or worse than expected.'' Movement needs a third teacher, \emph{error}. The hand missed the cup by two centimeters to the left: that is not ``worse,'' it is a vector that tells each output where and by how much it was wrong. An engineer would call this supervised learning.

The difference between the second and the third teacher is the difference between one wire for everyone and a separate wire for each output. In Proposition~\ref{prop:broadcastcost}, learning from one shared number slowed down with each neuron whose noise affected the result. If instead each output $i$ receives its own error $e_i$, the weights can be changed by the simplest rule:
\begin{empheq}[box=\keybox]{equation}
\frac{\dd w_{ij}}{\dd t} \;=\; -\eta\,e_i(t)\,x_j(t) .
\label{eq:lms}
\end{empheq}
This is the least-mean-squares rule, long known in adaptive-filter engineering~\cite{WidrowHoff1960}: a weight changes in proportion to the product of its input and the error of its output and, on average, follows the gradient of the squared error. No exploratory noise is needed, as it was in Chapter~\ref{sec:dofa}: the direction of the correction arrives directly over the error wire. And the speed does not drop with the number of outputs, because other outputs' errors do not reach the synapse.

\begin{keyblock}
\begin{proposition}[Three teachers]
\label{prop:teachers}
The Hebb rule teaches, without a teacher, what is frequent; the \nt{DOPA} signal teaches, with one number for the whole network, what pays off, and more slowly with each neuron; an error wire to each output teaches, by the rule~\eqref{eq:lms}, what is accurate, at a speed that does not depend on the number of outputs. The price of the third way is a separate wire for each output and someone who knows the right answer.
\end{proposition}
\end{keyblock}

\section{The error wire}

Who can know the right answer for a movement? The sensors themselves. If the hand has drifted to the left of the target, the eye and the stretch sensors of Chapter~\ref{sec:sensors} report it. What is needed is a circuit that delivers this error to the right outputs.

The brain has a region that seems to be built for exactly this~\cite{Marr1969,Albus1971}. Its circuitry is unusually regular, almost crystalline, and the rule~\eqref{eq:lms} is easy to recognize in it (Fig.~\ref{fig:errorwire}).

\emph{Input expander.} Signals about the command and about the state of the body arrive at a huge layer of small \nt{GLUT} neurons. There are about seventy billion of them, more than in all the rest of the brain combined~\cite{Azevedo2009}. Each of them mixes several inputs, and the signal is spread out into a very long vector. The trick is familiar to an engineer: raise the dimension of the input, and in the new space almost any desired dependence can be obtained by a simple weighted sum~\cite{Cover1965}.

\emph{Output neurons.} The weighted sum is computed by about thirty million large output neurons~\cite{Andersen1992cereb}, each with on the order of a hundred thousand inputs from the expander. And these are \nt{GABA} neurons: the result of learning is passed on not by excitation but by inhibition, as in the veto of Chapter~\ref{sec:gaba}.

\emph{The error wire.} Each output neuron receives one more input, a special one: exactly one \nt{GLUT} wire, which fires rarely, about once per second, but each time evokes a powerful burst in the output neuron~\cite{Ito2001}. This is $e_i$: the probability of the burst depends on the direction of the movement error~\cite{Herzfeld2018}. Coincidence of the error burst with activity of an input weakens that input, which is exactly the term $-\eta\,e_i x_j$ in~\eqref{eq:lms}.

\begin{figure}[t]
\centering
\begin{tikzpicture}[>=Stealth,
  gran/.style={circle,draw,thick,fill=blue!6,minimum size=0.32cm,inner sep=0pt},
  outn/.style={circle,draw,very thick,fill=red!10,minimum size=1.1cm,inner sep=0pt,font=\scriptsize},
  inh/.style={-{Bar[width=7pt]},very thick,red!70!black}]
% inputs
\foreach \y/\lab in {1.6/{command},0.4/{body state}}{
  \draw[->,thick,blue!60!black] (-0.6,\y) node[left,font=\scriptsize,black] {\lab} -- (0.35,\y);}
% expansion layer
\foreach \i in {0,...,11}{\node[gran] (g\i) at ({0.8+0.28*mod(\i,2)},{-0.3+0.23*\i}) {};}
\foreach \i in {0,...,11}{\pgfmathsetmacro{\yy}{mod(\i,3)>0 ? 1.6 : 0.4}\draw[gray!60!black] (0.35,\yy) -- (g\i);}
\node[font=\scriptsize,align=center] at (0.95,-1.05) {expander\\$\sim7\cdot10^{10}$ \nt{GLUT}};
% parallel wires to the output neuron
\foreach \i in {0,...,11}{\draw[->,gray!70!black] (g\i.east) -- (4.1,{1.0+0.07*(\i-5.5)});}
\node[outn] (P) at (4.6,1.0) {\nt{GABA}};
\node[font=\scriptsize,align=center,above=2pt] at (P.north) {output neuron\\$\sim10^5$ inputs $x_j$, weights $w_{ij}$};
\draw[inh] (P) -- (6.8,1.0) node[right,font=\scriptsize,align=left] {correction\\to movement};
% error wire
\draw[->,very thick,red!70!black,densely dashed] (4.6,-1.4) node[below,font=\scriptsize,black,align=center] {one error wire $e_i$\\\nt{GLUT}, $\sim1$ per s} -- (P.south);
\end{tikzpicture}
\caption{The supervised-learning circuit as the brain apparently implements it. A huge layer of \nt{GLUT} neurons spreads the command and body state out into a long vector $x_j$; the output \nt{GABA} neuron sums it with weights $w_{ij}$. The single error wire brings $e_i$; coincidence of the error with an active input weakens the weight of that input, as in the rule~\eqref{eq:lms}.}
\label{fig:errorwire}
\end{figure}

One difficulty is familiar from Chapter~\ref{sec:dofa}: the error arrives later than the input that caused it. A miss is visible tens to hundreds of milliseconds after the command. So the synapse must remember that it was active; an eligibility trace is needed, as in the rule~\eqref{eq:threefactor}. And the length of this trace, judging by experiments, is tuned to the feedback delay of each task: where the error arrives after a hundred milliseconds, the input that was active a hundred milliseconds before it is weakened most~\cite{Suvrathan2016}.

\begin{keyblock}
\begin{proposition}[The error wire]
\label{prop:errorwire}
In the circuit shown in Fig.~\ref{fig:errorwire}, each output neuron receives exactly one error wire, and coincidence of the error with an input weakens that input. This is the least-mean-squares rule~\eqref{eq:lms} applied to an input expanded to a huge dimension. The structure of the circuit and the weakening on coincidence have been measured; that the whole circuit works as supervised learning is the leading hypothesis.
\end{proposition}
\end{keyblock}

\section{The internal model as an adaptive filter}

What does such a circuit learn? The most natural answer is prediction. Let the input be a command $u(t)$ passed through a set of filters with different time constants, as in the expander: $x_j(t)=(g_j*u)(t)$. The output is their weighted sum, a prediction of what the sensors will report:
\begin{empheq}[box=\keybox]{equation}
\hat y(t) \;=\; \sum_j w_j\,(g_j*u)(t), \qquad e(t) \;=\; \hat y(t) - y(t) ,
\label{eq:adaptivefilter}
\end{empheq}
and the weights learn by~\eqref{eq:lms}. This is an \emph{adaptive filter}: a circuit that tunes its own transfer function so as to reproduce an unknown system~\cite{Fujita1982}. Here the unknown system is one's own body, with its mass, elasticity, and delays. The learned filter is the internal model of Chapter~\ref{sec:muscles}: it says what the sensors will report without waiting for them.

Such a model is useful twice over. First, its prediction can be used in place of the delayed feedback, and then the stability condition~\eqref{eq:delaystable} no longer ties the controller's hands: corrections can be made quickly, from the model. Second, the difference between what was predicted and what was received is a signal of the \emph{unexpected}. Everything the machine did itself, the model predicted and subtracted; what remains is what the world did. That, according to one hypothesis, is why we cannot tickle ourselves~\cite{Blakemore1998}.

\section{When the world changes: fast and slow}

Let us test the circuit with an experiment that is easy to set up. A person reaches for a target, and the world changes unexpectedly: for example, a sideways force acts on the arm. The first movements miss, and then the miss shrinks. If the force is removed, the arm at first misses to the \emph{other} side: the model has learned the force and keeps compensating for it. This is direct evidence that a model has been learned, not just that things ``started to work.''

Experiments also show something subtler: two processes learn at the same time, a fast one that learns quickly and forgets quickly, and a slow one that learns slowly and remembers for a long time~\cite{Smith2006}. The corrections are updated after each movement, event by event:
\begin{empheq}[box=\keybox]{equation}
e = p - (x_f + x_s), \qquad x_f \leftarrow A_f\,x_f + B_f\,e, \qquad x_s \leftarrow A_s\,x_s + B_s\,e ,
\label{eq:tworate}
\end{empheq}
where $p$ is the perturbation, $x_f$ and $x_s$ are the learned corrections of the two processes, $A$ is how much of a correction is retained until the next movement, and $B$ is how strongly it learns from the error. The fast process has a large $B$ and an $A$ noticeably below one; the slow process is the opposite.

\begin{figure}[t]
\centering
\begin{tikzpicture}[>=Stealth,x=0.034cm,y=1.9cm]
\draw[->,gray!70!black] (0,0) -- (355,0) node[right,font=\small,black] {movement};
\draw[->,gray!70!black] (0,-1.15) -- (0,1.25);
\foreach \v in {-1,1}{\draw[gray!60!black] (-3,\v) -- (3,\v); \node[left,font=\scriptsize] at (0,\v) {$\v$};}
\foreach \n in {100,200,300}{\draw[gray!60!black] (\n,-0.04) -- (\n,0.04); \node[below,font=\scriptsize] at (\n,0) {\n};}
\draw[thin,gray!70!black] plot file {figures/tworate_p.dat};
\draw[thick,orange!85!black] plot file {figures/tworate_fast.dat};
\draw[thick,blue!60!black] plot file {figures/tworate_slow.dat};
\draw[very thick,black] plot file {figures/tworate_net.dat};
\node[font=\scriptsize,gray!50!black] at (120,1.12) {perturbation $p$};
\node[font=\scriptsize,black,anchor=west] at (238,0.52) {sum: rebound};
\node[font=\scriptsize,blue!60!black,anchor=west] at (150,0.39) {slow $x_s$};
\node[font=\scriptsize,orange!85!black,anchor=west] at (60,0.08) {fast $x_f$};
\draw[densely dashed,gray!60!black] (226,-1.15) -- (226,1.25);
\node[font=\scriptsize,align=center,anchor=south west] at (228,-1.1) {no more\\errors};
\end{tikzpicture}
\caption{Two learning processes by the model~\eqref{eq:tworate}, $A_f=0.92$, $B_f=0.1$, $A_s=0.996$, $B_s=0.02$. Two hundred movements with perturbation $p=1$, then six with the reverse, $p=-1$, until the sum returns to zero. After the dashed line errors stop arriving, and the sum returns by itself toward the old correction: the fast process has forgotten the reverse perturbation, while the slow one remembers the original.}
\label{fig:tworate}
\end{figure}

The model~\eqref{eq:tworate} explains an experiment that is hard to understand without it (Fig.~\ref{fig:tworate}). In our calculation, after two hundred movements with the perturbation the sum of the corrections reaches $0.84$, with most of it held by the slow process, $0.63$. Then six movements with the reverse perturbation bring the sum back to zero: the fast process has gone negative, $-0.53$, while the slow one is still positive, $0.46$. If errors now stop being reported, the fast process forgets its correction within tens of movements and the slow one much more slowly, and the sum rises \emph{by itself} to $0.37$ within forty movements: the old skill returns without any training. Such spontaneous recovery has been measured in humans; a model with a single process cannot produce it, since without errors it simply decays to zero.

Why two processes? A plausible engineering answer comes from the optimal filter of Chapter~\ref{sec:acet}. The body changes for different reasons at different rates: muscles tire within minutes, and grow and weaken over months. If disturbances come both fast and slow, the best estimate consists of two filters with different time constants, each catching its own kind~\cite{Kording2007}. The fast process is a temporary correction, ``the arm is tired''; the slow one is ``the arm has changed.'' This is still a hypothesis, but it explains why a skill learned in one day is partly lost by the next and why it is learned faster the second time.

\section{Summary}

\begin{table}[t]
\centering
\footnotesize
\setlength{\tabcolsep}{4pt}
\renewcommand{\arraystretch}{1.15}
\begin{tabular}{>{\raggedright}p{3.0cm}>{\raggedright}p{4.6cm}>{\raggedright\arraybackslash}p{5.2cm}}
\toprule
What the circuit does & How & What is known \\
\midrule
learns with a teacher & an error wire to each output, the rule~\eqref{eq:lms} & circuit structure and weakening on coincidence measured; supervised learning is the leading hypothesis \\
expands the input & seventy billion \nt{GLUT} neurons & number of neurons measured; role of expansion is theory \\
accounts for the error delay & a trace tuned to the delay & measured \\
builds an internal model & adaptive filter~\eqref{eq:adaptivefilter} & well-founded hypothesis \\
learns at two rates & fast and slow processes~\eqref{eq:tworate} & aftereffect and spontaneous recovery measured; two processes are a model \\
\bottomrule
\end{tabular}
\caption{How the machine learns to move.}
\label{tab:motorlearning}
\end{table}

The machine now has three teachers (Table~\ref{tab:motorlearning}). Hebb compresses what is frequent; \nt{DOPA} teaches, with one number, to choose what pays off; the error wire teaches accuracy, and teaches it quickly, because each output gets its own error. The brain apparently uses all three, each where it is cheapest: the broadcast signal for a few decisions, and separate wires for precise tuning of the body, where the sensors themselves report the right answer.

\section{Exercises}

\begin{exercise}[\lvlA]
\label{ex:15:1}
\emph{Capacity of the error-wire circuit.} The circuit has $3\cdot10^7$ output neurons with $10^5$ inputs each, each with its own error wire ($1$~spike/s). A neuron with $n$ inputs can learn any labeling of $\approx2n$ vectors~\cite{Cover1965}. (a)~How many associations does the circuit store? (b)~Using~\eqref{eq:spikeentropy} with $\Delta t=10\ms$, estimate the shortest time to fill one neuron's capacity.
\end{exercise}

\begin{exercise}[\lvlB]
\label{ex:15:2}
\emph{Speed of the least-mean-squares rule.} The modes of the rule~\eqref{eq:lms} decay at rates $\eta\lambda_k(C)$. For five filters~\eqref{eq:adaptivefilter} driven by white noise, $C_{jk}=2\sqrt{\tau_j\tau_k}/(\tau_j+\tau_k)$. Find the ratio of the fastest to the slowest rate if the $\tau_j$ grow by a factor of $2$ ($10$--$160\ms$) and by a factor of $4$.
\end{exercise}

\begin{exercise}[\lvlB]
\label{ex:15:3}
\emph{Analysis of the two processes.} Write the model~\eqref{eq:tworate} with the parameters of Fig.~\ref{fig:tworate} as $\mathbf x_{n+1}=M\mathbf x_n+\mathbf b\,p$. (a)~From the eigenvalues of $M$, find the time constants in movements. (b)~Find the steady-state $x_f$, $x_s$, and their sum for constant $p=1$.
\end{exercise}

\begin{exercise}[\lvlB]
\label{ex:15:4}
\emph{Simulation: faster the second time.} Using the model~\eqref{eq:tworate} with the parameters of \texttt{models/\allowbreak motor\_learning.py}, find the number of movements until $x_f+x_s\ge0.8$ at $p=1$: (a)~for an untrained machine; (b)~after $200$ movements with $p=1$ and a reverse perturbation until the sum is zero, as in Fig.~\ref{fig:tworate}. Can a single-process model speed up like this?
\end{exercise}

\begin{exercise}[\lvlB]
\label{ex:15:5}
\emph{A controller with an internal model.} A plant $\dd x/\dd t=ku$ is observed with delay $d$; the controller $u=-G\hat x$ predicts $\hat x(t)=x(t-d)+\hat k\int_{t-d}^{t}u\,\dd s$. (a)~With $\hat k=k=1$ and $d=150\ms$, find $G$ for a time constant of $20\ms$ and compare it with the limit~\eqref{eq:delaystable}. (b)~Is the circuit stable with $\hat k=0.4k$?
\end{exercise}

\begin{exercise}[\lvlB]
\label{ex:15:6}
\emph{An adaptive filter learns a muscle.} The plant is the twitch~\eqref{eq:twitch} of unit area, $T_c=50\ms$; the filters~\eqref{eq:adaptivefilter} are RC circuits with $\tau_j=12.5$, $25$, $50$, $100$, $200\ms$; the input is a new normal random number every $10\ms$. How many seconds does the rule~\eqref{eq:lms} with $\eta=50$ and $200$ take to reduce the error to $0.5\%$? Why are the weights still far from optimal even after $300$~s?
\end{exercise}

\begin{exercise}[\lvlC]
\label{ex:15:7}
\emph{Two processes as an optimal filter.} The perturbation is the sum of processes $p^{f,s}_{n+1}=A_{f,s}\,p^{f,s}_n+w^{f,s}_n$ with noise variances $q_f$, $q_s$, and the error is observed with noise of variance $R$; the Kalman filter gives~\eqref{eq:tworate}. For what $q_f/R$ and $q_s/R$ do $A_f=0.92$, $A_s=0.996$ yield $B_f=0.1$, $B_s=0.02$? How do $B_f$ and $B_s$ change if $q_s$ is quadrupled?
\end{exercise}

\chapter{The Second Output: How the Brain Controls Body Chemistry}
\chaptermark{The Chemical Output}
\label{sec:chemoutput}

\section{The fast and the slow outputs}

In Chapter~\ref{sec:muscles} we saw the machine's fast output: the muscles. But it also has a second, slow one. The brain keeps body temperature, heart rate, and the amounts of water and sugar within the right limits, and it prepares the body for running or for sleep. To do this it does not move joints; it releases chemical signals. There are two routes. The first is the nerves that run directly to the internal organs: the heart, blood vessels, gut, and glands. The second is the blood: some neurons and glands release their substance not into the cleft to a neighbor but straight into the bloodstream.

The blood is the most broadcast bus of any we have met in this book. The broadcast channels of Chapters~\ref{sec:serocomputer}--\ref{sec:acet} sent a signal to large regions of the brain; the blood sends it to every cell of the body. Blood circulates through the body in about a minute, so a message reaches everyone within tens of seconds and lasts minutes to hours, until the substance is broken down. Such a broadcast has no address at all. As in Proposition~\ref{prop:receptor}, the receiver chooses the meaning of the message: only cells that have a receptor for the substance respond.

\section{Gas and brake: two wires to the organs}

The best example is the heart. The heart has its own pacemaker, like the \nt{SERO} neuron of Chapter~\ref{sec:serocomputer}: it sets the rhythm by itself, and with all nerves cut the heart beats about a hundred times a minute. The brain does not generate the rhythm; it only adjusts it, and it does so over two wires of opposite sign (Fig.~\ref{fig:gasbrake}). One carries \nt{NORA}: this is the \emph{gas}, which speeds the rhythm up. The other carries \nt{ACET}: this is the \emph{brake}, which slows it down. Roughly,
\begin{empheq}[box=\keybox]{equation}
f \;\approx\; f_0 \;+\; a_S\,S \;-\; a_P\,P ,
\label{eq:gasbrake}
\end{empheq}
where $f_0\approx100$ beats per minute is the pacemaker's intrinsic rhythm, and $S$ and $P$ are the activity of the gas and the brake. At rest the brake is on, and the heart usually beats about sixty to seventy times a minute; under load the brake is released and the gas is pressed.

\begin{figure}[t]
\centering
\begin{tikzpicture}[>=Stealth,
  blk/.style={draw,very thick,rounded corners=2pt,align=center,font=\small},
  pace/.style={circle,draw,very thick,double,double distance=1.2pt,minimum size=1.8cm,inner sep=0pt,font=\scriptsize,fill=red!8,align=center},
  inh/.style={-{Bar[width=7pt]},thick,red!70!black}]
\node[blk,fill=blue!5,text width=1.6cm] (B) at (0,0) {brain\\command};
\node[pace] (H) at (6.2,0) {pace-\\maker};
\draw[->,very thick,orange!85!black] (B.north east) to[bend left=20] node[above,font=\scriptsize,black,align=center] {\nt{NORA}: gas, seconds} (H.north west);
\draw[inh,very thick,blue!60!black] (B.south east) to[bend right=20] node[below,font=\scriptsize,black,align=center] {\nt{ACET}: brake, fractions of a second} (H.south west);
\draw[->,very thick] (H) -- (8.4,0) node[right,font=\small] {rate $f$};
\node[blk,fill=orange!10,text width=1.6cm] (A) at (0,-2.6) {\nt{ADRE}\\into blood};
\draw[->,thick,orange!85!black] (B) -- (A);
\foreach \y/\lab in {-2.0/{heart},-2.6/{vessels},-3.2/{liver, muscles}}{
  \draw[->,thick,densely dashed,orange!85!black] (A.east) -- (4.3,\y) node[right,font=\scriptsize,black] {\lab};}
\end{tikzpicture}
\caption{Two wires of opposite sign to the heart's pacemaker: the \nt{NORA} gas is slow, the \nt{ACET} brake is fast. Below, the same gas over the broadcast bus: \nt{ADRE} cells release adrenaline into the blood, and every organ with a matching receptor receives the signal (dashed arrows).}
\label{fig:gasbrake}
\end{figure}

This is the two-wire code of Chapter~\ref{sec:glutcomputer} and the muscle pair of Chapter~\ref{sec:muscles}, except that now the wires control not a joint but the tempo of a pacemaker. And the wires have different speeds. Both act as low-pass filters, but the brake passes changes several times faster than the gas~\cite{Berger1989}. \nt{ACET} takes a short path: the protein bound to the receptor opens a potassium channel by itself, with no second messenger inside the cell~\cite{Logothetis1987gbg}, whereas \nt{NORA} acts through a chain of reactions inside the cell. An engineer will recognize the trick of two actuators of different speed: the fast one does precise instant trimming, the slow one makes large, lasting changes.

This is also where \nt{ADRE} finds its role; Table~\ref{tab:types} said that its role in the brain is unclear. Outside the brain it is clear. Some cells of the gas pathway do not send a wire to an organ but release adrenaline straight into the blood. It is the same gas, but on the broadcast bus: within tens of seconds it reaches the heart, vessels, liver, and muscles all at once and lasts for minutes. The addressed \nt{NORA} gas trims one organ; the broadcast \nt{ADRE} gas switches the whole body into running mode.

\section{A cascade with feedback}

Many hormones are released not by a single gland but by a cascade. A small group of neurons at the base of the brain releases a first signal into the blood; it makes an intermediate gland release a second one; that one makes the final gland release the hormone that actually acts on the body. And the final hormone comes back to the first stages and inhibits them. The result is a three-stage amplifier enclosed in negative feedback: a level regulator, like the \nt{SERO} system in Eq.~\eqref{eq:feedback}.

We describe each stage as an RC circuit with time constant $\tau$ and gain $g$, and the feedback by a coefficient $k$:
\begin{equation}
\tau\,\frac{\dd x_1}{\dd t} = -x_1 + g\bigl[u - k\,x_3\bigr]_+ ,\qquad
\tau\,\frac{\dd x_2}{\dd t} = -x_2 + g\,x_1 ,\qquad
\tau\,\frac{\dd x_3}{\dd t} = -x_3 + g\,x_2 ,
\label{eq:cascade}
\end{equation}
where $u$ is the brain's command and $x_3$ is the level of the final hormone. The loop gain is $K=g^3k$. At equilibrium $x_3=g^3u/(1+K)$, and, as with any feedback regulator, a large loop gain makes the level insensitive to the spread in the stage gains. But each stage shifts the phase: at the frequency $\omega=\sqrt3/\tau$ the three identical stages together give a half-turn shift and an attenuation by a factor of eight. Hence the classical result of control theory:
\begin{empheq}[box=\keybox]{equation}
K \;<\; 8 ,
\label{eq:cascadestable}
\end{empheq}
otherwise the regulator starts to oscillate, with a period of about $2\pi\tau/\sqrt3\approx3.6\,\tau$.

\begin{keyblock}
\begin{proposition}[Accuracy versus stability]
\label{prop:cascade}
For a three-stage cascade with proportional feedback the level error is $1/(1+K)$, and the cascade is stable only for $K<8$. Such a regulator therefore cannot hold the level more accurately than about ten percent; trying to raise the gain turns it into an oscillator.
\end{proposition}
\end{keyblock}

We checked this on model~\eqref{eq:cascade} (Fig.~\ref{fig:cascade}). At $K=4$ the level oscillates a little after switch-on and settles. At $K=7$ it also settles, but slowly. At $K=12$ the oscillations neither decay nor grow without bound: a stage cannot output a negative signal, and the cascade enters steady pulsations between $0.83$ and $1.24$ of the mean level with a period of $3.7$--$3.9\,\tau$.

\begin{figure}[t]
\centering
\begin{tikzpicture}[>=Stealth,x=0.3cm,y=1.2cm]
\draw[->,gray!70!black] (0,0) -- (41.5,0) node[right,font=\small,black] {$t/\tau$};
\draw[->,gray!70!black] (0,0) -- (0,2.8) node[left,font=\small,black] {$x_3/x_3^*$};
\foreach \t in {0,10,20,30,40}{\draw[gray!60!black] (\t,-0.05) -- (\t,0.05); \node[below,font=\scriptsize] at (\t,0) {\t};}
\foreach \f in {1,2}{\draw[gray!60!black] (-0.3,\f) -- (0.3,\f); \node[left,font=\scriptsize] at (0,\f) {\f};}
\draw[densely dashed,gray!60!black] (0,1) -- (40,1);
\draw[thick,blue!60!black] plot file {figures/cascade_K4.dat};
\draw[thick,red!70!black] plot file {figures/cascade_K12.dat};
\node[font=\scriptsize,blue!60!black,anchor=west] at (16,0.55) {$K=4$: settles};
\node[font=\scriptsize,red!70!black,anchor=west] at (16,1.75) {$K=12$: pulsates};
\end{tikzpicture}
\caption{A three-stage cascade with feedback, Eq.~\eqref{eq:cascade}; the level of the final hormone as a fraction of its equilibrium value $x_3^*$. With a loop gain below eight the level settles; above eight the cascade generates steady pulsations on its own.}
\label{fig:cascade}
\end{figure}

Real cascade hormones do come out in pulses: the level of the stress hormone, for example, oscillates with a period of about an hour. According to one hypothesis, these pulsations are generated by the delayed feedback between the stages itself, and no separate rhythm generator is needed~\cite{Walker2010}. This is the same cause of oscillation as in the \nt{GLUT}--\nt{GABA} loop of Proposition~\ref{prop:ei} and in the muscle stretch loop of condition~\eqref{eq:delaystable}: negative feedback that comes late.

\section{The pulse code of hormones}

Pulses are not only a side effect; they can also be a code. The neurons that control reproduction release their signal into the blood in short pulses about once an hour. If the receiving gland is given the same signal continuously instead of in pulses, it does not work at full strength as it does with pulses; it falls silent. If hourly pulses are resumed, it works again~\cite{Belchetz1978}. So the receiver reads not the level but the \emph{frequency} of the pulses.

For an engineer there is no mystery here. A receptor that tires under prolonged exposure and stops responding is a high-pass filter, like the depressing synapse~\eqref{eq:depression} of Chapter~\ref{sec:code}. It suppresses a constant signal and passes changes. Such a receiver can be told anything only through pulses, and the information moves from the level into the frequency and shape of the pulses. Moreover, different pulse frequencies act differently on different cells of the same gland, so one chemical line can carry more than one quantity, just as the single \nt{DOPA} wire of Chapter~\ref{sec:dofaalt} carried a fast and a slow component.

\section{The daily rhythm: a slow mode generator}

The slowest rhythm of the body is the daily one. It is generated by delayed negative feedback inside cells: genes produce proteins that, with a delay of several hours, suppress their own production~\cite{TakahashiJS2017}. Delayed negative feedback once again, the fourth time in this book, and once again a rhythm, this time with a period of about a day. The master generator is a small group of tens of thousands of neurons at the base of the brain, whose rhythm synchronizes the other cells of the body.

In humans the intrinsic period of this generator is not exactly a day but $24.18$ hours on average~\cite{Czeisler1999}. Every day it is adjusted by light arriving from the sensors of the eye (Chapter~\ref{sec:sensors}). To an electronics engineer this is a \emph{phase-locked loop}: an oscillator with its own frequency $\omega_0$ locks to an external signal of frequency $\omega$, and the phase difference $\varphi$ obeys
\begin{empheq}[box=\keybox]{equation}
\frac{\dd\varphi}{\dd t} \;=\; (\omega-\omega_0) \;-\; K_\varphi\,\sin\varphi .
\label{eq:pll}
\end{empheq}
Lock holds if $|\omega-\omega_0|<K_\varphi$. A generator with a period of $24.18$ hours must shift by about eleven minutes every day; morning light is usually enough for such a shift. In the polar night, or in a cave without light, there is no locking, and the rhythm drifts to its own period.

The daily generator is not a computer's clock generator. It does not count computation steps and does not synchronize neuron spikes. It switches \emph{modes}: sleep and wakefulness, hormone levels, body temperature, readiness to eat. It is a slow broadcast register of the same kind as the registers of Chapters~\ref{sec:serocomputer}--\ref{sec:acet}, except that its value changes on a schedule set to the sun.

\begin{keyblock}
\begin{proposition}[The chemical output]
\label{prop:chemoutput}
The machine's slow output is built from the same parts as the fast one. The organs receive two wires of opposite sign and different speed, the \nt{NORA} gas and the \nt{ACET} brake, while the whole organism at once receives broadcast signals through the blood, including \nt{ADRE} adrenaline. Hormone levels are held by cascades with negative feedback, whose accuracy is limited by stability; some hormones are coded by pulse frequency; the daily rhythm is a slow oscillator phase-locked to light. The layout of the pathways and the experiments with pulses and with the period are measured; the explanation of cascade pulsations by the feedback itself is a hypothesis.
\end{proposition}
\end{keyblock}

\section{Summary}

\begin{table}[t]
\centering
\footnotesize
\setlength{\tabcolsep}{4pt}
\renewcommand{\arraystretch}{1.15}
\begin{tabular}{>{\raggedright}p{3.1cm}>{\raggedright}p{4.8cm}>{\raggedright\arraybackslash}p{4.9cm}}
\toprule
What the chemical output does & How & What is known \\
\midrule
trims the organs & \nt{NORA} gas and \nt{ACET} brake~\eqref{eq:gasbrake} & measured; the brake is faster than the gas \\
switches the whole body into running mode & \nt{ADRE} adrenaline into the blood & measured \\
holds hormone levels & cascade with feedback, $K<8$~\eqref{eq:cascadestable} & cascades and feedback measured; pulsations from the loop itself are a hypothesis \\
transmits by pulse frequency & a tiring receptor as a high-pass filter & dependence on pulses measured \\
switches modes over the day & oscillator phase-locked to light~\eqref{eq:pll} & period and light locking measured \\
\bottomrule
\end{tabular}
\caption{How the machine controls body chemistry.}
\label{tab:chemoutput}
\end{table}

The machine has two outputs (Table~\ref{tab:chemoutput}). The fast one is the muscles: milliseconds, addressed, to each joint. The slow one is body chemistry: seconds, minutes, and days, over two wires to the organs and through the blood to everyone at once. Both outputs use the same parts (two wires of opposite sign, pacemakers, feedback and its delay), and they are the same parts from which the machine is built inside.

\section{Exercises}

\begin{exercise}[\lvlA]
\label{ex:16:1}
\emph{Blood as a filter.} A hormone is cleared from the blood with half-life $t_{1/2}=10$~min and released in short pulses every $T=60$~min. By what factor does the peak concentration exceed the trough, and by what factor is the fundamental harmonic of the pulses attenuated? For what $t_{1/2}$ is the peak only twice the trough?
\end{exercise}

\begin{exercise}[\lvlA]
\label{ex:16:2}
\emph{Daily locking.} The intrinsic period of generator~\eqref{eq:pll} is $24.18$~h, and light shifts the phase by at most an hour per day: $K_\varphi=2\pi/24$~rad/day. Find the steady phase difference $\varphi^*$ in minutes and the relaxation time. How many days after a $\pm6$~h jump of the external signal until the mismatch is below an hour?
\end{exercise}

\begin{exercise}[\lvlB]
\label{ex:16:3}
\emph{Accuracy versus stability.} The linearized cascade~\eqref{eq:cascade} is lengthened to $n$ identical stages. Find the critical gain $K_c(n)$ and the smallest achievable error $1/(1+K_c)$ for $n=3$ and $n=6$. What does this error tend to as $n\to\infty$?
\end{exercise}

\begin{exercise}[\lvlB]
\label{ex:16:4}
\emph{Gas and brake.} In~\eqref{eq:gasbrake} the gas and brake contributions are outputs of RC circuits with $\tau_S=2$~s and $\tau_P=0.4$~s, with commands of at most $80$ and $60$ beats per minute. The rhythm is $f_0=100$; at rest the brake is $35$, the gas $0$, and $f=65$. How fast can $f=120$ be reached? And if the brake is kept on and only the gas is used?
\end{exercise}

\begin{exercise}[\lvlB]
\label{ex:16:5}
\emph{Simulation: a cascade with delay.} Add a delay $d$ to the feedback of the function \texttt{cascade} from \texttt{models/\allowbreak chem\_models.py} (copy it; do not run the file): the first stage sees $x_3(t-d)$. Find $K_c$ and the period at the boundary for $d/\tau=0$, $0.5$, $1$, $2$, and check by simulation at $0.9K_c$ and $1.1K_c$.
\end{exercise}

\begin{exercise}[\lvlB]
\label{ex:16:6}
\emph{A receiver that reads pulses.} The receptor tires: $y=[c-a]_+$, $\tau_a\,\dd a/\dd t=c-a$, $\tau_a=30$~min. The hormone arrives in $6$-min pulses with period $T$ at the same mean dose $\bar c$. Find the mean response for $T=15$, $60$, $120$~min and $T\to\infty$. What is it at a constant concentration $\bar c$?
\end{exercise}

\begin{exercise}[\lvlC]
\label{ex:16:7}
\emph{Oscillator or feedback?} Propose an experiment that distinguishes pulsations generated by the delayed feedback of cascade~\eqref{eq:cascade} from a separate pacemaker. Can the hypotheses be told apart by the shift of the period if $K$ can be changed only by a factor of one and a half and the period is measured to $5\%$?
\end{exercise}

\chapter{Debugging the Machine}
\chaptermark{Debugging the Machine}
\label{sec:debugging}
\begin{chaptergoals}
\item why a probe hearing a thousand out of $8.6\cdot10^{10}$ neurons can still decode movement;
\item why an implant must detect spikes on the chip and send only events, not the raw stream;
\item how a decoder reads group rates, and why interfaces extract only a few bits per second;
\item why writing into the brain is harder than reading, and what the probe cannot see.
\end{chaptergoals}

\section{How to debug a machine without a schematic}

An engineer handed a working board without a schematic debugs it with four techniques. Put on a \emph{probe} and watch what goes along a wire. \emph{Inject a test signal} and watch what changes. \emph{Disable} a piece of the circuit and watch what stops working. And \emph{read the control registers} to find out what mode the board is in. This whole book is an attempt to read the brain's schematic, and in this chapter we look at which of the four techniques engineers can already use and what the previous chapters tell them to expect.

The most visible example of such debugging today is brain–computer interfaces: devices that read neuron spikes and turn them into commands for external machines or, conversely, feed signals into the brain from outside. Most of the chapter is devoted to them, and above all to what the company Neuralink is doing.

\section{The probe: what an electrode hears}

The brain's probe is a thin electrode inserted into the tissue. It hears the spikes of neurons lying within a radius of about a hundred micrometers, usually one to several of them. A spike on the electrode is a blip of tens to hundreds of microvolts lasting about a millisecond, and its shape lets one tell neighboring neurons apart. Large \nt{GLUT} neurons are heard best: they are the majority in the cortex (Table~\ref{tab:types}), and their currents are stronger.

The main difficulty is obvious at once. A good probe today has hundreds or thousands of electrodes, while the brain has $8.6\cdot10^{10}$ neurons. We are eavesdropping on about one hundred-millionth of the machine. Why can anything be understood from such a tiny sample at all? Chapter~\ref{sec:code} gave the answer. Neighboring neurons are noisy in a correlated way, and averaging over a group hits a limit (Proposition~\ref{prop:pooling}): a hundred neurons carry almost as much as a thousand. Moreover, the task solved by the motor cortex is low-dimensional: the arm has few joints (Chapter~\ref{sec:muscles}), and the activity of the hundreds of neurons controlling it fits into a space of only one or two dozen shared variables~\cite{Gallego2017}. A probe on a hundred neurons hears almost all of these variables. If the brain stored an independent message in every neuron, such eavesdropping would be useless.

\section{The data stream: compression right on the probe}

A probe produces a huge stream. To see the spike shape, each electrode must be digitized about $20$ thousand times a second with a precision of about ten bits. A thousand electrodes give
\begin{empheq}[box=\keybox]{equation}
\begin{aligned}
R_{\text{raw}} \;&=\; N\,f_s\,b \;\approx\; 10^3\cdot2\cdot10^4\cdot10 \;\approx\; 2\cdot10^8\ \text{bit/s},\\
R_{\text{spk}} \;&\approx\; N\,r\,\log_2\frac{e}{r\,\Delta t} \;\approx\; 8\cdot10^4\ \text{bit/s}.
\end{aligned}
\label{eq:probedata}
\end{empheq}
The second estimate is the capacity of the spike stream~\eqref{eq:spikeentropy} at $r=10$ spikes per second and precision $\Delta t=1\ms$: everything the stream actually contains takes two and a half thousand times less. For an implanted device this difference is decisive: the raw stream cannot be sent by radio through the skin with the power that can be afforded inside the head without heating the tissue. So an implanted device must detect spikes right on the chip next to the electrodes and send out only events: the channel number and the time of the spike. The first description of the Neuralink system had not reached this point yet: the chips amplify and digitize the signals, the whole raw stream leaves over a cable, and spikes are detected by software outside; on-chip compression and wireless links are named there as a plan~\cite{Musk2019}.

This is a familiar trick. The eye compresses the signal a hundredfold before sending it down a long cable (Chapter~\ref{sec:sensors}); an event camera transmits changes, not frames. To fit into its wire, the probe is forced to do what the brain itself does: speak in spikes, and only about news.

\section{What Neuralink does}

The Neuralink device is a probe built around these constraints~\cite{Musk2019}. Instead of a rigid array of needles there are many flexible threads thinner than a hair, with electrodes along each: in the first description of the system, up to $3072$ electrodes on $96$ threads, $32$ per thread, and in the device implanted in humans, according to the company, about a thousand channels on $64$ threads. Flexible threads damage the tissue less and tolerate better the fact that the brain constantly shifts slightly inside the skull. A robot inserts the threads, avoiding blood vessels. In the first description, as said above, the raw stream~\eqref{eq:probedata} went out over a cable. The device for humans, according to the company, is built differently: the housing is fully covered by skin, spikes are detected inside, data leave by radio, and power arrives wirelessly.

The first device's task is reading. The threads are placed in the part of the cortex that plans arm movements, and a decoder turns spikes into cursor movement on a screen. A person with paralyzed arms controls a computer by imagining movement. The idea itself is not new: interfaces on rigid arrays of a hundred electrodes have long allowed controlling a cursor~\cite{Hochberg2006}, writing text by imagining handwriting movements~\cite{Willett2021}, and even turning attempts to speak into on-screen text at about sixty words per minute~\cite{Willett2023}. What is new at Neuralink is engineering: channel count, wireless operation, a sealed housing, and robotic insertion, that is, an attempt to make a probe that can be worn for years outside the lab. As far as we know, results of human trials have been published mostly by the company itself rather than in peer-reviewed papers; the company reported, in particular, that some threads shifted after insertion and the number of working channels dropped. For debugging this is a common nuisance: a probe that moves relative to the board is now hearing different wires.

The company's second direction is writing: a vision device that feeds signals into the visual cortex of a person who has lost sight. More on it below, in the section on writing.

\section{Reading: the decoder as a recipient}

An interface decoder is a recipient in the sense of Proposition~\ref{prop:reader}: it decides for itself which part of the message to read. Practically all interfaces read \emph{group rates}: they count the spikes of each channel in windows of a few tens of milliseconds and sum them with weights chosen so as to produce the intended movement. This is the group rate code of Section~\ref{sec:rate}, and it is chosen for a reason: with a few spikes per window the rate of a single neuron is undefined, whereas the sum over a hundred already works.

How much information can be extracted? Writing with the ``mental hand'' ran at about ninety characters per minute~\cite{Willett2021}, i.e., one and a half characters per second: even at five bits per character this is less than eight bits per second. Cursor control, according to Neuralink, gives about ten bits per second. Yet the upper bound~\eqref{eq:spikeentropy} for a \emph{single} neuron is tens of bits per second, and for a hundred, thousands. An interface extracts from the neurons it eavesdrops on a small fraction of what they could carry. The bottleneck is not the wire but how many independent variables the group carries (Proposition~\ref{prop:pooling}) and how well the decoder understands the code, which, as we saw in Chapter~\ref{sec:code}, has not been fully deciphered.

There is a second feature, familiar from Chapter~\ref{sec:motorlearning}. The decoder and the brain learn from each other. The brain sees where the cursor went, gets an error, and adjusts its commands: the same motor learning over an error wire. And engineers periodically refit the decoder weights, because the neurons under the threads change and the connections in the network relearn. The result is two learners adapting to each other; to keep such a pair from running away, it helps to separate their learning rates: let one learn fast and the other slowly.

\begin{keyblock}
\begin{proposition}[Why a probe on a thousand neurons works]
\label{prop:probe}
An interface that eavesdrops on a tiny fraction of neurons can control movement because group activity is low-dimensional and its noise is correlated: a few hundred neurons carry almost all the shared variables. For the same reason an interface extracts only a few bits per second, as many as there are independent variables in the task, not as many as the spike stream can hold. The performance of interfaces is measured; how much better the decoder will become once the code is understood is an open question.
\end{proposition}
\end{keyblock}

\section{Writing: harder than reading}

Feeding a signal into the machine is harder than reading one. An electrode that passes current excites not one neuron but all the neurons and wires around it, regardless of type and sign. It cannot write a code in which it matters \emph{which} neuron fired and \emph{when} (Chapter~\ref{sec:code}); it can only roughly set \emph{where} and \emph{how strongly}.

For vision this is partly enough. A person sees current through an electrode in the visual cortex as a spot of light at a particular place in the visual field; when spots were lit simultaneously, up to sixteen at once, a person who had lost sight recognized some letters and distinguished the edges of objects~\cite{Fernandez2021}. This is a place code, and it can be written. But recall Chapter~\ref{sec:sensors}: the eye sends the brain not pixels but a compressed description (edges, changes, brightening and darkening) over separate wires. Writing ``spots of light'' into the cortex is writing pixels into a machine that expects a completely different code at its input. That is why artificial vision is still cruder than the electrode count would suggest. There are also test signals that need no electrode: visual illusions feed an unusual input through the eye and push the machine out of its normal mode, and each error points to its own mechanism (Section~\ref{sec:illusions}).

\section{What the probe does not see}

An electrode hears the address bus: the spikes of \nt{GLUT} neurons and, less well, of small \nt{GABA} neurons. But in Chapters~\ref{sec:serocomputer}--\ref{sec:acet} we saw that the operating mode of the whole machine is set by broadcast registers: the concentrations of \nt{SERO}, \nt{DOPA}, \nt{NORA}, and \nt{ACET}. The electrode does not hear them: the source neurons are vanishingly few, and their signal is not a spike next to the probe but a concentration in a volume. A debugger who sees the data bus but not the control registers will always be puzzled: the same input will give different responses because somewhere $g$, $a$, or $\tau_\gamma$ has changed. Concentrations can be measured with chemical sensors right in the tissue~\cite{Bunin1998}, but such sensors are not yet used in interfaces. Entirely outside the probe's view is also the machine's second, slow output, the hormones in the blood (Chapter~\ref{sec:chemoutput}): they are measured in blood samples, not with an electrode.

Nor does the probe see the weights, and it is in them that almost all memory and everything learned is stored. They can only be guessed from how the responses change. And it does not see pain the way a person feels it: in Chapter~\ref{sec:pain} we saw that the strength of the alarm signal is set by the machine itself, by gates in the spinal cord and regulators from above, so one cannot read from a sensor how much it hurts.

\section{Summary: debugging and open questions}

\begin{table}[t]
\centering
\footnotesize
\setlength{\tabcolsep}{4pt}
\renewcommand{\arraystretch}{1.15}
\begin{tabular}{>{\raggedright}p{2.2cm}>{\raggedright}p{3.6cm}>{\raggedright}p{3.4cm}>{\raggedright\arraybackslash}p{4.0cm}}
\toprule
Debugging technique & What the interface does & What the book explains & What is known \\
\midrule
probe & hears hundreds to thousands of neurons & low dimensionality and correlated noise (Ch.~\ref{sec:code}) & measured \\
compression on the probe & transmits events instead of the raw signal & capacity of the spike stream~\eqref{eq:spikeentropy} & solved in engineering \\
reading & group rates $\to$ movement, text, speech & the decoder is a recipient; group rate code & works; a few bits per second \\
joint learning & brain and decoder adapt to each other & error wire (Ch.~\ref{sec:motorlearning}) & observed; the adaptation rules are under study \\
writing & current lights spots in the visual field & a place code can be written, a timing code cannot (Ch.~\ref{sec:sensors}) & spots measured; useful vision not yet \\
reading registers & hardly done & broadcast channels (Ch.~\ref{sec:serocomputer}--\ref{sec:acet}) & chemical sensors exist in experiments \\
\bottomrule
\end{tabular}
\caption{Debugging the machine: what brain–computer interfaces can do and what the book's chapters say about them.}
\label{tab:debugging}
\end{table}

Table~\ref{tab:debugging} sums up the chapter. Brain–computer interfaces can already put a probe on the address bus and read a few bits per second from it, and the fact that this works at all is explained by the low dimensionality and correlated noise of Chapter~\ref{sec:code}. They can write into the machine only crudely, because its input code is compressed and addressed to individual wires. And the control registers and the weights, which make the machine trainable and tunable, are still almost invisible to the probe.

Every open question of Chapter~\ref{sec:synthesis} is also a task for the debugger. Which code to read will be settled when probes become denser and decoders learn to use spike timing, not just rates. How a multilayer network learns can be tested by putting probes in several layers at once and watching how their responses change during learning. What the broadcast channels transmit will become visible when chemical probes join the electrical ones. This book is an attempt to read the machine's schematic from its behavior and from laws every engineer knows. The debuggers being built now will make it possible to check this schematic with a probe.

\section{Exercises}

\begin{exercise}[\lvlA]
\label{ex:17:1}
\emph{Radio link budget.} Transmission through the skin costs $1$~nJ per bit, and the transmitter may use at most $10$~mW. What power is needed for the raw stream~\eqref{eq:probedata} with $N=1000$ electrodes and for the spike stream? What energy per bit would the raw stream require?
\end{exercise}

\begin{exercise}[\lvlA]
\label{ex:17:2}
\emph{How many bits in a hundred neurons.} A decoder sums the spikes of $N$ neurons over $50\ms$ windows at $r=20$~spikes/s; the pairwise noise correlation is $c=0.1$, and the signal changes the group rate by $30\%$ (rms). Estimate the rate $\tfrac12\log_2(1+\text{SNR})$ per window for $N=10$, $100$, $1000$ and $N\to\infty$. And for $c=0$, $N=1000$?
\end{exercise}

\begin{exercise}[\lvlB]
\label{ex:17:3}
\emph{Speed of a keyboard interface.} One and a half times a second an interface selects one of $N=32$ characters: the intended one with probability $P$, errors equally likely. How many bits per second does it transmit at $P=0.99$, $0.94$, $0.8$? How many characters per second are needed at $P=0.94$ for $10$~bit/s?
\end{exercise}

\begin{exercise}[\lvlB]
\label{ex:17:4}
\emph{Simulation: a group decoder.} $N$ neurons with $\lambda_i=[b+m\,\mathbf v\cdot\mathbf p_i]_+$, $b=20$, $m=15$~spikes/s, $50\ms$ windows, $\mathbf v$ with standard deviation $0.5$ per component; all see a shared noise, $\mathbf v+\boldsymbol\xi$, $\sigma_\xi=0.2$. Simulate an unbiased population-vector decoder. At what $N$ are the two noises equal, and how many bits per component per window does $N\to\infty$ give?
\end{exercise}

\begin{exercise}[\lvlB]
\label{ex:17:5}
\emph{Two learners.} The brain issues the command $m=b\,v^*$, the decoder outputs $v=d\,m$, $\langle v^{*2}\rangle=1$. After each trial both take a gradient step on $\tfrac12(v-v^*)^2$ with rate $\eta$. Simulate from $b=1.2$, $d=1$ at $\eta=0.8$, $1.1$, $1.5$, $2$. Each alone is stable for $\eta<2$; for what $\eta$ is the pair stable?
\end{exercise}

\begin{exercise}[\lvlB]
\label{ex:17:6}
\emph{Design: a probe pipeline.} $1024$ channels at $20$~thousand samples per second, neurons at $10$~spikes/s, radio at $1$~nJ/bit. What threshold on white sample noise gives at most $0.1$~Hz of false spikes per channel? We send spike counts per $20\ms$ at $2$~bits each: what is the power, and how many times lower is it than for raw $10$-bit samples?
\end{exercise}

\begin{exercise}[\lvlC]
\label{ex:17:7}
\emph{The speed limit of an interface.} Estimate $R\le DW\log_2(1+\text{SNR})$ for $D=10$ variables with bandwidth $W=5$~Hz at $\text{SNR}\approx0.9$ (signal-like noise) and $90$ (independent noise of $1000$ neurons). About $10$~bit/s is measured. What experiment would distinguish the two explanations of the gap: signal-like noise and ignorance of the code?
\end{exercise}

\chapter{Neuron Types by Electrical Function}
\chaptermark{Neurons as Devices}
\label{sec:eetypes}

In Chapter~\ref{sec:neurontypes} we sorted neurons by the transmitter their output releases. An electronics engineer would sort them differently: by what the neuron does to its signal, that is, by what device it acts as. We have known the book's main device since Chapter~\ref{sec:glutcomputer}: the leaky integrate-and-fire neuron~\eqref{eq:glutneuron}, an RC circuit with a comparator. But the brain has others too. Table~\ref{tab:eetypes} collects them in four groups, and almost every one has already appeared in the book.

\begin{center}
\footnotesize
\setlength{\tabcolsep}{3pt}
\renewcommand{\arraystretch}{1.15}
\begin{longtable}{>{\raggedright}p{3.1cm}>{\raggedright}p{3.3cm}>{\raggedright}p{4.7cm}>{\raggedright\arraybackslash}p{2.4cm}}
\caption{Neurons as devices: the name, the electronic analog, what the device does, and where it appears in the book.}\label{tab:eetypes}\\
\toprule
Neuron & Device & What it does & Where in the book \\
\midrule
\endfirsthead
\toprule
Neuron & Device & What it does & Where in the book \\
\midrule
\endhead
\bottomrule
\endfoot
\multicolumn{4}{l}{\emph{Filters and integrators}} \\
leaky integrate-and-fire neuron & RC integrator with a comparator & sums inputs within a window $\tau$ and fires at threshold & Ch.~\ref{sec:glutcomputer}, \eqref{eq:glutneuron} \\
coincidence detector & AND with a time window & fires only if the inputs arrive almost together; very small $\tau$ & Ch.~\ref{sec:glutcomputer}, \ref{sec:code}, \eqref{eq:window} \\
phasic neuron & high-pass filter, edge detector & responds to a change of input and falls silent & Ch.~\ref{sec:sensors} \\
adapting neuron & high-pass filter with automatic gain control & rate falls under constant input & Ch.~\ref{sec:sensors}, \eqref{eq:nakarushton} \\
resonator & tank circuit, band-pass filter & responds best to input at its own frequency & Section~\ref{sec:resonator} \\
leak-free integrator & op-amp integrator & turns velocity into position and holds it & Section~\ref{sec:perfectint} \\
\midrule
\multicolumn{4}{l}{\emph{Converters}} \\
tonic neuron & voltage-to-frequency converter & rate follows the input linearly and hardly tires & Ch.~\ref{sec:code} \\
graded neuron (no spikes) & analog amplifier, buffer & passes a smooth voltage over a short distance & Ch.~\ref{sec:sensors} \\
rectifying neuron & half-wave rectifier & the rate is never negative, $[u]_+$ & Ch.~\ref{sec:glutcomputer}, \ref{sec:gaba} \\
ON--OFF pair & pair of rectifiers, two-wire code & one responds to ``more,'' the other to ``less'' & Ch.~\ref{sec:glutcomputer}, \eqref{eq:dualrail} \\
neuron with shunting inhibition & analog divider, gain control & divides its input rather than subtracting from it & Ch.~\ref{sec:gaba}, \eqref{eq:shunt} \\
\midrule
\multicolumn{4}{l}{\emph{Oscillators and time}} \\
pacemaker & relaxation oscillator, multivibrator & emits spikes at a constant rate on its own & Ch.~\ref{sec:serocomputer}, \ref{sec:dofa} \\
pacemaker with an input & voltage-controlled oscillator & its own rhythm, which the input speeds up or slows down & Ch.~\ref{sec:chemoutput} \\
bursting neuron & gated burst generator & emits bursts of rapid spikes, the ``dashes'' & Ch.~\ref{sec:code}, \eqref{eq:burst} \\
phase-locked neuron & phase detector, phase-locked loop & fires at a particular phase of the input wave & Ch.~\ref{sec:sensors}, \ref{sec:chemoutput} \\
axon with delay & delay line & delays the signal by a set time & Ch.~\ref{sec:glutcomputer}, Fig.~\ref{fig:itd} \\
\midrule
\multicolumn{4}{l}{\emph{Memory and switching}} \\
bistable neuron & Schmitt trigger, latch & a brief input puts it into a lasting state & Section~\ref{sec:bistable} \\
neuron with dendritic branches & multiplexer, small multilayer network & a branch passes its input only if an enabling input is present & Ch.~\ref{sec:synthesis} \\
\end{longtable}
\end{center}

One caveat matters more than the whole table: these are not different cells but different \emph{modes}. The same neuron can be an integrator under slow input and a coincidence detector when inhibition has shortened its $\tau$ (Chapter~\ref{sec:code}); a pacemaker in wakefulness and a silent receiver in sleep. Which device results is decided by the ion channels of the membrane and by the broadcast channels that tune them.

\section{Four devices not yet seen in the book}

\subsection{The resonator}
\label{sec:resonator}

Some neurons have a slow current that opposes any change in voltage: when the voltage rises, the current pulls it down; when it falls, the current pulls it up, but with a lag. To an electronics engineer such a current behaves like an inductance, and together with the membrane capacitance it forms a tank circuit. The neuron responds most strongly to input at a particular frequency, usually from a few to one or two dozen hertz, and more weakly to slower and faster input~\cite{HutcheonYarom2000}. This is a band-pass filter in a single cell: from a noisy input it picks out a rhythm at its own frequency, for example the local rhythm of Chapter~\ref{sec:gaba}.

\subsection{The leak-free integrator}
\label{sec:perfectint}

A leaky integrator remembers its input only for a time $\tau$, tens of milliseconds. But to stay still, the eye must remember its position for seconds: the command ``turn'' arrives as a brief velocity, while the eye must be held in the new position. The network solves this with the same feedback as the memory of Chapter~\ref{sec:glutcomputer}. If a group of neurons with time constant $\tau$ excites itself with weight $w$, then $\tau\,\dd r/\dd t=-r+w\,r+I$, and the group's time constant is
\begin{empheq}[box=\keybox]{equation}
\tau_{\text{eff}} \;=\; \frac{\tau}{1-w} .
\label{eq:perfectint}
\end{empheq}
With $\tau=0.1\,\text{s}$ and $w=0.995$ this gives $20\,\text{s}$: a nearly ideal integrator built from neurons each of which forgets on its own within a tenth of a second~\cite{Seung1996}. The price is fine tuning: with $w$ slightly above one the integrator runs into saturation; slightly below, it forgets quickly. Such an integrator therefore needs constant retuning by learning, like the learning described in Chapter~\ref{sec:motorlearning}.

\subsection{The bistable neuron}
\label{sec:bistable}

In Chapters~\ref{sec:glutcomputer} and~\ref{sec:gaba} a flip-flop was assembled from a group of neurons. But a flip-flop can also exist in a single cell. Some neurons have a current that, once switched on, sustains the excitation by itself: a brief input puts the neuron into lasting activity, a \emph{plateau}, and inhibition returns it to rest. This is a Schmitt trigger: the neuron has two stable states with hysteresis between them. The \nt{ACET} neurons that control muscles behave this way, for example, and this property is switched on by a broadcast channel: the plateau appears when serotonin reaches the neuron~\cite{Hounsgaard1988}. The same cell is a latch with \nt{SERO} and an ordinary integrator without it.

\subsection{Integrators and resonators}

Theory divides excitable cells into two classes~\cite{Izhikevich2007}. \emph{Integrators} resemble a voltage-controlled oscillator that starts from zero: just above threshold the neuron fires arbitrarily slowly, and its rate rises smoothly with the input. Such a neuron sums everything that arrives and conveys the size of the input well by its rate. \emph{Resonators} resemble an oscillator with a minimum frequency: they cannot fire slowly; at threshold input they immediately emit spikes at a finite rate, and they prefer input at their own frequency. The former are the natural devices of the rate code of Chapter~\ref{sec:code}, the latter of the timing code.

\begin{keyblock}
\begin{proposition}[One neuron, many devices]
\label{prop:eetypes}
By electrical function, neurons act as leaky and leak-free integrators, coincidence detectors, high-pass and band-pass filters, voltage-to-frequency converters, rectifiers, dividers, oscillators, delay lines, and flip-flops. Which device a given cell becomes is set by its ion channels and by the broadcast channels that tune them. The modes themselves are measured; the extent to which the brain uses such reconfiguration for computation is mostly a hypothesis.
\end{proposition}
\end{keyblock}

To an engineer this resembles a programmable logic array: the hardware is the same, and the configuration sets the device. Only here the configuration is changed not at programming time but on the fly, by the slow broadcast channels discussed in Chapters~\ref{sec:serocomputer}--\ref{sec:acet}.

\section{Exercises}

\begin{exercise}[\lvlA]
\label{ex:18:1}
\emph{Tolerance of a leak-free integrator.} Neurons with $\tau=0.1$~s hold eye position by~\eqref{eq:perfectint} for $T=1$~s with an error of at most $5\%$ either way. (a)~Find the allowed range of $w$. (b)~$w$ is the sum of the weights of $K$ synapses, each with an independent $10\%$ error. For what $K$ does the spread of the sum fit the tolerance?
\end{exercise}

\begin{exercise}[\lvlB]
\label{ex:18:2}
\emph{The resonator as a tank circuit.} Describe the slow current of Section~\ref{sec:resonator} by a variable $W$: $C\,\dd V/\dd t=-g_LV-g_wW+I$, $\tau_w\,\dd W/\dd t=V-W$. Show that this is an $RC$ circuit with a parallel branch $R_w=1/g_w$, $L_w=\tau_w/g_w$, and for $C/g_L=10\ms$, $\tau_w=100\ms$, $\alpha=g_w/g_L=4$ find the resonance frequency and the gain of $|Z|$ there over $|Z(0)|$.
\end{exercise}

\begin{exercise}[\lvlB]
\label{ex:18:3}
\emph{How an integrator starts firing.} (a)~A neuron $\tau\,\dd V/\dd t=-V+\mu$, $\tau=10\ms$, with threshold $\theta$ and reset to zero. For what $\mu/\theta-1$ is the rate $1$~Hz? (b)~What rate does the model $\tau\,\dd v/\dd t=v^2+\mu$ give (a spike is $v$ escaping to $+\infty$, reset to $-\infty$) at $\mu=0.01$?
\end{exercise}

\begin{exercise}[\lvlB]
\label{ex:18:4}
\emph{Simulation: integrator and resonator.} The model of Exercise~\ref{ex:18:2} with $g_L=1$, $\tau_m=10\ms$, $\tau_w=100\ms$, threshold $1$, and reset $V\to0$ ($W$ unchanged) receives $\mu=1.5\sin2\pi ft$, $f=1$--$40$~Hz. In what frequency band does the neuron spike at $\alpha=0$ and at $\alpha=4$? Compare with the condition $1.5\,|Z(f)|>1$.
\end{exercise}

\begin{exercise}[\lvlB]
\label{ex:18:5}
\emph{A single-cell latch.} A neuron with a plateau current: $\tau\,\dd r/\dd t=-r+F(wr+I)$, $F(x)=\min(1,\max(0,x))$, $\tau=10\ms$, $w=1.5$, background $I=-0.2$. (a)~What is the shortest pulse $I=0.3$ that moves the neuron from $r=0$ to the plateau? (b)~What is the smallest amplitude $B$ of a long pulse $I=-0.2-B$ that returns it to rest?
\end{exercise}

\begin{exercise}[\lvlB]
\label{ex:18:6}
\emph{Self-tuning of the integrator.} During gaze fixation the input is $I=0$, a slip sensor gives the drift velocity $e=\dd r/\dd t$, and $\dd w/\dd t=-\eta\,e\,r$. With $w=0.95$, $\tau=0.1$~s, $\langle r^2\rangle=1$, find the $\eta$ at which $|1-w|$ enters the tolerance of Exercise~\ref{ex:18:1} within $60$~s of fixations. What breaks if learning also runs during eye turns?
\end{exercise}

\begin{exercise}[\lvlC]
\label{ex:18:7}
\emph{Why reconfigure a device?} A broadcast channel switches $\alpha$ in the model of Exercise~\ref{ex:18:2} between $0$ and $4$. By what factor does the resonator beat the integrator in signal-to-noise ratio for a weak sine in white current noise at $11$~Hz and at $1$~Hz? Devise a task in which it is the mode switching itself that pays off.
\end{exercise}

\chapter{Neurons by Number of Dendrites}
\chaptermark{Number of Dendrites}
\label{sec:dendrites}

\section{The number of inputs as a circuit parameter}

In Chapter~\ref{sec:neurontypes} we sorted neurons by what their output releases, and in Chapter~\ref{sec:eetypes} by what device they act as. There is a third feature an engineer will notice at once: \emph{how many inputs a neuron has}. Inputs are collected by dendrites, the branches on which other neurons' axons land (Chapter~\ref{sec:dofa}, Section~\ref{sec:weightstore}). Some neurons have no dendrites at all, others one or two, still others a tree that collects a hundred thousand inputs. To an electronics engineer this is fan-in, and it almost always matches the task.

A simple estimate shows why the number and size of dendrites matter so much. The membrane is a capacitor with specific capacitance $c_m\approx1$~\textmu F/cm$^2$, and the larger the area $A$ of the neuron together with its dendrites, the more charge must be delivered to raise the voltage by $\Delta V$. Ignoring leakage, the time it takes an input current $I$ to do this is
\begin{empheq}[box=\keybox]{equation}
t \;\approx\; \frac{c_m\,A\,\Delta V}{I} .
\label{eq:dendriteload}
\end{empheq}
A small neuron with a few short dendrites has an area of about $2\cdot10^{-5}$~cm$^2$ and a capacitance of about $20$~pF. A single large synapse with a current of a few nanoamperes raises it by $20\mV$ in less than a tenth of a millisecond. A neuron with a large tree has about fifteen times the area and a capacitance of about $300$~pF; an ordinary synapse with a current of about $0.1$~nA would charge it in $60\ms$, much longer than the leak time constant, which means it would never charge it. Such a neuron needs dozens of simultaneous inputs. The numbers here are orders of magnitude, but the conclusion does not depend on them: \emph{few inputs means fast and precise; many inputs means slow and by the sum}.

\begin{figure}[t]
\centering
\begin{tikzpicture}[>=Stealth,
  soma/.style={circle,draw,very thick,fill=blue!6,minimum size=0.55cm,inner sep=0pt},
  ax/.style={->,very thick,red!70!black},
  den/.style={very thick,blue!60!black}]
% 0 dendrites: sensor on a line
\begin{scope}[xshift=0cm]
\draw[den,decorate,decoration={snake,amplitude=1pt,segment length=4pt}] (-0.9,0) -- (0,0);
\node[font=\scriptsize] at (-0.9,0.3) {sensor};
\draw[ax] (0,0) -- (1.0,0);
\draw[thick] (0.1,0) -- (0.1,0.45);
\node[soma,minimum size=0.4cm] at (0.1,0.65) {};
\node[font=\scriptsize,align=center] at (0.05,-0.75) {$0$\\line with sensor};
\end{scope}
% 1 dendrite
\begin{scope}[xshift=3.3cm]
\draw[den] (-0.9,0) -- (-0.28,0);
\node[soma] at (0,0) {};
\draw[ax] (0.28,0) -- (0.9,0);
\node[font=\scriptsize,align=center] at (0,-0.75) {$1$\\buffer, repeater};
\end{scope}
% 2 dendrites: one per ear
\begin{scope}[xshift=6.2cm]
\draw[den] (-0.28,0) -- (-0.9,0); \draw[den] (0.28,0) -- (0.9,0);
\node[soma] at (0,0) {};
\draw[ax] (0,-0.28) -- (0,-0.6);
\node[font=\scriptsize] at (-0.9,0.25) {left}; \node[font=\scriptsize] at (0.9,0.25) {right};
\node[font=\scriptsize,align=center] at (0,-1.0) {$2$\\AND in time};
\end{scope}
% ~4 short dendrites: mixer
\begin{scope}[xshift=9.0cm]
\foreach \a in {45,135,225,315}{\draw[den] (\a:0.28) -- (\a:0.7); \fill[blue!60!black] (\a:0.72) circle (0.06);}
\node[soma] at (0,0) {};
\draw[ax] (0.28,0) -- (1.0,0);
\node[font=\scriptsize,align=center] at (0,-1.0) {$\sim4$\\mixer};
\end{scope}
% many: summing tree
\begin{scope}[xshift=12.0cm]
\foreach \a in {60,90,120}{\draw[den] (0,0.28) -- ++(\a:0.55) coordinate (b); \foreach \d in {-20,20}{\draw[den] (b) -- ++(\a+\d:0.35);}}
\foreach \a in {-120,-60}{\draw[den] (0,-0.28) -- ++(\a:0.45);}
\node[soma] at (0,0) {};
\draw[ax] (0.28,0) -- (1.0,0);
\node[font=\scriptsize,align=center] at (0,-1.0) {$10^4$--$10^5$\\summer};
\end{scope}
\end{tikzpicture}
\caption{Five classes by number of inputs. Blue: dendrites (inputs); red: axon (output). The fewer the inputs, the faster and more precisely the neuron transmits an individual signal; the more inputs, the better it sums and classifies. A schematic, not a drawing of cell shapes.}
\label{fig:dendriteclasses}
\end{figure}

\section{Zero dendrites: a sensor at the end of a line}

The sensory neurons of touch, pain, and muscle stretch (Chapters~\ref{sec:sensors}, \ref{sec:pain}, \ref{sec:muscles}) have not a single dendrite on the cell body. The axon leaves the body as one twig and immediately splits in two: one branch goes to the skin or muscle, and its ending itself serves as the sensor; the other goes to the spinal cord. The spike runs from the sensor straight into the spinal cord, bypassing the cell body. To an engineer this is a communication line with the sensor at the far end and the cell body hanging off to the side like a power supply: it feeds the line but takes no part in transmission. The benefit is speed and reliability: there is not a single place on the path where the signal would have to be summed and a new decision made on whether to pass it on.

The light and sound sensors are an even more extreme case: they are not collecting neurons but the transducers themselves, whose input is not a synapse but a receptor protein (Fig.~\ref{fig:transducer}).

\section{One dendrite: a buffer}

A neuron with one dendrite and one axon receives one input line and passes it on. Olfactory neurons are built this way: a single dendrite reaches the surface of the nose and carries a single type of receptor~\cite{Mombaerts2004}; so are the relay cells of the retina between the light sensors and the output neurons of the eye; so are the neurons linking the ear's sensors to the brain. An engineer would call them a buffer or a repeater: they do not compute but deliver one signal, sometimes changing its form, for example turning a sensor's smooth voltage into spikes, as in Chapter~\ref{sec:sensors}.

\section{Two dendrites: AND in time}

The coincidence detectors that determine the direction of a sound (Fig.~\ref{fig:itd}) in mammals often have exactly two dendrites pointing in opposite directions: one receives inputs from the left ear, the other from the right~\cite{Grothe2010}. Why keep the inputs apart? Recall Eq.~\eqref{eq:shunt}: when many excitatory channels are open on one patch of membrane, the voltage there approaches the reversal potential, and each additional input adds less and less. So many inputs from one ear, gathered on one dendrite, saturate it and cannot on their own drive the neuron to threshold. But one input from each dendrite sum almost linearly. Two dendrites turn the neuron into a stricter AND: fire only if the signals came from both sides and together. Models show that such separation does improve coincidence detection~\cite{AgmonSnir1998}.

\section{A few short dendrites: mixer and relay}

\emph{Mixers.} In the motor learning circuit of Chapter~\ref{sec:motorlearning}, about seventy billion small \nt{GLUT} neurons expand the input into a very long vector. Each of them has only about four short dendrites, each ending in a ``claw'' that grips the terminal of one input. Each mixer thus sees only about four inputs out of many and works as a small AND element over its four. From $M$ input lines one can choose $\binom{M}{4}$ different quadruples: for $M=100$ this is almost four million. Why so many? A threshold element with $n$ inputs can correctly split into two classes at most
\begin{empheq}[box=\keybox]{equation}
P_{\max} \;\approx\; 2n
\label{eq:cover}
\end{empheq}
random patterns~\cite{Cover1965}. The output neuron of the same circuit receives about $10^5$ inputs from mixers, and by~\eqref{eq:cover} its upper bound is about $2\cdot10^5$ different associations. This is the limit for an ideal element: if all weights are positive, as for excitatory synapses, and a margin against noise is needed, the estimate for the same neuron is about $4\cdot10^4$ associations~\cite{Brunel2004}. Mixers with few inputs are needed precisely so that the large summer has many different, nearly independent inputs~\cite{Marr1969,Albus1971}.

\emph{Relays.} On the path from the ear to the direction detectors there are neurons with a few short dendrites that receive only a few huge synapses, and some receive essentially one. By~\eqref{eq:dendriteload} this is exactly what is needed: a small capacitance and a large current, so the neuron fires within a fraction of a millisecond after each input spike, losing almost no timing precision. One of these relays receives \nt{GLUT} and outputs \nt{GLYC}: an inverter with a precision of tens of microseconds that supplies fast inhibition to the direction detectors~\cite{BorstSoria2012}.

\section{Thousands of inputs: summer and classifier}

At the other end of the scale are the cortical \nt{GLUT} neurons with ten thousand inputs and the output \nt{GABA} neurons of the motor learning circuit with a hundred thousand. By~\eqref{eq:dendriteload} such a neuron is slow and cannot respond to an individual input: it sums. This is the integrating neuron of Chapter~\ref{sec:glutcomputer} and the summer from which classifiers are built. A large tree adds something new as well: its branches act as separate subcircuits with their own thresholds, so that one neuron behaves like a small multilayer network~\cite{Beniaguev2021,Gidon2020}.

\section{Dendrites that are also an output}

Some neurons have no axon at all, and their dendrites serve as both input and output: they receive signals and release transmitter themselves. The best-studied example is the retinal \nt{GABA} neurons involved in detecting the direction of motion (Chapter~\ref{sec:gaba}, Fig.~\ref{fig:direction}). In such a neuron each branch on its own responds to motion from the cell body toward its tip and delivers inhibition from that tip~\cite{Euler2002}. One neuron is several independent direction detectors, one per branch. To an engineer this is not one element but a chip with several channels in one package.

\section{Summary}

\begin{table}[t]
\centering
\footnotesize
\setlength{\tabcolsep}{3pt}
\renewcommand{\arraystretch}{1.15}
\begin{tabular}{>{\raggedright}p{1.9cm}>{\raggedright}p{3.4cm}>{\raggedright}p{2.6cm}>{\raggedright}p{1.8cm}>{\raggedright\arraybackslash}p{3.8cm}}
\toprule
Dendrites & Example & Device & Inputs & What is known \\
\midrule
$0$ & touch, pain, and stretch sensors & line with a sensor at the end & --- & established \\
$1$ & olfactory neurons, retinal relay cells, ear neurons & buffer, repeater & $1$ to a few & established \\
$2$ & sound direction detectors & AND in time & two bundles & structure established; the gain from separating inputs shown in models \\
$\sim4$ & mixers of the motor learning circuit & small AND element & $\sim4$ & established; the role of input expansion is a well-founded hypothesis \\
a few, short & relays on the path from the ear & relay, inverter & $1$ to a few huge & established \\
thousands & cortical \nt{GLUT} neurons, output neurons of motor learning & summer, classifier & $10^4$--$10^5$ & established; branches as subcircuits measured in individual cases \\
output dendrites & retinal direction \nt{GABA} neurons & multichannel chip & many & measured \\
\bottomrule
\end{tabular}
\caption{Neurons by number of dendrites.}
\label{tab:dendrites}
\end{table}

\begin{keyblock}
\begin{proposition}[The number of inputs follows the task]
\label{prop:dendrites}
Where the speed and precision of an individual signal are needed (in sensors, relays, and coincidence detectors), neurons have few dendrites, few inputs, and large synapses: the small capacitance~\eqref{eq:dendriteload} charges quickly. Where summing and classifying are needed, neurons have large trees with thousands of inputs. The structure of the classes is measured; the computational explanation is well founded but for many classes remains a hypothesis.
\end{proposition}
\end{keyblock}

Table~\ref{tab:dendrites} complements the two earlier sortings: by transmitter (Table~\ref{tab:types}) and by device (Table~\ref{tab:eetypes}). All three look at the same element from different sides: what it outputs, what it computes, and how much it listens to.

\section{Exercises}

\begin{exercise}[\lvlA]
\label{ex:19:1}
\emph{How many inputs a large neuron needs.} A neuron with $C=300$~pF and $\tau_m=20\ms$ receives synapses that are current sources of $0.1$~nA each; threshold is $20\mV$ above rest. (a)~How many synapses must act simultaneously to reach threshold at all; within $10\ms$; within $5\ms$? (b)~How long does a neuron with $C=20$~pF take to reach it from a $4$~nA synapse?
\end{exercise}

\begin{exercise}[\lvlA]
\label{ex:19:2}
\emph{Why four.} A mixer is an AND over $k$ lines out of $M=100$; half the lines are active in a pattern; there are $7\cdot10^{10}$ mixers. (a)~How many respond to a pattern for $k=2$; $4$; $40$? (b)~By what factor do mixers increase the number of associations~\eqref{eq:cover} of an output neuron with $10^5$ inputs over $M=100$ direct lines?
\end{exercise}

\begin{exercise}[\lvlB]
\label{ex:19:3}
\emph{Where $2n$ comes from.} A hyperplane through the origin splits $P$ points in general position in $n$-dimensional space into two classes in $C(P,n)=2\sum_{k=0}^{n-1}\binom{P-1}{k}$ ways. (a)~Prove that for $P=2n$ exactly half of the $2^P$ splits are separable. (b)~What fraction is separable for $n=100$ and $P=180$; $220$?
\end{exercise}

\begin{exercise}[\lvlB]
\label{ex:19:4}
\emph{Two dendrites as a strict AND.} A dendrite is a patch with leak $g_L$; each input opens a conductance $g_0=0.5\,g_L$ on it with reversal potential $E_E$; the soma sees $\kappa(V_1+V_2)$; threshold is $\theta=1.1\,\kappa E_E$. Why can no number of inputs on one dendrite fire the neuron, and how many inputs are needed on each?
\end{exercise}

\begin{exercise}[\lvlB]
\label{ex:19:5}
\emph{Designing a relay.} Spike jitter is $\sigma_t=\sigma_V/\dot V$, where $\dot V$ is the voltage slope at threshold. We need $\sigma_t\le20$~\textmu s with noise $\sigma_V=1\mV$; neglect leakage. How many synchronous $0.1$~nA synapses does a neuron with $C=300$~pF need for this, and what is the jitter of a neuron with $C=20$~pF and a single $4$~nA synapse?
\end{exercise}

\begin{exercise}[\lvlB]
\label{ex:19:6}
\emph{Simulation: charging a long dendrite.} Simulate a passive cable $\tau_m\partial_tV=\lambda^2\partial_x^2V-V+i/g$ of length $L=\ell/\lambda$ with sealed ends ($40$ compartments, Euler). A brief current pulse is injected at the far end: when, and to what fraction of the voltage from the same charge spread over the whole dendrite, does the near end rise for $L=0.2$; $1$; $2$?
\end{exercise}

\begin{exercise}[\lvlC]
\label{ex:19:7}
\emph{Research: the optimal number of inputs.} The capacitance is $C=C_0+nc_s$, $m$ inputs of current $I_s$ act simultaneously, and the neuron stores $P\approx2n$ associations~\eqref{eq:cover}. How does the response time depend on $P$ at constant $m$ and at a constant fraction $m=\varphi n$? Propose a cost criterion and find the optimal $n$ for a relay and for a classifier.
\end{exercise}

\chapter{What Neurons Do During Sleep}
\chaptermark{Sleep}
\label{sec:sleep}
\begin{chaptergoals}
\item why sleep is a set of modes switched by \nt{ACET}, \nt{NORA} and \nt{SERO}, not a shutdown;
\item how sleep pressure and two thresholds form a relay with hysteresis locked to the daily rhythm;
\item how fatigue turns the cortex into a relaxation oscillator, and how \nt{ACET} stops the swings;
\item how the day is replayed fast-forward, and why weights may be scaled down during sleep.
\end{chaptergoals}

\section{Sleep is not a shutdown but a different mode}

An engineer looking at a switched-off computer will say: it is doing nothing. That does not work for a sleeping brain. Neurons are not silent in sleep: in deep sleep the cortex emits spikes, and in REM sleep almost as many as during the day. What changes is not \emph{whether} neurons work but \emph{how} they work. Sleep is a set of machine modes, and they are switched by the same broadcast channels as during the day.

For each mode one can write down the ``register state'' (Table~\ref{tab:sleepmodes}). During the day \nt{ACET}, \nt{NORA}, and \nt{SERO} are active, the sensors are connected, and the muscles obey. In slow-wave, deep sleep all three channels quiet down, and the input from the senses is weakened~\cite{Hasselmo1999}: acetylcholine release in the cortex falls, and \nt{NORA} and \nt{SERO} neurons fire less often than during the day~\cite{Marrosu1995,AstonJonesBloom1981,McGintyHarper1976}. In REM sleep the picture is strange: \nt{ACET} rises again to its daytime level~\cite{Marrosu1995}, while \nt{NORA} and \nt{SERO} fall almost completely silent~\cite{AstonJonesBloom1981,McGintyHarper1976,JacobsAzmitia1992}. The body's muscles are switched off in REM sleep: the \nt{ACET} neurons that control them are strongly inhibited through \nt{GABA} and \nt{GLYC}~\cite{BrooksPeever2012}, by the same veto as in Chapter~\ref{sec:gaba}, only now imposed on the whole output at once.

\begin{table}[t]
\centering
\footnotesize
\setlength{\tabcolsep}{4pt}
\renewcommand{\arraystretch}{1.15}
\begin{tabular}{>{\raggedright}p{3.1cm}>{\raggedright}p{2.9cm}>{\raggedright}p{3.2cm}>{\raggedright\arraybackslash}p{3.4cm}}
\toprule
& wakefulness & slow-wave sleep & REM sleep \\
\midrule
\nt{ACET} (write, Ch.~\ref{sec:acet}) & high & low & high \\
\nt{NORA} (gain, Ch.~\ref{sec:nora}) & medium, bursts & low & nearly silent \\
\nt{SERO} (Ch.~\ref{sec:serocomputer}) & medium & low & nearly silent \\
input from sensors & connected & weakened & weakened \\
output to muscles & connected & weakened & switched off by inhibition \\
cortical activity & steady & slow swings: activity/silence & steady, as in the day \\
\bottomrule
\end{tabular}
\caption{The machine's register state in three modes. All rows are measured; the levels of \nt{ACET}, \nt{NORA}, and \nt{SERO}, by direct recordings across sleep stages~\cite{Marrosu1995,AstonJonesBloom1981,McGintyHarper1976}.}
\label{tab:sleepmodes}
\end{table}

\section{When to sleep: a relay with hysteresis}

What decides when the machine should switch into sleep? A simple two-part scheme works well~\cite{Borbely1982}. The first part is \emph{sleep pressure} $S$, which accumulates while we are awake and is discharged while we sleep. It is a capacitor that charges during the day and discharges at night:
\begin{empheq}[box=\keybox]{equation}
\frac{\dd S}{\dd t} \;=\; \frac{1-S}{\tau_r}\ \ \text{(wake)},\qquad
\frac{\dd S}{\dd t} \;=\; -\frac{S}{\tau_d}\ \ \text{(sleep)} .
\label{eq:sleeppressure}
\end{empheq}
The second part is two thresholds: the machine falls asleep when $S$ reaches the upper threshold $H(t)$ and wakes up when $S$ drops to the lower one $L(t)$. Both thresholds swing smoothly in step with the daily rhythm of Chapter~\ref{sec:chemoutput}. The result is a relay with hysteresis, the Schmitt trigger of Section~\ref{sec:bistable}, and together with the capacitor it forms a relaxation oscillator phase-locked by the daily rhythm~\eqref{eq:pll}.

\begin{figure}[t]
\centering
\begin{tikzpicture}[>=Stealth,x=0.16cm,y=4.2cm]
\foreach \a/\b in {0/4.3,21.2/29.3,45.4/53.4,69.5/72}{\fill[blue!8] (\a,0) rectangle (\b,1);}
\draw[->,gray!70!black] (0,0) -- (75,0) node[right,font=\small,black] {$t$, h};
\draw[->,gray!70!black] (0,0) -- (0,1.08) node[left,font=\small,black] {$S$};
\foreach \t in {0,24,48,72}{\draw[gray!60!black] (\t,-0.02) -- (\t,0.02); \node[below,font=\scriptsize] at (\t,0) {\t};}
\draw[thick,densely dashed,red!70!black] plot file {figures/sleep_H.dat};
\draw[thick,densely dashed,green!45!black] plot file {figures/sleep_L.dat};
\draw[very thick,blue!60!black] plot file {figures/sleep_S.dat};
\node[font=\scriptsize,red!70!black,anchor=west] at (73,0.72) {$H$: fall asleep};
\node[font=\scriptsize,green!45!black,anchor=west] at (73,0.24) {$L$: wake up};
\node[font=\scriptsize,blue!60!black] at (25.25,0.93) {sleep};
\node[font=\scriptsize,blue!60!black] at (49.4,0.93) {sleep};
\end{tikzpicture}
\caption{Sleep pressure $S$~\eqref{eq:sleeppressure} with $\tau_r=18.2$~h, $\tau_d=4.2$~h. During the day $S$ charges up to the upper threshold $H$; at night (shaded) it discharges to the lower threshold $L$; both thresholds swing with the daily rhythm. The machine settles by itself into a regime of about $16$ hours awake and $8$ hours asleep.}
\label{fig:twoprocess}
\end{figure}

In our model, with the usual values $\tau_r=18.2$~h and $\tau_d=4.2$~h, the machine settles by itself into $16.1$ hours awake and $8.0$ hours asleep (Fig.~\ref{fig:twoprocess}). The model also explains something that seems odd: after forty hours without sleep the pressure reaches $0.90$, yet recovery sleep in the model lasts only $8.8$ hours, not a day longer than usual. The discharge is exponential, and a high charge drains quickly; the ``debt'' is repaid not by the length of sleep but by its depth.

What physically plays the role of $S$? A strong candidate is adenosine, the ``load counter'' of Appendix~\ref{sec:othermediators}: its concentration rises during wakefulness and falls during sleep~\cite{PorkkaHeiskanen1997,Dunwiddie2001}. That sleep pressure is adenosine and nothing else is a hypothesis.

\section{Slow-wave sleep: the whole network as a relaxation oscillator}

In deep sleep the cortical neurons work in turns as a whole group: less than once a second they switch on together for a few hundred milliseconds (the \emph{active state}) and fall silent together (the \emph{silent state}); these swings run below $1$~Hz, most often at $0.3$--$0.4$~Hz under anesthesia~\cite{Steriade1993}. We can build these slow swings from parts we already have. Take a group of \nt{GLUT} neurons with strong feedback onto itself, the memory of Chapter~\ref{sec:glutcomputer} with its two stable states, and add slow fatigue, as in the rhythm generator of Chapter~\ref{sec:muscles}:
\begin{empheq}[box=\keybox]{equation}
\tau\,\frac{\dd r}{\dd t} = -r + F\bigl(w\,r + I - a\bigr),
\qquad
\tau_a\,\frac{\dd a}{\dd t} = -a + b\,r .
\label{eq:updown}
\end{empheq}
The group switches on, works, and tires; fatigue $a$ eats away the upper state, and the group drops into silence; in silence the fatigue wears off, the lower state disappears, and the group switches on again. The result is the same relaxation oscillator as sleep pressure, only about eighty thousand times faster.

\begin{figure}[t]
\centering
\begin{tikzpicture}[>=Stealth,x=2.9cm,y=1.0cm]
\begin{scope}[yshift=1.9cm]
\draw[->,gray!70!black] (0,0) -- (4.1,0);
\draw[->,gray!70!black] (0,0) -- (0,1.25) node[left,font=\small,black] {$r$};
\draw[thick,blue!60!black] plot file {figures/updown_sleep.dat};
\node[font=\scriptsize,anchor=west] at (4.15,0.5) {sleep: $b=5$};
\end{scope}
\draw[->,gray!70!black] (0,0) -- (4.1,0) node[right,font=\small,black] {$t$, s};
\draw[->,gray!70!black] (0,0) -- (0,1.25) node[left,font=\small,black] {$r$};
\draw[thick,red!70!black] plot file {figures/updown_wake.dat};
\node[font=\scriptsize,anchor=west] at (4.15,0.5) {day: $b=1$};
\foreach \t in {0,1,2,3,4}{\draw[gray!60!black] (\t,-0.05) -- (\t,0.05); \node[below,font=\scriptsize] at (\t,0) {\t};}
\end{tikzpicture}
\caption{The same group~\eqref{eq:updown} in two modes: $w=5$, $I=2$, $\theta=2.5$, $\tau=10\ms$, $\tau_a=0.3$~s. With strong fatigue ($b=5$, as in slow-wave sleep) the group swings between activity and silence with a period of $1.1$~s (a model value; in the cortex the period is longer than a second, about $3$~s under anesthesia). Weaken the fatigue ($b=1$), which is what acetylcholine does, and the same group works steadily.}
\label{fig:updown}
\end{figure}

What moves the network from swings to steady daytime work? Acetylcholine suppresses the slow currents in neurons that create fatigue~\cite{McCormick1992}. In our model, with strong fatigue, $b=5$, the group swings with a period of $1.1$~s (faster than the measured swings, but of the same order) and is active about $35\%$ of the time, averaged over a whole number of periods; with weak fatigue, $b=1$, the same group works steadily (Fig.~\ref{fig:updown}). One register, the \nt{ACET} level, switches an oscillator into an amplifier. On top of the slow swings, sleep also carries faster bursts of a $7$--$15$ cycles-per-second rhythm; they are generated by the \nt{GLUT}--\nt{GABA} loop between the cortex and an intermediate relay node~\cite{SteriadeMcCormickSejnowski1993}, a relative of the rhythm of Chapter~\ref{sec:gaba}.

\section{Replays: fast-forwarding the day}

In Chapter~\ref{sec:acet} we saw that in slow-wave sleep neurons that worked together during the day fire together again, in the same order. Let us add one engineering detail: the replays run \emph{fast-forward}. A sequence that took seconds during the day is played back in sleep in tenths of a second~\cite{Nadasdy1999}: in the hippocampus about twenty times faster~\cite{LeeWilson2002}, in the cortex six to seven times faster~\cite{Euston2007}. And it is not played back at random moments: the short bursts of replay in one region align with the swings of activity and silence and with the fast rhythm bursts in the cortex. If this alignment is artificially strengthened, memory improves~\cite{Maingret2016}; this is a causal link, not a coincidence.

Why would the machine fast-forward the day? Engineers know a difficulty called \emph{catastrophic forgetting}: a network that learns new things in sequence, one after another, corrupts the old. The cure is \emph{interleaving}: learn the new mixed with the old, in small portions, many times. The hypothesis of two learning systems~\cite{McClelland1995} is that a fast memory records the whole day, and during sleep it plays it back to the slow cortex little by little, interleaved, many times. This is the same ``fast buffer, slow storage'' pair as in Chapter~\ref{sec:acet}, and the same pair of fast and slow learners as in Chapter~\ref{sec:motorlearning}. Replays and their alignment are measured; that their purpose is interleaved training of the slow memory is a well-founded hypothesis.

\section{Renormalizing the weights}

During the day the Hebbian rules of Chapter~\ref{sec:dofa} mostly strengthen connections. If they only strengthen, the weights grow, and with them the energy cost grows and the sensitivity falls: a neuron whose inputs are all strong stops distinguishing them. According to the \emph{synaptic homeostasis} hypothesis~\cite{TononiCirelli2014}, sleep is needed, among other things, to bring the weights back to a working level: all connections are weakened by the same factor, so their \emph{ratios}, which hold what was learned, are preserved. To an engineer this is normalization, as in rule~\eqref{eq:oja}, only done not on the fly but at a separate time.

Direct measurements do not contradict this: in mice, the contact area of cortical synapses after sleep is about $18\%$ smaller than after wakefulness, the reduction is proportional to size, and the largest, most stable contacts hardly change~\cite{deVivo2017}. Weight scaling is measured; that it is the main job of sleep is a hypothesis.

\section{REM sleep: a machine disconnected from its inputs and outputs}

In REM sleep the cortex works almost as during the day, but the input from the senses is weakened and the output to the muscles is switched off by inhibition. To an engineer this is a familiar mode: \emph{a bench test}. The circuit runs at full power, but its inputs are tied to an internal generator and its outputs are disconnected from the actuators so that nothing breaks. According to one hypothesis, dreams are exactly such an internal run of the world model assembled during the day; what exactly the network does then, and whether it learns, is unknown. Only what is recorded in Table~\ref{tab:sleepmodes} is measured here: which channel is on, which is silent, and that the output is blocked.

\section{Maintenance and local sleep}

It was long believed that during sleep the extracellular gaps widen and the brain is flushed of metabolic waste faster~\cite{Xie2013}. Recent experiments that measured clearance by a different method found the opposite: it is slower during sleep~\cite{Miao2024}. The question is currently open, and we leave it open.

Another fact is more solid: sleep is a property of local networks, not of the whole machine at once. If an animal is kept awake for a long time, individual patches of its cortex briefly drop into silence, as in slow-wave sleep, while the animal is awake and moving, and at those moments the animal makes more errors~\cite{Vyazovskiy2011}. Each network tires on its own and switches on its own; a global mode arises when the broadcast channels switch all networks at once.

\begin{keyblock}
\begin{proposition}[Sleep as a set of modes]
\label{prop:sleep}
Neurons do not switch off in sleep: the broadcast channels move the machine into other modes. When to sleep is decided by a relay with hysteresis~\eqref{eq:sleeppressure}, locked to the daily rhythm. In slow-wave sleep the network becomes a relaxation oscillator~\eqref{eq:updown} and replays the day fast-forward; in REM sleep it works disconnected from its inputs and outputs. The modes, the replays, and the weight scaling are measured; their purposes, interleaved learning and renormalization, are well-founded hypotheses.
\end{proposition}
\end{keyblock}

\section{Summary}

\begin{table}[t]
\centering
\footnotesize
\setlength{\tabcolsep}{4pt}
\renewcommand{\arraystretch}{1.15}
\begin{tabular}{>{\raggedright}p{3.2cm}>{\raggedright}p{4.4cm}>{\raggedright\arraybackslash}p{5.2cm}}
\toprule
What happens & How & What is known \\
\midrule
mode switching & registers \nt{ACET}, \nt{NORA}, \nt{SERO} (Table~\ref{tab:sleepmodes}) & measured \\
when to sleep & capacitor $S$ and relay with hysteresis~\eqref{eq:sleeppressure} & the model describes experiments well; adenosine as $S$ is a hypothesis \\
slow swings & group with feedback and fatigue~\eqref{eq:updown} & swings measured; the model is the simplest circuit \\
fast-forwarding the day & replays $6$--$20$ times faster, aligned with the swings & measured; strengthening the alignment improves memory \\
why replay & interleaved training of the slow memory & well-founded hypothesis \\
renormalizing weights & scaling of connections during sleep & reduction of contacts measured; as the main role of sleep, a hypothesis \\
REM sleep & work with inputs and output disconnected & mode measured; purpose unknown \\
flushing & widening of extracellular gaps & conflicting results \\
local sleep & individual patches drop into silence & measured \\
\bottomrule
\end{tabular}
\caption{What neurons do during sleep.}
\label{tab:sleep}
\end{table}

Table~\ref{tab:sleep} sums up the chapter. Sleep turned out to be not a pause in the machine's work but maintenance on the run: mode switching by the same broadcast channels, an oscillator built from the same parts as the stepping rhythm, fast-forward replay of the day for the slow memory, and renormalization of the weights. By day the machine writes; by night it rewrites and puts in order what was written.

\section{Exercises}

\begin{exercise}[\lvlA]
\label{ex:20:1}
\emph{A relay without the daily rhythm.} Let the thresholds be constant: $H=0.67$, $L=0.17$ (the mean thresholds in Fig.~\ref{fig:twoprocess}). Derive from~\eqref{eq:sleeppressure} the durations of wake and sleep and the oscillator period for $\tau_r=18.2$~h, $\tau_d=4.2$~h. Why does the model with swinging thresholds still settle at $16.1+8.0=24$~h? Which device of Chapter~\ref{sec:chemoutput} does this effect belong to?
\end{exercise}

\begin{exercise}[\lvlA]
\label{ex:20:2}
\emph{Sleep debt.} Sleep begun at pressure $S_0$ lasts $t_s=\tau_d\ln(S_0/L)$, $\tau_d=4.2$~h, where $L$ is the lower threshold at waking. (a)~After $40$~h without sleep $S_0=0.90$, and sleep lasted $8.8$~h. What was $L$? How long would sleep last with the usual $L\approx0.092$? (b)~Overnight the synaptic contact area shrinks by $18\%$. By how much must the weights grow during the day?
\end{exercise}

\begin{exercise}[\lvlB]
\label{ex:20:3}
\emph{Designing the thresholds.} Choose constant thresholds $H$ and $L$ for which relay~\eqref{eq:sleeppressure} with $\tau_r=18.2$~h and $\tau_d=4.2$~h gives exactly $16$~h of wake and $8$~h of sleep. How is such a machine worse than one with swinging thresholds if one night it is woken after $4$~hours?
\end{exercise}

\begin{exercise}[\lvlB]
\label{ex:20:4}
\emph{The period of the slow swings.} In model~\eqref{eq:updown} $F(x)=1/\bigl(1+e^{-(x-\theta)/k}\bigr)$, $w=5$, $I=2$, $\theta=2.5$, $k=0.25$, $b=5$, $\tau\ll\tau_a=0.3$~s. (a)~Find the turning points $a_{\text{hi}}$ and $a_{\text{lo}}$ of the curve $r=F(wr+I-a)$ where activity and silence disappear. (b)~Taking $r\approx1$ during activity and $r\approx0$ during silence, find the period and the active fraction.
\end{exercise}

\begin{exercise}[\lvlB]
\label{ex:20:5}
\emph{When the swings switch off.} In the model of Exercise~\ref{ex:20:4} the fixed point lies at the intersection of the curve $a=wr+I-F^{-1}(r)$ with the line $a=br$. Oscillations run while the intersection is on the middle branch, between the turning points. For what $b$ is this so? How does \nt{ACET}, by changing $b$, switch the oscillator into an amplifier?
\end{exercise}

\begin{exercise}[\lvlB]
\label{ex:20:6}
\emph{Simulation: debt or phase.} In \texttt{sleep\_models.py} take the first waking after $48$~h and keep the machine awake for $D=24$, $32$, $40$, $48$, $64$~h (\texttt{stay\_awake}, \texttt{days=8}). How long does recovery sleep last for each $D$? What determines its duration?
\end{exercise}

\begin{exercise}[\lvlB]
\label{ex:20:7}
\emph{Simulation: forgetting and interleaving.} A linear neuron $y=\mathbf w\cdot\mathbf x$, $n=40$, learns by the rule $\mathbf w\leftarrow\mathbf w-0.5\,(y-y^*)\,\mathbf x$. Tasks A and B have $10$ patterns each, $x_j\sim\mathcal N(0,1/n)$, $y^*=\pm1$. What is the mean squared error on A after $200$ epochs of A followed by $200$ epochs of B? And after $200$ epochs of A and B interleaved?
\end{exercise}

\begin{exercise}[\lvlC]
\label{ex:20:8}
\emph{Research: interleaving or renormalization?} The neuron of Exercise~\ref{ex:20:7} with a threshold; two night procedures: (i)~all weights are multiplied by $0.82$; (ii)~old patterns are replayed interleaved with new ones. What does each do to the weight ratios and to the neuron's responses? How can the hypotheses be told apart by synapse sizes after sleep?
\end{exercise}

\chapter{A Machine Made of Living Neurons}
\chaptermark{A Machine Made of Living Neurons}
\label{sec:livingmachine}
\begin{chaptergoals}
\item why a culture with no broadcast channels fires network bursts, and what sets their interval;
\item how living networks change their responses under feedback, and what remains unproven;
\item how an organoid serves as a reservoir whose linear readout is trained outside, in silicon;
\item how cultures have been taught so far, and why the teacher still stands outside the dish.
\end{chaptergoals}

\section{A test bench with an open schematic}

In Chapter~\ref{sec:debugging} we debugged a machine that has no schematic: a probe hears a thousand neurons out of eighty billion, and everything else has to be guessed. An engineer in that position would do something simple: build a small piece of the circuit on a breadboard, where every part is visible. With neurons, this can be done too.

Cortical neurons are dissociated and grown on a chip with an electrode array, from a few dozen to tens of thousands of electrodes. Within a few weeks the cells sprout processes and wire up into a network of thousands or hundreds of thousands of neurons, mostly \nt{GLUT} with a smaller fraction of \nt{GABA} that depends on the preparation: in hippocampal cultures it is about $6\%$ of the neurons~\cite{Benson1994}. Every electrode can both listen, like the probe of Chapter~\ref{sec:debugging}, and inject current, like a test signal. The second kind of bench uses \emph{organoids}: three-dimensional clumps of neural tissue about a millimeter across, grown from human stem cells~\cite{Lancaster2013}. On such benches living networks run for months; there are facilities where experiments on organoids can be run remotely, over the internet, for more than a hundred days in a row~\cite{Jordan2024}. The field is called \emph{biological computing} or \emph{organoid intelligence}~\cite{Smirnova2023}.

The bench differs from the brain in two important ways. It is tiny: it has a hundred thousand times fewer neurons. And it has no broadcast channels: a cortical culture contains no \nt{DOPA}, \nt{SERO}, \nt{NORA}, or \nt{ACET} neurons. It is a machine with a data bus and flow control but no control registers. Let us see what it can do.

\section{What the network does on its own}

Left alone, a culture neither stays silent nor produces uniform noise. From time to time almost the whole network fires at once, in a \emph{network burst}, and then falls quiet again; the intervals between network bursts range from a second to minutes, depending on the culture~\cite{Wagenaar2006}. To an engineer this is a familiar picture: an amplifier with strong positive feedback and no load turns into an oscillator.

We already know the mechanism. Strong mutual connections between \nt{GLUT} neurons set off an avalanche, as in Chapter~\ref{sec:gaba} when the conditions of Proposition~\ref{prop:ei} are violated. The avalanche quickly depletes the transmitter store in the synapses~\eqref{eq:depression}, and the network goes silent. The store recovers with time constant $\tau_{\text{rec}}$, and once the fraction $x^*$ needed for a new avalanche has recovered, a new network burst follows. The interval between network bursts is
\begin{empheq}[box=\keybox]{equation}
T_{\text{burst}} \;\approx\; \tau_{\text{rec}}\,\ln\frac{1}{1-x^*} .
\label{eq:burstinterval}
\end{empheq}
With $\tau_{\text{rec}}=0.8\,\text{s}$ from Chapter~\ref{sec:code} and $x^*=0.9$ this gives about two seconds; with $x^*=0.99$, about four. This is a model estimate with the book's $\tau_{\text{rec}}$, and it lands at the short end of the measured intervals. Direct measurement in microcultures shows the transmitter store recovering over tens of seconds after a network burst~\cite{CohenSegal2011}, so $\tau_{\text{rec}}$ is larger there; with such a $\tau_{\text{rec}}$, formula~\eqref{eq:burstinterval} gives tens of seconds to minutes, as in slow cultures. This is a relaxation oscillator, a relative of the rhythm generator of Chapter~\ref{sec:muscles}: there it was recharged by the fatigue of neuron groups, here by the transmitter store.

Network bursts are what the network does when it has nothing to do. In the living brain a similar picture appears in deep sleep (Chapter~\ref{sec:sleep}) and in the developing brain, before real inputs arrive. The resemblance is suggestive, but that the mechanisms are the same is a hypothesis.

\section{Teaching a network without a dopamine teacher}

Since the culture has no \nt{DOPA}, the ``better or worse'' signal has to be supplied by us, through the electrodes. Experiments show that a network on the bench really does change its responses, and the most interesting thing about them is \emph{what} it is taught with.

\emph{A teacher who falls silent.} The network is stimulated through one electrode once every one to three seconds until the desired response appears on another electrode $50\pm10\ms$ after the stimulus. As soon as it appears, stimulation is switched off for five minutes. After several such cycles the desired response appears right away~\cite{Shahaf2001}. The reward here is silence: the stimulus shakes the network, and ending the stimulus leaves it in the state that produced the right answer. This is gradient descent by random perturbation from Chapter~\ref{sec:dofa} (Proposition~\ref{prop:gradient}), where the shaking itself plays the role of the error: keep shaking until it comes out right.

\emph{A network that plays tennis.} In an experiment on a bench with a dense electrode array, about a million neurons were connected to a computer game of one-paddle tennis~\cite{Kagan2022}. The ball position was delivered as current to ``sensory'' electrodes: where the ball is, by a place code; how far away, by frequency. Activity in two ``motor'' regions moved the paddle up or down. The learning rule was unusual. If the paddle hit the ball, the network got a short \emph{predictable} stimulus; if it missed, it got several seconds of \emph{unpredictable} random current. Over twenty minutes of play, runs of hits grew longer. A significant improvement within a session appeared only with full feedback; with silence in place of the stimulus the culture also played better by the end of the session than without feedback, but with a smaller effect, and there was no robust learning from session to session over several days. The experiment is analyzed in detail in Chapter~\ref{sec:dishbrain}.

The authors explained the result with the hypothesis that the network acts so as to make its input predictable~\cite{Friston2010}. This is the same idea as spike economy in Chapter~\ref{sec:code} and frame prediction in Chapter~\ref{sec:recognition}: the machine spends less when the input is expected. But both the experiment and, especially, the authors' talk of the culture's ``sentience'' drew sharp criticism: the effect is small, and it can be explained more simply~\cite{Balci2023}. An honest assessment is this: the culture on the bench changes its behavior under feedback, which is measured; that it learns specifically to predict its input is a hypothesis.

\emph{A network that drives a body.} In a similar way, a culture was connected to a simple virtual ``body'' and taught to move it toward a target by choosing the spatiotemporal pattern of training stimuli~\cite{Bakkum2008}. None of these experiments involves dopamine: the teacher is a current pattern chosen by the engineer. What changes when the teacher becomes a program, or dopamine itself, is the subject of Sections~\ref{sec:dishdopamine}--\ref{sec:cartpole}.

\begin{figure}[t]
\centering
\begin{tikzpicture}[>=Stealth,
  blk/.style={draw,very thick,rounded corners=2pt,align=center,font=\footnotesize,minimum height=0.8cm,fill=blue!5}]
% (a) closed loop
\node[font=\small] at (1.5,4.3) {(a) closed loop};
\node[blk,text width=2.6cm] (G) at (1.5,3.3) {game: ball and paddle};
\draw[very thick,fill=orange!10,rounded corners=3pt] (0.5,0.2) rectangle (2.5,1.6);
\foreach \x in {0.7,1.0,...,2.35}{\foreach \y in {0.4,0.7,1.0,1.3}{\fill[gray!60] (\x,\y) circle (0.04);}}
\draw[->,thick,blue!60!black] (G.west) -| (-0.5,0.9) -- (0.5,0.9);
\node[font=\scriptsize,blue!60!black,anchor=east] at (-0.6,2.1) {ball $\to$ current};
\draw[->,thick,red!70!black] (2.5,0.9) -- (3.5,0.9) |- (G.east);
\node[font=\scriptsize,red!70!black,anchor=west] at (3.6,2.1) {spikes $\to$ paddle};
\node[font=\scriptsize] at (1.5,-0.2) {neurons on an electrode array};
\node[font=\scriptsize,align=center] at (1.5,-0.85) {hit: predictable stimulus\\miss: random current};
% (b) reservoir
\begin{scope}[xshift=7.9cm]
\node[font=\small] at (1.6,4.3) {(b) reservoir};
\node[blk,text width=1.1cm] (X) at (-0.4,1.0) {input\\$u(t)$};
\draw[very thick,fill=orange!10] (1.6,1.0) circle (0.8);
\foreach \a/\r in {20/0.35,80/0.5,150/0.3,210/0.55,280/0.4,330/0.2,110/0.15}{\fill[red!60!black] ({1.6+\r*cos(\a)},{1.0+\r*sin(\a)}) circle (0.05);}
\node[font=\scriptsize] at (1.6,-0.2) {organoid: no teacher};
\node[blk,text width=1.8cm,fill=green!8] (W) at (4.0,1.0) {readout\\$W_{\text{out}}$};
\draw[->,thick] (X) -- (0.8,1.0);
\draw[->,thick] (2.4,1.0) -- node[above,font=\scriptsize] {$\mathbf r(t)$} (W);
\draw[->,thick] (W) -- (4.0,2.3) node[above,font=\scriptsize] {$y(t)$};
\node[font=\scriptsize,green!40!black] at (4.0,0.3) {supervised training};
\end{scope}
\end{tikzpicture}
\caption{Two ways to make a living network compute. (a) Closed loop: the ball position is delivered as current, spikes from the motor regions move the paddle, and a predictable or random stimulus after each hit or miss serves as the teacher. (b) Reservoir~\eqref{eq:reservoir}: the organoid gets no training signal; it splits the input into many nonlinear responses with memory, and a single linear readout outside learns from the error. The organoid's own connections change as well, without a teacher.}
\label{fig:livingmachine}
\end{figure}

\section{A living reservoir}

There is another way to use a living network: not to train it at all. In engineering this is called \emph{reservoir computing}~\cite{Maass2002,JaegerHaas2004}. Take a ready-made nonlinear system with memory, any system, as long as its response to the input is rich and depends on the recent past. The input $u(t)$ drives it, producing a response vector $\mathbf r(t)$, and only a linear readout is trained:
\begin{empheq}[box=\keybox]{equation}
y(t) \;=\; W_{\text{out}}\,\mathbf r(t) ,
\label{eq:reservoir}
\end{empheq}
with weights fitted by the least-mean-squares rule~\eqref{eq:lms} of Chapter~\ref{sec:motorlearning}. The reservoir does what the small \nt{GLUT} neurons of Chapter~\ref{sec:motorlearning} did: it splits the input into a long vector in which the desired classes become separable by a plane (Chapter~\ref{sec:dendrites}).

An organoid on an electrode array proved suitable as such a reservoir. Speech sounds delivered to it as current patterns were classified by speaker with about $78\%$ accuracy after training a single linear readout~\cite{Cai2023}. The $78\%$ figure is the one the authors report and the press repeats. The result is useful, but it matters where the ``intelligence'' in it lies: the organoid supplies nonlinearity and fading memory, while error-driven learning happens outside, in silicon (Fig.~\ref{fig:livingmachine}). The organoid does not stay unchanged, though: according to the authors, its own functional connectivity changed during training, and they call this unsupervised learning within the organoid itself~\cite{Cai2023}. How much of the accuracy comes from this change and how much from the readout has not been separated.

\section{What the bench tells us about the machine}

\emph{Energy.} By estimate~\eqref{eq:energy} a spike costs about $10^{-10}$~J. A culture of a million neurons firing at one spike per second spends on signaling
\begin{empheq}[box=\keybox]{equation}
P \;\approx\; 10^{6}\times1\,\text{s}^{-1}\times10^{-10}\,\text{J} \;=\; 10^{-4}\,\text{W},
\label{eq:dishpower}
\end{empheq}
a tenth of a milliwatt. The electronics that stimulate, record, and process these spikes spend orders of magnitude more. As in Chapter~\ref{sec:debugging}, the bottleneck is not the computation but the connection to it.

\emph{Missing parts.} The bench shows how much a \nt{GLUT} bus with \nt{GABA} flow control can do without control registers: self-organize, generate network bursts, change its responses under feedback, and serve as a good reservoir. What it cannot do without them is decide for itself what counts as success. On the bench that signal is supplied by the engineer; in the brain, as we saw in Chapters~\ref{sec:dofa}--\ref{sec:acet}, by the broadcast channels.

\emph{Ethics.} Organoids are grown from human cells, and the more complex they become, the more serious the question of whether such tissue can feel anything. The engineering view suggests a partial answer: pain in Chapter~\ref{sec:pain} is not the signal of a single sensor but a whole alarm system with gates, a volume control, and broadcast channels, none of which the bench has. But this is an argument, not a proof, and the field itself proposes carrying out ethical assessment alongside the experiments, from the very start~\cite{Smirnova2023}.

\begin{keyblock}
\begin{proposition}[A living network on the bench]
\label{prop:livingmachine}
A network of \nt{GLUT} and \nt{GABA} neurons on an electrode array generates network bursts on its own~\eqref{eq:burstinterval}, changes its responses under feedback, and can serve as a reservoir~\eqref{eq:reservoir} for a readout trained outside. All of this is measured. That such a network learns by some general rule, for example to predict its input, is a hypothesis, and most specialists do not accept loud claims about its ``sentience.''
\end{proposition}
\end{keyblock}

\section{Dopamine on a flash of light}
\label{sec:dishdopamine}

The only direct experiment we know of in which a culture was given dopamine as a teacher was run on a platform of organoids that researchers access remotely~\cite{Jordan2024}. Dopamine in a light-sensitive ``cage'' was added to the fluid bathing the organoids: the bound molecule does not act on receptors. The concentration of caged dopamine is $1$~mM. When the network is to be rewarded, it is illuminated with ultraviolet light: the cage breaks apart, and dopamine is released into the volume.

The loop works like this. The same stimulus is delivered over and over; if it evoked more spikes than before, a flash is triggered and dopamine is released. Over $240$ repetitions the number of evoked spikes grew overall. The authors themselves write that a single experiment of this kind cannot show that the growth was caused by dopamine~\cite{Jordan2024}.

For an engineer this is a first attempt to build the three-factor rule~\eqref{eq:threefactor} in a dish: sender and recipient activity leaves an eligibility trace, and a common signal decides whether to lock in the change. But here the signal can only be positive: reward is either present or absent, and the setup cannot deliver a negative error $\delta<0$. Dopamine also spreads out and is cleared slowly, like the concentration in~\eqref{eq:lowpass}, so frequent flashes may simply raise its overall level. To tell learning apart from this, two controls are needed: the same number of flashes at random times unrelated to the network's response, and the same flashes without caged dopamine. A test after rest is also needed: does the change persist once the dopamine is gone?

\section{Dopamine teaches silicon}
\label{sec:biohybrid}

In the reverse experiment, living cells teach electronics~\cite{Keene2020}. Dopamine-releasing cells are grown in a microchannel above an organic neuromorphic element, a transistor whose channel conductance plays the role of a synaptic weight. Dopamine released by the cells is oxidized at the element's electrode, and this changes the conductance. The device also mimics the mechanism that clears transmitter from the cleft, so the weight can not only be changed for a long time but also reset.

In circuit terms this is a leaky integrator with a chemical input: the weight accumulates the dopamine flux, and ``clearance'' sets how fast it is forgotten, the same RC circuit as in~\eqref{eq:lowpass}. What matters here is the direction of the link. The probe of Chapter~\ref{sec:debugging} cannot see the broadcast registers because they are chemical. This element reads exactly a chemical register and turns it into an electrical weight. For brain--computer interfaces this is a path toward a sensor that hears not only the data bus but also the control registers.

\section{Organoids with their own dopamine neurons}
\label{sec:midbrainorg}

Organoids resembling the brain region that houses the \nt{DOPA} neurons can be grown from human stem cells. They contain working dopamine neurons that release dopamine~\cite{Jo2016}. Such organoids are grown mainly to study disease. We found no experiments in which their dopamine served as a teacher for another network.

For a machine made of living neurons this is the missing part. The bench of Chapter~\ref{sec:livingmachine} is a data bus and flow control without control registers; here a source of the broadcast signal already exists. What is missing is a critic: a network that would compute the prediction error~\eqref{eq:ctd} and drive these neurons. Connecting a cortical organoid to such an organoid so that dopamine is released on error is an open problem, and for now it is a hypothesis, not an experiment.

\section{An organoid balances a pole without dopamine}
\label{sec:cartpole}

The most complete of the recent experiments does without dopamine altogether~\cite{Robbins2026}. Mouse cortical organoids were connected in a closed loop to the cart-pole task: organoid activity moves the cart, and the pole position is fed back by stimulation. Between trials the organoid receives short high-frequency training stimuli. Which ones is chosen by a reinforcement-learning program.

The results are measured:
\begin{itemize}
\item in most organoids, training signals chosen by the program gave better performance than randomly chosen signals or no signal;
\item after $45$ minutes of rest the improvement did not persist;
\item blocking both types of glutamate receptors, AMPA and NMDA, abolished the improvement.
\end{itemize}

The last point says where the learned change lives: in \nt{GLUT} synapses, and through the receptor that in Section~\ref{sec:weightstore} works as a detector of input--output coincidence. The second point shows what is missing: there is a fast change but no consolidation, like the fast learner of Chapter~\ref{sec:motorlearning} without the slow one. And the teacher's role has changed: the engineer choosing a current pattern has been replaced by a program, but the teacher still stands outside the dish.

Among earlier experiments on training cultures on electrode arrays we likewise found none that used dopamine: the training signal was always electrical stimulation.

\section{Who teaches the culture}

\begin{table}[t]
\centering
\footnotesize
\setlength{\tabcolsep}{4pt}
\renewcommand{\arraystretch}{1.15}
\begin{tabular}{>{\raggedright}p{2.7cm}>{\raggedright}p{2.5cm}>{\raggedright}p{3.2cm}>{\raggedright\arraybackslash}p{4.4cm}}
\toprule
Experiment & Who teaches & Training signal & What is known \\
\midrule
stimulation stops at the desired response & engineer & end of stimulation & the response is locked in: measured \\
DishBrain (ch.~\ref{sec:dishbrain}) & engineer & predictable or random current & significant growth of rallies within a session only with full feedback, a smaller effect with silence: measured; not robust from session to session; the cause is debated \\
cart-pole & reinforcement-learning program & high-frequency stimuli chosen by the program & better than random, but not retained after rest: measured \\
dopamine on a flash & rule ``more spikes means reward'' & dopamine released by light & growth in spikes: observation; role of dopamine not proven \\
biohybrid synapse & living cells & dopamine from cells & long-lasting change and reset of the element's weight: measured \\
organoids with \nt{DOPA} neurons & --- & --- & a dopamine source exists; not used as a teacher \\
\bottomrule
\end{tabular}
\caption{Who teaches living neurons on a chip, and with what.}
\label{tab:dishteachers}
\end{table}

In every experiment where the culture itself learns, the teacher is still outside (Table~\ref{tab:dishteachers}): an engineer, a program, or a rule set by an engineer. Dopamine has entered the dish only once, and its role remains to be proven. Building a real broadcast channel in a dish, with a critic, an error, and an eligibility trace as in Chapter~\ref{sec:dofa}, is the next step, and no one has taken it yet.


\section{Summary}

\begin{table}[t]
\centering
\footnotesize
\setlength{\tabcolsep}{4pt}
\renewcommand{\arraystretch}{1.15}
\begin{tabular}{>{\raggedright}p{2.8cm}>{\raggedright}p{4.2cm}>{\raggedright\arraybackslash}p{5.8cm}}
\toprule
Experiment & What the living network does & What is known \\
\midrule
network with no input & network bursts at intervals from a second to minutes~\eqref{eq:burstinterval} & measured; the mechanism via synaptic depletion is the accepted model \\
a teacher who falls silent & the desired response is locked in when stimulation is switched off & measured \\
playing tennis & runs of hits grow longer; significant improvement within a session only with full feedback & behavior change measured; learning from session to session not robust; effect size and the input-prediction explanation disputed \\
virtual body & movement toward a target with a chosen stimulus pattern & measured \\
reservoir & nonlinearity and memory for the readout~\eqref{eq:reservoir} & measured; error-driven learning is outside, in silicon; the organoid's connections also change \\
energy & about $10^{-4}$~W per million neurons~\eqref{eq:dishpower} & estimate from~\eqref{eq:energy} \\
\bottomrule
\end{tabular}
\caption{What machines made of living neurons have shown.}
\label{tab:livingmachine}
\end{table}

A machine made of living neurons (Table~\ref{tab:livingmachine}) is a bench that shows what a data bus and flow control can do without control registers. They can do a lot, but they cannot decide on their own what is good. On the bench that role is played by the engineer with a chosen current pattern; in the brain, by the few broadcast wires to which the middle of the book is devoted.

\section{Exercises}

\begin{exercise}[\lvlA]
\label{ex:21:1}
\emph{Interval between network bursts.} A network burst depletes the store to zero, then $\tau_{\text{rec}}\,\dd x/\dd t=1-x$, and a new network burst starts at $x=x^*$; $\tau_{\text{rec}}=0.8$~s. (a)~Find the interval for $x^*=0.9$ and $0.99$. (b)~The threshold $x^*$ varies randomly by $\pm0.005$ around $0.99$. What range do the intervals span?
\end{exercise}

\begin{exercise}[\lvlA]
\label{ex:21:2}
\emph{Energy of the bench.} (a)~What power does estimate~\eqref{eq:dishpower} give for the whole brain ($8.6\cdot10^{10}$ neurons)? (b)~The bench records $1024$ electrodes at $20$~kHz with $10$~bits per sample. At $1$~nJ per bit, how many times more does transmitting this stream cost than the power of the culture itself?
\end{exercise}

\begin{exercise}[\lvlB]
\label{ex:21:3}
\emph{Design: suppressing network bursts.} Sparse stimulation through many electrodes makes every synapse release transmitter at rate $r_b$; the store obeys $\tau_{\text{rec}}\,\dd x/\dd t=1-x-Ur_b\tau_{\text{rec}}x$, $U=0.5$, $\tau_{\text{rec}}=0.8$~s. What is the smallest $r_b$ at which the store never reaches $x^*=0.9$; $0.99$, and how much weaker does the synapse become?
\end{exercise}

\begin{exercise}[\lvlB]
\label{ex:21:4}
\emph{Reservoir memory is at most the number of neurons.} A reservoir with $N$ outputs $\mathbf r(t)$ receives a white input $u(t)$ with zero mean and unit variance; $m_k$ is the largest squared correlation of a readout $\mathbf w_k\cdot\mathbf r(t)$ with $u(t-k)$. Prove that the memory capacity $\mathrm{MC}=\sum_{k\ge0}m_k\le N$.
\end{exercise}

\begin{exercise}[\lvlB]
\label{ex:21:5}
\emph{Simulation: a reservoir in silicon.} Reservoir $\mathbf r(t+1)=\tanh(W\mathbf r(t)+\mathbf v\,u(t)+\mathbf c)$, $N=30$, $W$ random with spectral radius $0.9$, $v_i\sim U(-0.5,0.5)$, $c_i\sim U(-0.2,0.2)$, $u\sim U(-1,1)$, readout by ridge regression. Find the memory capacity of Exercise~\ref{ex:21:4} with and without $\tanh$. Which reservoir remembers more?
\end{exercise}

\begin{exercise}[\lvlB]
\label{ex:21:6}
\emph{A readout for a living reservoir.} An organoid provides $1024$ channels and $P=240$ training recordings of two classes. (a)~How many features $n$ can be kept so that, by the formula of Exercise~\ref{ex:19:3}, a random labeling is separable with probability at most $1\%$? (b)~How many sigmas above chance is an accuracy of $78\%$ on $100$ test recordings?
\end{exercise}

\begin{exercise}[\lvlC]
\label{ex:21:7}
\emph{Research: a teacher without a broadcast channel.} Propose a stimulation protocol that implements the three-factor rule of Chapter~\ref{sec:dofa} on the bench through $8$--$1024$ electrodes, and a frozen-readout control that distinguishes weight learning from a shift in network state. What result would refute the learning hypothesis?
\end{exercise}

\chapter{DishBrain}
\chaptermark{DishBrain}
\label{sec:dishbrain}
\begin{chaptergoals}
\item how DishBrain codes the ball by place and rate and reads the paddle as a count difference;
\item why a count in a $10\ms$ window is noisy, so the network can shift only the mean;
\item how shaking the network only after misses drifts it toward configurations that fail less;
\item what was measured, what is disputed, and how to rebuild the bench to test the mechanism.
\end{chaptergoals}

\section{A network in a dish plays tennis}

DishBrain is the name its authors gave to a bench on which a culture of living neurons, grown on a chip with electrodes, plays a simplified one-paddle tennis game~\cite{Kagan2022}. It is the most discussed experiment with a machine made of living neurons (such machines are surveyed in Chapter~\ref{sec:livingmachine}), and it is especially convenient for this book. Here a small network is accessible in its entirety: the experimenter sets every input, and every output is recorded. All the questions of the preceding chapters, how the signal is encoded, who teaches, what the probe sees, can be put to a network that has no eyes, no muscles, and no \nt{DOPA} neurons. Let us take the experiment apart the way an engineer takes apart someone else's bench: circuit, input, output, teacher, result, disputes.

\section{The bench}

The neurons are taken from embryonic mouse cortex, about $800\,000$ cells per chip, or grown from human stem cells, about a million per chip. For several weeks they grow on the chip, grow into one another, and begin to produce spikes and bursts on their own. Beneath them, on an $8$~mm$^2$ area, lie $26\,000$ platinum electrodes; up to $1024$ of them can be recorded at once, and in practice stimulation goes through eight.

A loop is closed around the culture (Fig.~\ref{fig:dishloop}). A computer simulates the game: the ball flies, bounces off the walls, and the paddle moves up and down. The ball position is turned into current on eight ``sensory'' electrodes. Spikes from two ``motor'' regions are counted in $10\ms$ windows and move the paddle. After each hit or miss the culture receives one more stimulus, the feedback. The $10\ms$ windows are the clock of the bench computer, not of the culture: the network itself still runs asynchronously, like every network in this book.

\begin{figure}[t]
\centering
\begin{tikzpicture}[>=Stealth,
  blk/.style={draw,very thick,rounded corners=3pt,align=center,font=\small,minimum height=1.2cm},
  lab/.style={font=\scriptsize,align=center}]
\node[blk,fill=blue!6,text width=3.4cm] (C) at (0,0) {culture on a chip\\\scriptsize $\sim10^6$ neurons, $26\,000$ electrodes};
\node[blk,fill=orange!10,text width=3.0cm] (G) at (7.2,0) {computer:\\tennis game};
\draw[->,very thick,blue!60!black] (G.north west) to[bend right=25] node[lab,above] {ball: \emph{place} = which of 8 electrodes,\\\emph{rate} = 4--40 spikes/s} (C.north east);
\draw[->,very thick,red!70!black] (C.south east) to[bend right=25] node[lab,below] {paddle: spike counts of the\\``up'' and ``down'' regions per $10\ms$} (G.south west);
\draw[->,thick,densely dashed,green!45!black] (G.east) -- ++(0.9,0) |- ++(0,-2.6) -| node[lab,pos=0.25,below] {hit: predictable stimulus; miss: unpredictable} (C.south);
\end{tikzpicture}
\caption{The DishBrain loop. Blue arrow: input; the ball position is encoded by a place code and a rate code (Chapter~\ref{sec:code}). Red: output; the difference between the counts of two regions moves the paddle. Green dashed: feedback after each rally, the bench's only ``teacher.''}
\label{fig:dishloop}
\end{figure}

\section{Input and output}

The \emph{input} is written in two codes of Chapter~\ref{sec:code} at once. Which of the eight electrodes is stimulated tells how high the ball is: this is a place code, like that of the sensors in Chapter~\ref{sec:sensors}. How often the stimulation pulses come tells how far away the ball is: $4$ per second when it is at the far wall, rising linearly to $40$ per second when it is at the paddle's wall. This is a rate code; the amplitude of the ball pulses is $75$~mV. The culture does not ``know'' either code in advance: the code is set by the experimenter, and the network has to learn to read it.

The \emph{output} is read the same way as in the interfaces of Chapter~\ref{sec:debugging}. Spikes of one motor region push the paddle up, those of the other push it down; in the simplest form, for each window,
\begin{empheq}[box=\keybox]{equation}
\Delta p \;=\; c\,\bigl(n_\uparrow - n_\downarrow\bigr),
\label{eq:dishreadout}
\end{empheq}
where $n_\uparrow$, $n_\downarrow$ are the spike counts of the two regions in $10\ms$, and $c$ is a bench constant. This is the two-wire code of Chapter~\ref{sec:glutcomputer}: the direction of movement is carried not by the sign of a signal but by which of the two wires is louder.

Consider the noise of such an output. If a region emits a total of $R$ spikes per second, a window $T=10\ms$ holds $RT$ of them on average, and by~\eqref{eq:rateerror} the relative spread is $1/\sqrt{RT}$. At $R=1000$ spikes per second this is about $30$ percent in each window; at $R=100$, more than a hundred. The paddle inevitably jitters, and the network can learn in only one way: by shifting the \emph{mean} count difference in the right direction. These are the same laws that limit any interface of Chapter~\ref{sec:debugging}, only now on both sides of the loop.

\section{A teacher without dopamine}

The most interesting part of the bench is the teacher. A culture of cortical neurons has no \nt{DOPA} neurons, and the bench sends no ``better than expected'' signal. The feedback is different. If the network hits the ball, a short \emph{predictable} stimulus goes to all eight sensory electrodes at once: $75$~mV pulses at $100$ per second for $0.1\,\text{s}$. If it misses, for four seconds it receives a stimulus the authors call \emph{unpredictable}: $150$~mV pulses at $5$ per second, at random times, on randomly chosen electrodes; then four seconds of pause without stimulation, and the ball starts again~\cite{Kagan2022}.

The authors started from the hypothesis that the network acts so as to make its input predictable~\cite{Friston2010}. For an engineer the meaning is simple: the culture has exactly one way to get rid of four seconds of random noise, which is to hit the ball. We met the same idea in the spike economy of Chapter~\ref{sec:code} and in frame prediction in Chapter~\ref{sec:recognition}.

But is predictability really what is needed here? A cruder mechanism is possible. Suppose the unpredictable stimulus after a miss simply \emph{scrambles} recent changes in the network, while the calm stimulus after a hit leaves them alone. Describe the network's configuration by a single vector $\boldsymbol\psi$, and let the network miss with probability $P_{\text{miss}}(\boldsymbol\psi)$ in configuration $\boldsymbol\psi$. If the kicks are symmetric, so that going from $\boldsymbol\psi$ to $\boldsymbol\psi'$ is as likely as the reverse, then in steady state the outflow from each configuration balances the inflow, and the fraction of time the network spends in configuration $\boldsymbol\psi$ is
\begin{empheq}[box=\keybox]{equation}
\rho(\boldsymbol\psi) \;\propto\; \frac{1}{P_{\text{miss}}(\boldsymbol\psi)} .
\label{eq:dishdensity}
\end{empheq}
A configuration that misses ten times less often is held ten times longer. No one tells the network which way to change its weights: good configurations are simply destroyed less often.

We tested this on a model that reduces the culture to two numbers: the paddle arrives at the point $a\,y+b$ when the ball arrives at height $y$, and it hits if it misses by less than $0.1$. After a miss both numbers get a random kick; after a hit they stay the same. The fraction of hits over $2000$ rallies grew from $0.21$ in the first two hundred to $0.38$ in the last five hundred (forty ``cultures,'' Fig.~\ref{fig:dishmodel}). If the kick is given after \emph{every} rally, the feedback carries no information, and the fraction stays around $0.12$; if there are no kicks at all, it stays at $0.19$. The bench experiment agrees with this: a significant improvement within a session appeared only with full feedback. With silence in place of the stimuli the culture still played better by the end of the session than without feedback, but with a smaller effect~\cite{Kagan2022}; note that in this condition, too, a miss is followed by a long pause and a hit by a short one, so the difference between outcomes does not vanish entirely.

\begin{figure}[t]
\centering
\begin{tikzpicture}[x=1cm,y=5cm]
\draw[->,gray!70!black] (0,0) -- (10.6,0);
\draw[->,gray!70!black] (0,0) -- (0,0.5) node[above,font=\small,black] {fraction of hits};
\foreach \v in {0.1,0.2,0.3,0.4}{\draw[gray!60!black] (-0.06,\v) -- (0.06,\v); \node[left,font=\scriptsize] at (0,\v) {\v};}
\foreach \x/\f/\l/\lab in {1.8/0.21/0.38/{noise after\\a miss},5.2/0.13/0.11/{noise after\\every rally},8.6/0.19/0.19/{no noise}}{
  \fill[blue!35] (\x-0.8,0) rectangle (\x-0.05,\f);
  \fill[red!60!black] (\x+0.05,0) rectangle (\x+0.8,\l);
  \node[below,font=\scriptsize,align=center] at (\x,-0.01) {\lab};
  \node[above,font=\scriptsize] at (\x-0.425,\f) {\f};
  \node[above,font=\scriptsize] at (\x+0.425,\l) {\l};}
\fill[blue!35] (7.4,0.47) rectangle (7.7,0.495); \node[right,font=\scriptsize] at (7.75,0.482) {first 200};
\fill[red!60!black] (7.4,0.41) rectangle (7.7,0.435); \node[right,font=\scriptsize] at (7.75,0.422) {last 500};
\end{tikzpicture}
\caption{The ``scramble on a miss'' model (\texttt{dishbrain\_model.py}): fraction of hits at the start and at the end of $2000$ rallies, averaged over forty model cultures. Learning occurs only if the kick follows a miss and not a hit, as in the experiment, where a significant improvement within a session appeared only with full feedback.}
\label{fig:dishmodel}
\end{figure}

\begin{keyblock}
\begin{proposition}[The scrambling teacher]
\label{prop:dishscramble}
If failure shakes the network and success leaves it alone, the network drifts by itself toward configurations that fail less often: the time spent in a configuration is inversely proportional to its probability of failure~\eqref{eq:dishdensity}. Such learning needs no reward signal, no sign of that signal, and no predictability as such; it is enough that the feedback after success differs from the feedback after failure. The change in the culture's behavior is measured; which mechanism works in the dish, input prediction or scrambling, has not been established.
\end{proposition}
\end{keyblock}

Compare this with Chapter~\ref{sec:dofa}. There, noise also served as search, but dopamine carried a sign: the error $\delta$ said which way to shift the weights, and the step followed the gradient~\eqref{eq:perturbation}. Here there is no sign, only ``keep'' or ``shake.'' Such a search is cruder and slower, but it asks less of the network: no \nt{DOPA} neurons, no critic, no eligibility trace stretched over seconds.

\section{What was measured and what is disputed}

Each game lasted $20$ minutes; the first $5$ minutes were compared with the last $15$. In cultures with full feedback, runs of hits grew longer and immediate misses became fewer; in controls with no cells, no game, or a computer-driven ``paddle,'' nothing of the kind happened. Cultures of human cells on average sustained rallies longer than mouse cultures. The improvement is statistically reliable, on nine mouse and eleven human cultures, with a hundred-odd games in each group, but it is small: the network did not become a good player; it became slightly better than chance. It is reliable within a session; the authors did not obtain robust learning from session to session over several days~\cite{Kagan2022}. A follow-up study by the same group found that during play the culture's activity approaches a \emph{critical} regime, the boundary between dying out and an avalanche, the same life on the edge as in Chapter~\ref{sec:gaba}, and the closer it gets, the better the play~\cite{Habibollahi2023}.

The main dispute was over words, not numbers. The paper's title claimed that neurons in a dish ``exhibit sentience,'' and a group of researchers objected: behavior change in a closed loop is not yet intelligence or sensation, and conclusions should be stated more modestly~\cite{Balci2023}. From an engineering standpoint two questions are open, and both concern the mechanism. First: does the network learn, that is, do its weights change for a long time, or does the feedback merely shift its current state? Second: is predictability specifically needed, or is any difference between the response to success and to failure enough, as in Proposition~\ref{prop:dishscramble}? A bench on which every electrode is accessible can answer both, and that is perhaps its main value.

\section{Bench specification: how to reproduce it}
\label{sec:dishspec}

Below is everything needed to build such a bench again: hardware, loop, input encoding, output readout, feedback, and experimental protocol. Each parameter is marked with its source. ``Paper'' means the value is taken from the text of~\cite{Kagan2022}. ``Choice'' means the paper does not give it, and to reproduce the bench you will have to set it yourself; we propose a default value and say how to check it.

\subsection*{Hardware and culture}

\begin{center}
\footnotesize
\setlength{\tabcolsep}{4pt}
\renewcommand{\arraystretch}{1.15}
\begin{tabular}{>{\raggedright}p{4.0cm}>{\raggedright}p{6.9cm}>{\raggedright\arraybackslash}p{1.6cm}}
\toprule
Parameter & Value & Source \\
\midrule
chip & MaxOne array: $26\,000$ platinum electrodes on an $8$~mm$^2$ area & paper \\
recording & up to $1024$ channels at once, $20\,000$ samples per second & paper \\
stimulation & up to $32$ electrodes technically, $8$ independent ones in the experiment & paper \\
cells & embryonic mouse cortex; human neurons from stem cells (two culture methods); HEK293T cells (not neurons) as a control & paper \\
plating & about $10^6$ cells per chip & paper \\
culture age at the experiment & mouse: from $\sim10$ days; human: $14$--$24$ days or $\sim80$ days depending on the culture method & paper \\
\bottomrule
\end{tabular}
\end{center}

\subsection*{Loop and timing}

The loop runs in $10\ms$ windows ($200$ samples): in each window the spikes on the motor-region electrodes are counted, the game is updated, and the next stimulation is delivered. The delay from a spike in the culture to the responding stimulus is about $5\ms$. During play each electrode's signal passes through a 2nd-order Bessel high-pass filter with a $100$~Hz cutoff; the signal's magnitude is smoothed by a 1st-order Bessel low-pass filter at $1$~Hz, and the spike threshold is proportional to this smoothed magnitude. The paper gives a threshold of six standard deviations of the background for separate activity measurements, not for play. To keep stimulation from being counted as a response of the culture, all signals are blanked when more than $15$ large spikes, above $75$~mV, arrive at once. All of this is from the paper.

\subsection*{Ball encoding}

\begin{center}
\footnotesize
\setlength{\tabcolsep}{4pt}
\renewcommand{\arraystretch}{1.15}
\begin{tabular}{>{\raggedright}p{3.4cm}>{\raggedright}p{7.5cm}>{\raggedright\arraybackslash}p{1.6cm}}
\toprule
Parameter & Value & Source \\
\midrule
sensory region & $8$ electrodes arranged so that neighboring electrodes mean neighboring ball positions & paper \\
place code & electrode number $k=1,\dots,8$ gives the ball position relative to the paddle center & paper \\
which coordinate & vertical ball position; for reproduction, relative to the paddle, $k=\lceil 8(y-p+\tfrac12)\rceil$ for $y-p\in[-\tfrac12,\tfrac12]$ & paper; formula is a choice \\
rate code & stimulation rate from $4$ to $40$~pulses/s carries the distance to the paddle & paper \\
scale direction & $4$~pulses/s when the ball is at the opposite wall, rising linearly to $40$~pulses/s at the paddle's wall: $f=4+36\,(1-x/L)$, $x$ is the distance to the paddle's wall, $L$ is the court width & paper \\
pulse shape & short rectangular biphasic: positive phase first, then negative & paper \\
amplitude & $75$~mV; chosen so as not to force inhibited cells to fire & paper \\
phase duration & not given; use the stimulation system's standard setting and keep it fixed throughout the experiment & choice \\
\bottomrule
\end{tabular}
\end{center}

\subsection*{Paddle readout}

The paper has two motor regions: spikes in the first move the paddle up, in the second down; the regions' contributions are weighted to account for different cell densities and activity~\cite{Kagan2022}. The exact formula is not given. For reproduction we take rule~\eqref{eq:dishreadout}, $\Delta p=c\,(n_\uparrow-n_\downarrow)$ for each $10\ms$ window, with a weight that equalizes the regions by their mean resting activity (choice): $n_\uparrow$ and $n_\downarrow$ are divided by their region's mean count per window, measured before play. We choose the coefficient $c$ so that a difference of one mean count moves the paddle by $1\%$ of the court height per window (choice); a constant advantage of one region then carries the paddle across the whole court in $1$~s.

The paper's text does not give the regions' layout on the chip in coordinates; there is only a diagram. For reproduction, three non-overlapping groups of electrodes suffice: eight sensory ones in the center and a group of recording electrodes on each side, one ``up'' and the other ``down'' (choice).

\subsection*{The game}

The court size, paddle length, and ball speed are not given in the paper; it says only that the ball flies at constant speed and bounces off the walls. For reproduction (choice): a $1\times1$ court, a paddle of length $0.2$ on the right wall (a hit when $|y-p|<0.1$, as in the model~\texttt{models/dishbrain\_model.py}), and a ball that crosses the court in $1$~s. A miss with no ball returned in the rally counts as an ``ace''; a rally of more than three consecutive hits counts as ``long'' (paper).

\subsection*{Feedback}

\begin{center}
\footnotesize
\setlength{\tabcolsep}{4pt}
\renewcommand{\arraystretch}{1.15}
\begin{tabular}{>{\raggedright}p{2.6cm}>{\raggedright}p{8.3cm}>{\raggedright\arraybackslash}p{1.6cm}}
\toprule
Event & What is delivered & Source \\
\midrule
hit & predictable stimulus: all $8$ sensory electrodes at once, $75$~mV, $100$~pulses/s, $100\ms$; the usual ball stimulation is interrupted for this time & paper \\
miss & unpredictable stimulus: $150$~mV, $5$~pulses/s, $4$~s, at random times on electrodes chosen at random from the $8$; then $4$~s without stimulation, and the ball starts in a random direction & paper \\
randomness law & not given; for reproduction, pulse times as a Poisson stream with mean rate $5$~pulses/s & choice \\
\bottomrule
\end{tabular}
\end{center}

\subsection*{Protocol and controls}

One session lasts $20$~min; the first $5$ and the last $15$~min are compared (paper). Three outcome measures: mean rally length, fraction of aces, and fraction of long rallies. Experimental conditions (paper):
\begin{itemize}
\item \emph{stimulus}: the full loop, as described above;
\item \emph{silence}: in place of the hit or miss stimulus, the same time with no stimulation at all, then the ball starts in a random direction;
\item \emph{no feedback}: after a miss the ball simply bounces and flies on, and ball-position stimulation continues;
\item \emph{rest}: the game runs and the paddle moves according to activity, but there is no stimulation at all;
\item \emph{controls}: a chip with medium only; HEK293T cells; an ``in silico culture,'' a simulation instead of cells.
\end{itemize}
Sample size in the main series: $9$ mouse cultures ($101$ sessions), $11$ human cultures ($138$), $6$ medium-only chips ($80$), $20$ cultures at rest ($42$), $3$ simulations ($38$), $399$ sessions in all. In the feedback series: $486$ sessions over $3$ days at $3$ sessions per day, conditions ``stimulus,'' ``silence,'' and ``no feedback'' ($n=27$, $17$, and $15$). Growth of the mean rally length from the first $5$ minutes to the last $15$ was obtained only in living cultures. In the feedback series, significant growth over time occurred only in the ``stimulus'' condition; at the end of the session the ``silence'' condition still outperformed ``no feedback,'' but with a smaller effect, and there was no robust learning from session to session over the three days~\cite{Kagan2022}.

\subsection*{The bench cycle step by step}

Every $10\ms$ the bench computer does the following.
\begin{enumerate}
\item Reads the spike counts $n_\uparrow$, $n_\downarrow$ in the two motor regions over the past window and normalizes them by the region's mean resting count.
\item Moves the paddle by $\Delta p=c\,(n_\uparrow-n_\downarrow)$ and clips it to the edges of the court.
\item Moves the ball one step; reflects it off the top, bottom, and left walls.
\item If the ball has reached the right wall: if $|y-p|<0.1$, it is a hit, the rally count grows by one, and the hit stimulus is delivered ($100\ms$); otherwise it is a miss, the rally is recorded, the miss stimulus is delivered ($4$~s), followed by a $4$~s pause, and the ball starts again.
\item Otherwise, chooses electrode $k$ from the ball's vertical position and rate $f$ from the distance to the paddle's wall, and delivers $75$~mV biphasic pulses at rate $f$ to electrode $k$ until the next window.
\end{enumerate}
This is enough to build a bench compatible with the published one and to test the chapter's main question: is the culture taught by predictability, or by a simple difference between the response to a hit and to a miss? For that, a fourth condition absent from the paper should be added to its three: after a miss, a \emph{predictable} stimulus of the same duration and energy as the unpredictable one. If the culture learns in this case too, the cause is the difference, not predictability.


\section{Summary}

\begin{table}[t]
\centering
\footnotesize
\setlength{\tabcolsep}{4pt}
\renewcommand{\arraystretch}{1.15}
\begin{tabular}{>{\raggedright}p{2.1cm}>{\raggedright}p{4.2cm}>{\raggedright}p{1.9cm}>{\raggedright\arraybackslash}p{4.6cm}}
\toprule
Bench unit & How it works & Chapter & What is known \\
\midrule
input & place (which of 8 electrodes) and rate (4--40 spikes/s) & \ref{sec:code}, \ref{sec:sensors} & set by the experimenter \\
output & count difference of two regions per $10\ms$~\eqref{eq:dishreadout} & \ref{sec:glutcomputer}, \ref{sec:debugging} & set by the experimenter; noise by~\eqref{eq:rateerror} \\
teacher & predictable stimulus after a hit, unpredictable after a miss; no dopamine & \ref{sec:dofa} & significant improvement within a session only with full feedback; smaller effect with silence; not robust from session to session \\
mechanism & input prediction or scrambling on a miss~\eqref{eq:dishdensity} & \ref{sec:dofa}, \ref{sec:recognition} & not established; both explain the result \\
network regime & approach to the critical regime during play & \ref{sec:gaba} & link to play quality measured \\
``sentience'' & the authors' claim & --- & not shown; disputed \\
\bottomrule
\end{tabular}
\caption{DishBrain through an engineer's eyes.}
\label{tab:dishbrain}
\end{table}

DishBrain is a machine made of living neurons in its simplest form: the input is set by a code, the output is read by counting, and the teacher is the difference between the response to success and to failure (Table~\ref{tab:dishbrain}). A network without eyes, muscles, or dopamine learned to play just a little, and that alone makes it a test bench for the ideas of this book. Whatever the mechanism turns out to be, prediction, scrambling, or something else, the answer will be found as Chapter~\ref{sec:debugging} advises: with a probe, a test signal, and switching parts off, on a machine that can at last be seen in its entirety.

\section{Exercises}

\begin{exercise}[\lvlA]
\label{ex:22:1}
\emph{How far must the count shift?} Two regions together emit $R_\uparrow+R_\downarrow=2000$ spikes per second, Poisson streams. (a)~What is the relative spread of one region's count in $10\ms$ at $R=1000$ and $R=100$? (b)~Over a $1$~s ball flight the displacements~\eqref{eq:dishreadout} add up. What is the smallest $R_\uparrow-R_\downarrow$ at which the mean displacement is twice the spread?
\end{exercise}

\begin{exercise}[\lvlA]
\label{ex:22:2}
\emph{How many rallies are needed to see learning?} Rallies are independent, and the fraction of hits is binomial. How many rallies are needed in each of two periods for the growth in the fraction of hits to exceed two standard deviations of the difference: (a)~from $0.20$ to $0.25$; (b)~from $0.21$ to $0.38$, as in the model?
\end{exercise}

\begin{exercise}[\lvlB]
\label{ex:22:3}
\emph{Deriving~\eqref{eq:dishdensity}.} The configuration $\boldsymbol\psi$ ranges over a finite set; on a miss (probability $P(\boldsymbol\psi)>0$) the network moves to $\boldsymbol\psi'$ with symmetric probability $K(\boldsymbol\psi,\boldsymbol\psi')$, and on a hit it stays. (a)~Prove that $\rho\propto1/P$ is the stationary distribution. (b)~What is the fraction of misses with two configurations with $P=0.9$ and $P=0.1$?
\end{exercise}

\begin{exercise}[\lvlB]
\label{ex:22:4}
\emph{The model's absorbing region.} The paddle arrives at $p=\min\bigl(1,\max(0,ay+b)\bigr)$, the ball at height $y\sim U(0,1)$, a hit if $|p-y|<h=0.1$. (a)~Prove that for $|b|<h$ and $|a-1+b|<h$ the network hits every ball, and find the area of this region. (b)~Find numerically the mean fraction of hits for $a\sim U(-1,1)$, $b\sim U(0,1)$.
\end{exercise}

\begin{exercise}[\lvlB]
\label{ex:22:5}
\emph{Simulation: a fence around the configurations.} Give $200$ cultures $20\,000$ rallies each in a copy of \texttt{dishbrain\_model.py}. What is the fraction of hits in the last $500$? And if after each kick the configuration is clipped to the rectangle $a\in[-1,2]$, $b\in[-1,1]$? Explain the difference using Exercise~\ref{ex:22:3}.
\end{exercise}

\begin{exercise}[\lvlB]
\label{ex:22:6}
\emph{Designing the input code.} The ball is hit if the height error is below $h=0.1$; the network reads the input over $T=0.25$~s, about $\sqrt{rT}$ levels up to rate $r$ are distinguishable, and stimulation is at most $40$ pulses per second. (a)~Is a place code on eight electrodes enough? (b)~Can the height be sent as a rate on a single electrode?
\end{exercise}

\begin{exercise}[\lvlC]
\label{ex:22:7}
\emph{Research: prediction or scrambling?} Generalize the model to $d$ tunable numbers ($p=\sum_{i<d}a_iy^i$). How does the number of rallies needed to reach a hit fraction of $0.5$ grow with $d$ under scrambling and under the signed-error rule~\eqref{eq:perturbation}? What bench experiment would separate the two hypotheses?
\end{exercise}

\appendix
\renewcommand{\chaptername}{\appendixname}
\chapter{Substances That Do Not Define a Neuron Type}
\label{sec:othermediators}

In chapter~\ref{sec:neurontypes} we sorted neurons by their primary transmitter and obtained ten types (table~\ref{tab:types}). Now it is time to add complications. A neuron releases more than its primary transmitter, and besides the primary transmitters other substances also work in the brain. They change how the network transmits and stores signals.

Besides the address and broadcast buses (section~\ref{sec:buses}) there is a third one: a \emph{back channel}. Several substances are released not by the sender but by the \emph{receiver}; they travel back across the cleft to the sender's output and change how much transmitter the sender will release next time. In communication networks this resembles an acknowledgment of receipt that the sender uses to tune its transmission. This is precisely the channel neurons use when they change the weights of their connections: through it the sender's output learns about the receiver's activity, that is, about the second factor of the learning rules of chapter~\ref{sec:dofa}.

The full list of known transmitters is long: more than a hundred substances, most of them short protein chains, the neuropeptides. But all of them fall into a few classes, and within each class the substances behave similarly. The substances that are never a primary transmitter are collected in table~\ref{tab:othermediators}, with the same three questions an engineer would ask: which bus carries the signal, does it act quickly, and what is its usual sign. They do not define the neuron type: they are released together with the primary transmitter, or released by cells other than neurons, or they travel in the reverse direction.

\begin{table}[t]
\centering
\footnotesize
\setlength{\tabcolsep}{4pt}
\renewcommand{\arraystretch}{1.12}
\begin{tabular}{>{\raggedright}p{2.25cm}>{\raggedright}p{2.8cm}>{\raggedright}p{1.8cm}>{\raggedright}p{1.5cm}>{\centering}p{0.8cm}>{\raggedright\arraybackslash}p{4.3cm}}
\toprule
Class & Substance & Bus & Speed & Sign & Role in computation \\
\midrule
Amino acids & D-serine & local & slow & AND & second key for the glutamate receptor: an ``AND'' condition \\
\midrule
Monoamines & trace amines & broadcast & slow & $\pm$ & fine-tuning of the monoamine channels \\
\midrule
\multirow{2}{2.25cm}{Purines}
 & ATP & address & fast and slow & $+$ & co-transmitter; signal between neurons and glia \\
 & adenosine & volume & slow & $-$ & accumulates with activity: a load counter~\cite{Dunwiddie2001} \\
\midrule
\multirow{8}{2.25cm}{Neuropeptides (more than a hundred)}
 & enkephalins, endorphins, dynorphin & volume & slow & $-$ & suppression of transmission \\
 & substance P & volume & slow & $+$ & slow amplification \\
 & neuropeptide Y & volume & slow & $-$ & slow inhibition \\
 & somatostatin & volume & slow & $-$ & slow inhibition, co-released with GABA \\
 & vasoactive intestinal peptide (VIP) & volume & slow & $+$ & disinhibition: inhibition of inhibitory neurons \\
 & cholecystokinin & volume & slow & $\pm$ & co-released with GABA in some inhibitory neurons \\
 & orexin & broadcast & slow & $+$ & stabilizer of the waking state \\
 & oxytocin, vasopressin & broadcast & slow & $\pm$ & global behavioral settings \\
\midrule
\multirow{2}{2.25cm}{Gases}
 & nitric oxide NO & back, volume & slow & $\pm$ & passes through membranes; back signal during weight changes~\cite{Garthwaite2008} \\
 & CO, H$_2$S & volume & slow & $\pm$ & role little studied \\
\midrule
Lipids & endocannabinoids (anandamide, 2-AG) & back & slow & $-$ & the receiver reduces the sender's release: weight feedback~\cite{WilsonNicoll2001} \\
\bottomrule
\end{tabular}
\caption{Substances that do not define a neuron type. \emph{Bus}, \emph{speed}, and \emph{sign} are as in table~\ref{tab:types}; in addition, a volume bus means the transmitter spreads through the extracellular space around the release site, a back bus means it travels from the receiver back to the sender, and a local bus means it acts near the release site. ``AND'' is an enabling signal: by itself it produces no current, but without it another input does not fire, like the second input of an AND gate. Neuropeptides are almost always released together with the primary transmitter, and mostly at high spike rates~\cite{vandenPol2012}.}
\label{tab:othermediators}
\end{table}

\chapter{Notation}
\label{sec:notation}

Neuron types are denoted by the first four letters of their primary transmitter: \nt{GLUT} is glutamate, \nt{GABA} is gamma-aminobutyric acid, \nt{GLYC} is glycine, \nt{ACET} is acetylcholine, \nt{DOPA} is dopamine, \nt{SERO} is serotonin, \nt{HIST} is histamine, \nt{NORA} is noradrenaline, \nt{ADRE} is adrenaline, and \nt{ASPA} is aspartate (Table~\ref{tab:types}).

There are fewer letters than quantities, and in different chapters the same letter sometimes means different things. In the table below each meaning has its own row; the right-hand column gives the formula where it is introduced.

\begin{center}
\footnotesize
\setlength{\tabcolsep}{4pt}
\renewcommand{\arraystretch}{1.12}
\begin{longtable}{>{\raggedright}p{2.0cm}>{\raggedright}p{9.6cm}>{\raggedright\arraybackslash}p{2.4cm}}
\toprule
Symbol & Meaning & Introduced in \\
\midrule
\endfirsthead
\toprule
Symbol & Meaning & Introduced in \\
\midrule
\endhead
\bottomrule
\endfoot
$a$ & slope of a linear transfer characteristic, $f(u)=au$ & \eqref{eq:feedback} \\
$a$ & \nt{ACET} level: $0$ for reading, $1$ for writing & \eqref{eq:acetnet} \\
$a_i$, $a$ & fatigue of a group: in the rhythm generator; in the slow-wave-sleep seesaw & \eqref{eq:halfcentre}, \eqref{eq:updown} \\
$a_S$, $a_P$ & sensitivity of the heart rate to the gas pedal and to the brake & \eqref{eq:gasbrake} \\
$A(r)$, $A_0$ & synaptic response to a spike at rate $r$; of a rested synapse & \eqref{eq:depression} \\
$A_1$ & recipient's response to a single quantum of transmitter & \eqref{eq:quantal} \\
$A$ & membrane area of a neuron including its dendrites & \eqref{eq:dendriteload} \\
$A_f$, $A_s$ & fraction of the correction retained until the next movement, for the fast and slow processes & \eqref{eq:tworate} \\
$b_i$ & intrinsic drive of a pacemaker & \eqref{eq:seroneuron} \\
$b$ & how fast work fatigues a group in the rhythm generator or in the sleep seesaw & \eqref{eq:halfcentre}, \eqref{eq:updown} \\
$b$ & number of bits per digitized sample & \eqref{eq:probedata} \\
$B_f$, $B_s$ & strength with which the correction learns from the error, for the fast and slow processes & \eqref{eq:tworate} \\
$c$ & transmitter concentration in the volume & \eqref{eq:seroneuron} \\
$c$ & correlation coefficient of the noise of neighboring neurons & \eqref{eq:correlated} \\
$c$ & coefficient with which the count difference moves the paddle & \eqref{eq:dishreadout} \\
$c_f$ & response of a $C$ neuron: the largest $s_{f,p}$ over all positions & \eqref{eq:scstage} \\
$c_m$ & specific membrane capacitance & \eqref{eq:dendriteload} \\
$C$ & cost of effort: an action taking time $\tau_a$ costs $C/\tau_a$ & \eqref{eq:vigor} \\
$d'$ & discriminability of two inputs: the difference divided by the noise & \eqref{eq:gainsnr} \\
$d$, $d_j$ & delay along the signal path & \eqref{eq:glutneuron}, \eqref{eq:vstar} \\
$d$ & delay of the feedback loop through a muscle & \eqref{eq:delaystable} \\
$D$ & diffusion coefficient & \eqref{eq:diffusionlength} \\
$D$ & denominator in the rates of the \nt{GLUT}--\nt{GABA} network & \eqref{eq:isn} \\
$D$ & waiting time for a delayed reward & \eqref{eq:patience} \\
$e$, $e_{ij}$ & eligibility trace in a synapse & \eqref{eq:threefactor} \\
$e$, $e_i$ & output or movement error arriving on the error wire & \eqref{eq:lms}, \eqref{eq:adaptivefilter}, \eqref{eq:tworate} \\
$E_{\text{spk}}$ & energy of one spike including the work of the synapses & \eqref{eq:energy} \\
$E$ & rate of the \nt{GLUT} group & \eqref{eq:ei} \\
$E_E$ & reversal potential of an excitatory synapse & \eqref{eq:shunt} \\
$f$, $F$ & transfer characteristic: rate as a function of input & \eqref{eq:ratenet}, \eqref{eq:gain} \\
$f$, $f_0$ & heart rate; intrinsic rhythm of the pacemaker & \eqref{eq:gasbrake} \\
$f_s$ & sampling rate of one electrode & \eqref{eq:probedata} \\
$F(t)$, $F_1$ & muscle force; force of a single twitch & \eqref{eq:forcesum} \\
$F_1$, $F_2$ & forces of two muscles pulling a joint in opposite directions & \eqref{eq:antagonists} \\
$g$ & receptor sensitivity to concentration & \eqref{eq:seroneuron} \\
$g_L$, $g_E$, $g_I$ & leak, excitatory, and inhibitory conductances & \eqref{eq:shunt} \\
$g$ & gain set by \nt{NORA} & \eqref{eq:gain}, \eqref{eq:machine} \\
$g$ & overall inhibition from \nt{GABA} in the memory network & \eqref{eq:acetnet} \\
$g$ & gain of the pain gate, set from above & \eqref{eq:gate} \\
$g$ & gain of one stage of the hormonal cascade & \eqref{eq:cascade} \\
$g_j$ & filter $j$ with its own time constant at the input of the adaptive filter & \eqref{eq:adaptivefilter} \\
$G$ & feedback loop gain & \eqref{eq:feedback} \\
$G$ & correction rate in a loop with delay & \eqref{eq:delaystable} \\
$h$, $h_I$ & external input to the \nt{GLUT} group and to the \nt{GABA} group & \eqref{eq:ei}, \eqref{eq:isn} \\
$h(t)$ & single muscle twitch & \eqref{eq:twitch} \\
$H(t)$ & upper threshold of sleep pressure: on reaching it, the machine falls asleep & \eqref{eq:sleeppressure} \\
$I$, $I_i$ & external drive to a neuron or group & \eqref{eq:ratenet}, \eqref{eq:wta}, \eqref{eq:updown} \\
$I$ & rate of the \nt{GABA} group & \eqref{eq:ei} \\
$I$ & brightness or intensity of a stimulus & \eqref{eq:photonnoise}, \eqref{eq:nakarushton} \\
$I$ & input current charging the membrane & \eqref{eq:dendriteload} \\
$k$ & number of spikes in a burst & \eqref{eq:burst} \\
$k$ & how many times the increment must exceed the noise & \eqref{eq:photonnoise} \\
$k$ & feedback coefficient of the hormonal cascade & \eqref{eq:cascade} \\
$k_s$ & strength with which pain drives sensitization & \eqref{eq:sensitization} \\
$K$ & number of outcome features & \eqref{eq:vectortd} \\
$K$ & joint stiffness & \eqref{eq:antagonists} \\
$K$ & loop gain of the hormonal cascade, $K=g^3k$ & \eqref{eq:cascade}, \eqref{eq:cascadestable} \\
$K_\varphi$ & strength with which the circadian oscillator locks to an external signal & \eqref{eq:pll} \\
$L$ & length of the delay line & \eqref{eq:itd} \\
$L$ & length of the wire from a pain sensor & \eqref{eq:paintime} \\
$L(t)$ & lower threshold of sleep pressure: on reaching it, the machine wakes up & \eqref{eq:sleeppressure} \\
$M$ & measure of the causal link between a predictor and the reward & \eqref{eq:retro} \\
$M_k$ & expectation of feature $k$ & \eqref{eq:vectortd} \\
$M_{ik}$ & strength of inhibition from \nt{GABA} neuron $k$ & \eqref{eq:machine} \\
$M$ & joint torque & \eqref{eq:antagonists} \\
$M$ & number of input lines to the mixers & \eqref{eq:cover} \\
$n$, $\bar n$ & number of spikes in a window; its mean & \eqref{eq:rateerror} \\
$n$ & number of inputs of a threshold element & \eqref{eq:cover} \\
$n_\uparrow$, $n_\downarrow$ & spike counts of the ``up'' and ``down'' regions per window & \eqref{eq:dishreadout} \\
$N$ & number of neurons & \eqref{eq:correlated}, \eqref{eq:energy} \\
$N$ & number of probe electrodes & \eqref{eq:probedata} \\
$N_s$ & number of release sites between two neurons & \eqref{eq:quantal} \\
$p$ & probability of transmitter release per spike & \eqref{eq:burst} \\
$p$ & fraction of active neurons in the memory network & \eqref{eq:acetnet} \\
$p$ & perturbation that is learned to be compensated & \eqref{eq:tworate} \\
$P$ & power spent on signaling & \eqref{eq:energy}, \eqref{eq:dishpower} \\
$P$ & uncertainty of the stored estimate & \eqref{eq:kalman} \\
$P$ & activity of the ``brake'': the \nt{ACET} wires to the heart & \eqref{eq:gasbrake} \\
$P_{\max}$ & how many random patterns a threshold element can split into two classes & \eqref{eq:cover} \\
$P_{\text{miss}}$ & probability of a miss & \eqref{eq:dishdensity} \\
$q$ & rate at which the world changes & \eqref{eq:kalman} \\
$q_k$ & rate of \nt{GABA} neuron $k$ & \eqref{eq:machine} \\
$q_0$ & background activity of the inhibitory interneuron & \eqref{eq:gate} \\
$q$ & rate of the inhibitory interneuron in the pain gate & \eqref{eq:gate} \\
$q$ & how many times stronger each motor unit is than the previous one & \eqref{eq:sizeprinciple} \\
$r$, $r_i$ & spike rate of a neuron & \eqref{eq:ratenet} \\
$r$, $r_t$ & reward (reward stream) & \eqref{eq:rpe}, \eqref{eq:ctd} \\
$\mathbf r(t)$ & response vector of the reservoir & \eqref{eq:reservoir} \\
$R$, $\bar R$ & reward received; its expectation or mean per unit time & \eqref{eq:perturbation}, \eqref{eq:vigor} \\
$R$ & noise variance at the filter input & \eqref{eq:kalman} \\
$R$, $r_0$ & delayed and immediate rewards & \eqref{eq:patience} \\
$R_{\max}$ & maximum transmission rate, bit/s & \eqref{eq:spikeentropy} \\
$r_{\max}$ & maximum spike rate & \eqref{eq:gain}, \eqref{eq:nakarushton} \\
$R_{\text{raw}}$, $R_{\text{spk}}$ & probe data rate: raw and spikes only, bit/s & \eqref{eq:probedata} \\
$s_j$ & sign of the outputs of neuron $j$ & \eqref{eq:dalesign} \\
$s_i$ & sign of the recipient's receptor & \eqref{eq:seroneuron} \\
$s$, $\hat s$ & state of the world; its estimate & \eqref{eq:rpe}, \eqref{eq:kalman} \\
$s_{f,p}$ & response of an $S$ neuron to feature $f$ at position $p$ & \eqref{eq:scstage} \\
$S$ & activity of the ``gas pedal'': the \nt{NORA} wires to the heart & \eqref{eq:gasbrake} \\
$S$ & sleep pressure & \eqref{eq:sleeppressure} \\
$t_1$ & time of the first spike & \eqref{eq:latency} \\
$t_k$ & time of the $k$-th spike & \eqref{eq:forcesum} \\
$t_{f,i}$ & template of feature $f$ & \eqref{eq:scstage} \\
$T$ & spike-counting window; photon integration time & \eqref{eq:rateerror}, \eqref{eq:photonnoise} \\
$T_c$ & rise time of a muscle twitch & \eqref{eq:twitch} \\
$T_{\text{burst}}$ & interval between network bursts of a culture & \eqref{eq:burstinterval} \\
$u$ & total input to a neuron & \eqref{eq:gain} \\
$u$, $u(t)$ & command from the brain: at the input of the adaptive filter; to the hormonal cascade & \eqref{eq:adaptivefilter}, \eqref{eq:cascade} \\
$u(t)$ & reservoir input & \eqref{eq:reservoir} \\
$U$ & level of a constant input & \eqref{eq:latency} \\
$U$ & fraction of the transmitter store released per spike & \eqref{eq:depression} \\
$v$ & signal speed in a wire; speed of a moving spot & \eqref{eq:itd}, \eqref{eq:vstar} \\
$V$ & membrane voltage & \eqref{eq:glutneuron}, \eqref{eq:shunt} \\
$V$ & expected future reward & \eqref{eq:rpe}, \eqref{eq:ctd} \\
$w$, $w_{ij}$, $W_{ij}$ & connection weight & \eqref{eq:glutneuron}, \eqref{eq:ratenet}, \eqref{eq:acetnet}, \eqref{eq:lms}, \eqref{eq:adaptivefilter} \\
$\mathbf w_k$ & input weights of $C$ neuron $k$ & \eqref{eq:traceinvariance} \\
$w_{q\mu}$, $w_{q\nu}$, $w_{z\mu}$ & connection weights in the pain gate & \eqref{eq:gate} \\
$W_{\text{out}}$ & weights of the reservoir's linear readout & \eqref{eq:reservoir} \\
$x$, $y$ & logic inputs and output & \eqref{eq:dualrail}, \eqref{eq:veto} \\
$x$, $y$ & sender and recipient activity & \eqref{eq:hebb} \\
$x$ & position of a detector on the delay line & \eqref{eq:itd} \\
$x_i$ & input from the sense organs & \eqref{eq:acetnet} \\
$x_p$ & input at position $p$ & \eqref{eq:scstage} \\
$\mathbf x$ & outputs of the $S$ neurons, input of a $C$ neuron & \eqref{eq:traceinvariance} \\
$x_j$ & input $j$ of the adaptive filter: the command after filter $g_j$ & \eqref{eq:lms}, \eqref{eq:adaptivefilter} \\
$x_f$, $x_s$ & corrections of the fast and slow processes & \eqref{eq:tworate} \\
$x_1$, $x_2$, $x_3$ & levels at the three stages of the hormonal cascade; $x_3$ is the final hormone & \eqref{eq:cascade} \\
$x^*$ & fraction of the recovered transmitter store at which a new avalanche starts & \eqref{eq:burstinterval} \\
$y$ & noisy filter input & \eqref{eq:kalman} \\
$y$, $\hat y$ & sensor signal; its prediction by the adaptive filter & \eqref{eq:adaptivefilter} \\
$y_k$, $\bar y_k$ & activity of $C$ neuron $k$; its trace & \eqref{eq:traceinvariance} \\
$y(t)$ & output of the reservoir readout & \eqref{eq:reservoir} \\
$z$ & output rate of the pain gate & \eqref{eq:gate} \\
\midrule
$\alpha_i^\pm$ & weight of positive and negative errors & \eqref{eq:asymmetric} \\
$\beta$ & strength of mutual inhibition & \eqref{eq:wta} \\
$\beta$ & strength of inhibition of the pain gate output & \eqref{eq:gate} \\
$\gamma$ & discount factor & \eqref{eq:rpe} \\
$\delta(\cdot)$ & delta function: an instantaneous jump & \eqref{eq:glutneuron} \\
$\delta$, $\delta_t$ & reward-prediction error & \eqref{eq:rpe}, \eqref{eq:ctd} \\
$\delta t$ & difference in the arrival times of sound at the two ears & \eqref{eq:itd} \\
$\Delta$, $\Delta_{\max}$ & interval between spikes; coincidence window & \eqref{eq:window} \\
$\Delta t$ & precision with which the recipient resolves time & \eqref{eq:spikeentropy} \\
$\Delta F$ & force increment from the next motor unit & \eqref{eq:sizeprinciple} \\
$\Delta I$ & discriminable brightness increment & \eqref{eq:photonnoise} \\
$\Delta p$ & paddle displacement per window & \eqref{eq:dishreadout} \\
$\Delta u$ & difference between the inputs of two groups & \eqref{eq:gainsnr} \\
$\Delta V$ & how much the membrane voltage must be raised & \eqref{eq:dendriteload} \\
$\Delta x$ & spacing between neighboring sensors & \eqref{eq:vstar} \\
$\varepsilon$ & relative difference between inputs & \eqref{eq:wtagain} \\
$\eta$ & learning rate & \eqref{eq:plasticity}, \eqref{eq:lms}, \eqref{eq:traceinvariance} \\
$\eta$ & noise in the network after the amplifier & \eqref{eq:softmax} \\
$\mu$ & touch input to the pain gate & \eqref{eq:gate} \\
$\nu$ & pain input & \eqref{eq:gate} \\
$\theta$ & firing threshold & \eqref{eq:window}, \eqref{eq:scstage} \\
$\kappa$ & amount of transmitter per spike & \eqref{eq:seroneuron} \\
$\kappa_i$ & ratio of the weights of positive and negative errors & \eqref{eq:expectile} \\
$\kappa_m$ & link between a muscle's force and its stiffness & \eqref{eq:antagonists} \\
$\ell$ & distance over which the transmitter spreads & \eqref{eq:diffusionlength} \\
$\ell$ & lever arm with which a muscle pulls on a joint & \eqref{eq:antagonists} \\
$\xi$ & random jitter of a neuron's output & \eqref{eq:perturbation} \\
$\rho(\boldsymbol\psi)$ & fraction of time spent in configuration $\boldsymbol\psi$ & \eqref{eq:dishdensity} \\
$\sigma$ & spread, noise scale & \eqref{eq:rateerror}, \eqref{eq:softmax} \\
$\sigma$ & brightness at which the response is half the maximum & \eqref{eq:nakarushton} \\
$\sigma_{\text{pre}}$, $\sigma_{\text{post}}$ & noise before and after the amplifier & \eqref{eq:gainsnr} \\
$\tau$ & membrane time constant & \eqref{eq:glutneuron} \\
$\tau$ & time constant of a stage of the hormonal cascade & \eqref{eq:cascade} \\
$\tau_{\text{eff}}$ & time constant of a self-exciting group & \eqref{eq:perfectint} \\
$\tau_c$ & lifetime of the transmitter in the volume & \eqref{eq:seroneuron} \\
$\tau_E$, $\tau_I$ & time constants of the \nt{GLUT} and \nt{GABA} groups & \eqref{eq:ei} \\
$\tau_e$ & lifetime of the eligibility trace & \eqref{eq:threefactor} \\
$\tau_{\text{tr}}$ & lifetime of the activity trace of a $C$ neuron & \eqref{eq:traceinvariance} \\
$\tau_s$ & time for sensitivity to return after injury & \eqref{eq:sensitization} \\
$\tau_r$, $\tau_d$ & time for sleep pressure to build during waking and to discharge during sleep & \eqref{eq:sleeppressure} \\
$\tau_{\text{rec}}$ & recovery time of the transmitter store & \eqref{eq:depression} \\
$\tau_\gamma$ & planning horizon & \eqref{eq:ctd} \\
$\tau_v$ & rate at which the expectation is refined & \eqref{eq:ramp} \\
$\tau_a$ & fatigue time in the rhythm generator or in the sleep seesaw & \eqref{eq:halfcentre}, \eqref{eq:updown} \\
$\tau_a^*$ & optimal time to perform an action & \eqref{eq:vigor} \\
$\Phi$ & right-hand side of the learning rule & \eqref{eq:plasticity} \\
$\Phi$ & photon flux & \eqref{eq:photonnoise} \\
$\varphi_k$ & feature $k$ of the obtained outcome & \eqref{eq:vectortd} \\
$\varphi$ & phase difference between the circadian oscillator and the external signal & \eqref{eq:pll} \\
$\boldsymbol\psi$ & network configuration in the DishBrain model & \eqref{eq:dishdensity} \\
$\omega$, $\omega_0$ & frequency of the external signal (light); intrinsic frequency of the circadian oscillator & \eqref{eq:pll} \\
\end{longtable}
\end{center}

\chapter{Answers}
\label{sec:answers}

Short answers to the end-of-chapter exercises. Full solutions are in a separate instructor's manual.

\answer{ex:01:1}{(a)~$8\cdot10^{10}$ \nt{GLUT} neurons are $10$ Earth populations; the broadcast neurons, $\approx1.9\cdot10^6$, make a city the size of Vienna. (b)~$\approx8\cdot10^{14}$ terminals, of which $\approx3\cdot10^{11}$ are broadcast, a share of $\approx4\cdot10^{-4}$: $\approx16$ times the share of these neurons ($2.2\cdot10^{-5}$).}
\answer{ex:01:2}{(a)~$\tau_c\ln(c_0/c_*)\approx4.6\ms$. (b)~The line clears when $T\ge\tau_c\ln101$: up to $\approx22$ releases per second instead of $\approx220$.}
\answer{ex:01:3}{$V(s_k)=\gamma^{K-1-k}r$; $\delta_{\text{lamp}}=\gamma^Kr=re^{-T/\tau_\gamma}\approx0.60\,r$ for any $\Delta$, $\tau_\gamma\approx1.95\,\text{s}$.}
\answer{ex:01:4}{$n^*=93$, $47$, $25$, $12$: $n^*\approx5/\alpha$ for small $\alpha$.}
\answer{ex:01:5}{\nt{GLUT}: $1\to10\to10^4\to10^7$, $\approx10^7$ neurons, $4$ links, $4\ms$. \nt{DOPA}: $1\to10\to3.2\cdot10^5$, $\approx3\cdot10^5$ neurons, $3$ links, $30$ times fewer; the same order as the number of \nt{DOPA} neurons.}
\answer{ex:01:6}{$\lambda=f_EJ_E-f_IJ_I$, whence $J_I=37.5$; the current ratio is $f_IJ_I/f_EJ_E=0.94$; an avalanche ($\lambda\ge1$) sets in when inhibition is weakened by only $6.7\,\%$.}
\answer{ex:01:7}{$n\ge57$ presentations per delay in each condition. The same decline would come, for example, from an imprecision of the internal time estimate that grows with $T$.}

\answer{ex:02:1}{A row $3.2$~mm long with a spacing of $25\um$: $\approx129$ detectors. A band $v\,\tau\ln2=1.7$~mm wide fires, $\approx69$ detectors, more than half the row; the direction is the center of the band.}
\answer{ex:02:2}{The window is $-\tau\ln\frac{w_2}{\theta-w_1}\le\delta\le\tau\ln\frac{w_1}{\theta-w_2}$, i.e., $[-0.235,\,0.178]\ms$; the center is $c=\frac\tau2\ln\frac{w_1(\theta-w_1)}{w_2(\theta-w_2)}=-28$~\textmu s. The direction shifts by $28$~\textmu s toward the stronger ear, almost three resolution steps.}
\answer{ex:02:3}{We need $s_i$ with $s_iw_{ij}s_j\ge0$; they exist if and only if the cycles are even. Mutual inhibition reduces (no stable oscillations); a \nt{GLUT}--\nt{GABA} loop does not (it can oscillate).}
\answer{ex:02:4}{Six neurons $g_k=[x+y+z\ge k]$ and $\bar g_k=[\bar x+\bar y+\bar z\ge4-k]$, $k=1,2,3$; $C=g_2$, $\bar C=\bar g_2$, $S=[g_1+\bar g_2+g_3\ge2]$, $\bar S=[\bar g_1+g_2+\bar g_3\ge2]$: eight neurons. From AND and OR: $12$.}
\answer{ex:02:5}{$T_{\min}(0.6)=0.718\,\tau$; $I_c=0.225$, and near the threshold $T_{\min}\propto(I-I_c)^{-1/2}$.}
\answer{ex:02:6}{(a)~$500$ links, $5\cdot10^4$ neurons; scatter $4.5\ms$ ($0.45\,\%$), relative error $\propto T^{-1/2}$. (b)~A random delay factor common to all links, with a $5\,\%$ scatter.}
\answer{ex:02:7}{Refresh once every $\tau_F\ln2=0.69\,\text{s}$: $4.3$ spikes per second per neuron versus $20$, $4.6$ times more economical ($4\cdot10^{-7}$ versus $2\cdot10^{-6}$~J at $10^{-10}$~J per spike). The flip-flop loses the bit; the capacitor keeps it if the silencing began no later than $0.49\,\text{s}$ after a refresh (probability $0.71$).}

\answer{ex:03:1}{$35\to17.5\mV$ at $g_E=g_L$; $56\to40\mV$ at $g_E=4g_L$, whereas subtraction would give $38.5\mV$: the shunt divides $g_E$ by $3$. The window narrows by half ($\tau_{\text{eff}}$ falls by a factor of $2$).}
\answer{ex:03:2}{$\operatorname{tr}=\frac{w_{EE}-1}{\tau_E}-\frac1{\tau_I}$, $\det=\frac{1-w_{EE}+w_{EI}w_{IE}}{\tau_E\tau_I}$; at the boundary $\omega=\sqrt{\det}$, $\tau_I=20\ms$, period $73\ms$.}
\answer{ex:03:3}{$d_c=\tau_E\arccos(-1/G)/\sqrt{G^2-1}$, period $2\pi\tau_E/\sqrt{G^2-1}\in(2d_c,4d_c)$. $G=2$: $d_c=12.1\ms$, period $36\ms$; $G=5$: $d_c=3.6\ms$, period $12.8\ms$, right in the band of fast rhythms.}
\answer{ex:03:4}{(a)~At $I_2=\beta I_1=2$ going up and $I_2=I_1/\beta=0.5$ going down, a hysteresis loop of width $1.5$. (b)~By $\tau\ln10/(\beta-1)=2.3\tau$ for each decade of $\varepsilon$.}
\answer{ex:03:5}{Three intermediaries, $d_k=5$, $10$, $15\ms$: the vetoes must cover $15\ms$, $5\ms$ each.}
\answer{ex:03:6}{$n+M+1=3+4+1=8$ neurons. Two: \nt{GABA} $=[x+y+z\ge2]$, output $=[x+y+z-2\,\nt{GABA}\ge1]$.}
\answer{ex:03:7}{$\binom84=70<100\le\binom94=126$: $9$ intermediaries (Sperner's theorem). Without direct connections, $100$: the rank of $\beta(J-I)$ is $n$.}
\answer{ex:03:8}{$h_c=\frac1{2w_{EE}}+\frac{w_{EI}w_{IE}-w_{EE}}{4w_{EE}^2}=\frac7{18}\approx0.39$; simulation gives a sign change between $h=0.38$ and $0.40$.}

\answer{ex:04:1}{(a)~$\approx1.1\cdot10^{11}$ bits/J. (b)~$\log_2\sqrt{10}\approx1.7$ bits versus $81$: $\approx50$ times less.}
\answer{ex:04:2}{$r^*=(P_0/E_{\text{spk}})\ln\bigl(1/(r^*\Delta t)\bigr)$, $P_0/E_{\text{spk}}=1.16\,\text{s}^{-1}$: $r^*\approx6$ spikes per second, $7.4\cdot10^{10}$ bits/J.}
\answer{ex:04:3}{$\mathcal I=\frac{N}{1+(N-1)c}=9.9$ (limit $10$) versus $\mathcal I=\frac{N}{1-c}=1111$: a factor of $\approx110$; correlation hurts only when it looks like the signal.}
\answer{ex:04:4}{$\theta\approx68.7$ and $19.9$ (in units of $w$), $m=19$ and $m=10$: the fast recipient needs almost half as many coincident inputs, although its mean input is five times smaller.}
\answer{ex:04:5}{$U\tau_{\text{rec}}\ge0.725\,\text{s}$ and $\tau_{\text{rec}}(1-2U)\le0.05\,\text{s}$, which requires $U\ge0.483$; the smallest $\tau_{\text{rec}}=0.725\,\text{s}$ at $U=1$ ($1.45\,\text{s}$ at $U=0.5$).}
\answer{ex:04:6}{(a)~$\theta\,e/(e-1)\approx1.58\,\theta$ for any $\tau$ and $\theta$. (b)~$14$ spikes.}
\answer{ex:04:7}{(a)~$1.62$ bits per word: $0.77$ in the spike count, $0.84$ in their arrangement. (b)~$1.13$ and $1.59$ bits: from $30$ words the estimate loses almost a third.}

\answer{ex:05:1}{$\ell\approx17\um$; $\approx100$~days: broadcasting across the brain is done by the branching wire, and diffusion only blurs each release site over $\sim20\um$.}
\answer{ex:05:2}{$r=ab/(1+G)$, $S=1/(1+G)$ exactly; $+1.7\,\%$ instead of $+20\,\%$.}
\answer{ex:05:3}{$T^*=1.21\,\text{s}$ for $G=2$ and $0.19\,\text{s}$ for $G=9$: realistically $G\sim2$--$4$, a suppression by a factor of $3$--$5$.}
\answer{ex:05:4}{$\mathrm{CV}=1/\sqrt{2\lambda\tau_c}\approx9\cdot10^{-4}$ for a total rate $\lambda$; a single neuron needs $\tau_c=12.5\,\text{s}$.}
\answer{ex:05:5}{$c\approx0.40$ in both cases; $\mathrm{CV}=0.050$ versus $0.158$, a factor of $\sqrt{10}$. The rates of individual neurons stay proportional to $a_i$: only the sum is regulated.}
\answer{ex:05:6}{$N_1=\overline{A\vee B}$, $N_2=\overline{A\vee N_1}$, $N_3=\overline{B\vee N_1}$, $N_4=\overline{N_2\vee N_3}$, XOR $=N_5=\overline{N_4}$; each switch takes $\tau_c\ln2$, and a path of $4$ gates takes $2.8\,\text{s}$ per XOR.}
\answer{ex:05:7}{(a)~$\tau_\gamma=6/\ln3\approx5.5\,\text{s}$. (b)~$\bar D<4\,\text{s}$ versus $D<2.2\,\text{s}$: an unpredictable delay makes the animal more patient.}
\answer{ex:05:8}{$k_2=1/\tau_d=10\,\text{s}^{-1}$, $k_1=1/\tau_d-1/\tau_\gamma=9.5\,\text{s}^{-1}$; $\tau_\gamma=1/(k_2-mk_1)$: doubling at $+2.6\,\%$ (and at $+5.3\,\%$ the horizon is infinite).}

\answer{ex:06:1}{$e=(1-e^{-T_a/\tau_e})e^{-d/\tau_e}$. (a)~$0.110$ versus $4.5\cdot10^{-5}$, a factor of $\approx2.5\cdot10^3$. (b)~$\tau_e^*=T_a/\ln(1+T_a/d)=0.59\,\text{s}$.}
\answer{ex:06:2}{The rate is $\nu_x\nu_y(A_+\tau_+-A_-\tau_-)=0.8A_+$ per second, $\approx2900A_+$ per hour; $\tau_-=\tau_+/0.6=33\ms$.}
\answer{ex:06:3}{$\mathbf e_1$ if $\mu^2+s_1^2>s_2^2$; here $1.01>0.25$, and the neuron learns $\mathbf e_1$ even though the spread along $\mathbf e_2$ is $25$ times larger: the rule finds the principal eigenvector of the correlation matrix, not of the covariance matrix.}
\answer{ex:06:4}{$V=Re^{-(t_c+L-t)/\tau_\gamma}$ between the lamp and the juice, so $\dot V=V/\tau_\gamma$; the burst has area $Re^{-L/\tau_\gamma}$: $0.61R$ and $0.78R$.}
\answer{ex:06:5}{Signal-to-noise ratio $1/(N+1)$, number of trials $\propto N+1$: $5\cdot10^5$ trials $\approx1$~month and $5\cdot10^{11}$ trials $\approx8\cdot10^4$~years.}
\answer{ex:06:6}{(a)~$k_2=1/\tau_d$, $k_1=1/\tau_d-1/\tau_\gamma$. (b)~The spurious error is $-(V/\tau_\gamma)\,\varepsilon/(1+\varepsilon)$, $\varepsilon=\tau_d/\tau_\gamma$, whence $\tau_d\le0.105\,\text{s}$.}
\answer{ex:06:7}{$0.46$, $0.90$, $0.99$, $0.95$, and $0.70$ for $d=3\,\text{s}$. The points fall on a single curve of the trace fraction $F=(1-e^{-T_{\text{cue}}/\tau_e})e^{-d/\tau_e}$: learning for $F\gtrsim0.1$, random choice for $F<10^{-2}$.}
\answer{ex:06:8}{$\eta^*=1/\bigl(2\sigma^2(N+2)\bigr)$, in $N+2$ trials; with $m$ of $N$ jittering, in $N(m+2)/m\ge N+2$: sparseness does not help.}

\answer{ex:07:1}{$V=\kappa p/\bigl(\kappa p+(1-\kappa)(1-p)\bigr)$: $0.059$, $0.20$, $0.50$; the median is $0$, while the point~\eqref{eq:expectile} changes smoothly with $\kappa$.}
\answer{ex:07:2}{$p=1/3$, $x=0.75$, $\sigma_r=0.35$, $V(0.2)=0.083$.}
\answer{ex:07:3}{$V=\sqrt\kappa/(\sqrt\kappa+\sqrt{1-\kappa})$, from $0.25$ to $0.75$; scatter $\propto\sqrt\eta$; for the drops $V=0.25+0.5\kappa$, the distributions differ by $0.05$ at the extreme neurons, so $\eta\lesssim0.01$ is needed.}
\answer{ex:07:4}{(a)~$J^*=2\sqrt{C\bar R}$. (b)~Overpayment $\bigl(k^{1/2}+k^{-1/2}\bigr)/2-1$: $6\,\%$ for $k=2$ and $25\,\%$ for $k=4$; accuracy ``within a factor of two'' is enough.}
\answer{ex:07:5}{A burst of area $R(e^{-1}-e^{-2})\approx0.23R$, equal to $2.3\,\text{s}$ of ramp; if dopamine reports $V$, there is no burst, and the level simply steps up by a factor of $e$.}
\answer{ex:07:6}{An event shifts the weight $\propto A\tau_f/(\tau_e+\tau_f)$, and the level gives an error $s\tau_f$: $9\,\text{s}\le\tau_f\le17\,\text{s}$.}
\answer{ex:07:7}{Initially $\Delta P=0.80$, $M=0.90$. (a)~$\Delta P=0.74$ ($-8\,\%$), $M=0.43$ ($-52\,\%$); (b)~$\Delta P=0.40$ ($-50\,\%$), $M=0.80$ ($-11\,\%$): a double dissociation.}
\answer{ex:07:8}{$N_{\text{eff}}\to2+\rho^2N\approx10^6$ trials: $10^4$ times fewer than with a single wire, but a $1\,\%$ leak still costs a million trials.}

\answer{ex:08:1}{(a)~$\approx1.7\cdot10^6$. (b)~$g=\frac{0.9}{\sqrt{0.19}}\,\sigma_{\text{post}}/\sigma_{\text{pre}}\approx16.5$: the answer depends only on the noise ratio.}
\answer{ex:08:2}{(a)~Eigenvalues $-(1\pm g\beta)/\tau$; the difference decays in $\tau/(1-g\beta)$: $40\ms$ at $g\beta=0.5$ (the input difference amplified threefold), $200\ms$ at $0.9$ ($19$-fold). (b)~$0.905$ and $1.105$; between the boundaries only the strong input wins.}
\answer{ex:08:3}{$P(k)=e^{gu_k/\sigma}/\sum_le^{gu_l/\sigma}$, that is,~\eqref{eq:softmax}; $g=10\ln9\approx22$.}
\answer{ex:08:4}{$\bar\Delta=1/3$, $g^*\approx3\ln(1/h)$: $21$, $14$, $9$ versus $16$--$23$, $11$ and $8$ in the model.}
\answer{ex:08:5}{$g^*$ from $16$--$23$ at $h=0.001$ to $6$--$8$ at $h=0.2$; the estimate $3\ln(1/h)$ holds within a factor of one and a half for $h\le0.05$; for $h\gtrsim\eta/10$ the rule $Q\leftarrow Q+\eta(R-Q)$ cannot absorb what exploration has found, and $g^*$ gets stuck near $8$.}
\answer{ex:08:6}{$T_p=\tau\ln9=44\ms$ for $g_{\text{lo}}=0$ (model: $44\ms$) and $70\ms$ for $g_{\text{lo}}=0.25$: a burst of $2$--$3\tau$ with gain near zero.}
\answer{ex:08:7}{(a)~Best constant: $0.925$ ($g=8$); oracle: $0.929$ ($16/8$); there is almost nothing to gain. (b)~For example, calm epochs with rare rewards ($p\in[0,0.3]$) and volatile ones with $h=0.1$: $0.850$ versus $0.872$; adjustment captures about a quarter at $\eta=0.2$ and almost all of it at $\eta=0.05$.}

\answer{ex:09:1}{(a)~$P_\infty=0.2$, memory $20\,\text{s}$. (b)~$P_\infty=1$, memory $1\,\text{s}$. (c)~Memory is proportional to the noise amplitude: sixfold.}
\answer{ex:09:2}{$P(t)=P_\infty\coth(t\sqrt{q/R})$. (a)~$\frac12\sqrt{R/q}=5\,\text{s}$, half the memory time. (b)~$t=\frac12\ln21\,\sqrt{R/q}\approx15\,\text{s}$: this is how long elevated \nt{ACET} must be sustained.}
\answer{ex:09:3}{Variance $qT/2+R/(2T)$; $T^*=\sqrt{R/q}$, variance $\sqrt{qR}=P_\infty$, the same as for the filter~\eqref{eq:kalman}; growth by a factor of $(\kappa+1/\kappa)/2$: $1.25$ and $5.05$.}
\answer{ex:09:4}{$a>a_w=1-\theta/(mw)$: $0.15$ for $w=0.041$, $0.22$ for $w=0.045$.}
\answer{ex:09:5}{Errors appear at $M\approx40$--$60$; at $M=80$, $1.9$ extra neurons; at $M=100$, a pattern fraction of $0.86$; the interference from other patterns has variance $\approx(M-1)p(1-p)/N$, so $M_{\max}\approx80$, $M_{\max}\propto N/p$.}
\answer{ex:09:6}{For example, $\eta(a)=\eta\min\bigl(1,\max(0,(a-0.3)/0.6)\bigr)$. With $\eta a$ all three foreign neurons entered the pattern; with $\eta(a)$, none, and writing at $a\ge0.9$ is unchanged (fraction $1.00$).}
\answer{ex:09:7}{(a)~$40$ replays; with $\eta(a)$ from Exercise~\ref{ex:09:6}, never. (b)~$200$ replays, $10$ nights.}

\answer{ex:10:1}{Bus $\approx10^{12}$ bit/s ($11.4$ bits per spike), control $\approx40$ bit/s; $2.5\cdot10^{10}\,\text{s}$, about eight hundred years.}
\answer{ex:10:2}{$a>1/3$ (at $a=0$, growing oscillations at $\approx67$~Hz). No: stability is set by the product $g(1-a)$.}
\answer{ex:10:3}{Mean $\sigma^2R'x_j$, variance $(N+1)\,x_j^2\sigma^4R'^2$; the ratio is $1/(N+1)$, so $n\approx z^2(N+1)\propto N$ trials are needed.}
\answer{ex:10:4}{$L_v\in[32.9,\,47.9]\ms$: delay $35.4\ms$, window $\pm7.5\ms$; false bindings $2\Delta_{\max}r_ar_v\approx0.38$ per second.}
\answer{ex:10:5}{$0.098$ versus $0.180$ at the best $\alpha=0.2$: $45\,\%$ smaller.}
\answer{ex:10:6}{$n=\ln\bigl(\ln1.5/(0.6g)\bigr)/\ln(1-\alpha)$: $24$ and $4$; at $g=2$, $11$ and $2$.}
\answer{ex:10:7}{$S'=S\bigl(1-4\eta\sigma^2+4\eta^2\sigma^4(G+2)\bigr)$, best $\eta\sigma^2=1/(2(G+2))$; $\approx10^6$ trials, about a month.}

\answer{ex:11:1}{(a)~$T=k^2/\bigl(\Phi(\Delta I/I)^2\bigr)=9\,\text{s}$. (b)~$45$ sensors, a spot of $\approx7\times7$, resolution worse by a factor of $\approx7$.}
\answer{ex:11:2}{(a)~$2$--$3.3$ bits per spike, $22$--$34\,\%$ of the limit ($9.1$--$9.8$ bits). (b)~$\log_2\binom{400}{20}\approx111$ bits; on average $1$ shared type.}
\answer{ex:11:3}{(a)~$\sigma/9\le I\le9\sigma$, $1.9$ orders of magnitude. (b)~$6$ and $7$.}
\answer{ex:11:4}{(a)~$f<343/0.44\approx780$~Hz. (b)~$625$ spikes, $125$ neurons.}
\answer{ex:11:5}{$21$ bits per event, $4.8\cdot10^5$ events/s, $7$ events per edge pixel: $136$ pixels/s. An ordinary camera: $1.25$ frames/s.}
\answer{ex:11:6}{Logarithmic: $\approx1.0\,\text{s}$ in both directions (theory $0.96\tau$); linear: $0.2\,\text{s}$ up and $3.0\,\text{s}$ down (theory $0.18\tau$ and $3.0\tau$).}
\answer{ex:11:7}{(a)~$\ln I$ is logistically distributed with median $\ln\sigma$ and scale $1$: a spread of $\pi/\sqrt3$ in natural logarithms, $0.79$ orders of magnitude. (b)~$r=r_{\max}\Phi_I(I)^2$.}

\answer{ex:12:1}{$0.60$ (without correlation it would be $0.18$); the limit is $\sqrt{c/(rT)}\approx0.58$: a hundred neurons distinguish only ``yes/no.''}
\answer{ex:12:2}{(a)~$10^8$ spikes, $\approx10\,\text{mJ}$: $0.3\,\%$ of the $3$~J of the whole brain. (b)~$3$~J, $\approx300$ times more.}
\answer{ex:12:3}{(a)~$\theta_A\in[2,3)$, $\theta_B\in[1,2)$; the other shape gives at most $2$ and $1$. (b)~$10\cdot96^2\approx9.2\cdot10^4$.}
\answer{ex:12:4}{(a)~$T_d=0.85\,\tau_a=0.34\,\text{s}$. (b)~$T_d\propto\tau_a$ and does not depend on $I$: $\tau_a\approx3.5\,\text{s}$ (simulation: $3.1\,\text{s}$).}
\answer{ex:12:5}{$p\approx0.17$. The fraction falls slowly: $0.90$, $0.88$, $0.86$, $0.84$, $0.82$, $0.79$ for $N=8$, $12$, $24$, $48$, $96$, $192$.}
\answer{ex:12:6}{All ten runs are error-free at $\tau_{\text{tr}}=20$, $50$, and $200\ms$ (at $30\ms$ one run in ten fails); at $5$--$10\ms$, about $0.5$; at $0.5$--$2\,\text{s}$ the quality drops to $0.86$. The upper edge: the trace must decay during the pause, $e^{-0.3\,\text{s}/\tau_{\text{tr}}}\ll1$.}
\answer{ex:12:7}{One delay is not enough: the $\pm5\ms$ window does not cover $\Delta x/v$ from $5$ to $20\ms$. Two inhibitory branches with $5<d_1\le10\ms$ and $15\le d_2\le d_1+10\ms$.}
\answer{ex:12:8}{\texttt{two\_stage}: $\Delta=2$, one flip gives a tie, two flips change the answer. The ones counter: one flip; $0.57$ at $p=0.1$.}

\answer{ex:13:1}{(a)~$q=100^{1/99}\approx1.048$, step $\approx4.8\%$; range $\approx2.2\cdot10^3$ ($67$~dB). (b)~A step of $22f_1$, i.e., $220\%$ at $10f_1$.}
\answer{ex:13:2}{$P_{\text{fail}}\approx1.8\cdot10^{-4}$ and $2.8\cdot10^{-16}$: about $300$ failures per day at $S=2$ and $5\cdot10^{-10}$ per day at $S=4$.}
\answer{ex:13:3}{(a)~$H(s)=eF_1T_c/(1+sT_c)^2$, a second-order filter with cutoff frequency $1/(2\pi T_c)\approx3.2$~Hz. (b)~$r\approx22$~spikes/s ($\approx20$ from the first harmonic).}
\answer{ex:13:4}{(a)~At $Gd=1/e$ the root $s=-1/d$ is double, and for no other $G$ does the rightmost root lie farther left. (b)~$d=30\ms$: $G\approx12$~s$^{-1}$, $\tau=30\ms$; $d=150\ms$: $G\approx2.5$~s$^{-1}$, $\tau=150\ms$. The time constant equals the delay.}
\answer{ex:13:5}{(a)~$E\approx0.427$, $T\approx1.70\,\tau_a=0.68$~s (the model with $\tau=20\ms$ gives $0.84$~s). (b)~The equation does not contain $I$; a rhythm exists for $b>\beta-1$.}
\answer{ex:13:6}{At $b=0.8$ there is no rhythm; at $b=1.05$, $1.2$, $1.5$, $2$, $4$, $T\approx2.73$, $1.74$, $1.19$, $0.84$, $0.45$~s; for any $I$, $T=0.840$~s: the input changes only the amplitude, not the tempo.}
\answer{ex:13:7}{The substitution $s=u-a$ gives the fastest decay rate $a+1/d$ at $G=e^{-1-ad}/d$; $a\ge33.3$~s$^{-1}$, $G=e^{-2}/d\approx4.5$~s$^{-1}$; the stronger the co-contraction (larger $a$), the smaller $G$ must be.}
\answer{ex:13:8}{An open problem. For example, a threshold in the fatigue drive, $b[r_i-r_0]_+$, gives $I_{\min}=\beta b r_0/(b-\beta+1)\approx0.53$ at $b=4$, $r_0=0.2$, and a period from $1.8$~s at $I=0.6$ to $0.52$~s at $I=3$.}

\answer{ex:14:1}{(a)~By $0.17$~s and by $1.5$~s. (b)~From distances shorter than $11$~cm.}
\answer{ex:14:2}{Touch without pain does not get through at any $\mu$ as long as $g\le\beta w_{q\mu}/w_{z\mu}=5$; touch $\mu=1$ gets through at $g>6$: strong sensitization turns an ordinary touch into pain.}
\answer{ex:14:3}{(a)~$g^*=(1-k_sC)/(1-k_sA)$, stable for $k_sA<1$. (b)~Without touch at $\nu\ge2$; with touch already at $\nu\ge1.7$.}
\answer{ex:14:4}{$V(t)=-Pe^{-(T-t)/\tau_\gamma}$. (a)~$-Pe^{-T/\tau_\gamma}\approx-0.37P$. (b)~$+Pe^{-(T-t_e)/\tau_\gamma}\approx+0.61P$; it grows with $t_e$ up to $+P$ and teaches the network to avoid pain.}
\answer{ex:14:5}{At $k_s=4$ there is no latching (a second stable state exists for $k_s\ge4.58$); $T_{\min}\approx4.35$, $1.40$, $0.65$, $0.32\,\tau_s$ for $k_s=4.6$, $5$, $6$, $8$.}
\answer{ex:14:6}{$w_{q\nu}=4/9\approx0.444$, $q_0=2/9\approx0.222$, $w_{q\mu}=49/45\approx1.089$; touch without pain is safe up to $g=49/9\approx5.4$.}
\answer{ex:14:7}{(a)~The threshold is $\nu_c(g)=\theta/g$ for $\nu_c\ge q_0/w_{q\nu}$, otherwise $(\theta+\beta q_0)/(g+\beta w_{q\nu})$; $g^*\approx2.45$, $1.0$, $0.25$. (b)~An open question.}

\answer{ex:15:1}{(a)~$2\cdot10^5$ per neuron, $6\cdot10^{12}$ for the circuit. (b)~$\approx8$~bit/s per wire, at least $2.5\cdot10^4$~s $\approx7$~h.}
\answer{ex:15:2}{Factor $2$: neighbor correlation $0.943$, $\lambda$ from $4.18$ to $7.3\cdot10^{-4}$, ratio $\approx5.7\cdot10^3$; factor $4$: correlation $0.8$, ratio $\approx94$.}
\answer{ex:15:3}{(a)~$\lambda=0.988$ and $0.808$: $82$ and $4.7$ movements. (b)~$x_s^*=0.690$, $x_f^*=0.172$, sum $0.862$.}
\answer{ex:15:4}{(a)~$116$ movements. (b)~$19$. A single-process model with a zero sum gives no savings.}
\answer{ex:15:5}{(a)~$G=50$~s$^{-1}$ versus the limit of $10.5$~s$^{-1}$ without a model. (b)~No: at $G=50$~s$^{-1}$ the critical delay is $d_c\approx103\ms$; at $d=150\ms$ one needs $\hat k>0.445k$ (for any $G$ and $d$, $\hat k>k/2$).}
\answer{ex:15:6}{The error falls below $0.5\%$ in $6$--$30$~s, with a floor of $\approx(6$--$9)\cdot10^{-4}$ against $7.5\cdot10^{-4}$ for the best weights; at $\eta=50$ the weights still differ from the best by up to $0.2$ along the ill-conditioned directions of $C$ (Exercise~\ref{ex:15:2}).}
\answer{ex:15:7}{$\mathbf B=A\mathbf K$; $q_f/R\approx0.035$, $q_s/R\approx8.6\cdot10^{-4}$; with $4q_s$, $B_f\approx0.088$, $B_s\approx0.047$: the slow process learns more than twice as fast, and the fast one more slowly.}

\answer{ex:16:1}{The blood is an RC circuit with $\tau=t_{1/2}/\ln2=14.4$~min; peak to trough $2^{T/t_{1/2}}=64$; the fundamental harmonic is attenuated to $0.55$; the ratio is $2$ for $t_{1/2}=T=60$~min.}
\answer{ex:16:2}{$\varphi^*=\arcsin\bigl((\omega-\omega_0)/K_\varphi\bigr)\approx41$~min ($\omega-\omega_0\approx0.047$~rad/day), relaxation time $1/(K_\varphi\cos\varphi^*)\approx3.9$~days; after a shift of $+6$ and $-6$~h, $7.7$ and $7.9$~days.}
\answer{ex:16:3}{$K_c(n)=\sec^n(\pi/n)$: $8$ and $2.37$, error $11\%$ and $30\%$; as $n\to\infty$, $K_c\to1$ and the error $\to50\%$.}
\answer{ex:16:4}{The extreme commands (gas $80$, brake $0$) give $f=120$ after $0.76$~s; ``gas only'' takes $2.33$~s, three times longer.}
\answer{ex:16:5}{$3\arctan(\omega\tau)+\omega d=\pi$, $K_c=(1+\omega^2\tau^2)^{3/2}$: $K_c=8$, $3.54$, $2.50$, $1.76$ and period $3.63$, $5.46$, $6.86$, $9.27\,\tau$; at $0.9K_c$ the oscillations decay, at $1.1K_c$ they persist.}
\answer{ex:16:6}{$0.60\bar c$, $0.87\bar c$, $0.90\bar c$, limit $0.91\bar c$; at a constant concentration the response is zero: such a recipient reads only pulses.}
\answer{ex:16:7}{Open problem. With feedback there is a rhythm only for $K>8$, and the period is proportional to the stage $\tau$ and nearly independent of $K$ ($3.64\,\tau$ at $K=8.5$, $4.23\,\tau$ at $K=40$): changing $K$ by a factor of one and a half shifts the period by $2$--$4\%$, below the measurement error. The decisive tests are opening the loop (a constant final hormone at the input of the first stage kills such a rhythm) or the response to a brief extra dose of hormone at different phases of the cycle.}

\answer{ex:17:1}{$0.2$~W versus $81$~\textmu W; the raw stream would need at most $50$~pJ/bit.}
\answer{ex:17:2}{At $c=0.1$: $5.6$, $8.7$, $9.2$~bit/s, limit $9.3$~bit/s; at $c=0$ and $N=1000$: $65$~bit/s.}
\answer{ex:17:3}{$B=\log_2N+P\log_2P+(1-P)\log_2\frac{1-P}{N-1}=4.87$, $4.37$, $3.29$~bits per character, i.e., $7.3$, $6.6$, $4.9$~bit/s; $10$~bit/s needs $2.3$~characters per second.}
\answer{ex:17:4}{Error variance $2b/(Nm^2T)+\sigma_\xi^2$ per component; the two contributions are equal at $N\approx90$; limit $\tfrac12\log_2(1+0.25/0.04)\approx1.4$~bits per component per window.}
\answer{ex:17:5}{The pair converges for $\eta_bd^2+\eta_db^2<2$, i.e., roughly for $\eta<1$: at $0.8$ it converges, at $1.1$ steady period-$2$ oscillations, at $1.5$ aperiodic ones, at $2$ runaway; the learners' rates must be separated.}
\answer{ex:17:6}{Threshold $\approx4.4\sigma$; $102$~kbit/s, $0.10$~mW versus $205$~mW for the raw stream, $2000$ times less; only $5.7\cdot10^{-5}$ of the bins saturate (more than $3$ spikes).}
\answer{ex:17:7}{Open problem. About $46$~bit/s at $\text{SNR}\approx0.9$ and about $330$~bit/s with independent noise. If signal-like noise is the obstacle, the extracted information saturates as the channel count grows; if ignorance of the code, it grows with a better decoder (for example, one that uses spike times).}

\answer{ex:18:1}{(a)~$0.9949\le w\le1.0049$, i.e., $|1-w|\lesssim0.005$. (b)~$K\ge400$ synapses.}
\answer{ex:18:2}{$(\omega_R\tau_w)^2=\sqrt{\alpha(\alpha+2+2\varrho)}/\varrho-1$, $\varrho=\tau_m/\tau_w$; $f_R\approx11.1$~Hz; $|Z(\omega_R)|/|Z(0)|=4.6$.}
\answer{ex:18:3}{(a)~$f=1/\bigl(\tau\ln\frac{\mu}{\mu-\theta}\bigr)$; $1$~Hz requires $\mu/\theta-1\approx3.7\cdot10^{-44}$. (b)~$T=\pi\tau/\sqrt\mu$, $f\approx3.2$~Hz: $f\propto\sqrt\mu$.}
\answer{ex:18:4}{The integrator fires at $f\lesssim17.7$~Hz (low-pass filter), the resonator only in the band $5.5$--$22$~Hz (band-pass filter).}
\answer{ex:18:5}{(a)~$T_{\min}=\frac{\tau}{w-1}\ln\frac{0.5}{0.3}\approx10.2\ms$. (b)~$B>w-1-0.2=0.3$.}
\answer{ex:18:6}{The tuning time constant is $\tau/(\eta\langle r^2\rangle)$, $\eta=\tau\ln10/60\,\text{s}\approx3.8\cdot10^{-3}$; with $I\neq0$ the weight drifts, so a copy of the command $I/\tau$ must be subtracted from $e$.}
\answer{ex:18:7}{Open problem. The resonator wins only by a factor of $1.35$ at $11$~Hz and loses by a factor of $17$ at $1$~Hz; the main gain from reconfiguration lies in the threshold selection of the band (Exercise~\ref{ex:18:4}).}

\answer{ex:19:1}{(a)~One synapse gives $6.7\mV$ ($R=66.7$~M$\Omega$), so reaching threshold at all needs at least $4$ (three suffice only asymptotically); $8$ within $10\ms$; $14$ within $5\ms$. (b)~$0.1\ms\ll\tau_m$; the leak correction is $\sim0.25\%$.}
\answer{ex:19:2}{(a)~$1.75\cdot10^{10}$ (a quarter of the mixers respond to everything), $4.4\cdot10^9$, and $0.06$ (at $k=40$ they almost never respond). (b)~By a factor of $10^3$: $2\cdot10^5$ versus $200$.}
\answer{ex:19:3}{(a)~$C(2n,n)=2^{2n-1}$ by the symmetry of binomial coefficients. (b)~$0.93$ and $0.09$; the transition width is $\sim\sqrt{2n}$ in $P$, relative width $\sim1/\sqrt n$.}
\answer{ex:19:4}{From one dendrite $V_{\text{soma}}<\kappa E_E<\theta$ for any number of inputs; with $g_0=0.5g_L$, $n\ge3$ are needed on each dendrite.}
\answer{ex:19:5}{$I\ge C\sigma_V/\sigma_t=15$~nA, i.e., $150$ synchronous synapses, which is unrealistic; the small neuron would need only $1$~nA, and with a $4$~nA synapse its jitter is $5$~\textmu s.}
\answer{ex:19:6}{Peak at the near end: $0.03\,\tau_m$ and $0.97$ ($L=0.2$); $0.32\,\tau_m$ and $0.66$ ($L=1$); $0.79\,\tau_m$ and $0.33$ ($L=2$). The estimate~\eqref{eq:dendriteload} from the total capacitance holds only for $L\ll1$.}
\answer{ex:19:7}{Open problem. At fixed $m$ the response time grows linearly with the stored $P$; at a fixed fraction $m=\varphi n$ it no longer depends on $n$, and the slowness of a large neuron is a consequence of summing many inputs, not of size as such.}

\answer{ex:20:1}{$t_{\text{w}}=\tau_r\ln\frac{1-L}{1-H}=16.8$~h, $t_{\text{s}}=\tau_d\ln\frac HL=5.8$~h, period $22.5$~h; the daily swing of the thresholds pulls the oscillator to $24$~h: a phase-locked loop~\eqref{eq:pll}.}
\answer{ex:20:2}{(a)~$L=0.111$; with $L=0.092$ sleep would last $9.6$~h. (b)~By $22\%$ ($1/0.82=1.22$), not by $18\%$.}
\answer{ex:20:3}{$H=\frac{p-1}{p-1/q}=0.623$, $L=H/q=0.093$, where $p=e^{16/\tau_r}$, $q=e^{8/\tau_d}$. After being woken at $4$~h the phase shifts permanently (falling asleep $7.2$~h earlier); with swinging thresholds it returns within two to three days.}
\answer{ex:20:4}{(a)~$r_\pm=\frac12\bigl(1\pm\sqrt{1-4k/w}\bigr)=0.947,\ 0.053$; $a_{\text{hi}}=3.51$, $a_{\text{lo}}=0.49$. (b)~Activity $\approx0.33$~s, silence $\approx0.59$~s, period $\approx0.93$~s, active fraction $0.36$.}
\answer{ex:20:5}{$a_{\text{hi}}/r_+<b<a_{\text{lo}}/r_-$: $3.71<b<9.20$. \nt{ACET} lowers $b$ below $3.71$: the intersection moves to the upper branch, and the active state ($r\approx1$) is stable.}
\answer{ex:20:6}{$4.9$, $6.8$, $8.8$, $5.5$, $8.9$~h for $D=24$, $32$, $40$, $48$, $64$: the duration is set by the daily phase at which sleep begins, not by the debt.}
\answer{ex:20:7}{After B the error on A averages $0.51$ (from $0.19$ to $1.08$ over $20$ sets; a zero output gives $1$). Interleaved, both errors are below $10^{-4}$.}
\answer{ex:20:8}{Open problem. Uniform scaling preserves the weight ratios, but it preserves the responses of a threshold neuron only if the threshold is scaled by the same factor; interleaved replays change the ratios. That the large contacts stay unchanged contradicts purely uniform scaling.}

\answer{ex:21:1}{(a)~$T=\tau_{\text{rec}}\ln\frac{1}{1-x^*}$: $1.84$~s and $3.68$~s. (b)~From $3.36$ to $4.24$~s: the sensitivity is $\dd T/\dd x^*=\tau_{\text{rec}}/(1-x^*)=80$~s, so as $x^*\to1$ the network bursts become irregular.}
\answer{ex:21:2}{(a)~$8.6$~W, as in~\eqref{eq:energy}. (b)~$205$~Mbit/s, $0.2$~W: $2\cdot10^3$ times more than the culture ($10^{-4}$~W).}
\answer{ex:21:3}{$x_{\text{ss}}=1/(1+Ur_b\tau_{\text{rec}})<x^*$ for $r_b>(1/x^*-1)/(U\tau_{\text{rec}})=0.28$~Hz at $x^*=0.9$ and $0.025$~Hz at $x^*=0.99$; the synaptic strength drops to $x_{\text{ss}}$, that is, by $10\%$ and $1\%$.}
\answer{ex:21:4}{The $u(t-k)$ are orthonormal, $m_k=\|P_Vu(t-k)\|^2$, where $P_V$ is the projector onto the $N$-dimensional span of the outputs; by Bessel's inequality $\sum_k m_k\le N$.}
\answer{ex:21:5}{$\mathrm{MC}\approx9$--$11$ with $\tanh$ and $15$--$18$ for the linear reservoir (three seeds), both $\le30$: nonlinearity is paid for with memory.}
\answer{ex:21:6}{(a)~$n\le102$. (b)~By $5.6$.}
\answer{ex:21:7}{An open problem. The key control: freeze the readout and check whether the improvement persists after a pause without stimulation; if it does not, the state changed, not the weights.}

\answer{ex:22:1}{(a)~$32\%$ and $100\%$. (b)~$R_\uparrow-R_\downarrow\ge2\sqrt{2000\,\text{Hz}/1\,\text{s}}\approx89$~Hz, only $4.5\%$ of the total rate.}
\answer{ex:22:2}{(a)~$n\ge4(p_1q_1+p_2q_2)/\Delta^2\approx560$. (b)~$\approx56$.}
\answer{ex:22:3}{(a)~Detailed balance: $\rho PK$ is symmetric. (b)~The fraction of misses equals the harmonic mean $1/\langle1/P\rangle=0.18$, versus $0.5$ without feedback.}
\answer{ex:22:4}{(a)~A parallelogram of area $4h^2=0.04$ (with the clipping of $p$ taken into account, $0.045$ numerically). (b)~$\approx0.18$.}
\answer{ex:22:5}{$0.49$: some cultures drift into a region where they almost always miss and wander there. With clipping, $1.00$: on a bounded set the density $\propto1/P$ gathers in the absorbing region.}
\answer{ex:22:6}{(a)~At least $5$ levels are needed; $8$ electrodes suffice. (b)~No: this needs $r_{\max}\ge25/T=100$~Hz.}
\answer{ex:22:7}{An open problem. The time of random search grows with $d$ roughly as the volume ratio $(L/h)^d$; for search with a signed error, linearly in $d$. The hypotheses are separated by a swap experiment: after a miss, a predictable but strong stimulus; after a hit, an unpredictable but weak one; several dozen cultures are needed in each group.}


\chapter{Experimental Supplement}
\label{sec:supplement}

For each chapter, this appendix collects its main claims and what they rest on: the experiment in which each was measured, what exactly was measured, and where it was published. The status in the right-hand column says how firm a claim is: \emph{measured} means there is a direct experiment; \emph{theory}, a mathematical consequence derived in the chapter or in a classic paper; \emph{model}, a result computed with the book's model (scripts in the \texttt{models/} folder); \emph{estimate}, an order-of-magnitude number without direct measurement; \emph{hypothesis}, a plausible explanation not yet proven by experiment. Where there is no experiment, we say so.

\section{Chapter~\ref{sec:neurontypes}. Neuron Types}

The chapter's numbers rest on direct cell counts in the human brain where such counts exist; the ``one neuron, one transmitter'' rule rests on cell atlases and on experiments that have also found the exceptions; the role of \nt{DOPA} rests on spike recordings; and the roles of the other broadcast channels rest on hypotheses.

{\footnotesize
\setlength{\tabcolsep}{3pt}\renewcommand{\arraystretch}{1.15}
\begin{longtable}{>{\raggedright}p{0.5cm}>{\raggedright}p{4.0cm}>{\raggedright}p{5.6cm}>{\raggedright}p{2.3cm}>{\raggedright\arraybackslash}p{1.7cm}}
\toprule
No. & Claim & Experiment and what was measured & Source & Status \\
\midrule\endfirsthead
\toprule
No. & Claim & Experiment and what was measured & Source & Status \\
\midrule\endhead
1 & The human brain has about $8.6\cdot10^{10}$ neurons, and almost all \nt{GLUT} neurons are small cerebellar neurons (Table~\ref{tab:types}) & Isotropic fractionator: tissue is dissolved into a uniform suspension of nuclei, and nuclei carrying the neuronal marker NeuN are counted. Adult men have $86.1\pm8.1$ billion neurons; only $19\,\%$ of them are in the cerebral cortex, and the bulk is in the cerebellum & \cite{Azevedo2009} & measured \\
2 & Almost every neuron has one principal transmitter (Proposition~\ref{prop:dale}), but there are exceptions & Transcriptomic atlas of the mouse brain: about $4\cdot10^6$ cells, $5322$ clusters; each cluster is assigned a transmitter from its synthesis and packaging genes. The exceptions were measured directly: in slices, \nt{DOPA} neurons also release GABA through the same vesicular transporter and inhibit striatal neurons; in some \nt{DOPA} axons, glutamate and dopamine leave from different terminals of the same wire (electron microscopy and optogenetics) & \cite{Strata1999,Yao2023,Tritsch2012,Zhang2015} & measured \\
3 & The entire column of the weight matrix has one sign~\eqref{eq:dalesign} & Derived in the chapter from Proposition~\ref{prop:dale} and the fact that glutamate receptors are almost all excitatory and GABA receptors inhibitory. There is no direct test of ``all recipients of one neuron receive a weight of one sign'' as a general rule; counterexamples are GABA co-release by \nt{DOPA} neurons (row~2) and recipients with receptors of opposite sign (row~11) & derivation in the chapter & theory \\
4 & Three-quarters to four-fifths of cortical neurons are \nt{GLUT}, a fifth to a quarter are \nt{GABA} & GABA immunostaining in columns $50\um$ wide through the full depth of the cortex in $10$ areas of monkey cortex: in $8$ areas GABA neurons make up about $25\,\%$ of all neurons, in primary visual cortex less than $20\,\%$ & \cite{Hendry1987} & measured \\
5 & A \nt{GLUT} neuron has about $10^4$ outputs & Postmortem stereology: the neocortex of young men has $1.64\cdot10^{14}$ synapses (electron microscopy, disector) and about $2.3\cdot10^{10}$ neurons. The ratio gives $\sim7\cdot10^3$ synapses per neuron; on average the number of inputs equals the number of outputs & \cite{Tang2001synapses,Pakkenberg1997} & measured \\
6 & The output of a \nt{DOPA} neuron branches into $10^5$--$10^6$ terminals; this is an estimate from target coverage, not a count & Viral membrane labeling of single rat \nt{DOPA} neurons and full axon reconstruction: one neuron covers $0.45$ to $5.7\,\%$ (on average $2.7\,\%$) of the striatal volume. Coverage was measured, not the number of terminals: that is inferred from coverage to order of magnitude, with no direct terminal count. For \nt{SERO} and \nt{NORA} the terminal numbers are rough estimates & \cite{Matsuda2009} & measured \\
7 & Broadcast neurons are few: \nt{DOPA} $\sim5\cdot10^5$ (the main one of two groups, a lower bound), \nt{SERO} $\sim3\cdot10^5$ (sum of two partial counts), \nt{HIST} $\sim6\cdot10^4$ (Table~\ref{tab:types}, Fig.~\ref{fig:typepie}) & Stereological counts in humans: $550\,000$ pigmented neurons in the substantia nigra (without the ventral tegmental area); $235\,000$ neurons in the dorsal raphe nucleus (all neurons of the nucleus) and another $125\,000$ serotonin neurons in the pontine tegmentum; about $64\,000$ histamine neurons in the hypothalamus. For \nt{NORA}, \nt{GLYC}, \nt{ACET}, and \nt{ADRE} there is no direct count: in the table these are rough order-of-magnitude estimates & \cite{Pakkenberg1991,Baker1990,Baker1991,Airaksinen1991} & measured \\
8 & Glutamate in the gap shoots up to about a millimolar and decays within $\sim1\ms$ & From the displacement of a rapidly unbinding competitive blocker from NMDA receptors during synaptic transmission (cultured hippocampal neurons): peak $1.1$\,mM, decay time constant $1.2\ms$ & \cite{Clements1992} & measured \\
9 & In working cortex the excitatory and inhibitory currents are nearly balanced, and the neuron sits near threshold & Theory: a network with strong connections settles into the balanced regime on its own. Experiment: in ferret prefrontal cortex in vivo, during ``up'' states the reversal potential of the synaptic current stays near $-37$\,mV for hundreds of milliseconds, although the conductance changes by $\sim21$\,nS: excitation and inhibition rise and fall together & \cite{vanVreeswijk1996,Haider2006} & measured \\
10 & \nt{DOPA} neurons report the reward-prediction error~\eqref{eq:rpe} (Fig.~\ref{fig:rpe}) & Spikes of monkey \nt{DOPA} neurons: a burst for unexpected juice; after learning, a burst for the cue and no response to predicted juice; a dip in rate when promised juice is withheld. In a quantitative experiment, the rate is proportional to the difference between the current reward and a weighted average of past ones, but only for positive errors. Equation~\eqref{eq:rpe} itself is the temporal-difference algorithm & \cite{Schultz1997,Bayer2005,SuttonBarto2018} & measured \\
11 & The receptor sets the sign and speed: ionotropic, fractions of a millisecond; metabotropic, much slower (Proposition~\ref{prop:receptor}) & Brief pulses of glutamate on membrane patches from hippocampal neurons: the current through ionotropic receptors rises (from $20$ to $80\,\%$) in $0.2$--$0.6\ms$ and decays in $2$--$3\ms$. The slow inhibitory potential of pyramidal neurons is abolished by a selective blocker of the metabotropic GABA receptor. Two families of dopamine receptors with opposite action; in the striatum, neurons with the D1 receptor need dopamine to strengthen synapses, neurons with D2 to weaken them & \cite{Colquhoun1992,DutarNicoll1988,Beaulieu2011,Shen2008} & measured \\
12 & A broadcast channel is not one wire but a bundle of a few lines (Section~\ref{sec:buses}) & Viral-genetic labeling of serotonin neurons in the mouse dorsal raphe nucleus: neurons projecting to the amygdala and to the frontal cortex lie in different parts of the nucleus, branch into complementary regions, and respond oppositely to an aversive stimulus. Review: subgroups of locus coeruleus neurons (\nt{NORA}) encode different processes & \cite{Ren2018,Poe2020} & measured \\
13 & \nt{NORA} sets the gain $g$ & The adaptive gain hypothesis; a network model in which catecholamines increase the steepness of the transfer characteristic. There is no direct measurement of ``$g$ rises with noradrenaline'' across the whole network; what has been measured is that pupil diameter follows the spikes of locus coeruleus neurons in the monkey & \cite{AstonJones2005,ServanSchreiber1990,Joshi2016} & hypothesis \\
14 & \nt{ACET} switches ``write/read'' and sets the learning rate & Hypothesis. Support: in rat olfactory cortex slices, acetylcholine agonists ($100$\,\textmu M) weaken the internal, ``recalling'' synapses by more than $60\,\%$ and the input synapses by less than $15\,\%$ & \cite{Hasselmo2006,HasselmoBower1992} & hypothesis \\
15 & \nt{SERO} sets the discount factor $\gamma$; serotonin neurons fire steadily, a few spikes per second, and are almost silent in one of the phases of sleep & The ``serotonin $=\gamma$'' hypothesis. Strongest argument: optogenetic activation of serotonin neurons in the mouse dorsal raphe nucleus reduces the number of premature give-ups while waiting for a delayed reward. The firing rate and silence during rapid-eye-movement sleep have been measured (review of recordings) & \cite{Doya2002,Miyazaki2014,JacobsAzmitia1992} & hypothesis \\
\bottomrule
\end{longtable}}

\section{Chapter~\ref{sec:glutcomputer}. What \nt{GLUT} Neurons Compute}

The chapter's logical claims are theorems about circuits with nonnegative weights, derived in the chapter itself; experiment supports the neuron model and the fact that circuits built according to these theorems (the coincidence detector, the delay line, the self-sustaining group, the chain) really do occur in the brain.

{\footnotesize
\setlength{\tabcolsep}{3pt}\renewcommand{\arraystretch}{1.15}
\begin{longtable}{>{\raggedright}p{0.5cm}>{\raggedright}p{4.0cm}>{\raggedright}p{5.6cm}>{\raggedright}p{2.3cm}>{\raggedright\arraybackslash}p{1.7cm}}
\toprule
No. & Claim & Experiment and what was measured & Source & Status \\
\midrule\endfirsthead
\toprule
No. & Claim & Experiment and what was measured & Source & Status \\
\midrule\endhead
1 & A neuron is an RC circuit with a threshold and reset~\eqref{eq:glutneuron} & A noisy current resembling the synaptic bombardment in the living brain was injected into the somata of rat cortical pyramidal neurons (slices). The dependence of the mean rate on the mean and variance of the current is described by a leaky integrate-and-fire neuron with added slow rate adaptation. The ranges of $\tau$ and delays are given in the chapter without reference to an experiment & \cite{Rauch2003} & measured \\
2 & Model~\eqref{eq:glutneuron} is a simplification: the input branches generate local spikes of their own & Recording directly from dendrites of pyramidal neurons in mouse visual cortex in vivo: a visual stimulus evokes local dendritic spikes that depend on NMDA receptors; suppressing them broadens the neuron's orientation tuning & \cite{Smith2013} & measured \\
3 & The coincidence window is $\Delta_{\max}=\tau\ln\frac{w}{\theta-w}$~\eqref{eq:window}; for $\theta=1.5w$ it equals $0.7\tau$ & A consequence of model~\eqref{eq:glutneuron}. A direct measurement of a narrow summation window exists for neurons with fast inhibition (Chapter~\ref{sec:gaba}, row~3 of its table) & derivation in the chapter & theory \\
4 & A network of \nt{GLUT} neurons alone computes only monotone functions of its inputs (Proposition~\ref{prop:monotone}) & Proof in the chapter: iteration from silence with a nondecreasing right-hand side. This is a claim about the model; there is no experiment testing it on a living network without inhibition & derivation in the chapter & theory \\
5 & Sense organs provide a pair of wires: some cells respond to an increase of the stimulus, others to a decrease & Retina: already at the bipolar cells the cone signal splits into on and off channels. Pharmacological block of on bipolar cells in the monkey: the channels run separately up to the cortex; after the block, detection of light increments is impaired but not of decrements & \cite{Schiller1992} & measured \\
6 & With a dual-rail code, a network of \nt{GLUT} neurons computes any logical function asynchronously, without a common beat, at the price of twice as many neurons (Proposition~\ref{prop:dualrail}, \eqref{eq:dualrail}, Fig.~\ref{fig:dualxor}) & The dual-rail code and completion detection from the signal itself are a standard technique of delay-insensitive asynchronous circuits. The six-neuron exclusive-OR circuit is built in the chapter. There is no experiment showing that the brain computes logical functions specifically with a dual-rail code & \cite{Sparso2001} & theory \\
7 & The largest interaural arrival-time difference in humans is $\approx0.6\ms$, and the brain resolves about ten microseconds & Listeners in a two-alternative experiment: the time-difference discrimination threshold (75\,\% correct) is $9$\,\textmu s for noise in the $150$--$1700$\,Hz band, $11$\,\textmu s for a $1$\,kHz tone, $28$\,\textmu s for a click. The $0.6\ms$ limit is the chapter's calculation, $0.22\,\text{m}/343\,\text{m/s}$ & \cite{KlumppEady1956} & measured \\
8 & In birds, a time difference is turned into a place by delay lines and coincidence detectors~\eqref{eq:itd} (Fig.~\ref{fig:itd}) & Barn owl: axons from the two ears enter the nucleus of coincidence detectors from opposite sides; the arrival time of spikes along the axon changes systematically with depth, and the spread of delays across the nucleus ($160$\,\textmu s over $700\um$) matches the range of interaural delays & \cite{Carr1990} & measured \\
9 & In mammals no row of delays has been found; the same task is solved by coincidence detectors with precisely timed inhibition from \nt{GLYC} neurons & In vivo recording in the gerbil medial superior olive: the responses do not form a map of directions; iontophoresis of glycine and its blocker shows that precisely timed \emph{glycinergic} inhibition sets the working range of delays. The inhibition here comes from \nt{GLYC}, not from \nt{GABA} & \cite{Brand2002,Grothe2010} & measured \\
10 & A group with strong feedback is a flip-flop; self-sustaining activity persists for seconds after the stimulus (Fig.~\ref{fig:bistable}) & Monkey prefrontal cortex in a delayed eye-movement task toward a remembered location (delay $1$--$6$\,s): $87$ of $170$ task-related neurons change their rate during the delay, and in $79\,\%$ of them this depends on the direction of the remembered location. That this activity is held by feedback with two stable states is theory & \cite{Funahashi1989,Wang2001} & measured \\
11 & A bit is written if the drive lasts longer than $\approx0.7\tau$ (Fig.~\ref{fig:recurrentgroup}b) & Euler integration of $\tau\,\dd r/\dd t=-r+f(wr+I)$ with $w=1$, $I=0.6$; there is no script in \texttt{models/}. No experiment & calculation in the chapter & model \\
12 & Memory can be stored in synaptic facilitation, with almost no spikes & Theory: residual calcium in the terminals holds the record for about a second. Support: in ferret medial prefrontal cortex (simultaneous recording from several neurons) there is a subnetwork of pyramidal neurons connected by facilitating synapses with long aftereffects. A review of observations of ``quiet'' intervals during memory maintenance. Which way the cortex uses when is an open question & \cite{Mongillo2008,WangMarkram2006,Stokes2015} & hypothesis \\
13 & Memory is erased by fatigue: the stock of transmitter is depleted & Paired recordings of cortical pyramidal neurons: the synaptic response falls with frequent spikes, and the rate of the fall is set by the release probability. That this is precisely what turns off a self-sustaining group has not been shown directly & \cite{Tsodyks1997} & measured \\
14 & A chain of groups is an event-triggered timer; in songbirds neurons fire strictly in turn & Recording of identified neurons of the HVC song nucleus in the zebra finch during sleep and singing: each neuron projecting to the motor nucleus fires one short burst at one precise moment of the song, and the neurons fire in sequence & \cite{Hahnloser2002} & measured \\
\bottomrule
\end{longtable}}

\section{Chapter~\ref{sec:gaba}. \nt{GLUT} and \nt{GABA} Together}

The chapter's logical and dynamical claims are derived in the chapter itself from linear analysis and Ohm's law; experiment shows that the circuits they rely on (the delayed veto, direct inhibition following excitation, stabilization by inhibition, the rhythm of the \nt{GLUT}--\nt{GABA} loop) really do operate in the brain.

{\footnotesize
\setlength{\tabcolsep}{3pt}\renewcommand{\arraystretch}{1.15}
\begin{longtable}{>{\raggedright}p{0.5cm}>{\raggedright}p{4.0cm}>{\raggedright}p{5.6cm}>{\raggedright}p{2.3cm}>{\raggedright\arraybackslash}p{1.7cm}}
\toprule
No. & Claim & Experiment and what was measured & Source & Status \\
\midrule\endfirsthead
\toprule
No. & Claim & Experiment and what was measured & Source & Status \\
\midrule\endhead
1 & A fifth to a quarter of cortical neurons are \nt{GABA} & GABA immunostaining in $10$ areas of monkey cortex: about $25\,\%$ of neurons in $8$ areas, less than $20\,\%$ in primary visual cortex. A review of the diversity of cortical inhibitory neurons & \cite{Hendry1987,Markram2004} & measured \\
2 & What inhibition does is largely decided by its landing site: dendrite, cell body, spike initiation site, or the terminal of another neuron's wire & Review: classes of cortical \nt{GABA} neurons differ in their target on the recipient. Simultaneous recording from the soma and dendrite of a pyramidal neuron: feedforward inhibition is much stronger at the soma than at the dendrites. Inhibition at the terminal of another neuron's wire in the spinal cord reduces transmitter release from one input line (review of measurements) & \cite{Markram2004,Pouille2001,Rudomin1999} & measured \\
3 & A sign change through an intermediary costs one delay, about a millisecond; inhibition arriving right after excitation narrows the coincidence window & Rat hippocampal pyramidal neurons in slices: inhibition through one intermediary follows direct excitation after a short delay, and the window in which two excitatory inputs sum to threshold is less than $2\ms$. On the dendrites, where this inhibition is weaker, the window is wider & \cite{Pouille2001} & measured \\
4 & The veto $a\wedge\bar b$~\eqref{eq:veto} with OR and a constant one is functionally complete (Proposition~\ref{prop:gabacomplete}) & Proof in the chapter. A classical result: a network of threshold elements with inhibitory inputs implements any logical expression. A claim about the model; no experiment & \cite{McCullochPitts1943} & theory \\
5 & A delayed veto gives a direction-of-motion detector with optimal speed $v^*=\Delta x/d$~\eqref{eq:vstar} & Rabbit retina: direction selectivity is created by inhibition that cancels excitation for motion in the ``null'' direction. Separate recording of the excitatory and inhibitory inputs of selective cells showed that starburst amacrine cells (\nt{GABA}) inhibit them asymmetrically, only with processes pointing to the ``null'' side. The formula for $v^*$ is the chapter's derivation & \cite{Barlow1965,Fried2002} & measured \\
6 & Shunting inhibition divides the voltage rather than subtracting~\eqref{eq:shunt} (Fig.~\ref{fig:shunt}) & Ohm's law for a divider of three conductances with the chloride equilibrium potential near rest; for a neuron that does not fire & derivation in the chapter & theory \\
7 & On the spike rate the shunt acts by subtraction, and division comes from inhibition together with balanced background noise & Model: in a firing neuron the current through the shunt is nearly constant, and the effect on rate is subtractive. Experiment: background excitatory and inhibitory conductances were injected into rat somatosensory cortex pyramidal neurons (slices) with a dynamic clamp; a simultaneous balanced increase of both divides the slope of the response with no noticeable shift in rate. Review: division by total activity occurs in many parts of the brain & \cite{HoltKoch1997,Chance2002,Carandini2012} & measured \\
8 & For $\beta>1$ mutual inhibition selects a single winner (Proposition~\ref{prop:wta}, \eqref{eq:wta}) & The eigenvalues $-(1\pm\beta)/\tau$ are derived in the chapter. The rates $0.73$ and $0.53$ at $\beta=0.5$ are the fixed point $r_{1,2}=(I_{1,2}-\beta I_{2,1})/(1-\beta^2)$; there is no script in \texttt{models/}. No direct experiment is cited in the chapter & derivation in the chapter & theory \\
9 & Memory reset by a signal through \nt{GABA}; two groups with mutual inhibition form a latch & Follows from Fig.~\ref{fig:bistable} and Proposition~\ref{prop:wta}. No experiment & derivation in the chapter & theory \\
10 & The equilibrium of the \nt{GLUT}--\nt{GABA} loop is stable if the inhibition is strong enough and fast enough (Proposition~\ref{prop:ei}) & Two-population model~\eqref{eq:ei}; conditions (i) and (ii) are the determinant and trace of its matrix, derived in the chapter & \cite{WilsonCowan1972} & theory \\
11 & For $w_{EE}=1.5$ and $w_{EI}w_{IE}=2$, fast inhibition ($\tau_I=2\ms$) holds $E^*=0.67h$, and slow inhibition ($30\ms$) drives oscillations (Fig.~\ref{fig:ei}) & Numerical solution of~\eqref{eq:ei}, script \texttt{figures/make\_ei.py} & calculation & model \\
12 & The cortex works in an inhibition-stabilized regime; the signature: push the \nt{GABA} neurons, and their rate drops~\eqref{eq:isn} & Theory predicted the paradoxical response. Experiment: intracellular recordings in cat primary visual cortex show that when the response is suppressed by a surrounding stimulus, the inhibitory conductance does not rise but falls together with the excitatory one. Optogenetic excitation of PV neurons in mouse visual, somatosensory, and motor cortex lowers the rate of inhibitory neurons, including without a stimulus and under light anesthesia. The coefficient $-1/3$ for the numbers of Fig.~\ref{fig:ei} is the chapter's derivation & \cite{Tsodyks1997isn,Ozeki2009,Sanzeni2020} & measured \\
13 & The loop ``\nt{GLUT} excites \nt{GABA}, \nt{GABA} inhibits \nt{GLUT}'' generates a fast rhythm with a period of $12$--$50\ms$ & Optogenetic stimulation of fast \nt{GABA} neurons in mouse somatosensory cortex at $8$--$200$\,Hz selectively enhances the gamma rhythm ($20$--$80$\,Hz, period $12$--$50\ms$); stimulation of \nt{GLUT} neurons enhances only slower rhythms & \cite{Cardin2009,Buzsaki2012} & measured \\
\bottomrule
\end{longtable}}

\section{Chapter~\ref{sec:code}. Morse Code}

The chapter's limiting estimates are information theory and arithmetic; what has been measured is where noise arises in the channel, how fast the brain works, and what a spike costs, while the question of which code the cortex actually uses remains open.

{\footnotesize
\setlength{\tabcolsep}{3pt}\renewcommand{\arraystretch}{1.15}
\begin{longtable}{>{\raggedright}p{0.5cm}>{\raggedright}p{4.0cm}>{\raggedright}p{5.6cm}>{\raggedright}p{2.3cm}>{\raggedright\arraybackslash}p{1.7cm}}
\toprule
No. & Claim & Experiment and what was measured & Source & Status \\
\midrule\endfirsthead
\toprule
No. & Claim & Experiment and what was measured & Source & Status \\
\midrule\endhead
1 & Rates of cortical \nt{GLUT} neurons range from fractions to tens of spikes per second, and most of the time neurons work near the lower end & Cell-attached recording (neuron selection independent of its activity) in the auditory cortex of awake rats: at any moment fewer than $5\,\%$ of neurons fire above $20$ spikes per second; some neurons have a background rate below $0.01$ per second; the rate distribution is lognormal & \cite{Hromadka2008} & measured \\
2 & The capacity of a spike channel is $R_{\max}\approx r\log_2\frac{e}{r\Delta t}$~\eqref{eq:spikeentropy}, and almost all of it lies in spike timing (Proposition~\ref{prop:capacity}) & Entropy of a binary string of bins of length $\Delta t$. The chapter's numbers: $8$ bits per spike at $r=10$ per second and $\Delta t=1\ms$, about $5$ bits at $\Delta t=10\ms$, fewer than $2$ bits per second for a spike counter & \cite{MacKay1952} & theory \\
3 & Sensory neurons transmit one to several bits per spike, in the best cases up to half the limit & Stimulus reconstruction from the spikes of sensory neurons whose stimulus is known; a summary of measurements in the book & \cite{Rieke1997} & measured \\
4 & The spike generator is precise; precise timing is set by fast transients of the input & Rat cortical neurons in slices: for a repeated current with fast fluctuations, spikes repeat with a precision better than $1\ms$; for a constant current the times drift & \cite{Mainen1995} & measured \\
5 & The synapse is unreliable: more than half of the spikes do not reach the recipient because of random release & Minimal stimulation of inputs to hippocampal pyramidal neurons (slices): more than half of the spikes evoke no response. Fluctuations of the stimulation threshold, conduction failures at axon branch points, and low experimental temperature were ruled out; what remains is probabilistic release & \cite{AllenStevens1994} & measured \\
6 & The spike train in living cortex looks random: intervals are scattered as in a Poisson stream, and the variance of the spike count is close to the mean; part of the ``noise'' is a message about the animal's movements & Recordings in visual areas V1 and MT of awake monkeys: coefficient of variation of the intervals about $1$, ratio of spike-count variance to mean as for a random process; a model summing independent small inputs gives much less variability. In mouse visual cortex a considerable part of the variability is predicted from video of the animal's own movements & \cite{Softky1993,Stringer2019} & measured \\
7 & The rate code is slow: error $1/\sqrt{rT}$~\eqref{eq:rateerror}, while the brain recognizes a picture in $\sim150\ms$ & The formula is Poisson statistics, the chapter's derivation. Experiment: people decided whether an unfamiliar photograph flashed for $20\ms$ contained an animal; the brain potentials for ``yes'' and ``no'' diverge after about $150\ms$. Splitting this time into ten stages of $10$--$20\ms$ is the chapter's estimate, not a measurement & \cite{Thorpe1996} & measured \\
8 & Averaging over a group is limited by correlation: the error falls by no more than a factor of $1/\sqrt c$~\eqref{eq:correlated} (Proposition~\ref{prop:pooling}) & Pairs of monkey MT neurons recorded simultaneously: the spike-count correlation coefficient averages $0.12$, and theoretical analysis shows that even such a correlation limits the gain from a group. Theory: only the part of the correlation that looks like the signal sets the limit. A model of the opposing position: the rate of a group of $50$--$100$ neurons is read within one interspike interval, $10$--$50\ms$, and the timing of individual spikes is not used in the cortex & \cite{Zohary1994,MorenoBote2014,ShadlenNewsome1998} & measured \\
9 & A number can be written in the first-spike time~\eqref{eq:latency}, in the spike order of a group, or in the phase of a local rhythm & Salamander retina: the spatial structure of a briefly flashed image is recorded in the relative latencies of the first spikes of the output neurons, almost independently of contrast. Rat hippocampal place cells: while passing through ``their'' place, spikes shift earlier and earlier in the theta cycle ($7$--$12$\,Hz), a shift of $100^\circ$ to $355^\circ$. How much the cortex uses spike timing is an open question (row~8) & \cite{Gollisch2008,OKeefe1993} & measured \\
10 & Which code gets read is decided by the recipient's time constant~\eqref{eq:reader} (Proposition~\ref{prop:reader}); in active cortex neurons are closer to coincidence detectors & Formula~\eqref{eq:reader} is the chapter's derivation. Review of in vivo recordings: in active cortex the membrane conductance is several times higher than at rest, and the time constant correspondingly shorter. That the cortex switches codes in precisely this way has not been shown directly & \cite{Destexhe2003} & hypothesis \\
11 & A depressing synapse transmits the change in rate, essentially $\Delta r/r$~\eqref{eq:depression} (Fig.~\ref{fig:depression}) & Paired recordings of cortical pyramidal neurons: the rate of depression is set by the release probability; with slow depression the response reflects rate, with fast depression coincidences. A model based on measured depression in rat cortex: equal relative rate changes on fast and slow inputs give equal responses, that is, the synapse transmits $\Delta r/r$. The numbers $1.7\to2.2$, $6.7$, and $90\ms$ are the chapter's calculation & \cite{Tsodyks1997,Abbott1997depression} & measured \\
12 & A burst is a ``dash'': its probability of being lost, $(1-p)^k$~\eqref{eq:burst}, is small, and facilitation lowers it further & Review: at unreliable synapses single spikes are often lost, while bursts are transmitted reliably thanks to facilitation; a burst induces long-term synaptic changes. That the brain reads a burst as a separate symbol is a hypothesis & \cite{Lisman1997,AllenStevens1994} & hypothesis \\
13 & Fast \nt{GABA} neurons set the rhythm and ``punctuation''; one more class inhibits the inhibitors and opens the gates (Proposition~\ref{prop:punctuation}) & Review of the properties of classes of cortical \nt{GABA} neurons. Optogenetics: rhythmic excitation of fast \nt{GABA} neurons evokes the gamma rhythm, and the response to a stimulus depends on the phase of the cycle. In mouse visual cortex PV neurons inhibit one another, SST neurons inhibit all the other classes, and VIP neurons inhibit mainly SST. Excitation of VIP neurons in awake mice releases \nt{GLUT} neurons from inhibition; they are switched on by a reinforcement signal. The division of roles as a general rule is a hypothesis & \cite{Tremblay2016,Cardin2009,Pfeffer2013,Pi2013} & hypothesis \\
14 & Energy limits the mean rate to about one spike per second~\eqref{eq:energy} & Energy budget of rodent gray matter from anatomical and physiological data: spikes and synaptic currents take $47\,\%$ and $34\,\%$ of the signaling energy; no more than $\sim15\,\%$ of neurons can be active at once. For the human cortex: no more than $\sim1\,\%$ of neurons can be strongly active at once. The figures of $\sim10^9$ ATP per spike and $\sim10$\,W for signaling are the chapter's estimate based on these calculations & \cite{Attwell2001,Lennie2003} & theory \\
15 & The code is sparse and tuned for the greatest number of bits per joule; the cortex trades precision for energy (Proposition~\ref{prop:economy}) & The efficient-coding hypothesis. Support: in food-restricted mice, AMPA receptor conductance and synaptic ATP use in the visual cortex fall ($-29\,\%$), the firing rate is preserved, and orientation tuning broadens by $32\,\%$ & \cite{Barlow1961,Padamsey2022} & hypothesis \\
\bottomrule
\end{longtable}}

\section{Chapter~\ref{sec:serocomputer}. \nt{SERO} neurons}

The properties of \nt{SERO} neurons themselves (rhythm, autoinhibition, sign set by the receptor, subsystems) are measured directly; the chapter's computations (regulator, filter, logic) follow from its equations; $\tau_c$, the diffusion coefficient of serotonin and the number of release sites are estimates; the loop's role as a level regulator in the brain is a hypothesis (in the raphe nucleus autoinhibition is point to point and causes only brief pauses), and the role of \nt{SERO} as the horizon is a hypothesis with behavioral experiments in its favor.

{\footnotesize
\setlength{\tabcolsep}{3pt}\renewcommand{\arraystretch}{1.15}
\begin{longtable}{>{\raggedright}p{0.5cm}>{\raggedright}p{4.0cm}>{\raggedright}p{5.6cm}>{\raggedright}p{2.3cm}>{\raggedright\arraybackslash}p{1.7cm}}
\toprule
No. & Claim & Experiment and what was measured & Source & Status \\
\midrule\endfirsthead
\toprule
No. & Claim & Experiment and what was measured & Source & Status \\
\midrule\endhead
1 & The sign of serotonin's action is chosen by the recipient's receptor (Proposition~\ref{prop:receptor}) & Serotonin receptors are divided into families by pharmacology and mechanism: 5-HT$_{1A}$ opens potassium channels through a G protein and inhibits the cell; 5-HT$_{2A}$ and the ionotropic 5-HT$_3$ excite. One transmitter gives opposite responses in different cells. & \cite{BarnesSharp1999} & measured \\
2 & In the waking state a \nt{SERO} neuron steadily fires a few spikes per second & Recordings of dorsal raphe neurons in freely moving cats: slow rhythmic firing of $2.8\pm0.2$ spikes/s in quiet waking; in slow-wave sleep the rate falls by $34$--$68\,\%$, in REM sleep by $98\,\%$. In anesthetized rats, $1.7\pm0.5$ spikes/s, very regular. & \cite{TrulsonJacobs1979, GagnonParent2014, JacobsAzmitia1992} & measured \\
3 & A \nt{SERO} neuron is a pacemaker: it needs no push for each spike, but it needs a steady background drive (in a slice, noradrenaline through $\alpha_1$ receptors; in the organism, from \nt{NORA}) & In a brain slice the cells fire slowly and steadily, with an afterhyperpolarization of $200$--$800\ms$ after each spike. Fewer cells fire spontaneously in a slice than in the organism: the steady rhythm needs a tonic noradrenaline input through $\alpha_1$ receptors, and blocking them abolishes it. & \cite{VandermaelenAghajanian1983} & measured \\
4 & Almost all receptors are metabotropic; the response is slow, rising and decaying within about a second & All receptor families except 5-HT$_3$ are coupled to G proteins. In a slice, evoked serotonin release produces an inhibitory current through 5-HT$_{1A}$ that rises and decays within a second. & \cite{BarnesSharp1999, CourtneyFord2016} & measured \\
5 & In target regions serotonin is released into the volume, not only into the cleft; in the raphe nucleus itself, inhibition of neighbors is point to point & Fast-scan cyclic voltammetry in slices: release per spike is the same in a region with synapses and in the raphe nucleus, where there are almost no synapses; it does not depend on blocking uptake or receptors; there are far fewer molecules per terminal than transporters and receptors, and the extracellular concentration matches the affinity of the main receptor. In the raphe nucleus, inhibition through 5-HT$_{1A}$ is point to point: the transporter keeps serotonin from accumulating outside the cleft. & \cite{Bunin1998, CourtneyFord2016} & measured \\
6 & The concentration in~\eqref{eq:seroneuron} decays because of reuptake; $\tau_c$ of order a second is an estimate & Voltammetry in the living brain: extracellular serotonin is limited primarily by reuptake and degradation, not by synthesis. The cited papers contain no direct measurement of $\tau_c$; its order of magnitude is the chapter's estimate. & \cite{Hashemi2012serotonin} & measured; $\tau_c$ is an estimate \\
7 & NOT and NOR with a single pacemaker; completeness of \nt{SERO} logic (Proposition~\ref{prop:serocomplete}, Fig.~\ref{fig:serogates}) & Follows from~\eqref{eq:seroneuron}: an inhibitable pacemaker is a NOR gate, and NOR is functionally complete. There is no experiment in which logic has been built from \nt{SERO} neurons; the condition of separate volumes does not hold in the brain. & derived in the chapter & theory \\
8 & Serotonin inhibits the \nt{SERO} neuron itself through receptors on it & In anesthetized rats, 5-HT$_{1A}$ agonists given intravenously and iontophoretically inhibit the firing of dorsal raphe neurons with the same dose--response relation as serotonin itself; 5-HT$_{1B}$ agonists have almost no effect. In a slice, endogenous serotonin from neighboring cells produces a current through 5-HT$_{1A}$ and a brief pause in firing without changing the mean rate. & \cite{Sprouse1987, CourtneyFord2016} & measured \\
9 & In the model the loop through a common volume is a level regulator: variation and the time constant shrink by a factor of $1+G$~\eqref{eq:feedback}; that the loop works this way in the brain is a hypothesis & The classical formula of an amplifier with negative feedback, derived again in the chapter. The loop gain $G$ of the \nt{SERO} system has not been measured. Measured: in a slice, autoinhibition is point to point, causes a brief pause and does not change the mean rate. & \cite{Black1934feedback, CourtneyFord2016} & theory; in the brain, hypothesis \\
10 & The concentration is a moving average of the rate, a low-pass filter~\eqref{eq:lowpass}, Fig.~\ref{fig:lowpass} & Solution of the linear equation~\eqref{eq:seroneuron}; the figure is computed from this formula with $\tau_c=1\,\text{s}$. There is no separate model in \texttt{models/}. & derived in the chapter & theory \\
11 & The tortuosity of the gaps slows the diffusion of a small molecule by a factor of about $2.6$ & Point-source method: iontophoresis of tetramethylammonium and recording of its concentration with an ion-selective electrode in slices and in the living brain. Tortuosity $\lambda\approx1.6$, so the effective $D$ is $\lambda^2\approx2.6$ times smaller than the free one; the extracellular volume fraction is about $0.2$. The diffusion coefficient of serotonin itself in tissue was not measured in these studies; in the chapter it is an estimate. & \cite{Sykova2008} & measured \\
12 & The length of an independent ``wire'' is $\ell\sim\sqrt{D\tau}\approx20$~\textmu m~\eqref{eq:diffusionlength}; the number of wires is limited by the number of such regions (Proposition~\ref{prop:volumewires}) & The diffusion law with the estimates of $D$ and $\tau$ substituted (rows~6 and~11); Proposition~\ref{prop:volumewires} is a consequence. There is no direct experiment counting the independent channels of volume transmission. & derived in the chapter & theory \\
13 & The output of a single \nt{SERO} neuron is spread over huge parts of the brain; the number of release sites is estimated at $\sim10^5$ & Complete reconstruction of single dorsal raphe neurons in the rat: all ascending axons branch heavily, the total axon length of one cell reaches $18.7$~cm, and the branches of one cell reach the striatum, prefrontal cortex and amygdala. The number of release sites per cell has not been counted directly. & \cite{GagnonParent2014} & measured; $10^5$ is an estimate \\
14 & \nt{SERO} neurons fall into several subsystems with different outputs and responses & Viral-genetic labeling in mice: neurons projecting to the amygdala and to the frontal cortex hardly overlap in their branching, receive different inputs and respond oppositely to an aversive stimulus; activating and silencing each group changes behavior in different ways. & \cite{Ren2018} & measured \\
15 & Patience condition $D<\tau_\gamma\ln(R/r_0)$~\eqref{eq:patience} with exponential discounting; with the measured, hyperbolic discounting, $D<\tau_\gamma(R/r_0-1)$ & Follows from discounting by $e^{-D/\tau_\gamma}$ or $1/(1+D/\tau_\gamma)$; derived in the chapter. Monkeys choosing at delays of seconds: discounting is closer to a hyperbola than to an exponential; the response of \nt{DOPA} neurons to the cue of a delayed reward also falls hyperbolically with delay. & derived in the chapter; \cite{KobayashiSchultz2008} & theory \\
16 & \nt{SERO} sets the horizon $\tau_\gamma$ (Section~\ref{sec:horizon}) & Optogenetic activation of dorsal raphe \nt{SERO} neurons in mice lengthens waiting for a delayed reward and increases persistence in foraging for a reward that has become scarce. In humans, lowering serotonin (a tryptophan-free diet) makes the discounting of delayed reward steeper. How the concentration becomes $\tau_\gamma$ is unknown. & \cite{Doya2002, Miyazaki2014, Lottem2018, Schweighofer2008} & hypothesis \\
\bottomrule
\end{longtable}}

\section{Chapter~\ref{sec:dofa}. Learning with \nt{DOPA}}

The mechanisms of the synapse (packets, coincidence detector, calcium, plasticity windows) and the behavior of dopamine as a prediction error are measured directly; that the rules lead to the gradient, and what broadcasting costs, are theorems; the outcome of learning to choose is computed with the book's model; the circuit that computes the error is a hypothesis.

{\footnotesize
\setlength{\tabcolsep}{3pt}\renewcommand{\arraystretch}{1.15}
\begin{longtable}{>{\raggedright}p{0.5cm}>{\raggedright}p{4.0cm}>{\raggedright}p{5.6cm}>{\raggedright}p{2.3cm}>{\raggedright\arraybackslash}p{1.7cm}}
\toprule
No. & Claim & Experiment and what was measured & Source & Status \\
\midrule\endfirsthead
\toprule
No. & Claim & Experiment and what was measured & Source & Status \\
\midrule\endhead
1 & Backpropagation of error in the form used in artificial networks has not been found in the brain & There is no direct experiment: this is a claim of absence. Reviews analyze what the algorithm requires (feedback connections mirroring the forward ones, a separate phase) and which biological approximations to it have been proposed. & \cite{Lillicrap2020, WhittingtonBogacz2019} & hypothesis \\
2 & The weight factors into the number of release sites, the release probability and the response to one packet: $w=N_s\,p\,A_1$~\eqref{eq:quantal} & Frog neuromuscular junction with reduced release: the recipient's response fluctuates in steps that are multiples of the spontaneous miniature response to one packet; the number of packets per spike obeys the Poisson law. & \cite{delCastillo1954} & measured \\
3 & The recipient's receptor is a coincidence detector; calcium triggers the weight change & Patch clamp of mouse neurons: Mg$^{2+}$ ions block the glutamate-opened channel at the resting potential and leave it upon depolarization. Hippocampal slices: a calcium chelator in the recipient blocks long-term potentiation, while light-triggered release of calcium inside the cell by itself strengthens transmission. & \cite{Nowak1984, Malenka1988calcium} & measured \\
4 & Either side carries out the decision: $A_1$ grows in the recipient, $p$ changes in the sender & Glutamate uncaged by light over a single spine causes rapid and lasting growth of the spine and of the current through its receptors; potentiation is accompanied by insertion of AMPA receptors. Depolarizing the recipient temporarily suppresses release by the sender; the effect is carried by endocannabinoids traveling back across the cleft (shown at \nt{GABA} inputs). Short-term depression depends on the sender's release probability. & \cite{Matsuzaki2004, Malinow2002, WilsonNicoll2001, Tsodyks1997} & measured \\
5 & A synapse strengthens when the sender and the recipient are active together~\eqref{eq:hebb} & Anesthetized rabbit: after brief trains of spikes along the hippocampal input pathway, the response is enhanced for $30$~min to $10$~h. In a hippocampal slice, the recipient's spikes strengthen only the synapses that are active at the same moment. & \cite{Bliss1973LTP, Kelso1986hebb} & measured \\
6 & Rule~\eqref{eq:oja} finds the principal eigenvector of the input correlations (Proposition~\ref{prop:oja}) & A theorem; the proof is in the chapter. There is no experiment in which a neuron has learned the principal component of its inputs; the mechanism that bounds weight growth in the synapse has not been established. & \cite{Oja1982} & theory \\
7 & Timing-based Hebb: a window of about $\pm20\ms$ (Fig.~\ref{fig:stdp}); chains grow from repeated sequences & Pairs of hippocampal neurons in culture: a recipient spike within $20\ms$ after a sender spike gives lasting strengthening, within $20\ms$ before, weakening; both depend on NMDA receptors. The ratio $A_-=0.6A_+$ in the figure is illustrative. That chains grow this way in a network has not been measured directly. & \cite{BiPoo1998} & measured \\
8 & Eligibility trace~\eqref{eq:threefactor}: dopamine arriving $1$--$2\,\text{s}$ after joint activity strengthens the connection, later dopamine does not (Proposition~\ref{prop:threefactor}) & Striatal slices, separate optogenetic stimulation of glutamate and dopamine inputs: the spine grows only if dopamine arrives $0.3$--$2\,\text{s}$ after the glutamate input; the window is set by the fast rise and decay of cAMP. In the cortex, pairing leaves separate traces for strengthening and weakening, which monoamine receptors convert into long-term changes within a window of order seconds. & \cite{Yagishita2014, He2015eligibility} & measured \\
9 & Windows lasting seconds also occur without dopamine: the third signal is strong excitation of the recipient & Awake mouse, area CA1: a single plateau event in the recipient strengthens inputs that arrived seconds before and after it, and creates a place field in a single pass. & \cite{Bittner2017} & measured \\
10 & The mean step of the three-factor rule follows the reward gradient~\eqref{eq:perturbation} (Proposition~\ref{prop:gradient}) & The gradient theorem for learning from random perturbations; derived in the chapter for small noise. & \cite{Williams1992} & theory \\
11 & Noise is search: the network learns from its own random deviations & Juvenile zebra finches: inactivating the output nucleus of the basal ganglia circuit abruptly and reversibly removes song variability. Adult zebra finches: noise is played on renditions of a syllable with pitch on one side of the mean, and the birds quickly shift the pitch the other way, learning from their own variability. & \cite{Olveczky2005, TumerBrainard2007} & measured \\
12 & A two-way choice is learned with $\tau_e=1\,\text{s}$ and $\delta=R-\bar R$; it is not learned with $\tau_e=50\ms$; without subtraction the weights grow without bound, and at a high learning rate the network can lock in the worse choice & \texttt{models/dofa\_learning.py}, $\eta=0.05$, $500$ trials, $6$ seeds. $\tau_e=1\,\text{s}$: fraction of $A$ choices $0.65$--$0.79\to0.96$--$1.00$, weights $0.85$--$1.33$. $\tau_e=50\ms$: $0.38$--$0.55$, weights unchanged. $\delta=R$: winner's weight $860$--$1190$. $\delta=R$, $\eta=0.5$, $3$ seeds: one locked onto $B$ ($0.04\to0$). The script prints all four cases; the rerun matches the chapter. & \texttt{models/\allowbreak dofa\_\allowbreak learning.py} & model \\
13 & The learning time from one shared signal grows in proportion to the number of noisy neurons (Proposition~\ref{prop:broadcastcost}) & Learning curves of a linear network: node perturbation is slower than the exact gradient by a factor equal to the number of output neurons, and weight perturbation is slower by a further factor equal to the number of inputs. & \cite{Werfel2005} & theory \\
14 & \nt{DOPA} outputs go primarily to action-selection regions & Complete axon reconstruction of single nigrostriatal \nt{DOPA} neurons in the rat: the branches of one cell cover $0.45$--$5.7\,\%$ of the striatal volume ($2.7\,\%$ on average), with almost no branches outside the striatum. & \cite{Matsuda2009} & measured \\
15 & The cortex learns to predict its input, and the error is computed on the spot & Review: sensory cortex shows responses to mismatch between expected and actual input. That this is a general learning rule of the cortex has not been proven. & \cite{Keller2018} & hypothesis \\
16 & Dopamine behaves like the error~\eqref{eq:ctd}: burst, transfer to the cue, dip (Fig.~\ref{fig:ctdcases}) & Monkey \nt{DOPA} neurons: a burst to unexpected juice; after learning, a burst to the cue and no response to the promised juice; a dip when the juice is omitted. The continuous form~\eqref{eq:ctd} is theory. & \cite{Schultz1997, Doya2000} & measured \\
17 & A circuit that computes $\dd V/\dd t$ and subtracts the expectation & Mice, recording and optogenetics in the ventral tegmental area: \nt{DOPA} neurons subtract the expectation from the reward; neighboring \nt{GABA} neurons inhibit them when reward is expected, and exciting or inhibiting these \nt{GABA} neurons changes the error. How the derivative is computed has not been shown. & \cite{Eshel2015arithmetic} & hypothesis \\
18 & The \nt{DOPA} neuron is a pacemaker; dips are shallower than bursts & In a slice, \nt{DOPA} neurons fire spontaneously and very regularly, riding on a slow pacemaker potential. In the monkey, the rate follows a positive error quantitatively but conveys a negative error only weakly. & \cite{GraceOnn1989, Bayer2005} & measured \\
19 & Actor and critic (Fig.~\ref{fig:actorcritic}) & The algorithm comes from reinforcement-learning theory. In humans (fMRI), the ventral striatum behaves like a critic both with and without choice, while the dorsal striatum does so only when actions must be chosen. & \cite{SuttonBarto2018, ODoherty2004actor} & hypothesis \\
20 & The sign of $\delta$ in the rule is flipped in neurons with the second receptor family & Striatal slices: in neurons with D1 receptors, pairing produces strengthening that requires D1; in neurons with D2, weakening that requires D2. & \cite{Shen2008} & measured \\
\bottomrule
\end{longtable}}

\section{Chapter~\ref{sec:dofaalt}. Other theories of \nt{DOPA}}

Each of the extended pictures rests on its own direct measurements of dopamine (the spread of reversal points, responses to unpleasant and novel events, the slow ramp, responses to a flavor switch), but the formulas that explain them are theories, and between two of them the experiments so far contradict one another.

{\footnotesize
\setlength{\tabcolsep}{3pt}\renewcommand{\arraystretch}{1.15}
\begin{longtable}{>{\raggedright}p{0.5cm}>{\raggedright}p{4.0cm}>{\raggedright}p{5.6cm}>{\raggedright}p{2.3cm}>{\raggedright\arraybackslash}p{1.7cm}}
\toprule
No. & Claim & Experiment and what was measured & Source & Status \\
\midrule\endfirsthead
\toprule
No. & Claim & Experiment and what was measured & Source & Status \\
\midrule\endhead
1 & The asymmetric rule~\eqref{eq:asymmetric} converges to the point~\eqref{eq:expectile}; for $\kappa=1/2$ it is the mean, and for a drop with probability $1/2$, $V=\kappa$ (Fig.~\ref{fig:expectiles}) & The point~\eqref{eq:expectile} is an expectile, the solution of a least-squares problem with different weights for positive and negative deviations; for the chapter's example the derivation is in the chapter itself. & \cite{NeweyPowell1987}; derived in the chapter & theory \\
2 & Different \nt{DOPA} neurons have different reversal points, ordered by asymmetry (Proposition~\ref{prop:expectile}) & Mice, optogenetically tagged \nt{DOPA} neurons of the ventral tegmental area, rewards of different sizes: the point where a burst turns into a dip differs between neurons; it is higher in neurons with greater response asymmetry; the shape of the reward distribution can be reconstructed from the group's responses. That this is the main function of the diversity is a hypothesis. & \cite{Dabney2020} & measured \\
3 & Some \nt{DOPA} neurons respond to unpleasant events with a burst & Monkeys: some \nt{DOPA} neurons are excited by a reward cue and inhibited by a cue for an air puff to the eye, others are excited by both; the groups lie in different parts of the midbrain. & \cite{Matsumoto2009, BrombergMartin2010} & measured \\
4 & The magnitude of the error, or novelty, sets the learning rate & The cited papers contain no direct experiment in which a salience signal in \nt{DOPA} neurons changes the learning rate; the review proposes the role of an ``attention, this matters'' signal. & \cite{BrombergMartin2010} & hypothesis \\
5 & A separate wire for the danger axis & Mice: \nt{DOPA} neurons projecting to the tail of the striatum respond to novel and intense stimuli, not to reward; ablating them weakens avoidance of a novel object. & \cite{Menegas2018} & measured \\
6 & The optimal action time $\tau_a^*=\sqrt{C/\bar R}$~\eqref{eq:vigor} & Minimum of the sum $C/\tau_a+\bar R\tau_a$; derived in the chapter. In the original model, the speed of actions is set by the cost of time, equal to the average reward. & \cite{Niv2007}; derived in the chapter & theory \\
7 & The slow dopamine level carries $\bar R$ and sets the speed of actions (Proposition~\ref{prop:multiplex}) & Rats, microdialysis and voltammetry in the nucleus accumbens: the dopamine level tracks the reward rate and the speed of actions from minute to minute. In humans, L-DOPA strengthens the dependence of reaction time on the average reward rate. That the slow level equals precisely $\bar R$ has not been proven. & \cite{Hamid2016, Beierholm2013vigor} & hypothesis \\
8 & The slow level combines two sources: release changes at the outputs without a change in spike rate, but a slow ramp also occurs in the rate itself & Rats, one task: release in the nucleus accumbens follows the changing reward expectation, while the spike rate of identified \nt{DOPA} neurons in the ventral tegmental area does not. Mice in a virtual corridor (row~10): a smooth ramp on approaching reward is present in the spikes of identified \nt{DOPA} neurons as well. & \cite{Mohebi2019, Kim2020} & measured \\
9 & Dopamine rises smoothly while the animal approaches a distant reward & Rats in a maze, voltammetry in the striatum: the dopamine concentration builds up over seconds as the goal approaches; the slope depends on the distance and on the reward size. & \cite{Howe2013ramp, Hamid2016} & measured \\
10 & The ramp is the error $\delta$ as the estimate is refined~\eqref{eq:ramp}; a jump in the corridor gives a burst (Fig.~\ref{fig:ramp}) & The formula and the value $\delta=0.75\,V$ per second are derived in the chapter. Experiment: mice in a virtual corridor, instantaneous teleportation closer to the goal; somatic spikes, somatic and axonal calcium, and the concentration in the striatum respond like an error, not like the expectation $V$. Which of the two explanations holds at all outputs is debated. & \cite{Kim2020} & measured \\
11 & Looking back: dopamine reports a measure of causal link~\eqref{eq:retro} & Mice, dopamine release in the nucleus accumbens in experiments where this theory and the error~\eqref{eq:ctd} predict different things: in several experiments release followed the former. & \cite{Jeong2022} & hypothesis \\
12 & The dopamine response gradually shifts from reward to cue during learning & Mice, signal from \nt{DOPA} neuron outputs in the ventral striatum over the course of learning: the burst gradually moves from the moment of reward to the moment of the cue, as learning by~\eqref{eq:ctd} requires. & \cite{Amo2022} & measured \\
13 & \nt{DOPA} neurons respond to a change in the reward's flavor at the same value & Rats, recordings from the ventral tegmental area: replacing a juice with another of equal value evokes a response, even though the error~\eqref{eq:ctd} is zero. & \cite{Takahashi2017} & measured \\
14 & Artificial bursts teach an association between two neutral stimuli & Rats, sensory preconditioning of two stimuli without reward: optogenetic activation of \nt{DOPA} neurons at the transition from the first stimulus to the second creates the association, and inhibition prevents it from being learned. & \cite{Sharpe2017} & measured \\
15 & A vector of errors per feature~\eqref{eq:vectortd} & A theory of feature-prediction errors; there is no direct test of the vector form. & \cite{Gardner2018} & hypothesis \\
16 & Different \nt{DOPA} neurons in the same task carry different quantities & Mice in a virtual corridor, two-photon imaging of \nt{DOPA} neurons: some are related to reward, others to speed and turning, still others to cues and choice accuracy. & \cite{Engelhard2019} & measured \\
17 & A bundle instead of a wire (Proposition~\ref{prop:bundle}) & Summary of rows~2--16: each decomposition is measured separately; there is no experimental evidence for a single picture in which all of them are parts of one signal. & --- & hypothesis \\
\bottomrule
\end{longtable}}

\section{Chapter~\ref{sec:nora}. \nt{NORA}}

The organization of the \nt{NORA} system, its two modes, its link to the pupil (not only through \nt{NORA}), the rise in signal-to-background ratio and the inverted U in behavior are measured directly; the multiplicative gain $g$ is a model; that gain switches the choice mode and acts as an inverse temperature are conclusions of the chapter; the benefit of adjusting the gain is computed with the book's model; ``\nt{NORA} as the exploration knob'' and ``\nt{NORA} as an interrupt'' are hypotheses.

{\footnotesize
\setlength{\tabcolsep}{3pt}\renewcommand{\arraystretch}{1.15}
\begin{longtable}{>{\raggedright}p{0.5cm}>{\raggedright}p{4.0cm}>{\raggedright}p{5.6cm}>{\raggedright}p{2.3cm}>{\raggedright\arraybackslash}p{1.7cm}}
\toprule
No. & Claim & Experiment and what was measured & Source & Status \\
\midrule\endfirsthead
\toprule
No. & Claim & Experiment and what was measured & Source & Status \\
\midrule\endhead
1 & A human has a few tens of thousands of \nt{NORA} neurons, almost all in one group & Serial sections of the human brainstem stained for tyrosine hydroxylase: $53\,900$ cells in the locus coeruleus and another $6\,260$ in the adjacent subnucleus. & \cite{Baker1989LC} & measured \\
2 & The outputs spread through almost the whole brain, but parts of the group partly serve different regions & Viral-genetic labeling in mice: \nt{NORA} neurons of the locus coeruleus receive inputs from wide areas of the brain and send outputs to wide areas of the brain and spinal cord. In rats: a group projecting to the amygdala enhances fear learning, and a separate group projecting to the prefrontal cortex enhances its extinction. & \cite{Schwarz2015LC, Uematsu2017LC, Poe2020} & measured \\
3 & Tonic mode: steady spikes at rates from a fraction of a spike to a few per second, changing with wakefulness & Rats, recordings of locus coeruleus neurons across the sleep--wake cycle: the rate is highest in waking, lower in slow-wave sleep, and almost zero in REM sleep; rate changes precede state transitions. & \cite{AstonJonesBloom1981} & measured \\
4 & Phasic mode: a short burst to a task-relevant event, at roughly the moment of decision & Monkeys, a task of detecting a rare target: all locus coeruleus neurons respond with a burst only to the target, $\sim90\ms$ after it and $\sim200\ms$ before the hand response; bursts are weaker when performance is poor. In a two-choice task the burst is more tightly locked to the response than to the stimulus, and it occurs on both correct and error trials. & \cite{AstonJones1994vigilance, Clayton2004LC} & measured \\
5 & Pupil diameter is an imprecise test point for \nt{NORA}: it follows \nt{NORA}, \nt{ACET} and other regions & Monkeys: locus coeruleus spikes, down to single spikes of one cell, predict pupil dilation; a similar but later relation exists in the colliculi and the cingulate cortex. Mice: pupil fluctuations follow the activity of both noradrenergic and cholinergic outputs in the cortex. & \cite{Joshi2016, Reimer2016} & measured \\
6 & Noradrenaline suppresses background activity more than evoked activity; the multiplicative gain~\eqref{eq:gain} is a model that captures this effect & Awake monkeys, auditory cortex, iontophoresis of noradrenaline: background activity of neurons is suppressed more than the response to conspecific calls, so the signal-to-noise ratio rises. The multiplicative sigmoid~\eqref{eq:gain} is taken from a network model; a literal multiplication of the slope has not been measured. & \cite{Foote1975NE, ServanSchreiber1990} & measured; $g$ is a model \\
7 & Gain raises $d'$ only against noise after the amplifier~\eqref{eq:gainsnr} (Proposition~\ref{prop:gainnoise}) & Derived in the chapter; in the network model, gain improves signal detection. There is no experiment that separates noise before and after the amplifier and tests this distinction. & derived in the chapter; \cite{ServanSchreiber1990} & theory \\
8 & The boundary $g\beta=1$ switches the choice network from ``deliberating'' to ``decided''~\eqref{eq:wtagain} (Proposition~\ref{prop:gainwta}, Fig.~\ref{fig:pitchfork}) & Linear analysis of two mutually inhibiting groups; derived in the chapter. There are no direct measurements of this threshold in the brain. & derived in the chapter & theory \\
9 & The gain is the inverse temperature of choice~\eqref{eq:softmax} & With logistic noise after the amplifier, the choice probability is softmax with $T=\sigma/g$; derived in the chapter. & derived in the chapter & theory \\
10 & \nt{NORA} sets how much to explore: high tonic activity means small effective $g$ (exploration), a phasic burst at the moment of decision means large $g$ (decision) & Humans: the baseline pupil is wider before a random choice than before choosing the best known option; people with wider pupils are more inclined to explore. Rats: switching into a random-choice mode is controlled by \nt{NORA} input to the anterior cingulate cortex. That the burst raises the effective $g$ has not been measured directly. & \cite{JepmaNieuwenhuis2011, Tervo2014stochastic, AstonJones2005} & hypothesis \\
11 & Reward depends on $g$ as an inverted U; the best $g$ is smaller in a fast-changing world (Proposition~\ref{prop:explore}, Fig.~\ref{fig:invertedU}) & \texttt{models/nora\_explore.py}, $60$ runs of $4000$ trials. A change once per $1000$ trials: best $g=16$, fraction $0.976$; once per $20$: $g=8$, fraction $0.886$; at $g=0.5$, $0.77$. The rerun matches the chapter. & \texttt{models/\allowbreak nora\_\allowbreak explore.py} & model \\
12 & The inverted U in behavior: best performance at moderate tonic activity with sharp bursts & Monkeys, the same task as in row~4: during periods of good performance the tonic rate is moderate and bursts to the target are sharp; at a high tonic rate there are no bursts and more false responses. Changes in tonic and evoked activity follow fluctuations in performance. & \cite{Usher1999LC, AstonJones2005} & measured \\
13 & \nt{NORA} reports unexpected uncertainty, \nt{ACET} expected uncertainty & A theory of Bayesian inference with two kinds of uncertainty; the cited papers contain no direct experiment separating these signals by transmitter. & \cite{YuDayan2005} & hypothesis \\
14 & After an abrupt change people learn faster, and the pupil dilates & Humans, a prediction task with sudden changes: brief pupil changes track the estimated probability of a change and, together with the baseline pupil, predict how strongly new data change the estimate (the learning rate); pupil changes unrelated to light and task alter this rate. That the pupil here reflects \nt{NORA} specifically is an inference from the indirect data of row~5. & \cite{Nassar2012} & measured \\
15 & The phasic burst works as an interrupt: the network resets its state & There is no direct experiment in which a \nt{NORA} burst resets the network's state; only bursts to important events have been measured (row~4). & \cite{BouretSara2005} & hypothesis \\
16 & Surprise-driven adjustment of $g$ or of the learning rate is hardly better than the best constant settings & \texttt{models/nora\_explore.py}, a world with epochs of $1000$ trials (a change once per $1000$ and once per $20$). Best constant $g=8$: $0.925$; adjusting $g$: up to $0.927$; adjusting the learning rate: up to $0.926$; constant learning rate $0.4$: $0.928$. The rerun matches the chapter. & \texttt{models/\allowbreak nora\_\allowbreak explore.py} & model \\
\bottomrule
\end{longtable}}

\section{Chapter~\ref{sec:acet}. Adding \nt{ACET}}

The three actions of acetylcholine are measured in slices and in the living brain, the conflict between writing and reading is shown with the book's model, and the idea that \nt{ACET} serves as a write-enable line and reports expected uncertainty is a hypothesis with partial experimental support.

{\footnotesize
\setlength{\tabcolsep}{3pt}\renewcommand{\arraystretch}{1.15}
\begin{longtable}{>{\raggedright}p{0.5cm}>{\raggedright}p{4.0cm}>{\raggedright}p{5.6cm}>{\raggedright}p{2.3cm}>{\raggedright\arraybackslash}p{1.7cm}}
\toprule
No. & Claim & Experiment and what was measured & Source & Status \\
\midrule\endfirsthead
\toprule
No. & Claim & Experiment and what was measured & Source & Status \\
\midrule\endhead
1 & The brain's \nt{ACET} neurons are few, gathered in a few groups, and their outputs spread throughout the cortex: a broadcast channel & Antibody staining for the acetylcholine-synthesizing enzyme and retrograde labeling in the rat: the cortex, hippocampus, amygdala, olfactory bulb and thalamus receive acetylcholine mainly from six compact cell groups of the basal forebrain and brainstem & \cite{Mesulam1983} & measured \\
2 & The acetylcholine level in the cortex is high in waking and low in slow-wave sleep & Microdialysis in the cortex and hippocampus of freely moving cats with EEG recording: acetylcholine release is $\sim100\%$ higher in quiet waking and $\sim175\%$ higher in active waking than in slow-wave sleep; in REM sleep, cortical release is as in waking & \cite{Marrosu1995,Hasselmo1999} & measured \\
3 & The meaning of the signal is set by the receptor: acetylcholine has fast channel receptors and slow ones that trigger reactions in the cell (Proposition~\ref{prop:receptor}) & Review of measurements: the action of acetylcholine on excitability, transmission and plasticity depends on the release site, the receptor subtype (nicotinic are ion channels, muscarinic act through G proteins) and the recipient cell & \cite{Picciotto2012} & measured \\
4 & First action: weaken the internal connections while hardly touching external inputs (factor $1-a$ in~\eqref{eq:acetnet}) & Rat olfactory cortex slices: at $100$\,\textmu M acetylcholine agonists, potentials of internal connections (layer Ib) fall by more than $60\%$, those of the external input (layer Ia) by less than $15\%$, in the same cell; the suppression is presynaptic. Hippocampal slices (CA1): $89\%$ versus $40\%$ for the internal and input pathways & \cite{HasselmoBower1992,HasselmoSchnell1994} & measured \\
5 & Second action: strengthen the input (factor $\frac{1+a}{2}$) & Neocortical slices: nicotinic receptors strengthen synapses of the thalamic input and do not act on intracortical ones; muscarinic receptors weaken both types. Acetylcholine raises neuronal excitability (review) & \cite{Gil1997,Hasselmo2006} & measured \\
6 & Third action: make weight changes easier (rate $\eta a$) & Rat: electrical stimulation of the nucleus basalis (the source of cortical acetylcholine) paired with a tone produces, in an adult animal, a strong reorganization of the auditory cortex map toward the presented sound & \cite{KilgardMerzenich1998,Hasselmo2006} & measured \\
7 & The Hebb rule for sparse patterns in the second equation of~\eqref{eq:acetnet} gives an associative memory that completes a pattern from a part & Capacity calculation for a network with sparse patterns and a covariance rule: at a small fraction $p$ of active neurons, the network stores considerably more patterns than at $p=1/2$ & \cite{Tsodyks1988} & theory \\
8 & Writing at low $a$ corrupts memory: the network writes down what it recalled, not what it saw; at $a=0.1$ the corruption is no weaker than at $a=0$ & \texttt{models/\allowbreak acet\_memory.py}: $200$ neurons, five patterns of $20$, a new pattern shares $10$ neurons with one of them, $10$ runs. At $a=0$ the new pattern is not memorized, and the old one drags along $5.9$ foreign neurons out of $10$; at $a=0.1$ the new pattern is recalled ($0.97$), but the two merge: the new one drags along $7.3$ foreign neurons, the old one $7.9$; from $a=0.2$ on, fewer than one & derived in the chapter & model \\
9 & Proposition~\ref{prop:writeread}: there is no level $a$ at which writing and reading are equally good (Fig.~\ref{fig:acetsim}) & The same script, learning rate $\eta a$: a new pattern is memorized in full after one presentation for $a\ge0.9$, to $0.8$ at $a=0.75$, not at all for $a\le0.5$. Reading from half a pattern: $1.0$ for $a\le0.2$, $0.96$ at $a=0.3$, $0.04$ at $a=0.4$ & derived in the chapter & model \\
10 & In the brain a single broadcast wire, \nt{ACET}, switches between writing and reading & The idea comes from models built on slice measurements. The most direct experiment is in humans: the muscarinic blocker scopolamine impairs learning of new word pairs, most of all those overlapping with old ones, but not recall of pairs already learned. There is no direct measurement of two network modes & \cite{HasselmoAndersonBower1992,Atri2004} & hypothesis \\
11 & The filter~\eqref{eq:kalman} is optimal; the input weight and the learning rate are one quantity, $P/R$; in steady state $P=\sqrt{qR}$ and the memory is $\sqrt{R/q}$ (Proposition~\ref{prop:kalman}) & Requires no experiment: the optimality of the Kalman--Bucy filter is a classical result of estimation theory; the steady state is derived in the chapter & \cite{KalmanBucy1961} & theory \\
12 & \nt{ACET} tells the network a quantity like $P$: the expected uncertainty & Proposed in a probabilistic model of attention. There is no direct measurement showing that the acetylcholine level tracks the uncertainty of an estimate & \cite{YuDayan2005} & hypothesis \\
13 & Some \nt{ACET} neurons respond to reward and punishment within twenty milliseconds or so, more strongly the more unexpected the reinforcement & Mouse, optogenetically identified cholinergic neurons of the basal forebrain: response to reward and punishment after $18\pm3$\,ms; the size grows with the surprise of the reinforcement; responses are similar in neurons of two nuclei & \cite{Hangya2015} & measured \\
14 & \nt{NORA} notices an unexpected change in the world, and \nt{ACET} keeps the network in write mode & The division of labor is proposed in a model. Indirectly: in humans, pupil diameter, linked to \nt{NORA}, tracks changes and the learning rate. There is no direct test of the \nt{NORA}--\nt{ACET} pair & \cite{YuDayan2005,Nassar2012} & hypothesis \\
15 & In slow-wave sleep the network, in read mode, replays what was written during the day, and these replays rewrite it into long-term memory & Recordings of rat hippocampal ensembles: cells that worked together during the day fire together more often in subsequent slow-wave sleep, and in the same order (measured). For rewriting: suppressing hippocampal ripples after learning impairs spatial memory, and lengthening spontaneous ripples improves it; that the cause is a low \nt{ACET} level has not been proven & \cite{WilsonMcNaughton1994,Skaggs1996,Girardeau2009,FernandezRuiz2019,Hasselmo1999} & hypothesis \\
\bottomrule
\end{longtable}}

\section{Chapter~\ref{sec:synthesis}. The Whole Machine}

This summary chapter brings in no new experiments: each row rests on the chapter where the claim is worked out and on the same sources. The \nt{GLUT} bus and flow control through \nt{GABA} are firmly established, while the roles of the broadcast channels and their joint operation are mostly hypothesis.

{\footnotesize
\setlength{\tabcolsep}{3pt}\renewcommand{\arraystretch}{1.15}
\begin{longtable}{>{\raggedright}p{0.5cm}>{\raggedright}p{4.0cm}>{\raggedright}p{5.6cm}>{\raggedright}p{2.3cm}>{\raggedright\arraybackslash}p{1.7cm}}
\toprule
No. & Claim & Experiment and what was measured & Source & Status \\
\midrule\endfirsthead
\toprule
No. & Claim & Experiment and what was measured & Source & Status \\
\midrule\endhead
1 & There are only four kinds of control signal for $8.6\cdot10^{10}$ neurons~\eqref{eq:machine} & Counting cell nuclei in brain homogenate (isotropic fractionator): an adult man has $86.1\pm8.1$ billion neurons & \cite{Azevedo2009} & measured \\
2 & Proposition~\ref{prop:architecture}, the first two layers: data travel along the \nt{GLUT} address bus, the sign of a connection is set by the sender, and \nt{GABA} neurons steer and normalize the flow (Chapters~\ref{sec:neurontypes}, \ref{sec:gaba}) & One main transmitter per neuron (review of measurements); feedforward inhibition narrows the window in which a pyramidal neuron sums its inputs; division by the total activity has been measured in many areas & \cite{Strata1999,Pouille2001,Carandini2012} & measured \\
3 & The elements are unreliable: a synapse loses most spikes (Table~\ref{tab:computer}, Chapter~\ref{sec:code}) & Hippocampal slices, minimal stimulation: more than half of the spikes evoke no response in CA1; threshold failures and axonal conduction failures are ruled out, and the cause is probabilistic transmitter release & \cite{AllenStevens1994} & measured \\
4 & The brain runs on $\sim20$ watts, on average about one spike per second per neuron (Chapter~\ref{sec:code}); engineers have not yet built such a machine & The estimate~\eqref{eq:energy} from measured costs: about $10^9$ ATP molecules per spike, including synaptic work (rodent cortex); a detailed estimate for the cortex gives an even lower rate. The state of the art is from a review & \cite{Attwell2001,Lennie2003,MehonicKenyon2022} & theory \\
5 & Learning at most synapses is the product of a local trace and one shared number $\delta$ (the second equation of~\eqref{eq:machine}, Chapter~\ref{sec:dofa}) & Monkey \nt{DOPA} neurons respond to the reward-prediction error; in striatal slices dopamine enlarges a spine only if it arrives $0.3$--$2$\,s after the glutamate input, so the trace has been measured & \cite{Schultz1997,Yagishita2014} & measured \\
6 & A separate error wire to each output exists only in the circuit for learning movements (Chapter~\ref{sec:motorlearning}) & Cerebellum: coincidence of the climbing-fiber signal with an input weakens that input; in the monkey, complex spikes of Purkinje cells carry the direction of the movement error and cause a correction & \cite{Ito2001,Herzfeld2018} & measured \\
7 & Each broadcast channel is a bundle of several lines (Proposition~\ref{prop:bundle}) & For \nt{DOPA}: in one task different neurons carry reward, movement, and stimulus features; the fast error and the slow dopamine level are distinct. For \nt{SERO}, \nt{NORA}, and \nt{ACET} such a breakdown is less well studied & \cite{Engelhard2019,Hamid2016,Mohebi2019} & measured \\
8 & \nt{SERO}: a level regulator and, by hypothesis, the horizon $\tau_\gamma$ (Chapter~\ref{sec:serocomputer}) & Self-inhibition: agonists of their own 5-HT$_{1A}$ receptors inhibit the serotonin neurons of the raphe (measured). The horizon was proposed in theory; the strongest experiment is that optogenetic activation of these neurons in the mouse lengthens waiting for a delayed reward & \cite{Sprouse1987,Doya2002,Miyazaki2014} & hypothesis \\
9 & \nt{NORA}: the gain $g$ chooses between searching and deciding (Chapter~\ref{sec:nora}) & Rise in signal-to-noise ratio: in the awake monkey, noradrenaline in the auditory cortex suppresses background activity more strongly than the response to conspecific calls (measured); multiplicative gain is a model; the role in choosing between search and decision is theory & \cite{Foote1975NE,ServanSchreiber1990,AstonJones2005} & hypothesis \\
10 & \nt{ACET}: switches between writing and reading (Chapter~\ref{sec:acet}) & Three actions (weakening internal connections, amplifying the input, facilitating plasticity) are measured; mode switching is indirectly supported by an experiment with scopolamine in humans & \cite{HasselmoBower1992,Gil1997,Atri2004} & hypothesis \\
11 & Formula~\eqref{eq:machine} describes the whole machine & A summary of the models of Chapters~\ref{sec:glutcomputer}--\ref{sec:acet}, each part resting on its own chapter; as a whole it has not been tested experimentally & derived in the chapter & hypothesis \\
12 & The knobs work in concert without central control: error, surprise, writing, return (Fig.~\ref{fig:timeline}) & No experiment: the sketch is qualitative. Individual links are the theory of the division of labor between \nt{NORA} and \nt{ACET} and the link between pupil size and learning rate in humans & \cite{YuDayan2005,Nassar2012} & hypothesis \\
13 & Learning from one shared signal is slower the more neurons affect the result (Proposition~\ref{prop:broadcastcost}) & Calculation of the learning curves of a linear network: learning by output perturbation is slower than gradient descent by a factor equal to the number of outputs & \cite{Werfel2005} & theory \\
14 & How the cortex tunes multilayer circuits is unknown; the candidates are bursts of spikes, computation in the neuron's branches, and self-supervised prediction & Bursts as carriers of error: model; only a $5$--$8$-layer network could reproduce the response of a detailed model of a cortical neuron: model; calcium spikes in the branches of human cortical neurons that give exclusive OR: measured in slices; self-supervised learning: a review of the evidence & \cite{Payeur2021,Beniaguev2021,Gidon2020,Keller2018} & hypothesis \\
\bottomrule
\end{longtable}}

\section{Chapter~\ref{sec:sensors}. The Machine's Inputs}

The structure of the sensors has been measured directly, by recording currents from single cells and vibrations of the cochlea and by counting retinal cells. The sensitivity limit and the shot-noise formula follow from physics, while the claim that the sensors' codes are tuned for the greatest information is a hypothesis with strong support.

{\footnotesize
\setlength{\tabcolsep}{3pt}\renewcommand{\arraystretch}{1.15}
\begin{longtable}{>{\raggedright}p{0.5cm}>{\raggedright}p{4.0cm}>{\raggedright}p{5.6cm}>{\raggedright}p{2.3cm}>{\raggedright\arraybackslash}p{1.7cm}}
\toprule
No. & Claim & Experiment and what was measured & Source & Status \\
\midrule\endfirsthead
\toprule
No. & Claim & Experiment and what was measured & Source & Status \\
\midrule\endhead
1 & A sensor is a transducer: the receptor gives a current, and the output of the sensory cells of the eye and ear is glutamate onto the \nt{GLUT} bus; the first cells produce no spikes (Fig.~\ref{fig:transducer}) & Turtle rods and cones release an excitatory amino acid, detected by glutamate channels on an excised membrane patch. Inner hair cells of the rat cochlea: currents in the nerve-fiber endings flow through glutamate AMPA receptors, and the release rate grows with smooth depolarization of the cell up to $150$\,Hz & \cite{CopenhagenJahr1989,GlowatzkiFuchs2002} & measured \\
2 & A light sensor in the dark responds to a single photon with a current of about a picoampere & Suction electrode on the outer segment of a toad rod: responses to dim flashes are quantized, an event of $\approx1$\,pA ($5\%$ of the dark current) with a spread of $0.2$\,pA, efficiency $0.5$. Humans report the arrival of a single photon more often than chance & \cite{Baylor1979,Tinsley2016} & measured \\
3 & A sound sensor detects a displacement of less than a nanometer & M\"ossbauer technique, basilar membrane of the guinea-pig cochlea at the $16$--$19$\,kHz place: at the auditory-nerve response threshold the membrane velocity is about $0.04$\,mm/s, i.e., a displacement of $\approx0.35$\,nm (our conversion). The membrane was measured, not the sensor's bundle & \cite{Sellick1982,Hudspeth1989} & measured \\
4 & Precision in the dark is limited by photon shot noise~\eqref{eq:photonnoise}; in dim light integration is longer & The formula follows from Poisson statistics. In a rod the noise in dim light matches the sum of independent quantal events; against a bright background the response to a flash is smaller and shorter. The ``tenths of a second'' figure is not checked separately here & \cite{Baylor1979} & theory \\
5 & The input range is ten orders of magnitude for light and twelve for sound, while the spike rate spans less than three & For sound: the basilar-membrane response at the characteristic frequency grows with a slope as low as $0.2$\,dB/dB over $40$--$80$\,dB, and compression is seen even above $100$\,dB. The order-of-magnitude figures themselves are standard estimates, unsourced in the chapter & \cite{Ruggero1997} & measured \\
6 & The sensor response is~\eqref{eq:nakarushton}, and $\sigma$ follows the mean brightness, shifting the curve along the logarithmic axis (Fig.~\ref{fig:adaptation}) & The formula was first fitted to S-potentials of the fish retina. Turtle cones: responses obey $V=I V_m/(I+\sigma)$ and span $\sim3.5$ orders of magnitude of brightness; after $2$--$3$ minutes of background the curve shifts horizontally, keeping its width. The link to division by total activity is from a review & \cite{NakaRushton1966,NormannPerlman1979,Carandini2012} & measured \\
7 & Sensors transmit changes, and their codes are tuned for the greatest information per spike (Proposition~\ref{prop:sensors}) & The principle was proposed theoretically. The strongest evidence: in the fly, the ``contrast $\to$ response'' curve of first-layer neurons matches the cumulative distribution of contrasts in natural scenes, which yields the greatest information & \cite{Barlow1961,Laughlin1981} & hypothesis \\
8 & On and off sensors form a two-wire code; the event camera reproduces it in silicon & A review of experiments: the on and off channels of the retina separate already at the first synapse and can be switched off separately. Camera: $128\times128$ pixels without frames, $120$\,dB range, $15$\,\textmu s latency & \cite{Schiller1992,Lichtsteiner2008,Gallego2022} & measured \\
9 & The eye has $\sim10^8$ light sensors and $\sim10^6$ output neurons: compression by $\sim100$ & Counts on whole-mount human retinas: on average $92$ million rods and $4.6$ million cones; $0.7$ to $1.5$ million output (ganglion) cells in different eyes & \cite{Curcio1990,CurcioAllen1990} & measured \\
10 & The mouse eye has more than thirty output channels & Mouse: two-photon calcium imaging of more than $11\,000$ output cells and clustering of the responses gives clearly more than $30$ functional types. The number for humans has not been measured & \cite{Baden2016} & measured \\
11 & The guinea-pig eye transmits $\sim875\,000$ bits/s; scaling to the human eye gives $\sim10^7$ bits/s & Guinea pig, multielectrode array, motion in natural scenes: $6$--$13$ bits/s per cell, $\sim875\,000$ bits/s for $\sim10^5$ output cells. For humans ($\sim10^6$ cells), $\sim10^7$ bits/s is the authors' extrapolation & \cite{Koch2006} & measured \\
12 & The ear is a bank of bandpass filters with a nearly logarithmic frequency map (Fig.~\ref{fig:filterbank}) & The place of maximum basilar-membrane vibration depends on frequency; the frequency--place function is nearly exponential and a single form describes the cochleas of humans and living animals & \cite{Greenwood1990,Robles2001} & measured \\
13 & The resonators are active: some sensors amplify weak vibrations thousands of times in amplitude ($66$--$76$\,dB), and the gain falls as loudness grows & Laser vibrometry at the base of the chinchilla cochlea: for weak tones at the characteristic frequency the gain is $66$--$76$\,dB relative to the stapes; death lowers sensitivity by $60$--$81$\,dB (a factor of $10^3$--$10^4$ in amplitude) and removes the compression. The source of the force is the outer hair cells & \cite{Ruggero1997,Robles2001} & measured \\
14 & The ear's amplifier sits at the edge of self-oscillation & Calculation: near a Hopf bifurcation point, range compression, sharp tuning, and combination tones, all the known nonlinearities of hearing, are unavoidable. That the cochlea is tuned in exactly this way has not been proved directly & \cite{Eguiluz2000} & hypothesis \\
15 & At low frequencies spikes are phase-locked to the sound (in the guinea pig, locking weakens from $\sim600$\,Hz and is noticeable up to $\sim3.5$\,kHz), and direction is found from them; at high frequencies only the place code remains & Guinea-pig auditory nerve: phase locking begins to fall at about $600$\,Hz and disappears above $3.5$\,kHz; in the cat and the monkey the limit is about an octave higher. The interaural time difference is from a review & \cite{PalmerRussell1986,Grothe2010} & measured \\
16 & The other sensors (Table~\ref{tab:sensors}): smell has $\sim400$ receptor types, one per neuron, and a combinatorial code; taste works partly through ATP and serotonin; touch has fast- and slow-adapting sensors; warm and cold are separate & Mouse: one receptor recognizes many odors, one odor activates many receptors, and different odors activate different sets. Without ATP receptors the taste nerve does not respond; taste buds release serotonin. Single-fiber recordings from the skin of the human hand: four types by adaptation speed. The cold channel is distinct from the heat channels & \cite{Mombaerts2004,Malnic1999,Finger2005,Huang2005,JohanssonVallbo1979,McKemy2002} & measured \\
\bottomrule
\end{longtable}}

\section{Chapter~\ref{sec:recognition}. Recognition}

The speed of recognition, the structure of the first stages, the linear readout of the response, and learning from the continuity of time are measured in humans and monkeys; the reduction of the whole chain to two operations and learning from video are tested on the book's models, and the explanations of illusions range from well founded to disputed.

{\footnotesize
\setlength{\tabcolsep}{3pt}\renewcommand{\arraystretch}{1.15}
\begin{longtable}{>{\raggedright}p{0.5cm}>{\raggedright}p{4.0cm}>{\raggedright}p{5.6cm}>{\raggedright}p{2.3cm}>{\raggedright\arraybackslash}p{1.7cm}}
\toprule
No. & Claim & Experiment and what was measured & Source & Status \\
\midrule\endfirsthead
\toprule
No. & Claim & Experiment and what was measured & Source & Status \\
\midrule\endhead
1 & Pixel-by-pixel template matching under shift does no better than a coin; a pair of stages~\eqref{eq:scstage} recognizes the shape at any position & \texttt{models/\allowbreak recognition\_models.py}, a ``retina'' of $24$ points, $4000$ trials: template, $0.49$, $0.51$, $0.52$ at flipped-point fractions $0$, $0.1$, $0.2$; the $S$ and $C$ pair, $1.00$, $0.86$, $0.70$ & derived in the chapter & model \\
2 & An animal in a picture is recognized in $\sim150$\,ms, a single pass through about ten stages of $10$--$20$\,ms each & Humans, ``animal or not'' task on photographs shown for $20$\,ms: evoked potentials for the two answers diverge after $\sim150$\,ms. Monkeys: the latency of the first response grows from stage to stage (V1 earliest, then V2 and V4). The number of stages and ``zero or one spike per stage'' are inferred from these numbers & \cite{Thorpe1996,Schmolesky1998} & measured \\
3 & Stage $S$ is AND over parts, stage $C$ is the maximum over positions (Proposition~\ref{prop:conveyor}) & Cat, primary visual cortex: simple cells respond to an edge of a given orientation at a given place, complex cells to the same orientation over a larger region. Intracellular recording from complex cells: for many cells the response to two bars is close to the larger of the two responses, on average closest to a max operation. The maximum via mutual inhibition is Proposition~\ref{prop:wta}; division by the total activity, a review & \cite{HubelWiesel1962,Lampl2004,Carandini2012} & measured \\
4 & Further up the chain the combinations are larger and more complex, and the response depends less and less on position & Monkeys, stimuli simplified down to the ``critical feature'': cells of V4 and posterior inferotemporal cortex respond to complex combinations of shape, color, and texture with small fields; in anterior inferotemporal cortex the fields are larger. Alternating $S$ and $C$ in a model reproduces this picture & \cite{KobatakeTanaka1994,Riesenhuber1999} & measured \\
5 & Convolutional networks repeat this alternation and predict the responses of the upper stages best; the object is read out linearly from the rates of a group of neurons & A network trained for recognition, with no fitting to the brain: the top layer predicts the responses of monkey inferotemporal neurons, the middle layers the responses of V4. A linear classifier on the activity of $\sim100$ randomly chosen cells, in windows from $12.5$\,ms, recognizes object and category across positions and sizes & \cite{Fukushima1980,Yamins2014,Hung2005} & measured \\
6 & The Hebb rule with a trace~\eqref{eq:traceinvariance} learns invariance from unlabeled video if the trace is at least $\sim20$\,ms & The idea of slow features and the trace rule is theory. \texttt{models/\allowbreak recognition\_models.py}, two $C$ neurons, $16$ inputs, $10$ runs, a $0.3$\,s pause between objects. At $\tau_{\text{tr}}=5$\,ms the match with the shape is $0.52$ on average, at $10$\,ms $0.56$; at $20$, $50$, and $200$\,ms, $1.00$ in all runs (at $30$\,ms one run in ten errs) & \cite{Foldiak1991,WiskottSejnowski2002} & model \\
7 & In the brain, invariance is learned from the continuity of time & Monkeys: the object was covertly swapped during an eye jump to a particular place in the field; the position invariance of inferotemporal neurons changed in line with the swap, significantly after just an hour. That this is the main rule of visual learning is a hypothesis & \cite{LiDiCarlo2008} & measured \\
8 & Motion: direction detectors, then the same AND/OR pair over large regions & Direction detectors in the rabbit retina; in monkey area MT all $110$ recorded neurons are direction selective, and the response to motion is stronger than in V1 & \cite{Barlow1965,Albright1984} & measured \\
9 & The upper stages send a prediction down, and the error goes up & The model explains extra-classical receptive field effects in primary visual cortex; individual observations agree (review), but it is not proven as the general design of the pipeline & \cite{RaoBallard1999,Keller2018} & hypothesis \\
10 & Hard pictures require feedback and several rounds & Monkeys: for ``hard'' pictures the decision in inferotemporal cortex forms $\sim30$\,ms later; these late responses are poorly predicted by a feedforward network and better by networks with feedback or by deeper ones & \cite{Kar2019} & measured \\
11 & An invisible pixel forgery fools a network; human vision is far more robust, but not immune & In networks: a small, specially chosen perturbation changes the answer; the ``panda $\to$ gibbon'' example is from the second paper. In humans, with brief presentation, forgeries that transfer well between networks shift the choice & \cite{Szegedy2014,Goodfellow2015,Elsayed2018} & measured \\
12 & Illusions of expectation: an unreliable input yields to expectation; faint things move ``more slowly,'' the dress is seen differently (Section~\ref{sec:illusions}) & Lower-contrast gratings appear to move more slowly. A survey of $1401$ people: the dress is seen in three stable variants, and a lighting cue switches the color; in $\sim13\,000$ observers the assumption about lighting decides. A Bayesian model with an expectation of slow motion explains a range of motion illusions & \cite{Thompson1982,LaferSousa2015,Wallisch2017,Weiss2002,Purves2011} & hypothesis \\
13 & Gain control: the waterfall, tilt, and contrast; the Hermann grid is not inhibition from neighbors & Direction detectors in the rabbit retina lower their rate under prolonged motion and fall below baseline after it stops. The tilt illusion is reproduced by a model with division by the activity of neighbors. Grid modifications that do not change the response of retinal output neurons abolish the spots, so the neighbor-inhibition explanation is rejected & \cite{BarlowHill1963,Schwartz2009,SchillerCarvey2005} & measured \\
14 & Delay compensation: a moving object appears ahead of a flash & The illusion is measured. Psychophysics contradicts both forward extrapolation of motion and a latency difference; a retrospective explanation has been proposed, based on events $\sim80$\,ms after the flash. The debate is not settled & \cite{Nijhawan1994,EaglemanSejnowski2000} & hypothesis \\
15 & Still ``snakes'' rotate: the latency difference~\eqref{eq:latency} is read by direction detectors & The illusion works without color too; pairs of neighboring elements produce it on their own; direction-selective neurons in monkey cortex give a directional response to these still pairs. In humans, rotation is triggered by microsaccades and blinks & \cite{Conway2005,OteroMillan2012} & measured \\
16 & Ambiguous pictures: mutual inhibition, fatigue, and noise; switches are caused mainly by noise & A model of competing neuron groups: switches in it are caused by \emph{noise}, and without noise there are none; it explains why the switching rate grows with stimulus strength. Fatigue plays a smaller role here & \cite{MorenoBote2007} & hypothesis \\
17 & Some illusions follow from the task the machine learns & A network trained to predict the next video frame ``sees'' rotation of the still rings and does not see it in control pictures; networks for image denoising and sharpening fall for color and brightness illusions, as humans do & \cite{Watanabe2018,GomezVilla2020} & measured \\
\bottomrule
\end{longtable}}

\section{Chapter~\ref{sec:muscles}. Output to muscles}

The design of the output (motor units, the synaptic safety factor, recruitment by size, the muscle pair, the stretch loop, and rhythm generators) is measured directly; the stability condition for a loop with delay is a theorem; the internal model of the body is a well-founded hypothesis; the numbers in the figures are computed with the book's model.

{\footnotesize
\setlength{\tabcolsep}{3pt}\renewcommand{\arraystretch}{1.15}
\begin{longtable}{>{\raggedright}p{0.5cm}>{\raggedright}p{4.0cm}>{\raggedright}p{5.6cm}>{\raggedright}p{2.3cm}>{\raggedright\arraybackslash}p{1.7cm}}
\toprule
No. & Claim & Experiment and what was measured & Source & Status \\
\midrule\endfirsthead
\toprule
No. & Claim & Experiment and what was measured & Source & Status \\
\midrule\endhead
1 & Motor unit: one \nt{ACET} neuron controls a group of fibers, from several hundred to a couple of thousand; in muscles for fine adjustment the units are smaller & Human muscles, counts of motor axons and muscle fibers: on average about $340$ fibers per neuron in an interosseous muscle of the hand, about $410$ in the brachioradialis, about $1930$ in the medial gastrocnemius. The chapter gives no exact number for the smallest units. & \cite{Feinstein1955mu} & measured \\
2 & The synapse on a muscle releases several times more transmitter than is needed to make the fiber fire & Review of measurements at neuromuscular synapses: the safety factor (ratio of released transmitter to threshold) in adult mammals is usually $3$--$5$; in humans the margin rests more on the folds of the receiving membrane than on the size of the release. & \cite{WoodSlater2001} & measured \\
3 & The twitch~\eqref{eq:twitch} rises over a time $T_c$ of order $50\ms$ & Single motor units of the human hand during voluntary effort, force averaged over the unit's spikes: rise time $30$--$100\ms$, under $70\ms$ for $80\,\%$ of units. The function $(t/T_c)e^{1-t/T_c}$ is an approximation from a motor-unit pool model. & \cite{MilnerBrown1973rec, Fuglevand1993} & measured \\
4 & Sum of twitches~\eqref{eq:forcesum}: at $5$, $15$, and $40$ spikes/s the force is $0.2$--$1.1F_1$, $1.8$--$2.2F_1$, and about $5.4F_1$ (Fig.~\ref{fig:tetanus}) & Computed with \texttt{models/\allowbreak muscle\_\allowbreak models.py}, $T_c=50\ms$: $0.21$--$1.10$, $1.76$--$2.19$, and $5.28$--$5.49\,F_1$ over the last $0.3\,\text{s}$. Linear summation is a simplification: the model has no saturation of a real muscle. & \texttt{models/\allowbreak muscle\_\allowbreak models.py} & model \\
5 & Units are recruited by size; the weakest are a hundred times weaker than the strongest & Cat ventral roots: with increasing reflex input, the neurons with small spikes in the root, that is, the small ones, fire first. Human hand interosseous muscle: unit twitch forces range from $0.1$ to $10$~g and grow almost linearly with the force at which the unit is recruited. & \cite{Henneman1957, MilnerBrown1973rec} & measured \\
6 & The force step is proportional to the force, $\Delta F/F\approx q-1$~\eqref{eq:sizeprinciple}: floating point & Derived in the chapter from the geometric series of forces. Consistent with experiment: at large efforts, far fewer new units are recruited for the same force increment. & derived in the chapter; \cite{MilnerBrown1973rec} & theory \\
7 & The scatter of force grows in proportion to the force itself & Voluntary isometric effort of a thumb muscle: the standard deviation of force grows linearly with mean force; for the same effort evoked by electrical stimulation, the scatter is constant. Pool model: the linear growth comes from recruiting units in order of force. & \cite{Jones2002sdn} & measured \\
8 & The smoothness of movements is a consequence of signal-dependent noise & Model: the trajectory that minimizes end-point scatter under such noise reproduces measured eye and arm trajectories. There is no direct experiment that separates this cause from others. & \cite{HarrisWolpert1998} & hypothesis \\
9 & The difference of the forces of a muscle pair sets the torque, the sum sets the stiffness~\eqref{eq:antagonists} & Theory of impedance control by co-contraction. Human arm: the stiffness of each joint, measured with small perturbations, is roughly proportional to the torque at it; different co-contraction ratios change the shape and orientation of the hand stiffness ellipse. & \cite{Hogan1984, GomiOsu1998stiff} & measured \\
10 & The stretch sensor excites its own \nt{ACET} neuron and, through an inhibitory intermediary, switches off the opponent (Fig.~\ref{fig:antagonists}) & Cat spinal cord: a single intermediary excited by stretch-sensor fibers produces monosynaptic inhibitory potentials in the \nt{ACET} neurons of the opposing muscle, synaptic delay $0.28$--$0.42\ms$. The intermediary's transmitter (glycine) was not identified in this experiment. & \cite{Jankowska1972Ia} & measured \\
11 & Loop delay: spinal $25$--$30\ms$, through the cortex one and a half to three times longer, through the eye $100$--$200\ms$ & Human arm, a perturbation of the joint: short-latency muscle response at $20$--$45\ms$, long-latency at $45$--$105\ms$, voluntary at $120$--$180\ms$. A shift in the seen position of the finger during a movement: correction after $160\ms$ on average. & \cite{Pruszynski2008rapid, SaundersKnill2003vis} & measured \\
12 & $\dd x/\dd t=-Gx(t-d)$ is stable for $Gd<\pi/2$~\eqref{eq:delaystable}; at the boundary the frequency is $1/(4d)$ & Theorem on the roots of a delay equation. Check with \texttt{models/\allowbreak muscle\_\allowbreak models.py}, $d=30\ms$: by $3\,\text{s}$ the amplitude is $3\cdot10^{-6}$ at $G=40$, $0.13$ at $G=50$, $38$ at $G=55$ per second (boundary $52$). & \cite{Hayes1950} & theory \\
13 & Physiological tremor at $8$--$12$~Hz; the stretch loop is one of its sources & Review of measurements: fine tremor around $10$~Hz is present in healthy people; its origin is mixed (central oscillations, properties of unit discharges, mechanical resonance, and resonance of the reflex loop), and different sources dominate under different conditions. & \cite{McAuleyMarsden2000} & hypothesis \\
14 & The brain predicts the result of a command from an internal model of the body & A person holds a load in a pinch grip and moves the arm: under inertial, viscous, and elastic loads, grip force changes simultaneously with the load, not with the loop delay, so the load is predicted. A review collects such experiments; the model itself has not been observed directly in neurons. & \cite{FlanaganWing1997grip, WolpertGhahramani2000} & hypothesis \\
15 & The rhythm of walking, breathing, and chewing is generated by local networks under constant input & Review of experiments: an isolated nervous system (spinal cord, ganglia) produces a rhythmic motor pattern without rhythmic input and without feedback from muscles. & \cite{MarderBucher2001} & measured \\
16 & Mutual inhibition with fatigue~\eqref{eq:halfcentre} gives a rhythm with period $0.84\,\text{s}\approx2\tau_a$ (Fig.~\ref{fig:halfcentre}) & Theorem: a network with mutual inhibition and adaptation that has no stable state under constant input must oscillate. Computed with \texttt{models/\allowbreak muscle\_\allowbreak models.py} ($I=1$, $\beta=2$, $b=2$, $\tau=20\ms$, $\tau_a=0.4\,\text{s}$): period $0.84\,\text{s}$. That real generators are built exactly this way has not been shown. & \cite{Matsuoka1985osc} & model \\
\bottomrule
\end{longtable}}

\section{Chapter~\ref{sec:pain}. Pain}

Pain sensors, the two wires, the gate in the spinal cord, the descending volume knob, sensitization, and the capture of attention have been measured directly; the gate and sensitization formulas are simplified models of the book; the rule by which the machine chooses the volume of the alarm is a hypothesis.

{\footnotesize
\setlength{\tabcolsep}{3pt}\renewcommand{\arraystretch}{1.15}
\begin{longtable}{>{\raggedright}p{0.5cm}>{\raggedright}p{4.0cm}>{\raggedright}p{5.6cm}>{\raggedright}p{2.3cm}>{\raggedright\arraybackslash}p{1.7cm}}
\toprule
No. & Claim & Experiment and what was measured & Source & Status \\
\midrule\endfirsthead
\toprule
No. & Claim & Experiment and what was measured & Source & Status \\
\midrule\endhead
1 & High threshold: heat thresholds of different pain sensors range from $41^\circ$ to $46$--$48^\circ$ & Single-fiber recordings. Monkey skin: median threshold of the slow (C) sensors $41^\circ$, of fast type~II sensors $46^\circ$, of type~I above $53^\circ$. Human skin: C sensors that also respond to pressure $40.7^\circ$, those that do not respond to pressure $48^\circ$. & \cite{Treede1995heat, Weidner1999cfib} & measured \\
2 & Pain sensors adapt far less than touch sensors and become more sensitive after mild injury; a weakening after strong heating is measured in one type & A mild skin burn ($50^\circ$, $100\,\text{s}$): after $5$--$10$~min the pain threshold in humans and the threshold of C sensors in the monkey are lowered by $2$--$6^\circ$; other skin sensors show no such shift. In fast type~II sensors the response after strong heating is suppressed. The comparison of adaptation with touch sensors is a general estimate, not backed here by a separate experiment. & \cite{LaMotte1982hyper, Treede1995heat} & measured \\
3 & Two wires: fast $5$--$30$~m/s, slow $0.5$--$2$~m/s; arrival time $t=L/v$~\eqref{eq:paintime} & Conduction velocity: fast pain sensors of the monkey average $15$ and $25$~m/s (two types); human C sensors $0.8$--$1.0$~m/s. For $L=1$~m this gives tens of milliseconds and about a second. & \cite{Treede1995heat, Weidner1999cfib} & measured \\
4 & Pain arrives twice: first fast and precise, then dull and lasting & Reaction time to heating of the human forearm: from $1100$ to $400\ms$ as the temperature rises; strong heating of hairy skin gives double pain, of the palm does not. First pain can be carried only by fibers faster than $6$~m/s; by latency it corresponds to fast type~II sensors, which are absent from the palm. The division of roles (``where'' and ``how bad'') is an interpretation. & \cite{CampbellLaMotte1983first, Treede1995heat} & measured \\
5 & The gate: an output neuron of the spinal cord receives pain and touch, and an inhibitory interneuron is excited by touch (Fig.~\ref{fig:gate}) & The gate scheme was proposed as a theory. Mice, genetic labeling and ablation of neuron groups: excitatory somatostatin neurons are needed for mechanical pain; inhibitory dynorphin neurons keep input from touch sensors from reaching them, and without these neurons a light touch causes pain. & \cite{MelzackWall1965, Duan2014} & measured \\
6 & Touch dampens pain & Humans, laser pain pulses on the arm: simultaneous touch lowers the pain rating and the ability to discriminate pulse energy; the effect weakens with the distance between the touch point and the pain point within one skin segment. & \cite{Mancini2014touch} & measured \\
7 & Assumption of gate theory: strong pain itself switches off the inhibitory interneuron, and the gate swings open & Marked in the chapter as an assumption of the original theory. The mouse work describes how the gate keeps touch out of the pain channel; switching off of the interneuron by pain is not shown in it. & \cite{MelzackWall1965} & hypothesis \\
8 & The gate~\eqref{eq:gate}: threshold $0.18$ without touch, $0.86$ with touch, $0.11$ at $g=2$; at $\nu=1$ touch lowers the output from $1$ to $0.25$, at $\nu=2$ it has no effect; touch alone does not get through (Fig.~\ref{fig:gatecurves}) & A calculation with \texttt{models/\allowbreak pain\_\allowbreak gate.py} using the weights of the text reproduces all these numbers. & \texttt{models/\allowbreak pain\_\allowbreak gate.py} & model \\
9 & Descending pathways from the brainstem change the gain of the gate in both directions & Rat, medulla, recordings during tail heating: some neurons fire before withdrawal (``on''), others fall silent (``off''); weak stimulation of this area suppresses the reflex. Review: the same circuits both facilitate and inhibit pain transmission, and opioids act through them. & \cite{Fields1983rvm, Fields2004} & measured \\
10 & Expectation of relief dampens pain through the opioid channel & People with acute pain: in those whose pain was reduced by the expectation of relief, blocking the opioid receptors brought the pain back to the level of the others; in the others the blockade did not change the pain. & \cite{Levine1978} & measured \\
11 & The brain chooses the volume of the alarm by the benefit of interrupting & There is no experiment in which the rule for choosing the volume has been measured. & --- & hypothesis \\
12 & Sensitization: after injury the gate output grows and the threshold falls, partly because of the spinal cord; it persists after the damaging stimulus has stopped & Rat: after skin injury the threshold of the flexion reflex falls and the response grows; electrophysiological analysis shows that part of the growth arises in the spinal cord itself. A damaging stimulus causes long-lasting changes in sensitivity that outlast the stimulus. The chapter gives no exact decay time. & \cite{Woolf1983} & measured \\
13 & Sensitization~\eqref{eq:sensitization} is positive feedback; with a strong loop the alarm can latch after healing & Follows from the model~\eqref{eq:sensitization} (an exercise at the end of the chapter). There is no experiment showing that persistent pain after healing is a latched loop of exactly this kind. & derived in the chapter & hypothesis \\
14 & Pain is a negative reward, and learning from it follows the temporal-difference error & Scanner, humans, chains of pain predictors: activity of the ventral striatum and the anterior insula matches temporal-difference learning signals. Monkey: some \nt{DOPA} neurons are excited by an aversive air puff and its predictor, others are inhibited. & \cite{Seymour2004, Matsumoto2009} & measured \\
15 & Pain interrupts the current work & Humans, an electrical pain stimulus and auditory probes: the response to a probe $100\ms$ after pain onset is markedly slowed, and after $1.5\,\text{s}$ there is no longer a difference from control. A review combines such experiments into a model of attentional capture. & \cite{Crombez1996disrupt, EcclestonCrombez1999} & measured \\
16 & Withdrawal is wired in the spinal cord & Rat: neurons of the deep layers of the dorsal horn reproduce the spatial pattern of the withdrawal reflex of individual muscles; they cannot be excited antidromically from the cervical cord, so they are spinal interneurons. & \cite{Schouenborg1995wr} & measured \\
\bottomrule
\end{longtable}}

\section{Chapter~\ref{sec:motorlearning}. Motor Learning}

The least-mean-squares rule and the advantage of a separate error wire are theorems; the structure of the error-wire circuit, weakening on coincidence, tuning of the trace to the delay, the aftereffect, and spontaneous recovery of a skill have been measured; that the whole circuit works as supervised learning and builds an internal model is the leading hypothesis; the numbers of the two-process model are a calculation with the book's model.

{\footnotesize
\setlength{\tabcolsep}{3pt}\renewcommand{\arraystretch}{1.15}
\begin{longtable}{>{\raggedright}p{0.5cm}>{\raggedright}p{4.0cm}>{\raggedright}p{5.6cm}>{\raggedright}p{2.3cm}>{\raggedright\arraybackslash}p{1.7cm}}
\toprule
No. & Claim & Experiment and what was measured & Source & Status \\
\midrule\endfirsthead
\toprule
No. & Claim & Experiment and what was measured & Source & Status \\
\midrule\endhead
1 & The rule~\eqref{eq:lms} on average follows the gradient of the squared error & A result of adaptive-filter theory: a correction proportional to the input and the output error is stochastic gradient descent on the mean squared error. & \cite{WidrowHoff1960} & theory \\
2 & With an error wire for each output, the speed does not depend on the number of outputs; with one shared number, it falls with the number of noisy neurons (Proposition~\ref{prop:teachers}) & Learning curves of a linear network: learning by random perturbation of the neurons is slower than the exact gradient by a factor equal to the number of perturbed neurons. & \cite{Werfel2005} & theory \\
3 & The error-wire circuit works as supervised learning (Proposition~\ref{prop:errorwire}) & Theories derived from the structure of the circuit. The experiments of rows 7--10 agree with them; that the whole circuit works exactly this way has not been proved. & \cite{Marr1969, Albus1971} & hypothesis \\
4 & The expander: about $7\cdot10^{10}$ small \nt{GLUT} neurons, more than in all the rest of the brain & Counting nuclei in a suspension of dissolved tissue: the human cerebellum has $69\cdot10^9$ neurons out of $86\cdot10^9$ in the whole brain. Stereology: $101\cdot10^9$ small neurons in the cerebellum of elderly men. & \cite{Azevedo2009, Andersen1992cereb} & measured \\
5 & Raising the dimension of the input makes the desired dependence linear & The function-counting theorem: a weighted sum of $n$ inputs with a threshold realizes an arbitrary labeling of about $2n$ vectors in general position. That the cerebellum uses exactly this is part of the hypothesis of row~3. & \cite{Cover1965} & theory \\
6 & About $3\cdot10^7$ output \nt{GABA} neurons, each with on the order of $10^5$ inputs & Stereology of the human cerebellum: $30.5\cdot10^6$ output neurons. Electron microscopy, rat: about $1.75\cdot10^5$ inputs from the expander per output neuron. The inhibitory transmitter of the output neurons is from a review. & \cite{Andersen1992cereb, NapperHarvey1988, Ito2001} & measured \\
7 & Exactly one error wire per output neuron, about once per second, each time a powerful burst & A review of recordings: each output neuron receives one climbing \nt{GLUT} input, which evokes a complex spike at a rate of about $1$~Hz. & \cite{Ito2001} & measured \\
8 & The probability of the burst depends on the direction of the movement error & Monkey, oculomotor cerebellum, saccade adaptation: the direction of the visual error sets the probability of a complex spike, and its size sets the timing; the burst changes the simple spikes on the next trial and shifts the movement along the cell's preferred error. & \cite{Herzfeld2018} & measured \\
9 & Coincidence of the error burst with an active input weakens that input ($-\eta e_i x_j$) & A review of experiments: joint stimulation of an input from the expander and the error wire causes long-term weakening of that input; stimulation of the input alone does not. & \cite{Ito2001} & measured \\
10 & The trace length is tuned to the error delay: with a $100\ms$ delay, the input that was active $100\ms$ earlier is weakened most & Slices and living mice: in the region responsible for oculomotor learning, weakening is narrowly tuned to an interval of about $120\ms$ between input and error, the time the error signal takes in this task; in another region different cells are tuned to different intervals. & \cite{Suvrathan2016} & measured \\
11 & The circuit is an adaptive filter~\eqref{eq:adaptivefilter} that learns an internal model of the body & A model derived from the structure of the circuit. There is no direct experiment in which the learned filter has been measured as the transfer function of the body. & \cite{Fujita1982} & hypothesis \\
12 & Prediction subtracts the consequences of one's own movements, so one cannot tickle oneself & Humans: one's own touch feels less ticklish than the same touch by someone else; in the scanner the somatosensory cortex responds more weakly to one's own touch, and the cerebellum is less active when the movement touches the skin. The explanation by subtraction of a prediction is a hypothesis. & \cite{Blakemore1998} & hypothesis \\
13 & Under an external force the miss shrinks, and after the force is removed the arm misses to the other side & People reach with a robot handle in a force field: the trajectories return to straight lines; after the field is removed, the aftereffect is a near mirror image of the initial distortions, and it transfers to untrained parts of the workspace. & \cite{ShadmehrMussaIvaldi1994} & measured \\
14 & Two processes learn~\eqref{eq:tworate}; after a reverse perturbation the skill returns by itself & Humans, a force field, then a short reverse field and trials with the error clamped: the correction returns by itself toward the old one; a two-process model reproduces this, a one-process model does not. & \cite{Smith2006} & measured \\
15 & The numbers of Fig.~\ref{fig:tworate}: $0.84$ (slow $0.63$) after $200$ movements; $6$ reverse ones; fast $-0.53$, slow $0.46$; rebound to $0.37$ within $40$ movements & Calculation with \texttt{models/\allowbreak motor\_\allowbreak learning.py}, $A_f=0.92$, $B_f=0.1$, $A_s=0.996$, $B_s=0.02$: $0.840$ ($0.634$ and $0.205$), $6$ movements, $-0.531$ and $0.457$, a peak of $0.371$ at the $40$th movement. & \texttt{models/\allowbreak motor\_\allowbreak learning.py} & model \\
16 & Two processes are needed because the body changes at different rates & Theory: optimal estimation with disturbances of several lifetimes yields several filters with different time constants and explains spontaneous recovery and faster relearning. & \cite{Kording2007} & hypothesis \\
\bottomrule
\end{longtable}}

\section{Chapter~\ref{sec:chemoutput}. The chemical output}

The two wires to the heart and their different speeds, negative feedback in hormone cascades, the pulse-frequency code, and the daily oscillator and its locking to light are measured directly; the bound $K<8$ is a theorem of control theory derived in the chapter; pulsations generated by the cascade feedback itself are a hypothesis with strong experimental support.

{\footnotesize
\setlength{\tabcolsep}{3pt}\renewcommand{\arraystretch}{1.15}
\begin{longtable}{>{\raggedright}p{0.5cm}>{\raggedright}p{4.0cm}>{\raggedright}p{5.6cm}>{\raggedright}p{2.3cm}>{\raggedright\arraybackslash}p{1.7cm}}
\toprule
No. & Claim & Experiment and what was measured & Source & Status \\
\midrule\endfirsthead
\toprule
No. & Claim & Experiment and what was measured & Source & Status \\
\midrule\endhead
1 & The heart's pacemaker beats by itself about $100$ times a minute; at rest the brake is on, and the rate is usually about $60$--$70$ & Humans, complete block of both wires to the heart: the intrinsic rate is above the resting rate and declines with age, about $100$ beats per minute in adults. So the brake dominates at rest. The resting rate of $60$--$70$ is a typical value, not cited separately. & \cite{JoseCollison1970ihr} & measured \\
2 & The \nt{NORA} gas and the \nt{ACET} brake are low-pass filters, the brake is faster~\eqref{eq:gasbrake} & Dog, frequency-modulated stimulation of the vagal or sympathetic cardiac nerve and the transfer function to atrial rate: both pathways are low-pass filters; the sympathetic one has a much lower corner frequency plus a pure delay of about $1.7\,\text{s}$. The characteristics depend on the mean tone, so the linear formula~\eqref{eq:gasbrake} is an approximation. & \cite{Berger1989} & measured \\
3 & The brake is fast because \nt{ACET} opens the potassium channel without a second messenger, while \nt{NORA} acts through a chain of reactions & Membrane patches of atrial cells: the potassium channel opened by acetylcholine is opened by the $\beta\gamma$ subunits of the receptor-coupled G protein, with no second messenger inside the cell. The \nt{NORA} pathway via cAMP is not cited here. & \cite{Logothetis1987gbg} & measured \\
4 & Blood circulates through the body in about a minute; \nt{ADRE} from the bloodstream reaches every organ with a receptor & A general order-of-magnitude estimate; stated as such in the chapter, no source cited. & --- & estimate \\
5 & A hormone cascade is enclosed in negative feedback: the final hormone inhibits the first stages & Review of experiments: adrenocortical hormones suppress ACTH release by the pituitary and releasing-hormone release by the hypothalamus in three time domains: fast, intermediate, and slow. & \cite{KellerWood1984fb} & measured \\
6 & Three identical stages are stable for $K<8$~\eqref{eq:cascadestable}; at the boundary the period is $2\pi\tau/\sqrt3\approx3.6\tau$ & Derived in the chapter: at the frequency $\sqrt3/\tau$ three stages give a $180^\circ$ phase shift and an attenuation by $8$. & derived in chapter & theory \\
7 & The level error is $1/(1+K)$, so such a regulator cannot do better than about $10\,\%$ (Proposition~\ref{prop:cascade}) & A consequence of row~6 for proportional feedback. The real axis has feedback at several stages and in several time domains (row~5), so the bound applies to model~\eqref{eq:cascade} and is not measured. & derived in chapter & theory \\
8 & At $K=4$ and $7$ the level settles; at $K=12$, steady pulsations with period $3.7$--$3.9\,\tau$ (Fig.~\ref{fig:cascade}) & Computed with \texttt{models/\allowbreak chem\_\allowbreak models.py}: at $K=12$ successive periods are $3.9$, $3.7$, $3.8$, $3.7\,\tau$; at $K=4$ the swing at the end of the run is $0.003$. & \texttt{models/\allowbreak chem\_\allowbreak models.py} & model \\
9 & The stress hormone is released in pulses about once an hour; the pulses are generated by the feedback itself, without a separate oscillator & Rats, constant infusion of releasing hormone: ACTH and corticosterone oscillate at the physiological, nearly hourly frequency, so no rhythmic input from the hypothalamus is needed for the pulses. A pituitary--adrenal model with delayed feedback produces such oscillations. & \cite{Walker2012glc, Walker2010} & hypothesis \\
10 & The receiver reads the pulse frequency, not the level & Monkeys with no release of their own gonadotropin-releasing hormone: constant infusion does not restore the release of pituitary hormones, hourly pulses do; switching back to constant infusion silences the response again. Receptor fatigue as a high-pass filter is an interpretation. & \cite{Belchetz1978} & measured \\
11 & Different pulse frequencies act differently on different cells of one gland & The same monkeys: raising the rate to $2$--$5$ pulses per hour lowers both pituitary hormones; lowering it to one per $3$~hours lowers one (LH) and raises the other (FSH), so their ratio is set by the rate. & \cite{Wildt1981gnrh} & measured \\
12 & The daily rhythm is generated by intracellular negative feedback with a delay of several hours & Review: clock-gene proteins suppress the transcription of their own genes with a delay; the transcription--translation loop gives a period of about a day. & \cite{TakahashiJS2017} & measured \\
13 & The master oscillator is a group of tens of thousands of neurons that synchronizes the rest of the body & Review: the suprachiasmatic nucleus is the master daily oscillator; its individual neurons in culture are autonomous oscillators, and the network synchronizes them; in the mouse, about $10^4$ neurons on each side. & \cite{Welsh2010scn} & measured \\
14 & The intrinsic period in humans is $24.18$~hours & Humans in dim light on a schedule decoupled from the day: the period of the melatonin, temperature, and cortisol rhythms averages $24.18$~hours in young and older people, with a narrow spread. & \cite{Czeisler1999} & measured \\
15 & Light adjusts the oscillator like a phase-locked loop~\eqref{eq:pll}; a shift of about $11$~min per day is needed & Humans, a single $6.7$-hour bright-light exposure at different times of day: a type~1 phase response curve with a $5$-hour range, delay before the body-temperature minimum and advance after it. Equation~\eqref{eq:pll} is the standard model; $11$~min $=0.18$~hours is arithmetic. & \cite{Khalsa2003prc} & measured \\
16 & Without light there is no locking, and the rhythm drifts to its own period & Totally blind people with no light perception living in ordinary society: in about half of them the melatonin and cortisol rhythms run at their own period, which does not match the day. & \cite{Sack1992blind} & measured \\
\bottomrule
\end{longtable}}

\section{Chapter~\ref{sec:debugging}. Debugging the machine}

What an electrode hears, the low dimensionality of activity, the performance of interfaces on rigid arrays, the joint learning of brain and decoder, and the visual spots produced by stimulation are measured in peer-reviewed experiments; the data-rate estimates are the chapter's arithmetic; for the Neuralink device implanted in humans, as far as we could check, data have been published mostly by the company itself.

{\footnotesize
\setlength{\tabcolsep}{3pt}\renewcommand{\arraystretch}{1.15}
\begin{longtable}{>{\raggedright}p{0.5cm}>{\raggedright}p{4.0cm}>{\raggedright}p{5.6cm}>{\raggedright}p{2.3cm}>{\raggedright\arraybackslash}p{1.7cm}}
\toprule
No. & Claim & Experiment and what was measured & Source & Status \\
\midrule\endfirsthead
\toprule
No. & Claim & Experiment and what was measured & Source & Status \\
\midrule\endhead
1 & An electrode hears the spikes of neurons within a radius of about a hundred micrometers, usually one to several; they are told apart by shape & Rat, area CA1, simultaneous intracellular and extracellular recording: the extracellular spike shape follows the width and amplitude of the intracellular one. Estimate: a tetrode could separate $60$--$100$ neurons, but in practice about six are isolated. & \cite{Henze2000extra} & measured \\
2 & The brain has $8.6\cdot10^{10}$ neurons; a probe hears about one hundred-millionth & Counting nuclei in a suspension of dissolved tissue: $86\cdot10^9$ neurons in the human brain. The fraction is the chapter's arithmetic. & \cite{Azevedo2009} & measured \\
3 & The activity of hundreds of motor-cortex neurons fits into one or two dozen shared variables & Review of recordings in monkey motor cortex: population activity is described by a few latent variables shared by many neurons. & \cite{Gallego2017} & measured \\
4 & Neighbors' noise is correlated, and a hundred neurons carry almost as much as a thousand (Proposition~\ref{prop:pooling}) & Monkey visual cortex: the noise correlation coefficient of neighboring neurons is about $0.1$, and the gain from averaging saturates at a hundred neurons. Theory of information-limiting correlations. & \cite{Zohary1994, MorenoBote2014} & measured \\
5 & Raw stream of $2\cdot10^8$~bit/s versus $8\cdot10^4$~bit/s in spikes~\eqref{eq:probedata} & The chapter's arithmetic from the capacity of the spike stream~\eqref{eq:spikeentropy}. The digitization parameters are real: the Neuralink chip samples at $19.3$~kHz, $10$~bits per channel. & derived in chapter; \cite{Rieke1997, Musk2019} & theory \\
6 & The first Neuralink system: chips amplify and digitize, the raw stream goes over a cable, spikes are detected by software outside; on-chip detection, a sealed housing, radio, and wireless power belong to the device for humans, according to the company & Company paper: the full raw stream over a USB-C cable, spikes detected by real-time software off the chip; on-board compression and wireless power and transmission are named as a plan for clinical devices. For the device implanted in humans, only company reports. That on-chip spike detection is necessary is the chapter's conclusion from~\eqref{eq:probedata}. & \cite{Musk2019}; Neuralink reports, no peer-reviewed paper & measured (company paper); for humans, company report \\
7 & Up to $3072$ electrodes on $96$ threads of $32$; a robot inserts the threads, avoiding vessels & Company paper: $3072$ electrodes on $96$ threads of $32$; the robot inserts $6$ threads ($192$ electrodes) per minute; up to $70\,\%$ of electrodes in rats with a chronically implanted array record spikes. & \cite{Musk2019} & measured (company paper) \\
8 & In humans, about a thousand channels on $64$ threads; some threads shifted; cursor at about $10$~bit/s & Company reports only. A PubMed search found no peer-reviewed paper with human results; this agrees with the caveat in the chapter text. & Neuralink reports; no peer-reviewed paper & measured (company report) \\
9 & An interface on an array of a hundred electrodes controls a cursor & A person with paralysis, $96$ electrodes in primary motor cortex: the intention to move the arm modulates spikes three years after the injury; the decoder gives a ``neural cursor'' (email, television), control of a prosthetic hand and a robotic arm. & \cite{Hochberg2006} & measured \\
10 & Writing with the ``mental hand'': about $90$ characters per minute, less than $8$~bit/s & A person with a paralyzed hand, attempted handwriting, recurrent decoder: $90$ characters per minute at $94.1\,\%$ accuracy in real time, above $99\,\%$ after autocorrection. The estimate in bits is the chapter's arithmetic. & \cite{Willett2021} & measured \\
11 & Attempts to speak are turned into text at about $60$ words per minute & A person who lost intelligible speech, arrays in the cortex: $62$ words per minute; $9.1\,\%$ word errors with a $50$-word vocabulary, $23.8\,\%$ with a $125\,000$-word vocabulary. & \cite{Willett2023} & measured \\
12 & An interface extracts a small fraction of the capacity: tens of bits per second for one neuron, thousands for a hundred (Proposition~\ref{prop:probe}) & The upper bound~\eqref{eq:spikeentropy} is theory; the comparison with rows~8 and 10 is the chapter's conclusion. That the bottleneck is low dimensionality and ignorance of the code rather than the wire has not been tested directly. & \cite{Rieke1997} & theory \\
13 & The brain and the decoder learn from each other & Monkeys control a cursor; the decoder adapts in closed loop while the neurons relearn at the same time: accuracy rises, the skill is retained and suffers less from movements of the monkey's own arm. The advice to separate the learning rates is an engineering rule; the chapter has no separate experiment for it. & \cite{Orsborn2014coad} & measured \\
14 & Current through an electrode excites the surrounding neurons and wires regardless of type & Mouse and cat cortex, two-photon calcium imaging: even a weak current excites sparse neurons scattered over up to millimeters, through direct excitation of axons in a volume tens of micrometers across. So the excitation is nonselective but not solid. & \cite{Histed2009stim} & measured \\
15 & Stimulation of the visual cortex lights spots; up to $16$ spots at once allow recognizing some letters and object edges & A blind person, $96$ electrodes in the visual cortex for $6$ months: single-electrode threshold $66.8\pm36.5$~\textmu A; simultaneous stimulation of groups (up to $16$ electrodes, the stimulator's limit) lowers the threshold and gives distinguishable patterns; some letters and object edges are recognized. & \cite{Fernandez2021} & measured \\
16 & Neuralink's second direction is a device that feeds signals into the visual cortex & Only company reports of development; no peer-reviewed data on its performance were found. & Neuralink reports & hypothesis \\
17 & Concentrations of broadcast transmitters can be measured with a chemical sensor in the tissue & Rat brain slices, fast-scan cyclic voltammetry at a carbon fiber: serotonin release and reuptake measured; the substance spills beyond the synapse. & \cite{Bunin1998} & measured \\
\bottomrule
\end{longtable}}

\section{Chapter~\ref{sec:eetypes}. Neurons as devices}

The operating modes of single cells are measured directly, with current injected through an electrode and voltage recording; the mathematics of the integrator and of the division into classes is classical theory, while the idea that the brain uses mode reconfiguration for computation remains a hypothesis.

{\footnotesize
\setlength{\tabcolsep}{3pt}\renewcommand{\arraystretch}{1.15}
\begin{longtable}{>{\raggedright}p{0.5cm}>{\raggedright}p{4.0cm}>{\raggedright}p{5.6cm}>{\raggedright}p{2.3cm}>{\raggedright\arraybackslash}p{1.7cm}}
\toprule
No. & Claim & Experiment and what was measured & Source & Status \\
\midrule\endfirsthead
\toprule
No. & Claim & Experiment and what was measured & Source & Status \\
\midrule\endhead
1 & Resonator: a slow current opposing voltage changes makes the neuron a band-pass filter peaking at a few to tens of hertz (Section~\ref{sec:resonator}) & In slices, sinusoidal current of smoothly rising frequency was injected into \nt{GLUT} neurons and the impedance $|Z(f)|$ was measured: a peak at $2$--$5$~Hz at $33^\circ$C (about $7$~Hz at $38^\circ$C); the resonance is created by slow potassium and cation currents and amplified by a persistent sodium current. A review describes the same technique for many cell classes & \cite{HuVervaekeStorm2002,HutcheonYarom2000} & measured \\
2 & The slow current behaves like an inductance and, together with the membrane capacitance, forms a tank circuit & The subthreshold response of an axon to current was measured and described by linearized conductance equations; the equivalent circuit contains a ``phenomenological'' inductance whose value depends on voltage & \cite{Mauro1970} & theory \\
3 & The time constant of a group with feedback is $\tau_{\text{eff}}=\tau/(1-w)$~\eqref{eq:perfectint}; with $\tau=0.1$~s and $w=0.995$ it is $20$~s & Derived in the chapter from the linear group equation; proposed in a theoretical paper as the mechanism for holding eye position & derived in chapter; \cite{Seung1996} & theory \\
4 & The eye-position integrator holds its input by the network, not by a property of a single cell & Intracellular recording in fish in the nucleus that holds eye position: after a fast eye turn the neuron's rate changes in a step and holds; a brief current into one cell produces only a transient shift, while the voltage steps persist below threshold as well, sustained by synaptic inputs & \cite{Aksay2001} & measured \\
5 & A leak-free integrator needs constant retuning by learning; when mistuned, it either forgets or runs into saturation & Fish were trained with visual feedback proportional to $\pm$ eye position: holding became unstable or leaky, and the neurons' time constant fell by more than an order of magnitude, to $<1$~s; ordinary visual feedback gradually restored stability & \cite{Major2004} & measured \\
6 & Bistable neuron: a brief input switches on a lasting plateau, inhibition switches it off; the property is enabled by serotonin (Section~\ref{sec:bistable}) & Intracellular recording from the \nt{ACET} neurons that control muscles in the cat: brief synaptic excitation puts the neuron into lasting activity, brief inhibition resets it; in the spinal cat the state appears after a serotonin precursor is given & \cite{Hounsgaard1988} & measured \\
7 & Two classes: integrators (rate rises from zero) and resonators (a finite rate right at threshold) & The classes are derived from bifurcation theory. Experiment: in cortical slices, \nt{GLUT} neurons start firing at arbitrarily low rates (type~1), fast \nt{GABA} neurons jump in at a finite rate (type~2) & \cite{Izhikevich2007,Tateno2004} & measured \\
8 & One neuron is an integrator or a coincidence detector: inhibition shortens the summation window & In hippocampal slices, the window within which excitatory inputs sum to threshold is shorter than $2$~ms; it is shortened by feedforward inhibition arriving right after the excitation; in dendrites, where inhibition is weaker, the window is wider & \cite{Pouille2001} & measured \\
9 & A delay line made of axons (Table~\ref{tab:eetypes}) & In the barn owl, delays were measured in the axons leading to coincidence detectors: different axon lengths give different delays, and the detectors are tuned to the difference in sound arrival times & \cite{Carr1990} & measured \\
10 & A pacemaker in wakefulness and a silent cell in sleep & \nt{SERO} neurons fire almost regularly during the day, slow down in slow-wave sleep, and are nearly silent in REM sleep (more in Chapter~\ref{sec:sleep}) & \cite{JacobsAzmitia1992,McGintyHarper1976} & measured \\
11 & The device is set by ion channels and by the broadcast channels that tune them (Proposition~\ref{prop:eetypes}) & Shown in small networks: modulatory substances change the intrinsic properties of neurons and the strength of synapses, so the same circuit produces different rhythms & \cite{Marder2012} & measured \\
12 & The brain uses mode reconfiguration for computation (Proposition~\ref{prop:eetypes}) & There is no direct experiment in which reconfiguring a cell's mode was needed to solve a task; there are only experiments in which modulators change a circuit's output & \cite{Marder2012} & hypothesis \\
\bottomrule
\end{longtable}}

\section{Chapter~\ref{sec:dendrites}. Number of dendrites}

The structure of the classes and the numbers of inputs are measured anatomically and by electrical recordings; the estimate~\eqref{eq:dendriteload} is simple capacitor physics, while the computational explanation of the number of inputs rests on theory and models.

{\footnotesize
\setlength{\tabcolsep}{3pt}\renewcommand{\arraystretch}{1.15}
\begin{longtable}{>{\raggedright}p{0.5cm}>{\raggedright}p{4.0cm}>{\raggedright}p{5.6cm}>{\raggedright}p{2.3cm}>{\raggedright\arraybackslash}p{1.7cm}}
\toprule
No. & Claim & Experiment and what was measured & Source & Status \\
\midrule\endfirsthead
\toprule
No. & Claim & Experiment and what was measured & Source & Status \\
\midrule\endhead
1 & Specific membrane capacitance $c_m\approx1$~\textmu F/cm$^2$ & Membrane was pulled from a neuron's cell body into a pipette, a voltage step was applied, the total capacitance was found from the decay of the capacitive current and then divided by the area: $0.9$~\textmu F/cm$^2$ in three neuron classes & \cite{Gentet2000} & measured \\
2 & Charging time $t\approx c_mA\Delta V/I$~\eqref{eq:dendriteload}: a large neuron needs dozens of simultaneous inputs, a small one is served by a single large synapse & Estimate from the capacitor charging law; the numbers ($20$ and $300$~pF, $0.1$ and several nA) are orders of magnitude & derived in chapter & theory \\
3 & One huge synapse gives a current of several nanoamperes and drives a small neuron to threshold at once & Simultaneous recording of the terminal and the receiver in a slice: single-synapse current $2$--$13$~nA, every input spike evokes a suprathreshold response, delay about $0.4$~ms at $36^\circ$C; the receiver outputs \nt{GLYC} & \cite{Borst1995,BorstSoria2012} & measured \\
4 & Zero dendrites: in touch and pain sensory neurons the spike travels from the sensor to the spinal cord past the cell body & The structure (an axon splitting in two at the cell body) is established anatomically. That conduction does not require an excitable cell body is shown in a validated model of such a neuron: without an excitable body the spike passes the fork just as reliably & \cite{AmirDevor2003} & theory \\
5 & One dendrite: an olfactory neuron carries one receptor type and transmits one input line & Review of experiments with receptor genes: each olfactory neuron expresses one receptor gene from a large family & \cite{Mombaerts2004} & measured \\
6 & Two dendrites: in sound direction detectors, inputs from the left and right ears land on different dendrites & Anatomy and recordings in mammals collected in a review: inputs from the two ears are separated onto two opposite dendrites & \cite{Grothe2010} & measured \\
7 & Separating inputs onto two dendrites makes the AND stricter & Shown in a model: saturation~\eqref{eq:shunt} on one dendrite keeps the inputs from one side from reaching threshold, and coincidence detection improves. There is no direct experiment in which the input layout was changed & \cite{AgmonSnir1998} & theory \\
8 & About $7\cdot10^{10}$ small \nt{GLUT} neurons in the motor learning circuit, with $\sim4$ inputs each & Cell counting by isotropic fractionation: about $69$ billion neurons in this part of the human brain. Recording of one such cell in a live rat: a touch brings bursts of discrete currents from a few inputs, and the cell fires when they sum & \cite{Azevedo2009,Chadderton2004} & measured \\
9 & A threshold element with $n$ inputs separates at most $P_{\max}\approx2n$ random patterns~\eqref{eq:cover}; for the output neuron of the motor learning circuit the bound is $\sim2\cdot10^5$, and a realistic estimate is $\sim4\cdot10^4$ associations & The bound is a classical result of combinatorial geometry. The realistic estimate is capacity theory with positive weights only and a noise margin, applied to the measured number of inputs of this neuron & \cite{Cover1965,Brunel2004} & theory \\
10 & Mixers with few inputs are needed to expand the input; ``four'' is close to optimal & Theory: the dimensionality of the representation given by a layer of mixers is maximal at roughly the number of inputs observed anatomically; the original expansion hypothesis is in the classical papers & \cite{Marr1969,Albus1971,LitwinKumar2017} & hypothesis \\
11 & Large trees: $10^4$--$10^5$ inputs & Synapse counts from electron micrographs: $\approx1.7\cdot10^5$ synapses on one output \nt{GABA} neuron of the motor learning circuit in the rat; about $30\,000$ excitatory and $1\,700$ inhibitory inputs on a hippocampal \nt{GLUT} neuron & \cite{NapperHarvey1988,Megias2001} & measured \\
12 & The branches of a large tree are separate subcircuits with their own thresholds; the neuron is a small multilayer network & Focal stimulation of two sites on thin dendrites: inputs on one branch sum sigmoidally, on different branches linearly. Graded calcium spikes were found in the dendrites of human cortical neurons that let a single cell solve a linearly nonseparable task. Model: reproducing one neuron requires a deep network & \cite{Polsky2004,Gidon2020,Beniaguev2021} & measured \\
13 & Output dendrites: each branch of a retinal \nt{GABA} neuron is a separate direction detector & Two-photon calcium imaging in individual branches: the signal in a branch is selective for motion from the cell body toward the tip, while the somatic voltage is not & \cite{Euler2002} & measured \\
14 & The number of inputs follows the task (Proposition~\ref{prop:dendrites}) & The structure of the classes is measured (rows 3--13); that the number of inputs is chosen for speed or for classification is shown only by theory and models for individual classes & \cite{LitwinKumar2017,AgmonSnir1998} & hypothesis \\
\bottomrule
\end{longtable}}

\section{Chapter~\ref{sec:sleep}. Sleep}

The sleep modes, the replays, and the shrinking of synapses are measured by neuron recordings, microdialysis, and electron microscopy; the sleep schedule and the slow swings are described by the book's models, while the purposes of sleep are mostly hypotheses.

{\footnotesize
\setlength{\tabcolsep}{3pt}\renewcommand{\arraystretch}{1.15}
\begin{longtable}{>{\raggedright}p{0.5cm}>{\raggedright}p{4.0cm}>{\raggedright}p{5.6cm}>{\raggedright}p{2.3cm}>{\raggedright\arraybackslash}p{1.7cm}}
\toprule
No. & Claim & Experiment and what was measured & Source & Status \\
\midrule\endfirsthead
\toprule
No. & Claim & Experiment and what was measured & Source & Status \\
\midrule\endhead
1 & Cortical neurons are not silent in sleep; the operating mode changes & Long recordings of cortical neurons in rats: in slow-wave sleep the rate stays nonzero, and periods of activity alternate with periods of silence; after long wakefulness the rate is higher in all states, after sleep it is lower & \cite{Vyazovskiy2009} & measured \\
2 & \nt{ACET} is high in the day and in REM sleep, low in slow-wave sleep (Table~\ref{tab:sleepmodes}) & Microdialysis in the cortex of freely moving cats: in wakefulness acetylcholine release is $100$--$175\%$ higher than in slow-wave sleep, and in REM sleep it is about at the level of quiet wakefulness & \cite{Marrosu1995} & measured \\
3 & \nt{NORA} and \nt{SERO} quiet down in slow-wave sleep and are nearly silent in REM sleep & Recordings of single \nt{NORA} neurons in rats and \nt{SERO} neurons in cats across sleep stages: the rate is highest in wakefulness, lower in slow-wave sleep, and nearly zero in REM sleep & \cite{AstonJonesBloom1981,McGintyHarper1976} & measured \\
4 & In REM sleep the muscles are switched off by inhibition through \nt{GABA} and \nt{GLYC} & In rats the activity of the \nt{ACET} neurons of a jaw muscle was measured and blockers were applied directly to them: the paralysis is lifted only if both types of \nt{GABA} receptors and the \nt{GLYC} receptors are all blocked & \cite{BrooksPeever2012} & measured \\
5 & Sleep pressure $S$ is a capacitor with two thresholds~\eqref{eq:sleeppressure}, $\tau_r=18.2$~h, $\tau_d=4.2$~h & The two-process model; the time constants are obtained from how EEG slow-wave power rises with wakefulness and falls during sleep, and the shape of the thresholds from the timing of spontaneous waking & \cite{Borbely1982,Daan1984} & theory \\
6 & The machine settles by itself into $16.1$~h awake and $8.0$~h asleep (Fig.~\ref{fig:twoprocess}) & Computed from~\eqref{eq:sleeppressure} with swinging thresholds, script \texttt{models/sleep\_models.py}: $16.1$~h awake and $7.9$--$8.0$~h asleep & --- & model \\
7 & After $40$~h without sleep $S=0.90$, but sleep lasts only $8.8$~h: the debt is repaid by depth, not by duration & Model: \texttt{models/sleep\_models.py} gives $S=0.90$ and $8.8$~h. Experiment: after $40.5$~h without sleep, in eight people EEG slow-wave power in the first recovery night is significantly above normal & \cite{Borbely1981} & model \\
8 & Physically, $S$ is adenosine & Microdialysis in cats: extracellular adenosine in one of the regions that control wakefulness rises during wakefulness, rises further with sleep deprivation, and falls during sleep. That $S$ is adenosine alone has not been shown & \cite{PorkkaHeiskanen1997,Dunwiddie2001} & hypothesis \\
9 & In slow-wave sleep the cortex swings: activity for a few hundred ms, then silence, less than once a second ($<1$~Hz, under anesthesia most often $0.3$--$0.4$~Hz) & Intracellular recording in anesthetized cats: depolarization for $0.8$--$1.5$~s and long hyperpolarization, rhythm $<1$~Hz, most often $0.3$--$0.4$~Hz; the same alternation of activity and silence is seen in natural sleep (row~1) & \cite{Steriade1993,Vyazovskiy2009} & measured \\
10 & The swings are a group with feedback and fatigue~\eqref{eq:updown}; at $b=5$ the period is $1.1$~s (a model value) and activity takes about $35\%$ of the time; at $b=1$ activity is steady (Fig.~\ref{fig:updown}) & Computation, script \texttt{models/sleep\_models.py}: period $1.11$~s, active time fraction $0.35$ (averaged over a whole number of periods); at $b=1$ the active fraction is $1.0$. The model period is shorter than measured (row~9) & --- & model \\
11 & \nt{ACET} switches off the swings and moves the cortex into steady activity & Stimulation of a nucleus of \nt{ACET} neurons in the cat blocked the slow rhythm in $79\%$ of cortical neurons and moved them into steady activity, mainly because the long hyperpolarizations disappeared. Acetylcholine suppresses the slow potassium currents that create fatigue & \cite{SteriadeAmzica1993,McCormick1992} & measured \\
12 & The $7$--$15$~Hz rhythm bursts are generated by the \nt{GLUT}--\nt{GABA} loop between the cortex and a relay node & In slices of the ferret's visual relay node, the bursts are created by rebound bursts of relay cells in response to inhibition from a neighboring \nt{GABA} nucleus, which they themselves excite & \cite{vonKrosigk1993,SteriadeMcCormickSejnowski1993} & measured \\
13 & The day's replays run fast-forward: about $20$ times faster in the hippocampus, $6$--$7$ times in the cortex & In slow-wave sleep, repeating sequences of hippocampal neurons run compressed in time. After running on a track, the sequences of place neurons were replayed in the same order in short bursts of $\sim100$~ms, compressed about $20$ times; in the cortex the compression is $6$--$7$ times & \cite{Nadasdy1999,LeeWilson2002,Euston2007} & measured \\
14 & Strengthening the alignment of replays with the cortex improves memory (a causal link) & In rats after training, stimulation locked to replays strengthened their coupling to the slow waves and rhythm bursts in the cortex: the next day memory was high, while in controls it was at chance & \cite{Maingret2016} & measured \\
15 & The purpose of replays is interleaved training of the slow memory & The theory of two learning systems and computations on artificial networks: sequential learning corrupts the old, interleaved learning does not. There is no direct experiment in the brain showing that replays are needed for exactly this & \cite{McClelland1995} & hypothesis \\
16 & After sleep synapses shrink in proportion to their size; large ones hardly change & Three-dimensional electron microscopy of $6920$ synapses in mouse cortex: the contact area after sleep is $\sim18\%$ smaller than after wakefulness, the reduction is proportional to size, and large stable synapses do not change & \cite{deVivo2017} & measured \\
17 & The main job of sleep is renormalizing the weights & The synaptic homeostasis hypothesis; rows 1 and 16 agree with it but do not prove that this is the main function & \cite{TononiCirelli2014} & hypothesis \\
18 & REM sleep is a bench test of the world model & The mode (Table~\ref{tab:sleepmodes}) is measured; that the network runs a world model is a theoretical proposal with no experiment & \cite{Hobson2009} & hypothesis \\
19 & Flushing of the brain in sleep: the question is open & One experiment: in sleep the extracellular space is $60\%$ wider and clearance is faster. Another, with a different measurement method: in sleep and under anesthesia clearance is markedly slower. The results contradict each other & \cite{Xie2013,Miao2024} & hypothesis \\
20 & Local sleep: patches of an awake animal's cortex briefly drop into silence & In rats after long wakefulness, the neurons of an individual cortical patch briefly fall silent, as in slow-wave sleep, with a slow wave in the local EEG; at such moments the rat makes more errors in a task & \cite{Vyazovskiy2011} & measured \\
\bottomrule
\end{longtable}}

\section{Chapter~\ref{sec:livingmachine}. A Machine Made of Living Neurons}

The behavior of cultures on electrode arrays is measured in direct experiments with recording and stimulation; the burst mechanism rests on the synaptic-depletion model and on experiments in microcultures, and the claims about dopamine in a dish rest on single experiments, some of which the authors themselves regard only as demonstrations.

{\footnotesize
\setlength{\tabcolsep}{3pt}\renewcommand{\arraystretch}{1.15}
\begin{longtable}{>{\raggedright}p{0.5cm}>{\raggedright}p{4.0cm}>{\raggedright}p{5.6cm}>{\raggedright}p{2.3cm}>{\raggedright\arraybackslash}p{1.7cm}}
\toprule
No. & Claim & Experiment and what was measured & Source & Status \\
\midrule\endfirsthead
\toprule
No. & Claim & Experiment and what was measured & Source & Status \\
\midrule\endhead
1 & A culture on a chip: mostly \nt{GLUT}, a smaller fraction of \nt{GABA} that depends on the preparation; about $6\%$ in hippocampal cultures & A network of thousands of neurons on an electrode array is a standard experiment (row~3). The fraction of \nt{GABA} neurons in hippocampal cultures was counted by GABA staining: about $6\%$ of neurons. For cortical cultures we give no verified number & \cite{Benson1994} & measured \\
2 & Organoids: three-dimensional neural tissue about a millimeter across, grown from human stem cells & Three-dimensional organoids with several brain regions were grown from stem cells, including cortex with progenitor layers and mature neurons & \cite{Lancaster2013} & measured \\
3 & With no input a culture fires in bursts at intervals from a second to minutes, depending on the culture & $58$ cortical cultures with densities from $3\,000$ to $50\,000$ neurons were recorded for five weeks ($963$ half-hour recordings): after two weeks, whole-network bursts dominate in almost all cultures; intervals between bursts range from $1$ to $300$~s and depend strongly on the preparation & \cite{Wagenaar2006} & measured \\
4 & A burst ends by transmitter depletion, interval $T_{\text{burst}}\approx\tau_{\text{rec}}\ln\frac{1}{1-x^*}$~\eqref{eq:burstinterval}; $2$--$4$~s is a model estimate with $\tau_{\text{rec}}=0.8$~s, measured recovery is slower & The formula is derived in the chapter. Model: a network with depressing synapses generates whole-network bursts on its own. Experiment in microcultures of $4$--$30$ neurons: burst duration depends on the time since the previous burst, and depleting the transmitter vesicles abolishes bursts; the authors conclude that a burst is limited by the transmitter store. Store recovery there takes tens of seconds, which with the same formula gives intervals of tens of seconds to minutes & derived in the chapter; \cite{Tsodyks2000,CohenSegal2011,Tsodyks1997} & theory \\
5 & ``A teacher who falls silent'': stimulation is switched off when the desired response appears, and the response is locked in & The culture was stimulated at $0.3$--$1$~Hz until the selected response appeared $50\pm10$~ms after the stimulus, then stimulation was switched off for $5$~min; after several cycles the desired response was evoked right away; learning curves were plotted & \cite{Shahaf2001} & measured \\
6 & A culture plays tennis; significant growth of runs within a session only with full feedback & Details in chapter~\ref{sec:dishbrain} and its supplement: with silence in place of the stimulus the culture is also better by the end of the session than without feedback, but with a smaller effect; learning from session to session is not robust & \cite{Kagan2022} & measured \\
7 & That the culture learns to predict its input is a hypothesis; loud claims of ``sentience'' are disputed & The free-energy principle is a theoretical proposal; there is no direct experiment that distinguishes it from simpler explanations. Thirty researchers pointed out in a response article that behavior change in a loop shows neither intelligence nor sentience & \cite{Friston2010,Balci2023} & hypothesis \\
8 & A culture drives a virtual body toward a target with a chosen stimulus pattern & A program chose training stimulation patterns based on the current performance; within tens of minutes the network reached the target states, and plasticity lasted $80$~min after $2$~h of training. The same stimuli delivered without a link to behavior had no effect & \cite{Bakkum2008} & measured \\
9 & Reservoir: only the linear readout $y=W_{\text{out}}\mathbf r$~\eqref{eq:reservoir} is trained & Theory of reservoir computing: any nonlinear system with fading memory plus a trained linear output & \cite{Maass2002,JaegerHaas2004} & theory \\
10 & An organoid as a reservoir identifies speakers with about $78\%$ accuracy (as reported by the authors); the readout learns from the error, and the organoid's connections change without a teacher & Speech sounds from several speakers were delivered to an organoid as current patterns on an electrode array; the trained readout identified the speaker with about $78\%$ accuracy, as the authors report. They also report that the organoid's functional connectivity changed during training and call this unsupervised learning within the organoid itself & \cite{Cai2023} & measured \\
11 & A million neurons spend about $10^{-4}$~W on spikes~\eqref{eq:dishpower} & Estimate: a spike cost of $10^{-10}$~J from the cortical energy budget~\eqref{eq:energy}, times the number of spikes; there is no direct measurement of a culture's power & derived in the chapter; \cite{Attwell2001} & theory \\
12 & Dopamine on a flash of light: $1$~mM caged dopamine, $240$ repetitions, the number of evoked spikes grew & On a remote organoid platform, dopamine in a light-sensitive cage ($1$~mM) was released by ultraviolet light ($365$~nm, $0.8$~s) whenever the stimulus evoked more spikes; over $240$ steps (about $5$~h) the number of spikes grew overall. The authors write that a single experiment cannot establish a link to dopamine; there are no controls & \cite{Jordan2024} & hypothesis \\
13 & Dopamine from living cells changes the weight of an organic transistor lastingly and reversibly & Dopamine-releasing cells were grown above an organic neuromorphic element; oxidation of dopamine at the electrode changes the channel conductance; long-term weight change and its reset were shown & \cite{Keene2020} & measured \\
14 & Organoids with working \nt{DOPA} neurons exist; their dopamine has not been used as a teacher & Electrically active mature \nt{DOPA} neurons and dopamine production were found in organoids grown from human stem cells. We found no experiments in the literature where this dopamine teaches another network & \cite{Jo2016} & measured \\
15 & An organoid balances a pole: training stimuli from a program beat random ones, without consolidation, via \nt{GLUT} synapses & Mouse cortical organoids in a closed loop with the cart-pole task: in most organoids, signals chosen by reinforcement learning gave better performance than random signals or none; after $45$~min of rest the improvement did not persist; blocking AMPA and NMDA receptors abolished it & \cite{Robbins2026} & measured \\
16 & Without control registers a network on the bench cannot decide on its own what counts as success (proposition~\ref{prop:livingmachine}) & In all experiments of rows 5--15 the training signal is set by the experimenter or a program; there is no experiment in which a culture developed a success criterion on its own. This is an inference from absence, not an experiment & --- & hypothesis \\
\bottomrule
\end{longtable}}

\section{Chapter~\ref{sec:dishbrain}. DishBrain}

The bench design and the game results come from one experimental study and its follow-up; the formulas for noise and for drift toward good configurations are derived in the chapter and checked with the book's model, while the learning mechanism in the dish has not been established.

{\footnotesize
\setlength{\tabcolsep}{3pt}\renewcommand{\arraystretch}{1.15}
\begin{longtable}{>{\raggedright}p{0.5cm}>{\raggedright}p{4.0cm}>{\raggedright}p{5.6cm}>{\raggedright}p{2.3cm}>{\raggedright\arraybackslash}p{1.7cm}}
\toprule
No. & Claim & Experiment and what was measured & Source & Status \\
\midrule\endfirsthead
\toprule
No. & Claim & Experiment and what was measured & Source & Status \\
\midrule\endhead
1 & Bench: about $800\,000$ mouse cortical cells or about $10^6$ human cells per chip; $26\,000$ electrodes on $8$~mm$^2$, recording up to $1024$, stimulation through $8$ & Methods of the study: MaxOne chip, $26\,000$ platinum electrodes on $8$~mm$^2$, recording of $1024$ channels at $20\,000$ samples per second; up to $32$ electrodes can technically be stimulated, $8$ could be stimulated independently; mouse cultures were used from about $10$~days & \cite{Kagan2022} & measured \\
2 & A loop with a $10$~ms clock: spikes from the motor regions move the paddle (fig.~\ref{fig:dishloop}) & Spikes are counted in $10$~ms windows; the delay from a spike to the responding stimulus is about $5$~ms, set by the hardware buffer & \cite{Kagan2022} & measured \\
3 & Input: place code (which of $8$ electrodes) and rate code, $4$--$40$~pulses/s & In the main experiment the stimulation rate rises linearly from $4$~Hz when the ball is at the far wall to $40$~Hz at the paddle; the amplitude of the ball pulses is $75$~mV; the pulses are biphasic and rectangular & \cite{Kagan2022} & measured \\
4 & Output $\Delta p=c\,(n_\uparrow-n_\downarrow)$~\eqref{eq:dishreadout}, count noise per window $1/\sqrt{RT}$ & The two-region difference rule is described in the study (with region weights by activity); the exact formula is not given. The noise estimate follows from Poisson counting and is derived in the chapter & \cite{Kagan2022}; derived in the chapter & theory \\
5 & Teacher: after a hit, a predictable stimulus of $75$~mV, $100$~pulses/s, $0.1$~s; after a miss, an unpredictable one, $150$~mV, $5$~pulses/s, $4$~s at random places and times, then a $4$~s pause & Methods: after a hit, $75$~mV, $100$~Hz, $100$~ms to all $8$ electrodes at once; after a miss, $150$~mV, $5$~Hz to random electrodes at random times for $4$~s, then $4$~s without stimulation. The culture contains no dopamine & \cite{Kagan2022} & measured \\
6 & The network acts so as to make its input predictable & The free-energy principle is a theoretical proposal from which the authors derived the protocol; the experiment agrees with it but does not distinguish it from simpler mechanisms (rows~7--8) & \cite{Friston2010,Kagan2022} & hypothesis \\
7 & If failure shakes the network and success does not, the time in a configuration is $\rho\propto1/P_{\text{miss}}$~\eqref{eq:dishdensity} (proposition~\ref{prop:dishscramble}) & Follows from flux balance in a Markov chain with symmetric kicks; derived in the chapter. There is no direct test on a culture & derived in the chapter & theory \\
8 & The ``scramble on a miss'' model learns: fraction of hits $0.21\to0.38$; with a kick after every rally $\approx0.12$, with no kicks $0.19$ (fig.~\ref{fig:dishmodel}) & Computation, script \texttt{models/dishbrain\_model.py}, $40$ model cultures with $2000$ rallies each: $0.21\to0.38$; $0.13\to0.11$; $0.19\to0.19$ & --- & model \\
9 & Runs of hits grow longer only in living cultures in the full loop; the controls do not learn & $399$ sessions of $20$~min: $9$ mouse cultures ($101$ sessions), $11$ human ($138$), $6$ medium-only chips ($80$), $20$ cultures at rest ($42$), $3$ simulations ($38$). The mean rally length over the last $15$~min exceeds that over the first $5$ only in living cultures; non-neuronal cells (HEK293T) and medium-only chips show no effect & \cite{Kagan2022} & measured \\
10 & Significant growth within a session only with full feedback; with silence in place of the stimulus, play is better by the end of the session than without feedback, but the effect is smaller & $486$ sessions over $3$ days: ``stimulus'' ($n=27$), ``silence'' ($17$), ``no feedback'' ($15$). Growth over time is significant only in the ``stimulus'' condition; at the end of the session ``silence'' outperformed ``no feedback'' with a smaller effect & \cite{Kagan2022} & measured \\
11 & The effect is small; whether the network learns for the long term is an open question & Within a session the improvement is reliable, but, as the authors themselves say, learning across sessions over several days was not robustly observed & \cite{Kagan2022} & measured \\
12 & During play the network approaches a critical regime, and the closer it gets, the better the play & Analysis of activity avalanches in the same cultures: structured input brings the network closer to the critical regime, and play quality correlates with proximity to it; but criticality alone, without information about the consequences of actions, is not enough for learning & \cite{Habibollahi2023} & measured \\
13 & The culture's ``sentience'' has not been shown & A response article by thirty researchers: behavior change in a closed loop proves neither intelligence nor sentience; no experiment shows it & \cite{Balci2023} & hypothesis \\
14 & Proposed fourth control: a predictable stimulus after a miss, of the same duration and energy (section~\ref{sec:dishspec}) & No such experiment has been done; this is the book's proposal & --- & hypothesis \\
\bottomrule
\end{longtable}}


\begin{thebibliography}{99}
\bibitem{Azevedo2009} F. A. C. Azevedo et al., \emph{Equal numbers of neuronal and nonneuronal cells make the human brain an isometrically scaled-up primate brain}, J. Comp. Neurol. \textbf{513}, 532 (2009).
\bibitem{Schultz1997} W. Schultz, P. Dayan, and P. R. Montague, \emph{A neural substrate of prediction and reward}, Science \textbf{275}, 1593 (1997).
\bibitem{SuttonBarto2018} R. S. Sutton and A. G. Barto, \emph{Reinforcement Learning: An Introduction}, 2nd ed. (MIT Press, Cambridge, MA, 2018).
\bibitem{Doya2002} K. Doya, \emph{Metalearning and neuromodulation}, Neural Networks \textbf{15}, 495 (2002).
\bibitem{Strata1999} P. Strata and R. Harvey, \emph{Dale's principle}, Brain Res. Bull. \textbf{50}, 349 (1999).
\bibitem{vanVreeswijk1996} C. van Vreeswijk and H. Sompolinsky, \emph{Chaos in neuronal networks with balanced excitatory and inhibitory activity}, Science \textbf{274}, 1724 (1996).
\bibitem{AstonJones2005} G. Aston-Jones and J. D. Cohen, \emph{An integrative theory of locus coeruleus-norepinephrine function: adaptive gain and optimal performance}, Annu. Rev. Neurosci. \textbf{28}, 403 (2005).
\bibitem{Beaulieu2011} J.-M. Beaulieu and R. R. Gainetdinov, \emph{The physiology, signaling, and pharmacology of dopamine receptors}, Pharmacol. Rev. \textbf{63}, 182 (2011).
\bibitem{Sparso2001} J. Spars{\o} and S. Furber (eds.), \emph{Principles of Asynchronous Circuit Design: A Systems Perspective} (Kluwer, Boston, 2001).
\bibitem{Carr1990} C. E. Carr and M. Konishi, \emph{A circuit for detection of interaural time differences in the brain stem of the barn owl}, J. Neurosci. \textbf{10}, 3227 (1990).
\bibitem{Wang2001} X.-J. Wang, \emph{Synaptic reverberation underlying mnemonic persistent activity}, Trends Neurosci. \textbf{24}, 455 (2001).
\bibitem{Hahnloser2002} R. H. R. Hahnloser, A. A. Kozhevnikov, and M. S. Fee, \emph{An ultra-sparse code underlies the generation of neural sequences in a songbird}, Nature \textbf{419}, 65 (2002).
\bibitem{Barlow1965} H. B. Barlow and W. R. Levick, \emph{The mechanism of directionally selective units in rabbit's retina}, J. Physiol. \textbf{178}, 477 (1965).
\bibitem{Carandini2012} M. Carandini and D. J. Heeger, \emph{Normalization as a canonical neural computation}, Nat. Rev. Neurosci. \textbf{13}, 51 (2012).
\bibitem{Pouille2001} F. Pouille and M. Scanziani, \emph{Enforcement of temporal fidelity in pyramidal cells by somatic feed-forward inhibition}, Science \textbf{293}, 1159 (2001).
\bibitem{Tsodyks1997isn} M. V. Tsodyks, W. E. Skaggs, T. J. Sejnowski, and B. L. McNaughton, \emph{Paradoxical effects of external modulation of inhibitory interneurons}, J. Neurosci. \textbf{17}, 4382 (1997).
\bibitem{Buzsaki2012} G. Buzs\'aki and X.-J. Wang, \emph{Mechanisms of gamma oscillations}, Annu. Rev. Neurosci. \textbf{35}, 203 (2012).
\bibitem{Rieke1997} F. Rieke, D. Warland, R. de Ruyter van Steveninck, and W. Bialek, \emph{Spikes: Exploring the Neural Code} (MIT Press, Cambridge, MA, 1997).
\bibitem{Mainen1995} Z. F. Mainen and T. J. Sejnowski, \emph{Reliability of spike timing in neocortical neurons}, Science \textbf{268}, 1503 (1995).
\bibitem{Softky1993} W. R. Softky and C. Koch, \emph{The highly irregular firing of cortical cells is inconsistent with temporal integration of random EPSPs}, J. Neurosci. \textbf{13}, 334 (1993).
\bibitem{Thorpe1996} S. Thorpe, D. Fize, and C. Marlot, \emph{Speed of processing in the human visual system}, Nature \textbf{381}, 520 (1996).
\bibitem{Zohary1994} E. Zohary, M. N. Shadlen, and W. T. Newsome, \emph{Correlated neuronal discharge rate and its implications for psychophysical performance}, Nature \textbf{370}, 140 (1994).
\bibitem{MorenoBote2014} R. Moreno-Bote et al., \emph{Information-limiting correlations}, Nat. Neurosci. \textbf{17}, 1410 (2014).
\bibitem{Gollisch2008} T. Gollisch and M. Meister, \emph{Rapid neural coding in the retina with relative spike latencies}, Science \textbf{319}, 1108 (2008).
\bibitem{OKeefe1993} J. O'Keefe and M. L. Recce, \emph{Phase relationship between hippocampal place units and the EEG theta rhythm}, Hippocampus \textbf{3}, 317 (1993).
\bibitem{Destexhe2003} A. Destexhe, M. Rudolph, and D. Par\'e, \emph{The high-conductance state of neocortical neurons in vivo}, Nat. Rev. Neurosci. \textbf{4}, 739 (2003).
\bibitem{Tsodyks1997} M. V. Tsodyks and H. Markram, \emph{The neural code between neocortical pyramidal neurons depends on neurotransmitter release probability}, Proc. Natl. Acad. Sci. USA \textbf{94}, 719 (1997).
\bibitem{Lisman1997} J. E. Lisman, \emph{Bursts as a unit of neural information: making unreliable synapses reliable}, Trends Neurosci. \textbf{20}, 38 (1997).
\bibitem{Attwell2001} D. Attwell and S. B. Laughlin, \emph{An energy budget for signaling in the grey matter of the brain}, J. Cereb. Blood Flow Metab. \textbf{21}, 1133 (2001).
\bibitem{Lennie2003} P. Lennie, \emph{The cost of cortical computation}, Curr. Biol. \textbf{13}, 493 (2003).
\bibitem{Yagishita2014} S. Yagishita et al., \emph{A critical time window for dopamine actions on the structural plasticity of dendritic spines}, Science \textbf{345}, 1616 (2014).
\bibitem{Werfel2005} J. Werfel, X. Xie, and H. S. Seung, \emph{Learning curves for stochastic gradient descent in linear feedforward networks}, Neural Comput. \textbf{17}, 2699 (2005).
\bibitem{Doya2000} K. Doya, \emph{Reinforcement learning in continuous time and space}, Neural Comput. \textbf{12}, 219 (2000).
\bibitem{Oja1982} E. Oja, \emph{Simplified neuron model as a principal component analyzer}, J. Math. Biol. \textbf{15}, 267 (1982).
\bibitem{BiPoo1998} G.-Q. Bi and M.-M. Poo, \emph{Synaptic modifications in cultured hippocampal neurons: dependence on spike timing, synaptic strength, and postsynaptic cell type}, J. Neurosci. \textbf{18}, 10464 (1998).
\bibitem{Williams1992} R. J. Williams, \emph{Simple statistical gradient-following algorithms for connectionist reinforcement learning}, Mach. Learn. \textbf{8}, 229 (1992).
\bibitem{Bayer2005} H. M. Bayer and P. W. Glimcher, \emph{Midbrain dopamine neurons encode a quantitative reward prediction error signal}, Neuron \textbf{47}, 129 (2005).
\bibitem{Shen2008} W. Shen, M. Flajolet, P. Greengard, and D. J. Surmeier, \emph{Dichotomous dopaminergic control of striatal synaptic plasticity}, Science \textbf{321}, 848 (2008).
\bibitem{Dabney2020} W. Dabney, Z. Kurth-Nelson, N. Uchida, C. K. Starkweather, D. Hassabis, R. Munos, and M. Botvinick, \emph{A distributional code for value in dopamine-based reinforcement learning}, Nature \textbf{577}, 671 (2020).
\bibitem{Matsumoto2009} M. Matsumoto and O. Hikosaka, \emph{Two types of dopamine neuron distinctly convey positive and negative motivational signals}, Nature \textbf{459}, 837 (2009).
\bibitem{BrombergMartin2010} E. S. Bromberg-Martin, M. Matsumoto, and O. Hikosaka, \emph{Dopamine in motivational control: rewarding, aversive, and alerting}, Neuron \textbf{68}, 815 (2010).
\bibitem{Menegas2018} W. Menegas, K. Akiti, R. Amo, N. Uchida, and M. Watabe-Uchida, \emph{Dopamine neurons projecting to the posterior striatum reinforce avoidance of threatening stimuli}, Nat. Neurosci. \textbf{21}, 1421 (2018).
\bibitem{Niv2007} Y. Niv, N. D. Daw, D. Joel, and P. Dayan, \emph{Tonic dopamine: opportunity costs and the control of response vigor}, Psychopharmacology \textbf{191}, 507 (2007).
\bibitem{Mohebi2019} A. Mohebi et al., \emph{Dissociable dopamine dynamics for learning and motivation}, Nature \textbf{570}, 65 (2019).
\bibitem{Hamid2016} A. A. Hamid et al., \emph{Mesolimbic dopamine signals the value of work}, Nat. Neurosci. \textbf{19}, 117 (2016).
\bibitem{Kim2020} H. R. Kim, A. N. Malik et al., \emph{A unified framework for dopamine signals across timescales}, Cell \textbf{183}, 1600 (2020).
\bibitem{Jeong2022} H. Jeong et al., \emph{Mesolimbic dopamine release conveys causal associations}, Science \textbf{378}, eabq6740 (2022).
\bibitem{Amo2022} R. Amo et al., \emph{A gradual temporal shift of dopamine responses mirrors the progression of temporal difference error in machine learning}, Nat. Neurosci. \textbf{25}, 1082 (2022).
\bibitem{Takahashi2017} Y. K. Takahashi et al., \emph{Dopamine neurons respond to errors in the prediction of sensory features of expected rewards}, Neuron \textbf{95}, 1395 (2017).
\bibitem{Sharpe2017} M. J. Sharpe et al., \emph{Dopamine transients are sufficient and necessary for acquisition of model-based associations}, Nat. Neurosci. \textbf{20}, 735 (2017).
\bibitem{Gardner2018} M. P. H. Gardner, G. Schoenbaum, and S. J. Gershman, \emph{Rethinking dopamine as generalized prediction error}, Proc. R. Soc. B \textbf{285}, 20181645 (2018).
\bibitem{Engelhard2019} B. Engelhard et al., \emph{Specialized coding of sensory, motor and cognitive variables in VTA dopamine neurons}, Nature \textbf{570}, 509 (2019).
\bibitem{Clements1992} J. D. Clements, R. A. J. Lester, G. Tong, C. E. Jahr, and G. L. Westbrook, \emph{The time course of glutamate in the synaptic cleft}, Science \textbf{258}, 1498 (1992).
\bibitem{Hasselmo2006} M. E. Hasselmo, \emph{The role of acetylcholine in learning and memory}, Curr. Opin. Neurobiol. \textbf{16}, 710 (2006).
\bibitem{JacobsAzmitia1992} B. L. Jacobs and E. C. Azmitia, \emph{Structure and function of the brain serotonin system}, Physiol. Rev. \textbf{72}, 165 (1992).
\bibitem{Schiller1992} P. H. Schiller, \emph{The ON and OFF channels of the visual system}, Trends Neurosci. \textbf{15}, 86 (1992).
\bibitem{Grothe2010} B. Grothe, M. Pecka, and D. McAlpine, \emph{Mechanisms of sound localization in mammals}, Physiol. Rev. \textbf{90}, 983 (2010).
\bibitem{Markram2004} H. Markram et al., \emph{Interneurons of the neocortical inhibitory system}, Nat. Rev. Neurosci. \textbf{5}, 793 (2004).
\bibitem{Joshi2016} S. Joshi, Y. Li, R. M. Kalwani, and J. I. Gold, \emph{Relationships between pupil diameter and neuronal activity in the locus coeruleus, colliculi, and cingulate cortex}, Neuron \textbf{89}, 221 (2016).
\bibitem{ServanSchreiber1990} D. Servan-Schreiber, H. Printz, and J. D. Cohen, \emph{A network model of catecholamine effects: gain, signal-to-noise ratio, and behavior}, Science \textbf{249}, 892 (1990).
\bibitem{YuDayan2005} A. J. Yu and P. Dayan, \emph{Uncertainty, neuromodulation, and attention}, Neuron \textbf{46}, 681 (2005).
\bibitem{Nassar2012} M. R. Nassar et al., \emph{Rational regulation of learning dynamics by pupil-linked arousal systems}, Nat. Neurosci. \textbf{15}, 1040 (2012).
\bibitem{BouretSara2005} S. Bouret and S. J. Sara, \emph{Network reset: a simplified overarching theory of locus coeruleus noradrenaline function}, Trends Neurosci. \textbf{28}, 574 (2005).
\bibitem{Hebb1949} D. O. Hebb, \emph{The Organization of Behavior: A Neuropsychological Theory} (Wiley, New York, 1949).
\bibitem{MacKay1952} D. M. MacKay and W. S. McCulloch, \emph{The limiting information capacity of a neuronal link}, Bull. Math. Biophys. \textbf{14}, 127 (1952).
\bibitem{WilsonCowan1972} H. R. Wilson and J. D. Cowan, \emph{Excitatory and inhibitory interactions in localized populations of model neurons}, Biophys. J. \textbf{12}, 1 (1972).
\bibitem{AllenStevens1994} C. Allen and C. F. Stevens, \emph{An evaluation of causes for unreliability of synaptic transmission}, Proc. Natl. Acad. Sci. USA \textbf{91}, 10380 (1994).
\bibitem{Barlow1961} H. B. Barlow, \emph{Possible principles underlying the transformations of sensory messages}, in \emph{Sensory Communication}, ed. W. A. Rosenblith (MIT Press, Cambridge, MA, 1961), pp.~217--234.
\bibitem{Cardin2009} J. A. Cardin et al., \emph{Driving fast-spiking cells induces gamma rhythm and controls sensory responses}, Nature \textbf{459}, 663 (2009).
\bibitem{Lillicrap2020} T. P. Lillicrap, A. Santoro, L. Marris, C. J. Akerman, and G. Hinton, \emph{Backpropagation and the brain}, Nat. Rev. Neurosci. \textbf{21}, 335 (2020).
\bibitem{Fremaux2016} N. Fr\'emaux and W. Gerstner, \emph{Neuromodulated spike-timing-dependent plasticity, and theory of three-factor learning rules}, Front. Neural Circuits \textbf{9}, 85 (2016).
\bibitem{Haider2006} B. Haider, A. Duque, A. R. Hasenstaub, and D. A. McCormick, \emph{Neocortical network activity in vivo is generated through a dynamic balance of excitation and inhibition}, J. Neurosci. \textbf{26}, 4535 (2006).
\bibitem{HoltKoch1997} G. R. Holt and C. Koch, \emph{Shunting inhibition does not have a divisive effect on firing rates}, Neural Comput. \textbf{9}, 1001 (1997).
\bibitem{Chance2002} F. S. Chance, L. F. Abbott, and A. D. Reyes, \emph{Gain modulation from background synaptic input}, Neuron \textbf{35}, 773 (2002).
\bibitem{Ozeki2009} H. Ozeki, I. M. Finn, E. S. Schaffer, K. D. Miller, and D. Ferster, \emph{Inhibitory stabilization of the cortical network underlies visual surround suppression}, Neuron \textbf{62}, 578 (2009).
\bibitem{HasselmoBower1992} M. E. Hasselmo and J. M. Bower, \emph{Cholinergic suppression specific to intrinsic not afferent fiber synapses in rat piriform (olfactory) cortex}, J. Neurophysiol. \textbf{67}, 1222 (1992).
\bibitem{HasselmoAndersonBower1992} M. E. Hasselmo, B. P. Anderson, and J. M. Bower, \emph{Cholinergic modulation of cortical associative memory function}, J. Neurophysiol. \textbf{67}, 1230 (1992).
\bibitem{Hasselmo1999} M. E. Hasselmo, \emph{Neuromodulation: acetylcholine and memory consolidation}, Trends Cogn. Sci. \textbf{3}, 351 (1999).
\bibitem{Tsodyks1988} M. V. Tsodyks and M. V. Feigel'man, \emph{The enhanced storage capacity in neural networks with low activity level}, Europhys. Lett. \textbf{6}, 101 (1988).
\bibitem{KalmanBucy1961} R. E. Kalman and R. S. Bucy, \emph{New results in linear filtering and prediction theory}, J. Basic Eng. \textbf{83}, 95 (1961).
\bibitem{WilsonMcNaughton1994} M. A. Wilson and B. L. McNaughton, \emph{Reactivation of hippocampal ensemble memories during sleep}, Science \textbf{265}, 676 (1994).
\bibitem{BarnesSharp1999} N. M. Barnes and T. Sharp, \emph{A review of central 5-HT receptors and their function}, Neuropharmacology \textbf{38}, 1083 (1999).
\bibitem{Bunin1998} M. A. Bunin and R. M. Wightman, \emph{Quantitative evaluation of 5-hydroxytryptamine (serotonin) neuronal release and uptake: an investigation of extrasynaptic transmission}, J. Neurosci. \textbf{18}, 4854 (1998).
\bibitem{Sprouse1987} J. S. Sprouse and G. K. Aghajanian, \emph{Electrophysiological responses of serotoninergic dorsal raphe neurons to 5-HT$_{1A}$ and 5-HT$_{1B}$ agonists}, Synapse \textbf{1}, 3 (1987).
\bibitem{Sykova2008} E. Syková and C. Nicholson, \emph{Diffusion in brain extracellular space}, Physiol. Rev. \textbf{88}, 1277 (2008).
\bibitem{WilsonNicoll2001} R. I. Wilson and R. A. Nicoll, \emph{Endogenous cannabinoids mediate retrograde signalling at hippocampal synapses}, Nature \textbf{410}, 588 (2001).
\bibitem{Garthwaite2008} J. Garthwaite, \emph{Concepts of neural nitric oxide-mediated transmission}, Eur. J. Neurosci. \textbf{27}, 2783 (2008).
\bibitem{vandenPol2012} A. N. van den Pol, \emph{Neuropeptide transmission in brain circuits}, Neuron \textbf{76}, 98 (2012).
\bibitem{Dunwiddie2001} T. V. Dunwiddie and S. A. Masino, \emph{The role and regulation of adenosine in the central nervous system}, Annu. Rev. Neurosci. \textbf{24}, 31 (2001).
\bibitem{Skaggs1996} W. E. Skaggs and B. L. McNaughton, \emph{Replay of neuronal firing sequences in rat hippocampus during sleep following spatial experience}, Science \textbf{271}, 1870 (1996).
\bibitem{Baylor1979} D. A. Baylor, T. D. Lamb, and K.-W. Yau, \emph{Responses of retinal rods to single photons}, J. Physiol. \textbf{288}, 613 (1979).
\bibitem{Hudspeth1989} A. J. Hudspeth, \emph{How the ear's works work}, Nature \textbf{341}, 397 (1989).
\bibitem{NakaRushton1966} K. I. Naka and W. A. H. Rushton, \emph{S-potentials from colour units in the retina of fish (Cyprinidae)}, J. Physiol. \textbf{185}, 536 (1966).
\bibitem{Lichtsteiner2008} P. Lichtsteiner, C. Posch, and T. Delbruck, \emph{A $128\times128$ 120 dB 15 $\mu$s latency asynchronous temporal contrast vision sensor}, IEEE J. Solid-State Circuits \textbf{43}, 566 (2008).
\bibitem{Koch2006} K. Koch et al., \emph{How much the eye tells the brain}, Curr. Biol. \textbf{16}, 1428 (2006).
\bibitem{Robles2001} L. Robles and M. A. Ruggero, \emph{Mechanics of the mammalian cochlea}, Physiol. Rev. \textbf{81}, 1305 (2001).
\bibitem{Mombaerts2004} P. Mombaerts, \emph{Genes and ligands for odorant, vomeronasal and taste receptors}, Nat. Rev. Neurosci. \textbf{5}, 263 (2004).
\bibitem{WoodSlater2001} S. J. Wood and C. R. Slater, \emph{Safety factor at the neuromuscular junction}, Prog. Neurobiol. \textbf{64}, 393 (2001).
\bibitem{Henneman1957} E. Henneman, \emph{Relation between size of neurons and their susceptibility to discharge}, Science \textbf{126}, 1345 (1957).
\bibitem{HarrisWolpert1998} C. M. Harris and D. M. Wolpert, \emph{Signal-dependent noise determines motor planning}, Nature \textbf{394}, 780 (1998).
\bibitem{Hogan1984} N. Hogan, \emph{Adaptive control of mechanical impedance by coactivation of antagonist muscles}, IEEE Trans. Autom. Control \textbf{29}, 681 (1984).
\bibitem{McAuleyMarsden2000} J. H. McAuley and C. D. Marsden, \emph{Physiological and pathological tremors and rhythmic central motor control}, Brain \textbf{123}, 1545 (2000).
\bibitem{WolpertGhahramani2000} D. M. Wolpert and Z. Ghahramani, \emph{Computational principles of movement neuroscience}, Nat. Neurosci. \textbf{3}, 1212 (2000).
\bibitem{MarderBucher2001} E. Marder and D. Bucher, \emph{Central pattern generators and the control of rhythmic movements}, Curr. Biol. \textbf{11}, R986 (2001).
\bibitem{Miyazaki2014} K. W. Miyazaki et al., \emph{Optogenetic activation of dorsal raphe serotonin neurons enhances patience for future rewards}, Curr. Biol. \textbf{24}, 2033 (2014).
\bibitem{Fuglevand1993} A. J. Fuglevand, D. A. Winter, and A. E. Patla, \emph{Models of recruitment and rate coding organization in motor-unit pools}, J. Neurophysiol. \textbf{70}, 2470 (1993).
\bibitem{Curcio1990} C. A. Curcio, K. R. Sloan, R. E. Kalina, and A. E. Hendrickson, \emph{Human photoreceptor topography}, J. Comp. Neurol. \textbf{292}, 497 (1990).
\bibitem{Hayes1950} N. D. Hayes, \emph{Roots of the transcendental equation associated with a certain difference-differential equation}, J. London Math. Soc. \textbf{25}, 226 (1950).
\bibitem{MelzackWall1965} R. Melzack and P. D. Wall, \emph{Pain mechanisms: a new theory}, Science \textbf{150}, 971 (1965).
\bibitem{Fields2004} H. Fields, \emph{State-dependent opioid control of pain}, Nat. Rev. Neurosci. \textbf{5}, 565 (2004).
\bibitem{Levine1978} J. D. Levine, N. C. Gordon, and H. L. Fields, \emph{The mechanism of placebo analgesia}, Lancet \textbf{312}, 654 (1978).
\bibitem{Woolf1983} C. J. Woolf, \emph{Evidence for a central component of post-injury pain hypersensitivity}, Nature \textbf{306}, 686 (1983).
\bibitem{Seymour2004} B. Seymour et al., \emph{Temporal difference models describe higher-order learning in humans}, Nature \textbf{429}, 664 (2004).
\bibitem{EcclestonCrombez1999} C. Eccleston and G. Crombez, \emph{Pain demands attention: a cognitive-affective model of the interruptive function of pain}, Psychol. Bull. \textbf{125}, 356 (1999).
\bibitem{WidrowHoff1960} B. Widrow and M. E. Hoff, \emph{Adaptive switching circuits}, IRE WESCON Convention Record, part 4, 96 (1960).
\bibitem{Marr1969} D. Marr, \emph{A theory of cerebellar cortex}, J. Physiol. \textbf{202}, 437 (1969).
\bibitem{Albus1971} J. S. Albus, \emph{A theory of cerebellar function}, Math. Biosci. \textbf{10}, 25 (1971).
\bibitem{Cover1965} T. M. Cover, \emph{Geometrical and statistical properties of systems of linear inequalities with applications in pattern recognition}, IEEE Trans. Electron. Comput. \textbf{EC-14}, 326 (1965).
\bibitem{Ito2001} M. Ito, \emph{Cerebellar long-term depression: characterization, signal transduction, and functional roles}, Physiol. Rev. \textbf{81}, 1143 (2001).
\bibitem{Suvrathan2016} A. Suvrathan, H. L. Payne, and J. L. Raymond, \emph{Timing rules for synaptic plasticity matched to behavioral function}, Neuron \textbf{92}, 959 (2016).
\bibitem{Fujita1982} M. Fujita, \emph{Adaptive filter model of the cerebellum}, Biol. Cybern. \textbf{45}, 195 (1982).
\bibitem{Blakemore1998} S.-J. Blakemore, D. M. Wolpert, and C. D. Frith, \emph{Central cancellation of self-produced tickle sensation}, Nat. Neurosci. \textbf{1}, 635 (1998).
\bibitem{Smith2006} M. A. Smith, A. Ghazizadeh, and R. Shadmehr, \emph{Interacting adaptive processes with different timescales underlie short-term motor learning}, PLoS Biol. \textbf{4}, e179 (2006).
\bibitem{Kording2007} K. P. K\"ording, J. B. Tenenbaum, and R. Shadmehr, \emph{The dynamics of memory as a consequence of optimal adaptation to a changing body}, Nat. Neurosci. \textbf{10}, 779 (2007).
\bibitem{Yao2023} Z. Yao et al., \emph{A high-resolution transcriptomic and spatial atlas of cell types in the whole mouse brain}, Nature \textbf{624}, 317 (2023).
\bibitem{Sanzeni2020} A. Sanzeni, B. Akitake, H. C. Goldbach, C. E. Leedy, N. Brunel, and M. H. Histed, \emph{Inhibition stabilization is a widespread property of cortical networks}, eLife \textbf{9}, e54875 (2020).
\bibitem{Padamsey2022} Z. Padamsey, D. Katsanevaki, N. Dupuy, and N. L. Rochefort, \emph{Neocortex saves energy by reducing coding precision during food scarcity}, Neuron \textbf{110}, 280 (2022).
\bibitem{Stringer2019} C. Stringer, M. Pachitariu, N. Steinmetz, C. B. Reddy, M. Carandini, and K. D. Harris, \emph{Spontaneous behaviors drive multidimensional, brainwide activity}, Science \textbf{364}, eaav7893 (2019).
\bibitem{Bittner2017} K. C. Bittner, A. D. Milstein, C. Grienberger, S. Romani, and J. C. Magee, \emph{Behavioral time scale synaptic plasticity underlies CA1 place fields}, Science \textbf{357}, 1033 (2017).
\bibitem{WhittingtonBogacz2019} J. C. R. Whittington and R. Bogacz, \emph{Theories of error back-propagation in the brain}, Trends Cogn. Sci. \textbf{23}, 235 (2019).
\bibitem{Reimer2016} J. Reimer et al., \emph{Pupil fluctuations track rapid changes in adrenergic and cholinergic activity in cortex}, Nat. Commun. \textbf{7}, 13289 (2016).
\bibitem{Beniaguev2021} D. Beniaguev, I. Segev, and M. London, \emph{Single cortical neurons as deep artificial neural networks}, Neuron \textbf{109}, 2727 (2021).
\bibitem{Gidon2020} A. Gidon et al., \emph{Dendritic action potentials and computation in human layer 2/3 cortical neurons}, Science \textbf{367}, 83 (2020).
\bibitem{Gallego2022} G. Gallego et al., \emph{Event-based vision: a survey}, IEEE Trans. Pattern Anal. Mach. Intell. \textbf{44}, 154 (2022).
\bibitem{Berger1989} R. D. Berger, J. P. Saul, and R. J. Cohen, \emph{Transfer function analysis of autonomic regulation. I. Canine atrial rate response}, Am. J. Physiol. \textbf{256}, H142 (1989).
\bibitem{Walker2010} J. J. Walker, J. R. Terry, and S. L. Lightman, \emph{Origin of ultradian pulsatility in the hypothalamic--pituitary--adrenal axis}, Proc. R. Soc. B \textbf{277}, 1627 (2010).
\bibitem{Belchetz1978} P. E. Belchetz, T. M. Plant, Y. Nakai, E. J. Keogh, and E. Knobil, \emph{Hypophysial responses to continuous and intermittent delivery of hypothalamic gonadotropin-releasing hormone}, Science \textbf{202}, 631 (1978).
\bibitem{Czeisler1999} C. A. Czeisler et al., \emph{Stability, precision, and near-24-hour period of the human circadian pacemaker}, Science \textbf{284}, 2177 (1999).
\bibitem{TakahashiJS2017} J. S. Takahashi, \emph{Transcriptional architecture of the mammalian circadian clock}, Nat. Rev. Genet. \textbf{18}, 164 (2017).
\bibitem{Musk2019} E. Musk and Neuralink, \emph{An integrated brain-machine interface platform with thousands of channels}, J. Med. Internet Res. \textbf{21}, e16194 (2019).
\bibitem{Hochberg2006} L. R. Hochberg et al., \emph{Neuronal ensemble control of prosthetic devices by a human with tetraplegia}, Nature \textbf{442}, 164 (2006).
\bibitem{Willett2021} F. R. Willett et al., \emph{High-performance brain-to-text communication via handwriting}, Nature \textbf{593}, 249 (2021).
\bibitem{Willett2023} F. R. Willett et al., \emph{A high-performance speech neuroprosthesis}, Nature \textbf{620}, 1031 (2023).
\bibitem{Gallego2017} J. A. Gallego, M. G. Perich, L. E. Miller, and S. A. Solla, \emph{Neural manifolds for the control of movement}, Neuron \textbf{94}, 978 (2017).
\bibitem{Fernandez2021} E. Fern\'andez et al., \emph{Visual percepts evoked with an intracortical 96-channel microelectrode array inserted in human occipital cortex}, J. Clin. Invest. \textbf{131}, e151331 (2021).
\bibitem{Tritsch2012} N. X. Tritsch, J. B. Ding, and B. L. Sabatini, \emph{Dopaminergic neurons inhibit striatal output through non-canonical release of GABA}, Nature \textbf{490}, 262 (2012).
\bibitem{Zhang2015} S. Zhang et al., \emph{Dopaminergic and glutamatergic microdomains in a subset of rodent mesoaccumbens axons}, Nat. Neurosci. \textbf{18}, 386 (2015).
\bibitem{Mongillo2008} G. Mongillo, O. Barak, and M. Tsodyks, \emph{Synaptic theory of working memory}, Science \textbf{319}, 1543 (2008).
\bibitem{Stokes2015} M. G. Stokes, \emph{`Activity-silent' working memory in prefrontal cortex: a dynamic coding framework}, Trends Cogn. Sci. \textbf{19}, 394 (2015).
\bibitem{Smith2013} S. L. Smith, I. T. Smith, T. Branco, and M. H\"ausser, \emph{Dendritic spikes enhance stimulus selectivity in cortical neurons in vivo}, Nature \textbf{503}, 115 (2013).
\bibitem{Baden2016} T. Baden et al., \emph{The functional diversity of retinal ganglion cells in the mouse}, Nature \textbf{529}, 345 (2016).
\bibitem{Lottem2018} E. Lottem et al., \emph{Activation of serotonin neurons promotes active persistence in a probabilistic foraging task}, Nat. Commun. \textbf{9}, 1000 (2018).
\bibitem{Payeur2021} A. Payeur, J. Guerguiev, F. Zenke, B. A. Richards, and R. Naud, \emph{Burst-dependent synaptic plasticity can coordinate learning in hierarchical circuits}, Nat. Neurosci. \textbf{24}, 1010 (2021).
\bibitem{MehonicKenyon2022} A. Mehonic and A. J. Kenyon, \emph{Brain-inspired computing needs a master plan}, Nature \textbf{604}, 255 (2022).
\bibitem{Poe2020} G. R. Poe et al., \emph{Locus coeruleus: a new look at the blue spot}, Nat. Rev. Neurosci. \textbf{21}, 644 (2020).
\bibitem{Hangya2015} B. Hangya, S. P. Ranade, M. Lorenc, and A. Kepecs, \emph{Central cholinergic neurons are rapidly recruited by reinforcement feedback}, Cell \textbf{162}, 1155 (2015).
\bibitem{FernandezRuiz2019} A. Fern\'andez-Ruiz et al., \emph{Long-duration hippocampal sharp wave ripples improve memory}, Science \textbf{364}, 1082 (2019).
\bibitem{Herzfeld2018} D. J. Herzfeld, Y. Kojima, R. Soetedjo, and R. Shadmehr, \emph{Encoding of error and learning to correct that error by the Purkinje cells of the cerebellum}, Nat. Neurosci. \textbf{21}, 736 (2018).
\bibitem{Duan2014} B. Duan et al., \emph{Identification of spinal circuits transmitting and gating mechanical pain}, Cell \textbf{159}, 1417 (2014).
\bibitem{Tremblay2016} R. Tremblay, S. Lee, and B. Rudy, \emph{GABAergic interneurons in the neocortex: from cellular properties to circuits}, Neuron \textbf{91}, 260 (2016).
\bibitem{Ren2018} J. Ren et al., \emph{Anatomically defined and functionally distinct dorsal raphe serotonin sub-systems}, Cell \textbf{175}, 472 (2018).
\bibitem{Pfeffer2013} C. K. Pfeffer, M. Xue, M. He, Z. J. Huang, and M. Scanziani, \emph{Inhibition of inhibition in visual cortex: the logic of connections between molecularly distinct interneurons}, Nat. Neurosci. \textbf{16}, 1068 (2013).
\bibitem{Pi2013} H.-J. Pi et al., \emph{Cortical interneurons that specialize in disinhibitory control}, Nature \textbf{503}, 521 (2013).
\bibitem{Keller2018} G. B. Keller and T. D. Mrsic-Flogel, \emph{Predictive processing: a canonical cortical computation}, Neuron \textbf{100}, 424 (2018).
\bibitem{HubelWiesel1962} D. H. Hubel and T. N. Wiesel, \emph{Receptive fields, binocular interaction and functional architecture in the cat's visual cortex}, J. Physiol. \textbf{160}, 106 (1962).
\bibitem{Riesenhuber1999} M. Riesenhuber and T. Poggio, \emph{Hierarchical models of object recognition in cortex}, Nat. Neurosci. \textbf{2}, 1019 (1999).
\bibitem{Fukushima1980} K. Fukushima, \emph{Neocognitron: a self-organizing neural network model for a mechanism of pattern recognition unaffected by shift in position}, Biol. Cybern. \textbf{36}, 193 (1980).
\bibitem{Yamins2014} D. L. K. Yamins et al., \emph{Performance-optimized hierarchical models predict neural responses in higher visual cortex}, Proc. Natl. Acad. Sci. USA \textbf{111}, 8619 (2014).
\bibitem{Hung2005} C. P. Hung, G. Kreiman, T. Poggio, and J. J. DiCarlo, \emph{Fast readout of object identity from macaque inferior temporal cortex}, Science \textbf{310}, 863 (2005).
\bibitem{WiskottSejnowski2002} L. Wiskott and T. J. Sejnowski, \emph{Slow feature analysis: unsupervised learning of invariances}, Neural Comput. \textbf{14}, 715 (2002).
\bibitem{Foldiak1991} P. F\"oldi\'ak, \emph{Learning invariance from transformation sequences}, Neural Comput. \textbf{3}, 194 (1991).
\bibitem{LiDiCarlo2008} N. Li and J. J. DiCarlo, \emph{Unsupervised natural experience rapidly alters invariant object representation in visual cortex}, Science \textbf{321}, 1502 (2008).
\bibitem{RaoBallard1999} R. P. N. Rao and D. H. Ballard, \emph{Predictive coding in the visual cortex: a functional interpretation of some extra-classical receptive-field effects}, Nat. Neurosci. \textbf{2}, 79 (1999).
\bibitem{Kar2019} K. Kar, J. Kubilius, K. Schmidt, E. B. Issa, and J. J. DiCarlo, \emph{Evidence that recurrent circuits are critical to the ventral stream's execution of core object recognition behavior}, Nat. Neurosci. \textbf{22}, 974 (2019).
\bibitem{Szegedy2014} C. Szegedy et al., \emph{Intriguing properties of neural networks}, in \emph{Proc. ICLR} (2014), arXiv:1312.6199.
\bibitem{Weiss2002} Y. Weiss, E. P. Simoncelli, and E. H. Adelson, \emph{Motion illusions as optimal percepts}, Nat. Neurosci. \textbf{5}, 598 (2002).
\bibitem{Purves2011} D. Purves, W. T. Wojtach, and R. B. Lotto, \emph{Understanding vision in wholly empirical terms}, Proc. Natl. Acad. Sci. USA \textbf{108}, Suppl.~3, 15588 (2011).
\bibitem{LaferSousa2015} R. Lafer-Sousa, K. L. Hermann, and B. R. Conway, \emph{Striking individual differences in color perception uncovered by `the dress' photograph}, Curr. Biol. \textbf{25}, R545 (2015).
\bibitem{Wallisch2017} P. Wallisch, \emph{Illumination assumptions account for individual differences in the perceptual interpretation of a profoundly ambiguous stimulus in the color domain: ``The dress''}, J. Vis. \textbf{17}(4), 5 (2017).
\bibitem{Schwartz2009} O. Schwartz, T. J. Sejnowski, and P. Dayan, \emph{Perceptual organization in the tilt illusion}, J. Vis. \textbf{9}(4), 19 (2009).
\bibitem{SchillerCarvey2005} P. H. Schiller and C. E. Carvey, \emph{The Hermann grid illusion revisited}, Perception \textbf{34}, 1375 (2005).
\bibitem{Nijhawan1994} R. Nijhawan, \emph{Motion extrapolation in catching}, Nature \textbf{370}, 256 (1994).
\bibitem{Conway2005} B. R. Conway, A. Kitaoka, A. Yazdanbakhsh, C. C. Pack, and M. S. Livingstone, \emph{Neural basis for a powerful static motion illusion}, J. Neurosci. \textbf{25}, 5651 (2005).
\bibitem{OteroMillan2012} J. Otero-Millan, S. L. Macknik, and S. Martinez-Conde, \emph{Microsaccades and blinks trigger illusory rotation in the ``rotating snakes'' illusion}, J. Neurosci. \textbf{32}, 6043 (2012).
\bibitem{MorenoBote2007} R. Moreno-Bote, J. Rinzel, and N. Rubin, \emph{Noise-induced alternations in an attractor network model of perceptual bistability}, J. Neurophysiol. \textbf{98}, 1125 (2007).
\bibitem{Watanabe2018} E. Watanabe, A. Kitaoka, K. Sakamoto, M. Yasugi, and K. Tanaka, \emph{Illusory motion reproduced by deep neural networks trained for prediction}, Front. Psychol. \textbf{9}, 345 (2018).
\bibitem{GomezVilla2020} A. G\'omez-Villa, A. Mart\'in, J. Vazquez-Corral, M. Bertalm\'io, and J. Malo, \emph{Color illusions also deceive CNNs for low-level vision tasks: analysis and implications}, Vision Res. \textbf{176}, 156 (2020).
\bibitem{delCastillo1954} J. del Castillo and B. Katz, \emph{Quantal components of the end-plate potential}, J. Physiol. \textbf{124}, 560 (1954).
\bibitem{Nowak1984} L. Nowak, P. Bregestovski, P. Ascher, A. Herbet, and A. Prochiantz, \emph{Magnesium gates glutamate-activated channels in mouse central neurones}, Nature \textbf{307}, 462 (1984).
\bibitem{Malinow2002} R. Malinow and R. C. Malenka, \emph{AMPA receptor trafficking and synaptic plasticity}, Annu. Rev. Neurosci. \textbf{25}, 103 (2002).
\bibitem{Matsuzaki2004} M. Matsuzaki, N. Honkura, G. C. R. Ellis-Davies, and H. Kasai, \emph{Structural basis of long-term potentiation in single dendritic spines}, Nature \textbf{429}, 761 (2004).
\bibitem{Rudomin1999} P. Rudomin and R. F. Schmidt, \emph{Presynaptic inhibition in the vertebrate spinal cord revisited}, Exp. Brain Res. \textbf{129}, 1 (1999).
\bibitem{HutcheonYarom2000} B. Hutcheon and Y. Yarom, \emph{Resonance, oscillation and the intrinsic frequency preferences of neurons}, Trends Neurosci. \textbf{23}, 216 (2000).
\bibitem{Seung1996} H. S. Seung, \emph{How the brain keeps the eyes still}, Proc. Natl. Acad. Sci. USA \textbf{93}, 13339 (1996).
\bibitem{Hounsgaard1988} J. Hounsgaard, H. Hultborn, B. Jespersen, and O. Kiehn, \emph{Bistability of alpha-motoneurones in the decerebrate cat and in the acute spinal cat after intravenous 5-hydroxytryptophan}, J. Physiol. \textbf{405}, 345 (1988).
\bibitem{Izhikevich2007} E. M. Izhikevich, \emph{Dynamical Systems in Neuroscience: The Geometry of Excitability and Bursting} (MIT Press, Cambridge, MA, 2007).
\bibitem{AgmonSnir1998} H. Agmon-Snir, C. E. Carr, and J. Rinzel, \emph{The role of dendrites in auditory coincidence detection}, Nature \textbf{393}, 268 (1998).
\bibitem{BorstSoria2012} J. G. G. Borst and J. Soria van Hoeve, \emph{The calyx of Held synapse: from model synapse to auditory relay}, Annu. Rev. Physiol. \textbf{74}, 199 (2012).
\bibitem{Euler2002} T. Euler, P. B. Detwiler, and W. Denk, \emph{Directionally selective calcium signals in dendrites of starburst amacrine cells}, Nature \textbf{418}, 845 (2002).
\bibitem{AstonJonesBloom1981} G. Aston-Jones and F. E. Bloom, \emph{Activity of norepinephrine-containing locus coeruleus neurons in behaving rats anticipates fluctuations in the sleep-waking cycle}, J. Neurosci. \textbf{1}, 876 (1981).
\bibitem{BrooksPeever2012} P. L. Brooks and J. H. Peever, \emph{Identification of the transmitter and receptor mechanisms responsible for REM sleep paralysis}, J. Neurosci. \textbf{32}, 9785 (2012).
\bibitem{Borbely1982} A. A. Borb\'ely, \emph{A two process model of sleep regulation}, Hum. Neurobiol. \textbf{1}, 195 (1982).
\bibitem{PorkkaHeiskanen1997} T. Porkka-Heiskanen et al., \emph{Adenosine: a mediator of the sleep-inducing effects of prolonged wakefulness}, Science \textbf{276}, 1265 (1997).
\bibitem{Steriade1993} M. Steriade, A. N\'u\~nez, and F. Amzica, \emph{A novel slow ($<1$ Hz) oscillation of neocortical neurons in vivo: depolarizing and hyperpolarizing components}, J. Neurosci. \textbf{13}, 3252 (1993).
\bibitem{McCormick1992} D. A. McCormick, \emph{Neurotransmitter actions in the thalamus and cerebral cortex and their role in neuromodulation of thalamocortical activity}, Prog. Neurobiol. \textbf{39}, 337 (1992).
\bibitem{SteriadeMcCormickSejnowski1993} M. Steriade, D. A. McCormick, and T. J. Sejnowski, \emph{Thalamocortical oscillations in the sleeping and aroused brain}, Science \textbf{262}, 679 (1993).
\bibitem{Nadasdy1999} Z. N\'adasdy et al., \emph{Replay and time compression of recurring spike sequences in the hippocampus}, J. Neurosci. \textbf{19}, 9497 (1999).
\bibitem{Maingret2016} N. Maingret et al., \emph{Hippocampo-cortical coupling mediates memory consolidation during sleep}, Nat. Neurosci. \textbf{19}, 959 (2016).
\bibitem{McClelland1995} J. L. McClelland, B. L. McNaughton, and R. C. O'Reilly, \emph{Why there are complementary learning systems in the hippocampus and neocortex}, Psychol. Rev. \textbf{102}, 419 (1995).
\bibitem{TononiCirelli2014} G. Tononi and C. Cirelli, \emph{Sleep and the price of plasticity: from synaptic and cellular homeostasis to memory consolidation and integration}, Neuron \textbf{81}, 12 (2014).
\bibitem{deVivo2017} L. de Vivo et al., \emph{Ultrastructural evidence for synaptic scaling across the wake/sleep cycle}, Science \textbf{355}, 507 (2017).
\bibitem{Xie2013} L. Xie et al., \emph{Sleep drives metabolite clearance from the adult brain}, Science \textbf{342}, 373 (2013).
\bibitem{Miao2024} A. Miao et al., \emph{Brain clearance is reduced during sleep and anesthesia}, Nat. Neurosci. \textbf{27}, 1046 (2024).
\bibitem{Vyazovskiy2011} V. V. Vyazovskiy et al., \emph{Local sleep in awake rats}, Nature \textbf{472}, 443 (2011).
\bibitem{Lancaster2013} M. A. Lancaster et al., \emph{Cerebral organoids model human brain development and microcephaly}, Nature \textbf{501}, 373 (2013).
\bibitem{Jordan2024} F. D. Jordan et al., \emph{Open and remotely accessible Neuroplatform for research in wetware computing}, Front. Artif. Intell. \textbf{7}, 1376042 (2024).
\bibitem{Smirnova2023} L. Smirnova et al., \emph{Organoid intelligence (OI): the new frontier in biocomputing and intelligence-in-a-dish}, Front. Sci. \textbf{1}, 1017235 (2023).
\bibitem{Wagenaar2006} D. A. Wagenaar, J. Pine, and S. M. Potter, \emph{An extremely rich repertoire of bursting patterns during the development of cortical cultures}, BMC Neurosci. \textbf{7}, 11 (2006).
\bibitem{Shahaf2001} G. Shahaf and S. Marom, \emph{Learning in networks of cortical neurons}, J. Neurosci. \textbf{21}, 8782 (2001).
\bibitem{Kagan2022} B. J. Kagan et al., \emph{In vitro neurons learn and exhibit sentience when embodied in a simulated game-world}, Neuron \textbf{110}, 3952 (2022).
\bibitem{Friston2010} K. Friston, \emph{The free-energy principle: a unified brain theory?}, Nat. Rev. Neurosci. \textbf{11}, 127 (2010).
\bibitem{Balci2023} F. Balci et al., \emph{A response to claims of emergent intelligence and sentience in a dish}, Neuron \textbf{111}, 604 (2023).
\bibitem{Bakkum2008} D. J. Bakkum, Z. C. Chao, and S. M. Potter, \emph{Spatio-temporal electrical stimuli shape behavior of an embodied cortical network in a goal-directed learning task}, J. Neural Eng. \textbf{5}, 310 (2008).
\bibitem{Maass2002} W. Maass, T. Natschl\"ager, and H. Markram, \emph{Real-time computing without stable states: a new framework for neural computation based on perturbations}, Neural Comput. \textbf{14}, 2531 (2002).
\bibitem{JaegerHaas2004} H. Jaeger and H. Haas, \emph{Harnessing nonlinearity: predicting chaotic systems and saving energy in wireless communication}, Science \textbf{304}, 78 (2004).
\bibitem{Cai2023} H. Cai et al., \emph{Brain organoid reservoir computing for artificial intelligence}, Nat. Electron. \textbf{6}, 1032 (2023).
\bibitem{Habibollahi2023} F. Habibollahi, B. J. Kagan, A. N. Burkitt, and C. French, \emph{Critical dynamics arise during structured information presentation within embodied in vitro neuronal networks}, Nat. Commun. \textbf{14}, 5287 (2023).
\bibitem{Robbins2026} A. Robbins et al., \emph{Goal-directed learning in cortical organoids}, Cell Rep. \textbf{45}, 116984 (2026).
\bibitem{Keene2020} S. T. Keene et al., \emph{A biohybrid synapse with neurotransmitter-mediated plasticity}, Nat. Mater. \textbf{19}, 969 (2020).
\bibitem{Jo2016} J. Jo et al., \emph{Midbrain-like organoids from human pluripotent stem cells contain functional dopaminergic and neuromelanin-producing neurons}, Cell Stem Cell \textbf{19}, 248 (2016).
\bibitem{Hendry1987} S. H. Hendry, H. D. Schwark, E. G. Jones, and J. Yan, \emph{Numbers and proportions of GABA-immunoreactive neurons in different areas of monkey cerebral cortex}, J. Neurosci. \textbf{7}, 1503 (1987).
\bibitem{Tang2001synapses} Y. Tang, J. R. Nyengaard, D. M. G. De Groot, and H. J. G. Gundersen, \emph{Total regional and global number of synapses in the human brain neocortex}, Synapse \textbf{41}, 258 (2001).
\bibitem{Pakkenberg1997} B. Pakkenberg and H. J. G. Gundersen, \emph{Neocortical neuron number in humans: effect of sex and age}, J. Comp. Neurol. \textbf{384}, 312 (1997).
\bibitem{Matsuda2009} W. Matsuda, T. Furuta, K. C. Nakamura, H. Hioki, F. Fujiyama, R. Arai, and T. Kaneko, \emph{Single nigrostriatal dopaminergic neurons form widely spread and highly dense axonal arborizations in the neostriatum}, J. Neurosci. \textbf{29}, 444 (2009).
\bibitem{Pakkenberg1991} B. Pakkenberg, A. M{\o}ller, H. J. G. Gundersen, A. Mouritzen Dam, and H. Pakkenberg, \emph{The absolute number of nerve cells in substantia nigra in normal subjects and in patients with Parkinson's disease estimated with an unbiased stereological method}, J. Neurol. Neurosurg. Psychiatry \textbf{54}, 30 (1991).
\bibitem{Baker1990} K. G. Baker, G. M. Halliday, and I. T\"ork, \emph{Cytoarchitecture of the human dorsal raphe nucleus}, J. Comp. Neurol. \textbf{301}, 147 (1990).
\bibitem{Baker1991} K. G. Baker, G. M. Halliday, P. Halasz, J.-P. Hornung, L. B. Geffen, R. G. H. Cotton, and I. T\"ork, \emph{Cytoarchitecture of serotonin-synthesizing neurons in the pontine tegmentum of the human brain}, Synapse \textbf{7}, 301 (1991).
\bibitem{Airaksinen1991} M. S. Airaksinen, A. Paetau, L. Palj\"arvi, K. Reinikainen, P. Riekkinen, R. Suomalainen, and P. Panula, \emph{Histamine neurons in human hypothalamus: anatomy in normal and Alzheimer diseased brains}, Neuroscience \textbf{44}, 465 (1991).
\bibitem{Colquhoun1992} D. Colquhoun, P. Jonas, and B. Sakmann, \emph{Action of brief pulses of glutamate on AMPA/kainate receptors in patches from different neurones of rat hippocampal slices}, J. Physiol. \textbf{458}, 261 (1992).
\bibitem{DutarNicoll1988} P. Dutar and R. A. Nicoll, \emph{A physiological role for GABA$_B$ receptors in the central nervous system}, Nature \textbf{332}, 156 (1988).

\bibitem{Rauch2003} A. Rauch, G. La Camera, H.-R. L\"uscher, W. Senn, and S. Fusi, \emph{Neocortical pyramidal cells respond as integrate-and-fire neurons to in vivo-like input currents}, J. Neurophysiol. \textbf{90}, 1598 (2003).
\bibitem{KlumppEady1956} R. G. Klumpp and H. R. Eady, \emph{Some measurements of interaural time difference thresholds}, J. Acoust. Soc. Am. \textbf{28}, 859 (1956).
\bibitem{Brand2002} A. Brand, O. Behrend, T. Marquardt, D. McAlpine, and B. Grothe, \emph{Precise inhibition is essential for microsecond interaural time difference coding}, Nature \textbf{417}, 543 (2002).
\bibitem{Funahashi1989} S. Funahashi, C. J. Bruce, and P. S. Goldman-Rakic, \emph{Mnemonic coding of visual space in the monkey's dorsolateral prefrontal cortex}, J. Neurophysiol. \textbf{61}, 331 (1989).
\bibitem{WangMarkram2006} Y. Wang, H. Markram, P. H. Goodman, T. K. Berger, J. Ma, and P. S. Goldman-Rakic, \emph{Heterogeneity in the pyramidal network of the medial prefrontal cortex}, Nat. Neurosci. \textbf{9}, 534 (2006).

\bibitem{McCullochPitts1943} W. S. McCulloch and W. Pitts, \emph{A logical calculus of the ideas immanent in nervous activity}, Bull. Math. Biophys. \textbf{5}, 115 (1943).
\bibitem{Fried2002} S. I. Fried, T. A. M\"unch, and F. S. Werblin, \emph{Mechanisms and circuitry underlying directional selectivity in the retina}, Nature \textbf{420}, 411 (2002).

\bibitem{Hromadka2008} T. Hrom\'adka, M. R. DeWeese, and A. M. Zador, \emph{Sparse representation of sounds in the unanesthetized auditory cortex}, PLoS Biol. \textbf{6}, e16 (2008).
\bibitem{Abbott1997depression} L. F. Abbott, J. A. Varela, K. Sen, and S. B. Nelson, \emph{Synaptic depression and cortical gain control}, Science \textbf{275}, 221 (1997).
\bibitem{ShadlenNewsome1998} M. N. Shadlen and W. T. Newsome, \emph{The variable discharge of cortical neurons: implications for connectivity, computation, and information coding}, J. Neurosci. \textbf{18}, 3870 (1998).

\bibitem{TrulsonJacobs1979} M. E. Trulson and B. L. Jacobs, \emph{Raphe unit activity in freely moving cats: correlation with level of behavioral arousal}, Brain Res. \textbf{163}, 135 (1979).
\bibitem{GagnonParent2014} D. Gagnon and M. Parent, \emph{Distribution of VGLUT3 in highly collateralized axons from the rat dorsal raphe nucleus as revealed by single-neuron reconstructions}, PLoS One \textbf{9}, e87709 (2014).
\bibitem{VandermaelenAghajanian1983} C. P. Vandermaelen and G. K. Aghajanian, \emph{Electrophysiological and pharmacological characterization of serotonergic dorsal raphe neurons recorded extracellularly and intracellularly in rat brain slices}, Brain Res. \textbf{289}, 109 (1983).
\bibitem{CourtneyFord2016} N. A. Courtney and C. P. Ford, \emph{Mechanisms of 5-HT$_{1A}$ receptor-mediated transmission in dorsal raphe serotonin neurons}, J. Physiol. \textbf{594}, 953 (2016).
\bibitem{Hashemi2012serotonin} P. Hashemi et al., \emph{Brain dopamine and serotonin differ in regulation and its consequences}, Proc. Natl. Acad. Sci. USA \textbf{109}, 11510 (2012).
\bibitem{Black1934feedback} H. S. Black, \emph{Stabilized feedback amplifiers}, Bell Syst. Tech. J. \textbf{13}, 1 (1934).
\bibitem{KobayashiSchultz2008} S. Kobayashi and W. Schultz, \emph{Influence of reward delays on responses of dopamine neurons}, J. Neurosci. \textbf{28}, 7837 (2008).
\bibitem{Schweighofer2008} N. Schweighofer et al., \emph{Low-serotonin levels increase delayed reward discounting in humans}, J. Neurosci. \textbf{28}, 4528 (2008).

\bibitem{Malenka1988calcium} R. C. Malenka, J. A. Kauer, R. S. Zucker, and R. A. Nicoll, \emph{Postsynaptic calcium is sufficient for potentiation of hippocampal synaptic transmission}, Science \textbf{242}, 81 (1988).
\bibitem{Bliss1973LTP} T. V. P. Bliss and T. L{\o}mo, \emph{Long-lasting potentiation of synaptic transmission in the dentate area of the anaesthetized rabbit following stimulation of the perforant path}, J. Physiol. \textbf{232}, 331 (1973).
\bibitem{Kelso1986hebb} S. R. Kelso, A. H. Ganong, and T. H. Brown, \emph{Hebbian synapses in hippocampus}, Proc. Natl. Acad. Sci. USA \textbf{83}, 5326 (1986).
\bibitem{He2015eligibility} K. He et al., \emph{Distinct eligibility traces for LTP and LTD in cortical synapses}, Neuron \textbf{88}, 528 (2015).
\bibitem{Olveczky2005} B. P. \"Olveczky, A. S. Andalman, and M. S. Fee, \emph{Vocal experimentation in the juvenile songbird requires a basal ganglia circuit}, PLoS Biol. \textbf{3}, e153 (2005).
\bibitem{TumerBrainard2007} E. C. Tumer and M. S. Brainard, \emph{Performance variability enables adaptive plasticity of `crystallized' adult birdsong}, Nature \textbf{450}, 1240 (2007).
\bibitem{Eshel2015arithmetic} N. Eshel et al., \emph{Arithmetic and local circuitry underlying dopamine prediction errors}, Nature \textbf{525}, 243 (2015).
\bibitem{GraceOnn1989} A. A. Grace and S.-P. Onn, \emph{Morphology and electrophysiological properties of immunocytochemically identified rat dopamine neurons recorded in vitro}, J. Neurosci. \textbf{9}, 3463 (1989).
\bibitem{ODoherty2004actor} J. O'Doherty et al., \emph{Dissociable roles of ventral and dorsal striatum in instrumental conditioning}, Science \textbf{304}, 452 (2004).

\bibitem{NeweyPowell1987} W. K. Newey and J. L. Powell, \emph{Asymmetric least squares estimation and testing}, Econometrica \textbf{55}, 819 (1987).
\bibitem{Beierholm2013vigor} U. Beierholm et al., \emph{Dopamine modulates reward-related vigor}, Neuropsychopharmacology \textbf{38}, 1495 (2013).
\bibitem{Howe2013ramp} M. W. Howe, P. L. Tierney, S. G. Sandberg, P. E. M. Phillips, and A. M. Graybiel, \emph{Prolonged dopamine signalling in striatum signals proximity and value of distant rewards}, Nature \textbf{500}, 575 (2013).

\bibitem{Baker1989LC} K. G. Baker, I. T\"ork, J.-P. Hornung, and P. Halasz, \emph{The human locus coeruleus complex: an immunohistochemical and three dimensional reconstruction study}, Exp. Brain Res. \textbf{77}, 257 (1989).
\bibitem{Schwarz2015LC} L. A. Schwarz et al., \emph{Viral-genetic tracing of the input--output organization of a central noradrenaline circuit}, Nature \textbf{524}, 88 (2015).
\bibitem{Uematsu2017LC} A. Uematsu et al., \emph{Modular organization of the brainstem noradrenaline system coordinates opposing learning states}, Nat. Neurosci. \textbf{20}, 1602 (2017).
\bibitem{AstonJones1994vigilance} G. Aston-Jones, J. Rajkowski, P. Kubiak, and T. Alexinsky, \emph{Locus coeruleus neurons in monkey are selectively activated by attended cues in a vigilance task}, J. Neurosci. \textbf{14}, 4467 (1994).
\bibitem{Clayton2004LC} E. C. Clayton, J. Rajkowski, J. D. Cohen, and G. Aston-Jones, \emph{Phasic activation of monkey locus ceruleus neurons by simple decisions in a forced-choice task}, J. Neurosci. \textbf{24}, 9914 (2004).
\bibitem{Foote1975NE} S. L. Foote, R. Freedman, and A. P. Oliver, \emph{Effects of putative neurotransmitters on neuronal activity in monkey auditory cortex}, Brain Res. \textbf{86}, 229 (1975).
\bibitem{JepmaNieuwenhuis2011} M. Jepma and S. Nieuwenhuis, \emph{Pupil diameter predicts changes in the exploration--exploitation trade-off: evidence for the adaptive gain theory}, J. Cogn. Neurosci. \textbf{23}, 1587 (2011).
\bibitem{Tervo2014stochastic} D. G. R. Tervo et al., \emph{Behavioral variability through stochastic choice and its gating by anterior cingulate cortex}, Cell \textbf{159}, 21 (2014).
\bibitem{Usher1999LC} M. Usher, J. D. Cohen, D. Servan-Schreiber, J. Rajkowski, and G. Aston-Jones, \emph{The role of locus coeruleus in the regulation of cognitive performance}, Science \textbf{283}, 549 (1999).

\bibitem{Mesulam1983} M.-M. Mesulam, E. J. Mufson, B. H. Wainer, and A. I. Levey, \emph{Central cholinergic pathways in the rat: an overview based on an alternative nomenclature (Ch1--Ch6)}, Neuroscience \textbf{10}, 1185 (1983).
\bibitem{Marrosu1995} F. Marrosu et al., \emph{Microdialysis measurement of cortical and hippocampal acetylcholine release during sleep-wake cycle in freely moving cats}, Brain Res. \textbf{671}, 329 (1995).
\bibitem{Picciotto2012} M. R. Picciotto, M. J. Higley, and Y. S. Mineur, \emph{Acetylcholine as a neuromodulator: cholinergic signaling shapes nervous system function and behavior}, Neuron \textbf{76}, 116 (2012).
\bibitem{HasselmoSchnell1994} M. E. Hasselmo and E. Schnell, \emph{Laminar selectivity of the cholinergic suppression of synaptic transmission in rat hippocampal region CA1: computational modeling and brain slice physiology}, J. Neurosci. \textbf{14}, 3898 (1994).
\bibitem{Gil1997} Z. Gil, B. W. Connors, and Y. Amitai, \emph{Differential regulation of neocortical synapses by neuromodulators and activity}, Neuron \textbf{19}, 679 (1997).
\bibitem{KilgardMerzenich1998} M. P. Kilgard and M. M. Merzenich, \emph{Cortical map reorganization enabled by nucleus basalis activity}, Science \textbf{279}, 1714 (1998).
\bibitem{Atri2004} A. Atri et al., \emph{Blockade of central cholinergic receptors impairs new learning and increases proactive interference in a word paired-associate memory task}, Behav. Neurosci. \textbf{118}, 223 (2004).
\bibitem{Girardeau2009} G. Girardeau, K. Benchenane, S. I. Wiener, G. Buzs\'aki, and M. B. Zugaro, \emph{Selective suppression of hippocampal ripples impairs spatial memory}, Nat. Neurosci. \textbf{12}, 1222 (2009).

\input{supplement/bib_10}
\bibitem{CopenhagenJahr1989} D. R. Copenhagen and C. E. Jahr, \emph{Release of endogenous excitatory amino acids from turtle photoreceptors}, Nature \textbf{341}, 536 (1989).
\bibitem{GlowatzkiFuchs2002} E. Glowatzki and P. A. Fuchs, \emph{Transmitter release at the hair cell ribbon synapse}, Nat. Neurosci. \textbf{5}, 147 (2002).
\bibitem{Tinsley2016} J. N. Tinsley et al., \emph{Direct detection of a single photon by humans}, Nat. Commun. \textbf{7}, 12172 (2016).
\bibitem{Sellick1982} P. M. Sellick, R. Patuzzi, and B. M. Johnstone, \emph{Measurement of basilar membrane motion in the guinea pig using the M\"ossbauer technique}, J. Acoust. Soc. Am. \textbf{72}, 131 (1982).
\bibitem{Ruggero1997} M. A. Ruggero, N. C. Rich, A. Recio, S. S. Narayan, and L. Robles, \emph{Basilar-membrane responses to tones at the base of the chinchilla cochlea}, J. Acoust. Soc. Am. \textbf{101}, 2151 (1997).
\bibitem{NormannPerlman1979} R. A. Normann and I. Perlman, \emph{The effects of background illumination on the photoresponses of red and green cones}, J. Physiol. \textbf{286}, 491 (1979).
\bibitem{Laughlin1981} S. Laughlin, \emph{A simple coding procedure enhances a neuron's information capacity}, Z. Naturforsch. C \textbf{36}, 910 (1981).
\bibitem{CurcioAllen1990} C. A. Curcio and K. A. Allen, \emph{Topography of ganglion cells in human retina}, J. Comp. Neurol. \textbf{300}, 5 (1990).
\bibitem{Greenwood1990} D. D. Greenwood, \emph{A cochlear frequency-position function for several species --- 29 years later}, J. Acoust. Soc. Am. \textbf{87}, 2592 (1990).
\bibitem{Eguiluz2000} V. M. Egu\'iluz, M. Ospeck, Y. Choe, A. J. Hudspeth, and M. O. Magnasco, \emph{Essential nonlinearities in hearing}, Phys. Rev. Lett. \textbf{84}, 5232 (2000).
\bibitem{PalmerRussell1986} A. R. Palmer and I. J. Russell, \emph{Phase-locking in the cochlear nerve of the guinea-pig and its relation to the receptor potential of inner hair-cells}, Hear. Res. \textbf{24}, 1 (1986).
\bibitem{Malnic1999} B. Malnic, J. Hirono, T. Sato, and L. B. Buck, \emph{Combinatorial receptor codes for odors}, Cell \textbf{96}, 713 (1999).
\bibitem{Finger2005} T. E. Finger et al., \emph{ATP signaling is crucial for communication from taste buds to gustatory nerves}, Science \textbf{310}, 1495 (2005).
\bibitem{Huang2005} Y.-J. Huang et al., \emph{Mouse taste buds use serotonin as a neurotransmitter}, J. Neurosci. \textbf{25}, 843 (2005).
\bibitem{JohanssonVallbo1979} R. S. Johansson and A. B. Vallbo, \emph{Tactile sensibility in the human hand: relative and absolute densities of four types of mechanoreceptive units in glabrous skin}, J. Physiol. \textbf{286}, 283 (1979).
\bibitem{McKemy2002} D. D. McKemy, W. M. Neuhausser, and D. Julius, \emph{Identification of a cold receptor reveals a general role for TRP channels in thermosensation}, Nature \textbf{416}, 52 (2002).

\bibitem{Schmolesky1998} M. T. Schmolesky et al., \emph{Signal timing across the macaque visual system}, J. Neurophysiol. \textbf{79}, 3272 (1998).
\bibitem{Lampl2004} I. Lampl, D. Ferster, T. Poggio, and M. Riesenhuber, \emph{Intracellular measurements of spatial integration and the MAX operation in complex cells of the cat primary visual cortex}, J. Neurophysiol. \textbf{92}, 2704 (2004).
\bibitem{KobatakeTanaka1994} E. Kobatake and K. Tanaka, \emph{Neuronal selectivities to complex object features in the ventral visual pathway of the macaque cerebral cortex}, J. Neurophysiol. \textbf{71}, 856 (1994).
\bibitem{Albright1984} T. D. Albright, \emph{Direction and orientation selectivity of neurons in visual area MT of the macaque}, J. Neurophysiol. \textbf{52}, 1106 (1984).
\bibitem{Goodfellow2015} I. J. Goodfellow, J. Shlens, and C. Szegedy, \emph{Explaining and harnessing adversarial examples}, in \emph{Proc. ICLR} (2015), arXiv:1412.6572.
\bibitem{Elsayed2018} G. F. Elsayed et al., \emph{Adversarial examples that fool both computer vision and time-limited humans}, in \emph{Advances in Neural Information Processing Systems 31} (2018), arXiv:1802.08195.
\bibitem{Thompson1982} P. Thompson, \emph{Perceived rate of movement depends on contrast}, Vision Res. \textbf{22}, 377 (1982).
\bibitem{BarlowHill1963} H. B. Barlow and R. M. Hill, \emph{Evidence for a physiological explanation of the waterfall phenomenon and figural after-effects}, Nature \textbf{200}, 1345 (1963).
\bibitem{EaglemanSejnowski2000} D. M. Eagleman and T. J. Sejnowski, \emph{Motion integration and postdiction in visual awareness}, Science \textbf{287}, 2036 (2000).

\bibitem{Feinstein1955mu} B. Feinstein, B. Lindeg{\aa}rd, E. Nyman, and G. Wohlfart, \emph{Morphologic studies of motor units in normal human muscles}, Acta Anat. \textbf{23}, 127 (1955).
\bibitem{MilnerBrown1973rec} H. S. Milner-Brown, R. B. Stein, and R. Yemm, \emph{The orderly recruitment of human motor units during voluntary isometric contractions}, J. Physiol. \textbf{230}, 359 (1973).
\bibitem{Jones2002sdn} K. E. Jones, A. F. de C. Hamilton, and D. M. Wolpert, \emph{Sources of signal-dependent noise during isometric force production}, J. Neurophysiol. \textbf{88}, 1533 (2002).
\bibitem{GomiOsu1998stiff} H. Gomi and R. Osu, \emph{Task-dependent viscoelasticity of human multijoint arm and its spatial characteristics for interaction with environments}, J. Neurosci. \textbf{18}, 8965 (1998).
\bibitem{Jankowska1972Ia} E. Jankowska and W. J. Roberts, \emph{Synaptic actions of single interneurones mediating reciprocal Ia inhibition of motoneurones}, J. Physiol. \textbf{222}, 623 (1972).
\bibitem{Pruszynski2008rapid} J. A. Pruszynski, I. Kurtzer, and S. H. Scott, \emph{Rapid motor responses are appropriately tuned to the metrics of a visuospatial task}, J. Neurophysiol. \textbf{100}, 224 (2008).
\bibitem{SaundersKnill2003vis} J. A. Saunders and D. C. Knill, \emph{Humans use continuous visual feedback from the hand to control fast reaching movements}, Exp. Brain Res. \textbf{152}, 341 (2003).
\bibitem{FlanaganWing1997grip} J. R. Flanagan and A. M. Wing, \emph{The role of internal models in motion planning and control: evidence from grip force adjustments during movements of hand-held loads}, J. Neurosci. \textbf{17}, 1519 (1997).
\bibitem{Matsuoka1985osc} K. Matsuoka, \emph{Sustained oscillations generated by mutually inhibiting neurons with adaptation}, Biol. Cybern. \textbf{52}, 367 (1985).

\bibitem{Treede1995heat} R.-D. Treede, R. A. Meyer, S. N. Raja, and J. N. Campbell, \emph{Evidence for two different heat transduction mechanisms in nociceptive primary afferents innervating monkey skin}, J. Physiol. \textbf{483}, 747 (1995).
\bibitem{Weidner1999cfib} C. Weidner et al., \emph{Functional attributes discriminating mechano-insensitive and mechano-responsive C nociceptors in human skin}, J. Neurosci. \textbf{19}, 10184 (1999).
\bibitem{CampbellLaMotte1983first} J. N. Campbell and R. H. LaMotte, \emph{Latency to detection of first pain}, Brain Res. \textbf{266}, 203 (1983).
\bibitem{LaMotte1982hyper} R. H. LaMotte, J. G. Thalhammer, H. E. Torebj{\"o}rk, and C. J. Robinson, \emph{Peripheral neural mechanisms of cutaneous hyperalgesia following mild injury by heat}, J. Neurosci. \textbf{2}, 765 (1982).
\bibitem{Mancini2014touch} F. Mancini, T. Nash, G. D. Iannetti, and P. Haggard, \emph{Pain relief by touch: a quantitative approach}, Pain \textbf{155}, 635 (2014).
\bibitem{Fields1983rvm} H. L. Fields, J. Bry, I. Hentall, and G. Zorman, \emph{The activity of neurons in the rostral medulla of the rat during withdrawal from noxious heat}, J. Neurosci. \textbf{3}, 2545 (1983).
\bibitem{Crombez1996disrupt} G. Crombez, C. Eccleston, F. Baeyens, and P. Eelen, \emph{The disruptive nature of pain: an experimental investigation}, Behav. Res. Ther. \textbf{34}, 911 (1996).
\bibitem{Schouenborg1995wr} J. Schouenborg, H.-R. Weng, J. Kalliom{\"a}ki, and H. Holmberg, \emph{A survey of spinal dorsal horn neurones encoding the spatial organization of withdrawal reflexes in the rat}, Exp. Brain Res. \textbf{106}, 19 (1995).

\bibitem{Andersen1992cereb} B. B. Andersen, L. Korbo, and B. Pakkenberg, \emph{A quantitative study of the human cerebellum with unbiased stereological techniques}, J. Comp. Neurol. \textbf{326}, 549 (1992).
\bibitem{ShadmehrMussaIvaldi1994} R. Shadmehr and F. A. Mussa-Ivaldi, \emph{Adaptive representation of dynamics during learning of a motor task}, J. Neurosci. \textbf{14}, 3208 (1994).

\bibitem{JoseCollison1970ihr} A. D. Jose and D. Collison, \emph{The normal range and determinants of the intrinsic heart rate in man}, Cardiovasc. Res. \textbf{4}, 160 (1970).
\bibitem{Logothetis1987gbg} D. E. Logothetis, Y. Kurachi, J. Galper, E. J. Neer, and D. E. Clapham, \emph{The $\beta\gamma$ subunits of GTP-binding proteins activate the muscarinic K$^+$ channel in heart}, Nature \textbf{325}, 321 (1987).
\bibitem{KellerWood1984fb} M. E. Keller-Wood and M. F. Dallman, \emph{Corticosteroid inhibition of ACTH secretion}, Endocr. Rev. \textbf{5}, 1 (1984).
\bibitem{Walker2012glc} J. J. Walker et al., \emph{The origin of glucocorticoid hormone oscillations}, PLoS Biol. \textbf{10}, e1001341 (2012).
\bibitem{Wildt1981gnrh} L. Wildt et al., \emph{Frequency and amplitude of gonadotropin-releasing hormone stimulation and gonadotropin secretion in the rhesus monkey}, Endocrinology \textbf{109}, 376 (1981).
\bibitem{Welsh2010scn} D. K. Welsh, J. S. Takahashi, and S. A. Kay, \emph{Suprachiasmatic nucleus: cell autonomy and network properties}, Annu. Rev. Physiol. \textbf{72}, 551 (2010).
\bibitem{Khalsa2003prc} S. B. S. Khalsa, M. E. Jewett, C. Cajochen, and C. A. Czeisler, \emph{A phase response curve to single bright light pulses in human subjects}, J. Physiol. \textbf{549}, 945 (2003).
\bibitem{Sack1992blind} R. L. Sack, A. J. Lewy, M. L. Blood, L. D. Keith, and H. Nakagawa, \emph{Circadian rhythm abnormalities in totally blind people: incidence and clinical significance}, J. Clin. Endocrinol. Metab. \textbf{75}, 127 (1992).

\bibitem{Henze2000extra} D. A. Henze et al., \emph{Intracellular features predicted by extracellular recordings in the hippocampus in vivo}, J. Neurophysiol. \textbf{84}, 390 (2000).
\bibitem{Histed2009stim} M. H. Histed, V. Bonin, and R. C. Reid, \emph{Direct activation of sparse, distributed populations of cortical neurons by electrical microstimulation}, Neuron \textbf{63}, 508 (2009).
\bibitem{Orsborn2014coad} A. L. Orsborn et al., \emph{Closed-loop decoder adaptation shapes neural plasticity for skillful neuroprosthetic control}, Neuron \textbf{82}, 1380 (2014).

\bibitem{HuVervaekeStorm2002} H. Hu, K. Vervaeke, and J. F. Storm, \emph{Two forms of electrical resonance at theta frequencies, generated by M-current, h-current and persistent Na$^+$ current in rat hippocampal pyramidal cells}, J. Physiol. \textbf{545}, 783 (2002).
\bibitem{Mauro1970} A. Mauro, F. Conti, F. Dodge, and R. Schor, \emph{Subthreshold behavior and phenomenological impedance of the squid giant axon}, J. Gen. Physiol. \textbf{55}, 497 (1970).
\bibitem{Aksay2001} E. Aksay et al., \emph{In vivo intracellular recording and perturbation of persistent activity in a neural integrator}, Nat. Neurosci. \textbf{4}, 184 (2001).
\bibitem{Major2004} G. Major et al., \emph{Plasticity and tuning of the time course of analog persistent firing in a neural integrator}, Proc. Natl. Acad. Sci. USA \textbf{101}, 7745 (2004).
\bibitem{Tateno2004} T. Tateno, A. Harsch, and H. P. C. Robinson, \emph{Threshold firing frequency-current relationships of neurons in rat somatosensory cortex: type 1 and type 2 dynamics}, J. Neurophysiol. \textbf{92}, 2283 (2004).
\bibitem{Marder2012} E. Marder, \emph{Neuromodulation of neuronal circuits: back to the future}, Neuron \textbf{76}, 1 (2012).
\bibitem{McGintyHarper1976} D. J. McGinty and R. M. Harper, \emph{Dorsal raphe neurons: depression of firing during sleep in cats}, Brain Res. \textbf{101}, 569 (1976).

\bibitem{Gentet2000} L. J. Gentet, G. J. Stuart, and J. D. Clements, \emph{Direct measurement of specific membrane capacitance in neurons}, Biophys. J. \textbf{79}, 314 (2000).
\bibitem{Borst1995} J. G. G. Borst, F. Helmchen, and B. Sakmann, \emph{Pre- and postsynaptic whole-cell recordings in the medial nucleus of the trapezoid body of the rat}, J. Physiol. \textbf{489}, 825 (1995).
\bibitem{AmirDevor2003} R. Amir and M. Devor, \emph{Electrical excitability of the soma of sensory neurons is required for spike invasion of the soma, but not for through-conduction}, Biophys. J. \textbf{84}, 2181 (2003).
\bibitem{Chadderton2004} P. Chadderton, T. W. Margrie, and M. H\"ausser, \emph{Integration of quanta in cerebellar granule cells during sensory processing}, Nature \textbf{428}, 856 (2004).
\bibitem{LitwinKumar2017} A. Litwin-Kumar et al., \emph{Optimal degrees of synaptic connectivity}, Neuron \textbf{93}, 1153 (2017).
\bibitem{NapperHarvey1988} R. M. A. Napper and R. J. Harvey, \emph{Number of parallel fiber synapses on an individual Purkinje cell in the cerebellum of the rat}, J. Comp. Neurol. \textbf{274}, 168 (1988).
\bibitem{Megias2001} M. Meg\'{\i}as, Z. Emri, T. F. Freund, and A. I. Guly\'as, \emph{Total number and distribution of inhibitory and excitatory synapses on hippocampal CA1 pyramidal cells}, Neuroscience \textbf{102}, 527 (2001).
\bibitem{Polsky2004} A. Polsky, B. W. Mel, and J. Schiller, \emph{Computational subunits in thin dendrites of pyramidal cells}, Nat. Neurosci. \textbf{7}, 621 (2004).
\bibitem{Brunel2004} N. Brunel et al., \emph{Optimal information storage and the distribution of synaptic weights: perceptron versus Purkinje cell}, Neuron \textbf{43}, 745 (2004).

\bibitem{Vyazovskiy2009} V. V. Vyazovskiy et al., \emph{Cortical firing and sleep homeostasis}, Neuron \textbf{63}, 865 (2009).
\bibitem{Daan1984} S. Daan, D. G. Beersma, and A. A. Borb\'ely, \emph{Timing of human sleep: recovery process gated by a circadian pacemaker}, Am. J. Physiol. \textbf{246}, R161 (1984).
\bibitem{Borbely1981} A. A. Borb\'ely et al., \emph{Sleep deprivation: effect on sleep stages and EEG power density in man}, Electroencephalogr. Clin. Neurophysiol. \textbf{51}, 483 (1981).
\bibitem{SteriadeAmzica1993} M. Steriade, F. Amzica, and A. Nu\~nez, \emph{Cholinergic and noradrenergic modulation of the slow ($\approx0.3$ Hz) oscillation in neocortical cells}, J. Neurophysiol. \textbf{70}, 1385 (1993).
\bibitem{vonKrosigk1993} M. von Krosigk, T. Bal, and D. A. McCormick, \emph{Cellular mechanisms of a synchronized oscillation in the thalamus}, Science \textbf{261}, 361 (1993).
\bibitem{LeeWilson2002} A. K. Lee and M. A. Wilson, \emph{Memory of sequential experience in the hippocampus during slow wave sleep}, Neuron \textbf{36}, 1183 (2002).
\bibitem{Euston2007} D. R. Euston, M. Tatsuno, and B. L. McNaughton, \emph{Fast-forward playback of recent memory sequences in prefrontal cortex during sleep}, Science \textbf{318}, 1147 (2007).
\bibitem{Hobson2009} J. A. Hobson, \emph{REM sleep and dreaming: towards a theory of protoconsciousness}, Nat. Rev. Neurosci. \textbf{10}, 803 (2009).

\bibitem{CohenSegal2011} D. Cohen and M. Segal, \emph{Network bursts in hippocampal microcultures are terminated by exhaustion of vesicle pools}, J. Neurophysiol. \textbf{106}, 2314 (2011).
\bibitem{Tsodyks2000} M. Tsodyks, A. Uziel, and H. Markram, \emph{Synchrony generation in recurrent networks with frequency-dependent synapses}, J. Neurosci. \textbf{20}, RC50 (2000).
\bibitem{Benson1994} D. L. Benson, F. H. Watkins, O. Steward, and G. Banker, \emph{Characterization of GABAergic neurons in hippocampal cell cultures}, J. Neurocytol. \textbf{23}, 279 (1994).

\input{supplement/bib_22}
\end{thebibliography}


\end{document}
